% !TeX program = xelatex
% Standalone English volume. Use XeLaTeX until the references stabilize.
\documentclass[a4paper,11pt]{article}
\usepackage[T1]{fontenc}
\usepackage{lmodern}
\usepackage[english]{babel}
\usepackage[left=26mm,right=26mm,top=25mm,bottom=25mm,headheight=15pt,headsep=7mm,footskip=11mm]{geometry}
\usepackage{amsmath,amssymb,mathtools,amsthm}
\newtheorem{proposition}{Proposition}[section]
\newtheorem{lemma}[proposition]{Lemma}
\newtheorem{theorem}[proposition]{Theorem}
\newtheorem{corollary}[proposition]{Corollary}
\newtheorem{maintheorem}[proposition]{Main Theorem}
\theoremstyle{definition}
\newtheorem{definition}[proposition]{Definition}
\theoremstyle{remark}
\newtheorem{remark}[proposition]{Remark}
\makeatletter
\renewcommand*\l@section{\@dottedtocline{1}{0em}{2.4em}}
\renewcommand*\l@subsection{\@dottedtocline{2}{1.45em}{3.45em}}
\renewcommand*\l@subsubsection{\@dottedtocline{3}{3.8em}{4.55em}}
\makeatother
\usepackage{xcolor,fancyhdr,titlesec,needspace,xurl}
\usepackage{longtable,multicol}
\usepackage{hyperref}
\definecolor{ink}{HTML}{1A3542}
\definecolor{muted}{HTML}{4F606A}
\hypersetup{unicode=true,colorlinks=true,linkcolor=ink,filecolor=ink,urlcolor=ink,citecolor=ink,bookmarksnumbered=true,
 pdfdisplaydoctitle=true,pdflang={en},
 pdftitle={The generalized Riemann hypothesis for Dirichlet L-functions. Edition 930, Volume XXII},
 pdfsubject={A natural-language reconstruction of the proposed proof},
 pdfauthor={A mathematical reconstruction}}
\titleformat{\section}{\Large\bfseries\color{ink}}{\thesection}{.65em}{}
\titlespacing*{\section}{0pt}{.8em}{.65em}
\setlength{\parindent}{1.2em}
\linespread{1.045}
\setlength{\parskip}{.22em plus .05em minus .03em}
\setlength{\emergencystretch}{2em}
\setcounter{topnumber}{2}
\setcounter{tocdepth}{3}\setcounter{secnumdepth}{3}
\widowpenalty=4500
\clubpenalty=4500
\pagestyle{fancy}\fancyhf{}
\fancyhead[L]{\small\color{muted}Dirichlet GRH \textperiodcentered\ Volume I}
\fancyhead[R]{\small\color{muted}Dirichlet zeros and power intervals}
\fancyfoot[C]{\small\thepage}
\renewcommand{\headrulewidth}{.3pt}
\newcommand{\lead}[1]{\par\smallskip\noindent{\bfseries\color{ink}#1.}\enspace}
\newcommand{\sour}[2]{\unskip\penalty150\hspace{.35em}{\small\hyperlink{ref#1}{[#2]}}}
\titleformat{\subsection}{\large\bfseries}{\thesubsection}{.6em}{}
\titlespacing*{\subsection}{0pt}{1.15em plus .15em}{.4em}
\setlength{\abovedisplayskip}{9pt plus 2pt minus 3pt}
\setlength{\belowdisplayskip}{9pt plus 2pt minus 3pt}
\newcommand{\RR}{\mathbb R}\newcommand{\CC}{\mathbb C}\newcommand{\QQ}{\mathbb Q}
\newcommand{\ZZ}{\mathbb Z}\newcommand{\ord}{\operatorname{ord}}
\newcommand{\Frob}{\operatorname{Frob}}\newcommand{\Gal}{\operatorname{Gal}}
\newcommand{\len}{\operatorname{length}}\newcommand{\Res}{\operatorname*{Res}}
\newcommand{\wmean}{\operatorname{av}}\newcommand{\Ker}{\operatorname{ker}}
\newcommand{\Lh}{\mathcal L}\newcommand{\Gg}{\mathcal G}
\newcommand{\real}{\operatorname{Re}}\newcommand{\imag}{\operatorname{Im}}
\newcommand{\abs}[1]{\left|#1\right|}\newcommand{\eps}{\varepsilon}
\newcommand{\anchoritem}[2]{\par\Needspace{2.8\baselineskip}\hypertarget{ref#1}{}\noindent{\bfseries\color{ink}#2.}\enspace}
\newcommand{\srcloc}[2]{{\small\nolinkurl{#1.lean}}, line~#2}
\numberwithin{equation}{section}
\raggedbottom

\newcommand{\Sel}{\mathcal S}
\newcommand{\Cl}{\operatorname{Cl}}
\newcommand{\Aut}{\operatorname{Aut}}
\newcommand{\GL}{\operatorname{GL}}
\newcommand{\OO}{\mathcal O}
\hypersetup{bookmarksdepth=3}


\addto\captionsenglish{\renewcommand{\partname}{Volume}}


\makeatletter
\expandafter\def\csname r@cor:actual50\endcsname{{6.31}{64}{}{proposition.6.31}{Dirichlet-GRH-v930-Volume-I.pdf}}
\expandafter\def\csname r@cor:actualgram535\endcsname{{113.8}{89}{}{proposition.113.8}{Dirichlet-GRH-v930-Volume-XXI.pdf}}
\expandafter\def\csname r@cor:rel155\endcsname{{106.21}{73}{}{proposition.106.21}{Dirichlet-GRH-v930-Volume-XX.pdf}}
\expandafter\def\csname r@cor:trace235\endcsname{{110.7}{43}{}{proposition.110.7}{Dirichlet-GRH-v930-Volume-XXI.pdf}}
\expandafter\def\csname r@eq:cross132\endcsname{{86.2}{5}{}{equation.86.2}{Dirichlet-GRH-v930-Volume-XVII.pdf}}
\expandafter\def\csname r@eq:gap\endcsname{{6.3}{26}{}{equation.6.3}{Dirichlet-GRH-v930-Volume-I.pdf}}
\expandafter\def\csname r@eq:length12\endcsname{{{L}}{30}{}{AMS.3}{Dirichlet-GRH-v930-Volume-I.pdf}}
\expandafter\def\csname r@eq:orb133\endcsname{{98.1}{3}{}{equation.98.1}{Dirichlet-GRH-v930-Volume-XIX.pdf}}
\expandafter\def\csname r@eq:period94\endcsname{{121.7}{51}{}{equation.121.7}{Dirichlet-GRH-v930-Volume-XXIII.pdf}}
\expandafter\def\csname r@eq:seed127\endcsname{{92.16}{20}{}{equation.92.16}{Dirichlet-GRH-v930-Volume-XVIII.pdf}}
\expandafter\def\csname r@eq:shear225\endcsname{{109.12}{19}{}{equation.109.12}{Dirichlet-GRH-v930-Volume-XXI.pdf}}
\expandafter\def\csname r@eq:success135\endcsname{{86.4}{8}{}{equation.86.4}{Dirichlet-GRH-v930-Volume-XVII.pdf}}
\expandafter\def\csname r@lem:completedcolumn535\endcsname{{113.6}{88}{}{proposition.113.6}{Dirichlet-GRH-v930-Volume-XXI.pdf}}
\expandafter\def\csname r@lem:conductorstabilizer532\endcsname{{96.23}{89}{}{proposition.96.23}{Dirichlet-GRH-v930-Volume-XVIII.pdf}}
\expandafter\def\csname r@lem:coordinates225\endcsname{{109.20}{18}{}{proposition.109.20}{Dirichlet-GRH-v930-Volume-XXI.pdf}}
\expandafter\def\csname r@lem:fixedflags302\endcsname{{112.1}{74}{}{proposition.112.1}{Dirichlet-GRH-v930-Volume-XXI.pdf}}
\expandafter\def\csname r@lem:flagdeterminant534\endcsname{{113.1}{84}{}{proposition.113.1}{Dirichlet-GRH-v930-Volume-XXI.pdf}}
\expandafter\def\csname r@lem:forms203\endcsname{{110.10}{46}{}{proposition.110.10}{Dirichlet-GRH-v930-Volume-XXI.pdf}}
\expandafter\def\csname r@lem:markedtransport533\endcsname{{96.27}{92}{}{proposition.96.27}{Dirichlet-GRH-v930-Volume-XVIII.pdf}}
\expandafter\def\csname r@lem:translate163\endcsname{{86.19}{16}{}{proposition.86.19}{Dirichlet-GRH-v930-Volume-XVII.pdf}}
\expandafter\def\csname r@lem:weyl209\endcsname{{93.8}{27}{}{proposition.93.8}{Dirichlet-GRH-v930-Volume-XVIII.pdf}}
\expandafter\def\csname r@lem:whittakermatrix526\endcsname{{96.2}{74}{}{proposition.96.2}{Dirichlet-GRH-v930-Volume-XVIII.pdf}}
\expandafter\def\csname r@lem:zero133\endcsname{{98.1}{3}{}{proposition.98.1}{Dirichlet-GRH-v930-Volume-XIX.pdf}}
\expandafter\def\csname r@prop:amplitudes126\endcsname{{92.22}{17}{}{proposition.92.22}{Dirichlet-GRH-v930-Volume-XVIII.pdf}}
\expandafter\def\csname r@prop:class236\endcsname{{109.40}{34}{}{proposition.109.40}{Dirichlet-GRH-v930-Volume-XXI.pdf}}
\expandafter\def\csname r@prop:field12\endcsname{{6.4}{27}{}{proposition.6.4}{Dirichlet-GRH-v930-Volume-I.pdf}}
\expandafter\def\csname r@prop:hyperplanetrace534\endcsname{{113.4}{86}{}{proposition.113.4}{Dirichlet-GRH-v930-Volume-XXI.pdf}}
\expandafter\def\csname r@prop:partnerweights531\endcsname{{96.20}{87}{}{proposition.96.20}{Dirichlet-GRH-v930-Volume-XVIII.pdf}}
\expandafter\def\csname r@prop:repeatedtransport534\endcsname{{113.5}{87}{}{proposition.113.5}{Dirichlet-GRH-v930-Volume-XXI.pdf}}
\expandafter\def\csname r@prop:retainedtest535\endcsname{{113.9}{90}{}{proposition.113.9}{Dirichlet-GRH-v930-Volume-XXI.pdf}}
\expandafter\def\csname r@prop:scale151\endcsname{{106.13}{67}{}{proposition.106.13}{Dirichlet-GRH-v930-Volume-XX.pdf}}
\expandafter\def\csname r@prop:weight163\endcsname{{86.20}{17}{}{proposition.86.20}{Dirichlet-GRH-v930-Volume-XVII.pdf}}
\expandafter\def\csname r@prop:weyl235\endcsname{{110.6}{42}{}{proposition.110.6}{Dirichlet-GRH-v930-Volume-XXI.pdf}}
\expandafter\def\csname r@sec:geom11\endcsname{{6.5}{30}{}{subsection.6.5}{Dirichlet-GRH-v930-Volume-I.pdf}}
\expandafter\def\csname r@sec:missing\endcsname{{6.6}{32}{}{subsection.6.6}{Dirichlet-GRH-v930-Volume-I.pdf}}
\expandafter\def\csname r@sec:mixedordinary541\endcsname{{113.6}{100}{}{subsection.113.6}{Dirichlet-GRH-v930-Volume-XXI.pdf}}
\expandafter\def\csname r@sec:relative155\endcsname{{106.15}{71}{}{subsection.106.15}{Dirichlet-GRH-v930-Volume-XX.pdf}}
\expandafter\def\csname r@seriesguide\endcsname{{1.1}{4}{}{subsection.1.1}{Dirichlet-GRH-v930-Volume-I.pdf}}
\expandafter\def\csname r@thm:char125\endcsname{{92.18}{15}{}{proposition.92.18}{Dirichlet-GRH-v930-Volume-XVIII.pdf}}
\expandafter\def\csname r@thm:columns235\endcsname{{109.38}{32}{}{proposition.109.38}{Dirichlet-GRH-v930-Volume-XXI.pdf}}
\expandafter\def\csname r@thm:common111\endcsname{{9.13}{118}{}{proposition.9.13}{Dirichlet-GRH-v930-Volume-I.pdf}}
\expandafter\def\csname r@thm:common133\endcsname{{98.2}{4}{}{proposition.98.2}{Dirichlet-GRH-v930-Volume-XIX.pdf}}
\expandafter\def\csname r@thm:conductor232\endcsname{{19.15}{12}{}{proposition.19.15}{Dirichlet-GRH-v930-Volume-IV.pdf}}
\expandafter\def\csname r@thm:cross132\endcsname{{86.3}{4}{}{proposition.86.3}{Dirichlet-GRH-v930-Volume-XVII.pdf}}
\expandafter\def\csname r@thm:crossclasses230\endcsname{{110.15}{50}{}{proposition.110.15}{Dirichlet-GRH-v930-Volume-XXI.pdf}}
\expandafter\def\csname r@thm:everydiagonal536\endcsname{{113.10}{91}{}{proposition.113.10}{Dirichlet-GRH-v930-Volume-XXI.pdf}}
\expandafter\def\csname r@thm:flagisometry544\endcsname{{113.24}{103}{}{proposition.113.24}{Dirichlet-GRH-v930-Volume-XXI.pdf}}
\expandafter\def\csname r@thm:identitycoefficient551\endcsname{{113.26}{105}{}{proposition.113.26}{Dirichlet-GRH-v930-Volume-XXI.pdf}}
\expandafter\def\csname r@thm:localmass156\endcsname{{87.24}{48}{}{proposition.87.24}{Dirichlet-GRH-v930-Volume-XVII.pdf}}
\expandafter\def\csname r@thm:mixedordinary541\endcsname{{113.22}{101}{}{proposition.113.22}{Dirichlet-GRH-v930-Volume-XXI.pdf}}
\expandafter\def\csname r@thm:mixedtransport535\endcsname{{113.7}{89}{}{proposition.113.7}{Dirichlet-GRH-v930-Volume-XXI.pdf}}
\expandafter\def\csname r@thm:moment286\endcsname{{105.3}{34}{}{proposition.105.3}{Dirichlet-GRH-v930-Volume-XX.pdf}}
\expandafter\def\csname r@thm:orbitaverage532\endcsname{{96.24}{90}{}{proposition.96.24}{Dirichlet-GRH-v930-Volume-XVIII.pdf}}
\expandafter\def\csname r@thm:orbits204\endcsname{{110.16}{51}{}{proposition.110.16}{Dirichlet-GRH-v930-Volume-XXI.pdf}}
\expandafter\def\csname r@thm:powers236\endcsname{{110.21}{56}{}{proposition.110.21}{Dirichlet-GRH-v930-Volume-XXI.pdf}}
\expandafter\def\csname r@thm:prime209\endcsname{{93.9}{28}{}{proposition.93.9}{Dirichlet-GRH-v930-Volume-XVIII.pdf}}
\expandafter\def\csname r@thm:rel155\endcsname{{106.20}{72}{}{proposition.106.20}{Dirichlet-GRH-v930-Volume-XX.pdf}}
\expandafter\def\csname r@thm:success135\endcsname{{86.7}{8}{}{proposition.86.7}{Dirichlet-GRH-v930-Volume-XVII.pdf}}
\expandafter\def\csname r@thm:suffix237\endcsname{{110.9}{45}{}{proposition.110.9}{Dirichlet-GRH-v930-Volume-XXI.pdf}}
\expandafter\def\csname r@thm:tensor249\endcsname{{106.9}{62}{}{proposition.106.9}{Dirichlet-GRH-v930-Volume-XX.pdf}}
\expandafter\def\csname r@thm:totaldiagonal540\endcsname{{113.19}{99}{}{proposition.113.19}{Dirichlet-GRH-v930-Volume-XXI.pdf}}
\makeatother
\begin{document}
\setcounter{part}{21}
\fancyhead[L]{\small\color{muted}Dirichlet GRH \textperiodcentered\ Volume XXII}
\fancyhead[R]{\small\color{muted}Signed orbital integrals}
\thispagestyle{plain}\begin{center}
{\LARGE\bfseries The generalized Riemann hypothesis\\[2pt]for Dirichlet \(L\)-functions}\\[6pt]
{\small 27 November 2026 \quad\textperiodcentered\quad Edition 930}\end{center}
\part{Signed orbital integrals on the marked rows}\label{part:XXII}
\begin{abstract}\noindent The common flag lattice controls denominators at the two terminal frequencies. We truncate the Gram determinant sum by height, compute the local integrals of its transverse lifts and distinguish rational classes in the signed period.\end{abstract}
\setcounter{section}{114}
\noindent The \hyperref[seriesguide]{guide in Volume~I} lists all twenty-seven volumes.
\begingroup\small\setlength{\parskip}{0pt}\tableofcontents\endgroup
\clearpage

\section{Terminal frequencies and common flag lattices}\label{nav:chapter-ordinary-terminal-lattices}\label{nav:volume-XXII}

\subsection{A uniform lattice count on the norm-one rows}
\label{sec:unipotentpacking552}

Let \(E=E_x\), \(F=E^+\), \(f=[F:\QQ]\) and \(Q=g_1g_2\) be
the marked level on the intervals of Section~\ref{sec:missing}.
We treat any ordinary diagonal row of terminal norm \(-1\).
Use its primitive integral frame \(v_1=e_1,\ldots,v_s=e_s\),
\(v_{s+1}=z\), from Lemma~\ref{lem:coordinates225}.
Write \(n=\dim_E V,\ d=s+1\). Its integral Witt decomposition
has hyperbolic partners \(f_1,\ldots,f_s\) and middle lattice
\(J=\mathcal O_E z\perp L\). Indeed orthogonal projection to
\(E z\) is integral because \(h(z,z)=-1\). The lattice \(L\)
is positive definite at all infinite places; put \(k=\dim_E EL\).

At each real place use the positive Hermitian norm in which the
hyperbolic basis vectors and \(z\) are orthonormal, the other
summand is perpendicular to them, and \(L\) has its given
positive Hermitian form. This specifies the real norms used in
the estimates below. The original form in these norms has matrix
a unitary involution, so \(w(g)=\|g\|_{\rm op}=\|g^{-1}\|_{\rm op}
\ge1\) for an isometry. The measures and the section \(W\) are
those of the row. For its real decorated stabilizer \(B_\tau\)
put \(\nu_\tau(y)=\inf_{b\in B_\tau}w(by_\tau)\) and
\(V(y)=\prod_\tau\nu_\tau(y)\). Write \(N_\Lambda\) for the
integral unipotent radical in \(\Gamma=\operatorname{Aut}(\Lambda,h)\).

\begin{lemma}[Counting lattice points through a filtration]
\leavevmode\label{lem:gradedpacking552}
Suppose a Hermitian lattice over \(\mathcal O_E\), as a module,
is a direct sum of finitely many blocks in a specified order.
Give each block a positive definite integral Hermitian form and
the direct sum its orthogonal norm at every infinite place.
Let \(b_\tau\) preserve the filtration by partial sums, acting
unitarily on each graded block. For an affine \(E\)-subspace
of dimension \(j\) and positive real radii \(R_\tau\), the
number of its lattice points \(u\) satisfying
\(\|b_\tau u_\tau\|\le R_\tau\) at every place is at most
\[
                \left(1+2(\textstyle\prod_\tau R_\tau)^{1/f}\right)^{2fj}.
\]
The same estimate holds for exterior powers with their induced
block decompositions and norms.
\end{lemma}
\begin{proof}
In one block a nonzero lattice vector \(v\) satisfies
\(\prod_\tau\|v_\tau\|\ge1\), since its squared Hermitian length
is a nonzero totally positive algebraic integer in \(F\).
Consider points in a fixed affine subspace of dimension \(j'\)
in this block, with arbitrary centers for the balls of radii
\(R_\tau\). Project the centers orthogonally onto the corresponding
complex affine subspaces. After rescaling by \(R_\tau\), the
points lie in a unit ball for the maximum of the place norms.
Two distinct points have distance at least
\((\prod_\tau R_\tau)^{-1/f}\). Disjoint open balls of half
this radius fit inside the unit ball enlarged by that radius.
Their volume ratio in real dimension \(2fj'\) gives the bound
in the statement with \(j'\).

Now fix block coordinates of \(u\) starting at the last block.
After the later coordinates have been fixed, the particular
block of \(b_\tau u_\tau\) in question is a unitary transform of
the current coordinate plus a known translation. The possible
coordinates lie in an affine subspace whose dimension is the
drop in dimension on specifying this block. This drop depends
only on the linear direction of the original affine space. Its
sum over all blocks is \(j\). Multiply the estimates just proved
for these successive choices.

In an exterior power, decompose by the number of factors taken
from each original block. Order the resulting blocks by comparing
these numbers from the last original block backwards. The
transform still preserves the filtration and acts unitarily
on its graded pieces, since replacement of a factor by one in
an earlier block lowers its index. Each piece is a tensor of
exterior powers of the integral blocks. Its induced form is
positive and integral by the determinant formula for exterior
pairings. This verifies all hypotheses for the same argument.
\end{proof}

\begin{theorem}[Height separation and counting on every norm-one row]
\leavevmode\label{thm:uniformheight552}
For each admissible \(\gamma\in\Gamma\) outside the rational
decorated stabilizer and every \(y\in G_\infty\) one has
\[
                  Q^f\le2^f V(y)V(\gamma^{-1}y).
\]
Set \(D=\sum_{j=1}^d j(n-j+1)\) and
\[
       C_{f,n,d}=3^{(2fk)^2}
             \prod_{j=1}^d(1+2\cdot4^j)^{2f(n-j+1)} .
\]
For any integers \(m_\tau\ge0\), the number of cosets
\(\delta\in N_\Lambda\backslash\Gamma\) with
\(2^{m_\tau}\le\nu_\tau(\delta y)<2^{m_\tau+1}\) everywhere
is at most \(C_{f,n,d}\,2^{4D\sum_\tau m_\tau}\).
\end{theorem}
\begin{proof}
Order the blocks as the successive \(e_i\)-lines, the line
\(E z\), the positive space \(EL\), and the successive
\(f_s,\ldots,f_1\)-lines. The decorated stabilizer acts
triangularly in these blocks with unitary graded actions.
Indeed it fixes each \(e_i\) modulo earlier vectors, and \(z\)
modulo their full span. In the middle it fixes the negative
vector and acts isometrically on its positive perpendicular.
Preservation of the hyperbolic pairing gives the reversed
dual action on the last block lines with diagonal one.

Let \(X_i\) be the span of the first \(i\) isotropic vectors,
with \(X_0=0\). Take the first \(i\le d\) such that
\(\gamma v_i-v_i\notin X_{i-1}\). Let \(p\) be the last
nonzero block component of this difference, necessarily beyond
\(X_{i-1}\). Admissibility at the primes over \(Q\) gives
\(\gamma v_i-v_i\in X_{i-1}+Q\Lambda_{\mathfrak w}\) there.
The integral splitting shows \(p\) belongs globally to \(Q\)
times its block lattice. Consequently
\(\prod_\tau\|p_\tau\|\ge Q^f\) by positivity and integrality
in that block, including the sign-corrected \(z\)-line.

For any \(b_1,b_2\in B_\tau\) compare \(p_\tau\) with the
same block of \(b_1\gamma b_2^{-1}v_i-v_i\). The vector
\(b_2^{-1}v_i-v_i\) is in \(X_{i-1}\), which \(\gamma\)
carries into \(X_{i-1}\) by the choice of \(i\). Also
\(b_1v_i-v_i\in X_{i-1}\). Thus the block in question is
the unitary graded transform of \(p_\tau\). Its norm is
at most \(1+w(b_1\gamma b_2^{-1})\), since \(v_i\) is a
unit vector. Factor the isometry as
\((b_1 y_\tau)(b_2\gamma^{-1}y_\tau)^{-1}\) and take the
two infima independently. This gives
\(\|p_\tau\|\le2\nu_\tau(y)\nu_\tau(\gamma^{-1}y)\).
Multiplication proves the height separation.

For the count put \(\Phi_j=v_1\wedge\cdots\wedge v_j\) for
\(1\le j\le d\). All \(\Phi_j\) have norm one and are fixed
by the real decorated stabilizer. A group element fixing every
\(\Phi_j\) must fix \(v_j\) modulo its predecessors: equality
of the wedge with the nonzero fixed preceding wedge gives
this assertion inductively. The integral stabilizer of the
tuple, modulo \(N_\Lambda\), embeds into
\(\operatorname{Aut}(L,h)\) by its middle action. This group
is finite since it preserves the positive integral trace form.
Its reduction on the underlying \(\ZZ\)-lattice of rank \(2fk\)
is injective modulo three. For if \(1+3^b A\), \(b\ge1\), has
a unit entry in \(A\), raising to a prime different from
three preserves that first nonzero level. Cubing raises it
by exactly one by the binomial expansion. Thus the kernel has
no element of prime order. Its finite subgroup is trivial,
giving the order bound \(3^{(2fk)^2}\).

In a nonempty box choose one representative \(\delta_0\).
Choose \(b_{0,\tau}\in B_\tau\) with
\(w(b_{0,\tau}\delta_0 y_\tau)<2^{m_\tau+1}\).
For any other representative in the box the integral tuple
\[
                Z_j=\delta_0\delta^{-1}\Phi_j
                        \in\bigwedge\nolimits^j\Lambda
\]
satisfies \(\|b_{0,\tau}Z_{j,\tau}\|\le 4^j2^{2jm_\tau}\).
In fact write \(b_{0,\tau}Z_j\) as the product of the exterior
actions of \(b_{0,\tau}\delta_0 y_\tau\) and
\((b_{\delta,\tau}\delta y_\tau)^{-1}\) on \(\Phi_j\),
where \(b_{\delta,\tau}\) realizes the same strict bound.
The exterior norm is bounded by the \(j\)-th power of the
ordinary operator norm.

On fixing the preceding tuple, the possibilities for \(Z_j\)
lie in the rational subspace \(Z_{j-1}\wedge V\) of dimension
\(n-j+1\), with \(Z_0=1\). This follows because the preceding
wedge is decomposable and nonzero. Apply
Lemma~\ref{lem:gradedpacking552} with those radii in the
exterior lattice. For this step the bound is at most
\[
 (1+2\cdot4^j)^{2f(n-j+1)}
                  2^{4j(n-j+1)\sum_\tau m_\tau}.
\]
Multiply over \(j\). Equal tuples give representatives differing
on the left by an integral element fixing every \(\Phi_j\).
The stabilizer order just proved bounds the number of such
left \(N_\Lambda\)-cosets and completes the count.
\end{proof}

\begin{theorem}[The arithmetic factor in the full ordinary diagonal error]
\leavevmode\label{thm:ordinaryunit552}
Set \(L_0=8(f+1)(n+1)^3\) and \(a=(L_0-4fn^2)/2\).
On every norm-one diagonal row, for all sufficiently large \(x\)
on the common intervals, one has
\[
 \frac{\sum_i\|Z_i-Z_i^B\|_{\rm HS}}{\lambda}
 \ \le\ Q^{-fa/2}\,d_\rho\,
 \frac{\displaystyle
       \sup_g V(g)^{L_0}\|W(g)\|\,
       \int_{B_\infty\backslash G_\infty}V(g)^a\|W(g)\|\,d\dot g}
      {\displaystyle
       \int_{B_\infty\backslash G_\infty}\|W(g)\|^2\,d\dot g}.
\]
Here \(d_\rho\) is the column dimension, \(\lambda\) is the
absolute value of the identity coefficient, and \(W\ne0\) is
the ordinary section. The factor \(Q^{-fa/2}\) is uniform in
the lattice class and the word. The estimate includes the
signed error at the marked identity index.
\end{theorem}
\begin{proof}
Put \(U=\sup V^{L_0}\|W\|\) and \(F_a=V^a\|W\|\).
Both the supremum and the integral in the assertion are
finite by Proposition~\ref{prop:weight163}; changing to the
specified compatible norm preserves finiteness since matrix
norms at a fixed row are equivalent. The box estimate of
Theorem~\ref{thm:uniformheight552} gives
\[
 \sup_y\sum_{\delta\in N_\Lambda\backslash\Gamma} F_a(\delta y)
 \le U C_{f,n,d}\big(1-2^{-(L_0-a-4D)}\big)^{-f}.
\]
Indeed the summand is bounded by
\(U2^{-(L_0-a)\sum m_\tau}\) in each box. Summing the geometric
series is valid since \(D\le n(n+1)(n+2)/6\) and
\(L_0-a-4D>1\).

For a nonstabilizing admissible isometry the height separation
in that theorem bounds its unweighted integrand in norm by
\((2/Q)^{fa}F_a(y)F_a(\gamma^{-1}y)\). Sum over double cosets
and unfold the nonnegative majorant as in
Theorem~\ref{thm:totaldiagonal540}. The result is bounded by
the last display times
\(|c_{\rm row}|(2/Q)^{fa}\int_{N_\Lambda\backslash G_\infty}F_a\).
This already bounds the sum of the norms before grouping by
marked index. Now \(\log C_{f,n,d}\ll f^2 k^2+fD\ll fa\)
with absolute constants, and the last geometric-series factor
is at most \(2^f\). Since \(\log Q\) tends to infinity on the
intervals, their product times \((2/Q)^{fa}\) is at most
\(Q^{-fa/2}\), uniformly in all the integers here.

Finally Theorem~\ref{thm:flagisometry544} gives
\(\lambda=|c_{\rm row}|\int_{N_\Lambda\backslash G_\infty}
\|W\|^2/d_\rho\). In the ratio of the two integrals the volume
of \(N_\Lambda\backslash B_\infty\) cancels because their
integrands are left \(B_\infty\)-invariant. This gives the
displayed estimate with the original finite level. The
remaining ratio concerns the sizes of the actual real
sections as the types vary. The norm-\(T\) middle lattice
also has a gluing index at \(T\), so its integral splitting
requires a separate calculation.
\end{proof}


\subsection{Both terminal norms in the ordinary period}
\label{sec:bothordinary553}

Let \(Q=g_1g_2\), let \(T=\varpi\bar\varpi\) be the terminal
prime, and use the polynomial-conductor fields \(E_x\) and
intervals of Theorem~\ref{thm:columns235}. Write \(F=E_x^+\),
\(f=[F:\QQ]\), and \(R_E=|\operatorname{disc}(E_x)|^{1/(2f)}\).
Keep all norm-one rows and the norm-\(T\) rows satisfying the
primitive hyperplane bound of that theorem for a fixed
constant \(A>0\). At norm \(T\) the bound retains a fraction
at least \(1-\exp(-Af(k-1))\) of the refined coordinates,
including in each stipulated terminal character. In each row
write \(\Lambda=X_{\mathcal O}\oplus J\oplus X_{\mathcal O}^*\)
as in Lemma~\ref{lem:coordinates225}; the ordered outer frame
has length \(s\), the middle lattice contains \(z\) of norm
\(-T_j\), \(T_j\in\{1,T\}\), and \(L=J\cap z^\perp\) has
positive rank \(k\). Put \(n=2s+1+k\). In this section the
outer line lattices are generated by the distinguished vectors
of the primitive frame model of Lemma~\ref{lem:coordinates225}.
The extent of the following estimate is determined by this
integral hypothesis.

On the middle space use the positive Hermitian form obtained by
reversing the sign on \(Ez\):
\[
        h_z^+(u,v)=h(u,v)+2h(u,z)h(z,v)/T_j .
\]
Take the outer hyperbolic basis vectors orthonormal and perpendicular
to this space for the positive norm. This defines the adapted
height \(V_j\) of row \(j\) as in Section~\ref{sec:unipotentpacking552}.
For mixed rows one can view \(y\) as a real isometry from a
reference space to the target and \(\gamma^{-1}y\) as an isometry
to the source. Use the corresponding positive norms in their
operator norms. Inversion preserves these norms since the
original Hermitian forms are unitary involutions in both spaces.

\begin{lemma}[The product minimum on the retained middle lattices]
\leavevmode\label{lem:middleminimum553}
For all sufficiently large \(x\), uniformly in the retained rows,
every \(0\ne u\in J\) satisfies
\[
                   \prod_\tau\|u_\tau\|_{h_z^+}\ge1 .
\]
At norm \(T\), every \(0\ne p\in L^\sharp\) satisfies the
stronger positive-space inequality
\[
       \prod_\tau\|p_\tau\|^2>
                    \bigl(kR_E/(2\pi\exp(A+4))\bigr)^f .
\]
\end{lemma}
\begin{proof}
For a positive lattice \(M\) write
\(\delta(M)=[M^\sharp:M]^{1/2}\) if \(M\) is integral, and
extend this definition to fractional lattices by covolume.
In the product of the complex embeddings its Euclidean
covolume is \(c_E^{\operatorname{rank}M}\delta(M)\), where
\(c_E\) is the covolume of \(\mathcal O_E\). This follows from
the determinant of the Hermitian Gram matrix for a free lattice;
passing to a finite-index free sublattice proves the same
formula for a projective lattice. At norm \(T\) we have
\(\delta(L)=T^f\) by Lemma~\ref{lem:completedcolumn535}.

For a nonzero vector \(p\in L^\sharp\), the hyperplane
\(H=L\cap p^\perp\) is primitive and
\(\mathfrak b=h(p,L)\) is a nonzero integral ideal. The
orthogonal projection of \(L\) onto \(Ep\) has lattice
\(\mathfrak b p/h(p,p)\). The covolume of \(L\) is the
covolume of \(H\) times that of this projected lattice, by a
block triangular determinant on real lattice bases. The field
constants cancel with the ranks, giving
\[
       \prod_\tau\|p_\tau\|^2
          = N_{E/\QQ}(\mathfrak b)\,\frac{\delta(H)}{\delta(L)} .
\]
Apply the retained bound of Theorem~\ref{thm:columns235} to
\(H\) and use \(N(\mathfrak b)\ge1\). This proves the second
assertion.

For \(u\in J\), its projection \(p\) onto \(EL\) lies in
\(L^\sharp\) because \(h(p,L)=h(u,L)\subset\mathcal O_E\).
The positive norm of \(u\) dominates that of \(p\) at every
place. As \(k\) tends to infinity and \(R_E\ge1\), the second
assertion suffices if \(p\ne0\). If \(p=0\), primitivity gives
\(u\in Ez\cap J=\mathcal O_Ez\); its norm product is at least
\(T^{f/2}\). At norm one the positive form is already integral
on the direct sum in Section~\ref{sec:unipotentpacking552},
so its norm-product bound follows directly.
\end{proof}

\begin{proposition}[Vanishing when only the terminal character changes]
\leavevmode\label{prop:terminalvanishing553}
Let \(\gamma\) be an integral isometry from a source row to
a target row. Suppose it transports each outer isotropic
vector to its target modulo the preceding span. If it fails
to transport the terminal vector modulo the isotropic span,
its ordinary orbital integral is zero. This holds also when
the two terminal norms differ.
\end{proposition}
\begin{proof}
The outer flags agree after transport, so their real unipotent
radicals agree and \(\gamma\) identifies their integral
unipotent subgroups. In fact these radicals are defined by
acting trivially on the graded spaces of the outer flag and
its orthogonal flag, including the full middle quotient.
They are independent of the last vector. Write \(N_\infty\)
and \(N_\Lambda\) for the common groups in target coordinates.

In the last simple coordinate of \(N_\infty\) we can prescribe
an arbitrary complex linear functional \(\phi\) on the
middle quotient at any real place, with all other places
trivial. Its action fixes the outer span pointwise and sends
a middle vector \(w\) to \(w+\phi(w)e_s\) modulo the preceding
span. Such elements are the real shears of
\eqref{eq:shear225}, with the adjoint block and quadratic
correction supplied by the isometry equation. The earlier
simple characters vanish on them. In the integrand the two
terminal characters give the additive character at
\(\phi(\bar z_t-\bar\gamma\bar z_s)\), up to reversing the sign
by the pairing convention. Here bars on vectors and the map
denote their classes in the middle quotient. This description
incorporates the frequency since the last vector in
\eqref{eq:period94} is \(\beta e\). Transport of the last
outer graded line is the identity, so introduces no scalar.
The difference of the terminal classes is nonzero. The
nontrivial real additive character composed with the trace
from \(\CC\) is therefore nontrivial on these shears.

The real integration domain is \(N_\Lambda\backslash G_\infty\).
The product of the two real sections transforms on the left
by the ratio of their unipotent characters. For an admissible
orbital term this ratio descends to the unipotent quotient:
the finite factors on its common rational stabilizer cancel,
as in Proposition~\ref{prop:partnerweights531}. Equivalently
one can first average over the common adelic unipotent
quotient in the unfolding. Integrating over
\(N_\Lambda\backslash N_\infty\) in the real formula gives
zero, since right translation by an element on which this
ratio is nontrivial preserves the quotient measure. The
integrals converge by Proposition~\ref{prop:weight163}.
This proves the vanishing.
\end{proof}

\begin{theorem}[The arithmetic error on the retained primitive-frame rows]
\leavevmode\label{thm:bothordinary553}
Take \(L_0=8(f+1)(n+1)^3,\ a=(L_0-4fn^2)/2\), and set
\[
 U_j=\sup V_j^{L_0}\|W_j\|,\qquad
 I_j=\int_{N_j\backslash G_\infty}V_j^a\|W_j\| .
\]
For every pair of retained rows, with either terminal frequency
at either end, and for all sufficiently large \(x\) on the
specified intervals,
\[
       \sum_{i\in G}\|Z_i^{tu}-B_i^{tu}\|_{\rm HS}
       \ \le\ |c_{t,u}|\,Q^{-fa/2}
                        \min\{I_t U_u,I_u U_t\}.
\]
Here \(u,t\) are source and target row indices. All the orbital
coefficients have their original finite marking, signs, measures
and real sections.
\end{theorem}
\begin{proof}
When the integral lattices are not isometric the arrays are
zero. For an isometry with a nonzero term outside \(B^{tu}\),
Proposition~\ref{prop:terminalvanishing553} forces a first
failure at some outer index \(i\le s\). The corresponding two
outer vectors have norm one in their respective adapted norms.
Use the block filtration with the \(e_j\)-lines first, then
the single middle block \(J\), then the \(f_j\)-lines in
reverse order. The target stabilizer acts isometrically on
the middle block for \(h_{z_t}^+\). Take the last nonzero
block of \(\gamma e_{u,i}-e_{t,i}\). It lies beyond the
preceding target span and belongs to \(Q\) times its lattice
by the marking. In a line its norm product is at least \(Q^f\)
by integrality; in \(J\) this follows from
Lemma~\ref{lem:middleminimum553}. Multiplying \(\gamma\) on
the left and right by arbitrary real decorated stabilizers
preserves this block norm, exactly as in
Theorem~\ref{thm:uniformheight552}: the earlier source vectors
map into the earlier target span. Taking infima in the two
real operator norms gives
\[
           Q^f\le 2^f V_t(y)V_u(\gamma^{-1}y).
\]

We record the cost of the middle lattice in the independent
coset count. The form \(T_j h_z^+\) is integral on \(J\).
On every graded block of the \(j'\)-th exterior lattice the
positive pairing becomes integral on multiplication by
\(T^{j'}\), since \(T_j=1\) or \(T\). Hence the separation
estimate in Lemma~\ref{lem:gradedpacking552} still applies
with its radii multiplied by \(T^{j'/2}\). The outer wedges
have norm one and the last wedge has norm \(\sqrt{T_j}\).
Thus the packing proof of Theorem~\ref{thm:uniformheight552}
applies with the new constant
\[
 C_T=3^{(2fk)^2}
   \prod_{j'=1}^{s+1}(1+2\cdot4^{j'}T^{(j'+1)/2})^{2f(n-j'+1)}
\]
and the same box exponent \(4D\), where
\(D=\sum_{j'=1}^{s+1}j'(n-j'+1)\). The stabilizer acts on
the positive lattice \(L\), so the same bound modulo three
applies at both frequencies. For a mixed pair the cosets to
count are integral isometries from one fixed lattice to the
other modulo the target unipotent subgroup. In each nonempty
box \(\delta_0\delta^{-1}\) is still an automorphism of the
target lattice. Its exterior tuples therefore lie in exactly
the lattices used in that proof. This gives the count with
the same constant for either orientation.

Now \(\log C_T\ll f^2 k^2+fn(s+1)^2\log(2T)\), and \(a\gg fn^3\).
On the intervals of Theorem~\ref{thm:columns235} one has
\(\log T=O(m x^\varepsilon)\), \(s+1\le m=O(\sqrt x)\),
\(n\ge k\asymp x^{9/2+10\varepsilon}\), and
\(\log Q\ge16x^\varepsilon\). Thus \(\log C_T=o(fa\log Q)\):
the first term divided by \(fa\log Q\) is
\(O(1/(n\log Q))\), and the second is \(O(m^3/(f n^2))\).
Also \(L_0-a-4D>1\). The factor from the geometric series in
the box indices is at most \(2^f\), and
\[
              C_T\,2^f(2/Q)^{fa}\le Q^{-fa/2}
\]
eventually, independently of the pair of rows. Unfolding the
majorant with the displayed height separation gives this factor
times \(|c_{t,u}|I_tU_u\). Inverting isometries and changing
variables gives the same bound with \(I_uU_t\). The terms
whose real integral vanished were removed before taking the
majorant. The triangle inequality after collecting signed
terms at each actual index proves the result.
\end{proof}

In Theorem~\ref{thm:identitycoefficient551} this controls the
term \(R_{\iota_t}^{tu}\) as well as the sum at distinct
indices. The real-section norms, the ratios of the volume
factors to the integral period normalizations, and the total
over the retained rows still enter the interval comparison.


\subsection{The outer flag ideals in the ordinary period}
\label{sec:flagideals554}

On the fields and intervals of Theorem~\ref{thm:columns235},
write \(E=E_x,\ F=E^+,\ f=[F:\QQ]\), and let
\(\mathfrak d=\mathfrak d_{F/\QQ}\mathcal O_E\).
Number the outer vectors \(e_i\) from \(1\) to \(s\), set
\(X_i=\sum_{h\le i}Ee_h\), and put \(X=X_s\). The prescribed
outer part of a marked row is
\[
 A=\Lambda\cap X=\bigoplus_{i=1}^{s}\mathfrak a_i e_i,
 \qquad \mathfrak a_i=\mathfrak d^{-(s-i+1)}.
\]
These ideals are units at the two original level primes, which
split completely in \(E\). The row is Hermitian self-dual.
Its terminal vector \(z\in\Lambda\cap X^\perp\) has norm
\(-T_j\), with \(T_j=1\) or \(T_j=T=\varpi\bar\varpi\).
Its image \(\bar z\) in \(X^\perp/X\) is primitive: the
multipliers \(c\in E\) for which \(c\bar z\) belongs to the
quotient lattice form \(\mathcal O_E\). These are the defining
row conditions in
\srcloc{Formalism/SharpRestriction/CuspSpanNormBounds/BlockEnds}{116},
\srcloc{Formalism/SharpRestriction/CuspSpanNormBounds/FlagEEqHyperbolic}{93}
and
\srcloc{Formalism/SharpRestriction/CuspSpanNormBounds/IsPrimitiveSub}{99}.
We retain all unit-terminal rows and impose the hyperplane
inequality of Theorem~\ref{thm:columns235} at terminal norm \(T\).

\begin{lemma}[A Witt decomposition with the prescribed outer ideals]
\leavevmode\label{lem:projectivewitt554}
There are hyperbolic partners \(f_i\) and a self-dual middle
lattice \(J\) containing \(z\) such that
\[
 \Lambda=A\oplus J\oplus A',\qquad
 A'=\bigoplus_{i=1}^{s}\bar{\mathfrak a}_i^{-1} f_i,\qquad
 h(e_i,f_h)=\delta_{ih}.
\]
Here \(A,A'\) are isotropic, \(J\) is perpendicular to both,
and \(J\cap Ez=\mathcal O_Ez\).
\end{lemma}
\begin{proof}
The submodule \(A+\mathcal O_E z\) is saturated in \(\Lambda\).
Indeed a vector of \(\Lambda\) in \(X+Ez\) maps to
\(c\bar z\) in \(X^\perp/X\). Primitivity gives \(c\in\mathcal O_E\),
and subtraction of \(cz\) leaves an element of \(A\).
Over the Dedekind ring \(\mathcal O_E\) the torsion-free quotient
is projective, so the submodule is a direct summand. Integral
functionals on \(A\) may therefore be extended to \(\Lambda\)
with value zero on \(z\). Hermitian self-duality represents
these functionals by vectors perpendicular to \(z\).
Projectivity of the conjugate dual \(A^\vee\) gives an integral
linear section \(u:A^\vee\to\Lambda\cap z^\perp\) for the
pairing with \(A\).

Let \(H\) be the Gram pairing on the image of \(u\). The finite
extension \(E/F\) is unramified at all finite places on these
fields. The trace from \(\mathcal O_E\) to \(\mathcal O_F\) is
surjective: locally this follows by reduction to the trace of a
separable finite residue field extension, or by projection in
the split algebra. The trace ideal is consequently the unit
ideal globally. Choose \(t\in\mathcal O_E\) with \(t+\bar t=1\).
By perfect duality between \(A\) and \(A^\vee\), the integral
form \(-tH\) is represented by an integral correction from
\(A^\vee\) to \(A\), with the adjoint form \(-\bar tH\).
Adding this correction to \(u\) changes the Gram pairing by
\(-H\) and keeps the image perpendicular to \(z\). The image
and \(A\) now form a self-dual hyperbolic sublattice:
the cross pairing is perfect and both diagonal pairings vanish.
The integral orthogonal projection onto this sublattice splits
it off, and its perpendicular \(J\) is self-dual. Rational
partners dual to the specified \(e_i\) in the corrected image
identify its lattice with the displayed \(A'\). Primitivity in
the quotient gives the last equality for \(J\).
\end{proof}

For a nonzero fractional ideal \(\mathfrak b\), use its positive
absolute norm \(N\mathfrak b\) and put
\(\eta(\mathfrak b)=(N\mathfrak b)^{-1/(2f)}\).
At every real place give \(e_i,f_i\) respective lengths
\(c_i=\eta(\mathfrak a_i)\) and \(c_i^{-1}\). Use the positive
majorant \(h_z^+\) of Section~\ref{sec:bothordinary553} on \(J\).
All these blocks are mutually perpendicular for the positive
norm. The original Hermitian form is represented by a unitary
involution in the resulting orthonormal basis. Define \(V_j\)
as the product of the real operator norms minimized over the
left decorated stabilizer, with these target norms and a single
fixed compatible reference norm on the domain.

\begin{lemma}[Product minima with a fractional ideal]
\leavevmode\label{lem:twistedminimum554}
Let \(\Pi\) be any one of the outer line lattices or the middle
lattice in Lemma~\ref{lem:projectivewitt554}. If
\(0\ne v\in\mathfrak b\Pi\), then on the retained rows, for all
sufficiently large \(x\),
\[
                   \prod_{\tau}\|v_\tau\|\ \ge\sqrt{N\mathfrak b}.
\]
For a graded block \(\Pi'\) of \(\bigwedge^j\Lambda\) containing
\(b\) middle factors, a nonzero \(v\in\mathfrak b\Pi'\) satisfies
\[
              \prod_{\tau}\|v_\tau\|
                  \ \ge\sqrt{N\mathfrak b}\,T_j^{-bf/2}.
\]
\end{lemma}
\begin{proof}
An outer block has the form \(\mathfrak h v_0\) with
\(\|v_0\|_\tau=\eta(\mathfrak h)\). For
\(0\ne x\in\mathfrak b\mathfrak h\), the ideal containment
\((x)\subset\mathfrak b\mathfrak h\) gives
\(\prod_\tau|\tau(x)|\ge\sqrt{N\mathfrak b\,N\mathfrak h}\).
Multiplying by \(\eta(\mathfrak h)^f\) cancels the second factor.

For the middle lattice at norm \(T\) project \(v\) to
\(E L\), where \(L=J\cap z^\perp\), and denote the projection
by \(p\). If \(p\ne0\), the pairing ideal \(\mathfrak c\) of
\(p\) with \(L\) lies in \(\mathfrak b\) or its conjugate,
according to the pairing convention. In particular
\(N\mathfrak c\ge N\mathfrak b\). Put \(H=L\cap p^\perp\).
The covolume calculation of Lemma~\ref{lem:middleminimum553}
holds for fractional pairing ideals as well: the projected
line lattice is \(\mathfrak c p/h(p,p)\), up to conjugation
of the ideal. Hence
\(\prod_\tau\|p_\tau\|^2=N\mathfrak c\,\delta(H)/\delta(L)\).
The retained hyperplane bound and \(\delta(L)=T^f\) give
\(\delta(H)/\delta(L)>1\) uniformly for large \(x\).
Orthogonal projection decreases the positive norm at each
place. If \(p=0\), one has
\(\mathfrak bJ\cap Ez=\mathfrak b z\) by localization at each
prime and \(J\cap Ez=\mathcal O_Ez\). The product bound follows
from \(\|z\|_\tau=\sqrt T\). At terminal norm one the middle
majorant is integral. Its value on \(v\) belongs to
\(\mathfrak b\bar{\mathfrak b}\cap F\), so the norm to \(\QQ\)
from \(F\) of that value is at least \(N\mathfrak b\).

For the exterior assertion write the block as an outer tensor
line \(\mathfrak h v_0\) tensored with \(\bigwedge^b J\).
The normalized length of the line generator is
\(\eta(\mathfrak h)\) by multiplicativity. The form
\(T_jh_z^+\) is integral on \(J\), so its induced exterior
pairing is integral and equals \(T_j^b\) times the exterior
majorant. On a nonzero vector in
\((\mathfrak b\mathfrak h)\bigwedge^b J\), its value lies in
\((\mathfrak b\mathfrak h)\overline{(\mathfrak b\mathfrak h)}
\cap F\) and is totally positive. Its norm from \(F\) is at
least \(N(\mathfrak b\mathfrak h)\). Cancel the line factor
as before and take square roots.
\end{proof}

\begin{theorem}[The ordinary orbital error for the prescribed outer ideals]
\leavevmode\label{thm:idealordinary554}
Put \(n=2s+k+1,\ L_0=8(f+1)(n+1)^3,\ a=(L_0-4fn^2)/2\), and
\(Q=g_1g_2\). Let \(N_j\) be the integral unipotent subgroup of
row \(j\), and \(W_j\) its prescribed ordinary real section. Define
\[
 U_j=\sup_y V_j(y)^{L_0}\|W_j(y)\|,\qquad
 I_j=\int_{N_j\backslash G_\infty}V_j(y)^a\|W_j(y)\|\,dy .
\]
For every retained pair \(t,u\) on the same fields and intervals,
with the original two-prime marking and either terminal norm,
\[
 \sum_{i\in G}\|Z_i^{tu}-B_i^{tu}\|_{\rm HS}
 \ \le\ |c_{t,u}|\, Q^{-fa/2}\min\{ I_t U_u,I_u U_t\}
\]
for all sufficiently large \(x\). The bound uses the height
just defined for the fractional line ideals.
\end{theorem}
\begin{proof}
An isometry transporting all the outer vectors modulo earlier
vectors but changing the terminal class has zero orbital
integral by Proposition~\ref{prop:terminalvanishing553}.
The proof there uses the common rational unipotent radical,
its lattice, and the two characters on it; it applies with
the present ideals. For any other nonzero term outside
\(B^{tu}\), let \(i\) be the first failed outer index. Order
the blocks as in Theorem~\ref{thm:bothordinary553} and let \(p\)
be the highest nonzero component of \(\gamma e_{u,i}-e_{t,i}\).
It lies beyond \(X_{t,i-1}\). For all \(b\in\mathfrak a_i\)
both \(b e_{u,i}\) and \(b e_{t,i}\) are integral. Hence
\(p\in\mathfrak a_i^{-1}\Pi\) in its block lattice \(\Pi\).
At the primes over \(Q\), the marked upper unitriangular
condition kills this component modulo \(Q\), and
\(\mathfrak a_i\) is a unit. Intersecting the local
conditions gives \(p\in Q\mathfrak a_i^{-1}\Pi\).
The first bound of Lemma~\ref{lem:twistedminimum554} gives
\(\prod_\tau\|p_\tau\|\ge Q^f c_i^f\). Left and right decorated
stabilizers preserve the norm of the highest component, since
the earlier vectors have transported into the earlier span.
The triangle inequality used in Theorem~\ref{thm:uniformheight552}
has right side \(2^f c_i^f V_t(y)V_u(\gamma^{-1}y)\), as the
source and target distinguished vectors both have length \(c_i\).
Cancelling \(c_i^f\) gives the required height separation.

For the coset count set
\(\mathfrak h_j=\prod_{i\le\min(j,s)}\mathfrak a_i\).
The initial \(j\)-fold wedge \(\Phi_j\), with last vector \(z\)
if \(j=s+1\), belongs to
\(\mathfrak h_j^{-1}\bigwedge^j\Lambda\). This follows by
wedging the integral line modules and, at the last step,
the integral vector \(z\). Its length at each place is at most
\(\eta(\mathfrak h_j)\sqrt T\). In a nonempty height box the
tuples \(Z_j=\delta_0\delta^{-1}\Phi_j\) lie in that same
fractional exterior lattice because \(\delta_0\delta^{-1}\)
is an integral automorphism. After applying the fixed left
stabilizer element for the box, their local radii are bounded
by \(4^j2^{2jm_\tau}\eta(\mathfrak h_j)\sqrt T\).
In each graded block their pairwise differences, when nonzero,
have product norm at least
\(\eta(\mathfrak h_j)^f T^{-jf/2}\), by the exterior assertion.
The factor \(\eta(\mathfrak h_j)\) therefore cancels in the
packing ratio. More precisely, condition on the later blocks
before applying translated-ball packing as in
Lemma~\ref{lem:gradedpacking552}; the remaining dimension
drops sum to \(n-j+1\) for each fixed preceding nonzero wedge.
Thus the box count has exactly the upper bound
\[
  3^{(2fk)^2}
  \prod_{j=1}^{s+1}(1+2\cdot4^j T^{(j+1)/2})^{2f(n-j+1)}
  \,2^{4D\sum_\tau m_\tau},
  \qquad D=\sum_{j=1}^{s+1}j(n-j+1).
\]
Here the stabilizer of the tuple modulo its integral unipotent
subgroup embeds into the automorphism group of \(L\): its middle
action fixes \(z\), and acting trivially on \(L\) fixes the
whole rational middle quotient. Reduction modulo three on the
underlying positive \(\ZZ\)-lattice gives the first factor as
in Theorem~\ref{thm:uniformheight552}.

For each fixed row the displayed \(U_j,I_j\) are finite by
Proposition~\ref{prop:weight163} and equivalence of the two
positive norms at that row. Now use the height separation in
the orbital integral, unfold the absolute majorant in either
direction, and sum the boxes. The exponent \(D\) and all box
constants equal those in Theorem~\ref{thm:bothordinary553},
where their logarithms were estimated on these intervals.
Its last geometric-series bound gives \(Q^{-fa/2}\).
This proves the stated estimate with the actual fractional
ideals, including the coefficient at the identity marking.
\end{proof}


\subsection{The real section at the prime terminal frequency}
\label{sec:realheight555}

Keep the outer ideals, level and retained row coordinates of
Theorem~\ref{thm:idealordinary554}. Use the signed orthonormal
norm on the common reference space. For a row \(j\), write
\(V_j^0\) for its height when the outer \(e_i,f_i\) all have
length one, and continue to write \(V_j\) for its height with
the prescribed lengths \(c_i,c_i^{-1}\). Both use \(h_z^+\) on
the middle space. The two quantities in the period bound can
be expressed in the first metric by an explicit change of the
real Bessel character.

\begin{proposition}[The different in the real Bessel character]
\leavevmode\label{prop:archdifferent555}
Put \(R=|\operatorname{disc}(F)|^{1/f}\). Then
\(c_i=R^{s-i+1}\). In the real unitary group let \(C\) act
by \(c_i\) on \(e_i\), by \(c_i^{-1}\) on \(f_i\), and by the
identity on the middle space. At all places together put
\[
                   J_C=\prod_{i=1}^s c_i^{\,2f(n+1-2i)} .
\]
The section \(W^\circ(g)=J_C^{1/2}W(C^{-1}g)\) has the same
quotient \(L^2\)-norm as \(W\). In standard outer coordinates
its simple additive characters have all been divided by \(R\);
the last has coefficient \(\beta/R\) when its original
coefficient is \(\beta\). Its right representation and its
compact column type are preserved.

Write \(v_j=\operatorname{vol}(N_j\backslash B_{j,\infty})\),
where \(N_j\) is the integral unipotent subgroup. Use the
quantity \(U_j\) of Theorem~\ref{thm:idealordinary554} and write
\(I_j\) for its integral. With the original quotient measures,
\[
 \begin{aligned}
 U_j&=J_C^{-1/2}\sup_g V_j^0(g)^{L_0}\|W^\circ(g)\|,\\
 I_j&=v_jJ_C^{1/2}\!
       \int_{B_{j,\infty}\backslash G_\infty}
                  V_j^0(g)^a\|W^\circ(g)\|\,d\dot g .
 \end{aligned}
\]
Consequently the powers of \(J_C\) cancel in each of
\(I_tU_u\) and \(I_uU_t\) for two rows with the same \(E,n,s\).
\end{proposition}
\begin{proof}
The absolute norm from \(E\) of \(\mathfrak d_{F/\QQ}\mathcal O_E\)
is \(|\operatorname{disc}F|^2\), by extension of ideals and the
norm of the different. The formula for \(c_i\) follows.
Conjugation by \(C\) normalizes the decorated real stabilizer
\(B=N\rtimes H\) and centralizes \(H\). It multiplies each of
the first \(s-1\) simple coordinates of the radical by
\(c_i/c_{i+1}=R\) and the terminal coordinate by \(c_s=R\).
In the transformation law for \(W(C^{-1}g)\) the argument
of the character is \(C^{-1}n'C\), giving the claimed coefficients.
All these operations commute with right translation.

At one real place, the independent coordinates on the Lie
algebra of the radical are the strict upper entries on the
outer space, the block from the middle to the outer space,
and a skew-Hermitian \(s\)-by-\(s\) block on the far corner.
The first gives determinant \(\prod_{i<j}(c_i/c_j)^2\), as
each entry is complex. The second gives \(\prod_i c_i^{2(k+1)}\).
In the last, each real diagonal coordinate is multiplied by
\(c_i^2\) and each complex off-diagonal coordinate by \(c_ic_j\).
Its determinant is \(\prod_i c_i^{2s}\). Multiplication and
\(n=2s+k+1\) give \(J_C\) over all \(f\) places.

For a nonnegative function \(F\) on \(B\backslash G_\infty\)
the quotient measure satisfies
\(\int F(Cg)\,d\dot g=J_C\int F(g)\,d\dot g\).
Indeed for a compactly supported function on the group first
integrate over \(b\in B\); the substitution \(b'=CbC^{-1}\)
multiplies Haar measure by \(J_C\). Left translation on the
ambient group preserves Haar measure. The quotient integration
formula gives the assertion, which extends to nonnegative
functions by local quotient sections and monotone convergence.
This proves the \(L^2\) normalization. Moreover the length
of a vector in the weighted metric is its unweighted length
after application of \(C\). Taking the infimum over \(B\)
therefore gives \(V_j(g)=V_j^0(Cg)\). Section norms and both
heights are \(B\)-invariant. Integrating first over
\(N_j\backslash B\) supplies \(v_j\); it is finite since the
integral unipotent lattice is cocompact and \(H\) is compact.
The two claimed identities now follow by substitution.
\end{proof}

For a nonzero section of either frequency and a specified
height \(V\), form the dimensionless ratio
\[
 \mathcal M(W,V)=
 \frac{\displaystyle\sup_g V(g)^{L_0}\|W(g)\|
        \displaystyle\int_{B\backslash G_\infty}
                         V(g)^a\|W(g)\|\,d\dot g}
      {\displaystyle\int_{B\backslash G_\infty}\|W(g)\|^2\,d\dot g}.
\]
Set \(p=L_0+a\). If \(d\mu_W=\|W\|^2d\dot g/\|W\|_2^2\), then
\begin{equation}
                      \mathcal M(W,V)\ge\int V^p\,d\mu_W .
 \label{eq:momentcomparison555}
\end{equation}
This follows by bounding \(V^{L_0}\|W\|\) at each point by
its supremum and integrating \(V^a\|W\|\) times this bound.

\begin{lemma}[A constraint on two translates of a real section]
\leavevmode\label{lem:twomoments555}
Use one real Witt model with the same outer lengths at terminal
frequencies \(1,\varpi\). If \(W_\varpi\) is the normalized
left translate of \(W_1\) of Proposition~\ref{prop:scale151}, then
\[
            \mathcal M(W_1,V)\mathcal M(W_\varpi,V)
                                     \ \ge\ T^{fp/2}.
\]
\end{lemma}
\begin{proof}
Let \(a_\varpi\) multiply every \(e_i\) by \(\varpi_\tau\),
every \(f_i\) by \(\bar\varpi_\tau^{-1}\), and fix the middle
space at place \(\tau\). The decorated stabilizer fixes
\(e_1\). Thus \(b_1a_\varpi b_2^{-1}\) acts on that vector by
\(\varpi_\tau\) for any \(b_1,b_2\) in the stabilizer, and
its norm is at least \(\sqrt T\). Factor this operator as
\((b_1g)(b_2a_\varpi^{-1}g)^{-1}\). The norm of the inverse
of an isometry between compatible majorants equals the norm
of the isometry. Taking the two infima and then the product
over places gives \(V(g)V(a_\varpi^{-1}g)\ge T^{f/2}\).
The probability measure for \(W_\varpi\) is the pushforward
of \(\mu_{W_1}\) under \(g\mapsto a_\varpi^{-1}g\), by the
quotient measure calculation in Proposition~\ref{prop:scale151}.
Apply Cauchy--Schwarz to \(V(g)^{p/2}\) and
\(V(a_\varpi^{-1}g)^{p/2}\) under \(\mu_{W_1}\). The moment
comparison \eqref{eq:momentcomparison555} proves the assertion.
\end{proof}

\begin{theorem}[A lower bound for the prime-frequency majorant]
\leavevmode\label{thm:primemoment555}
For every retained prime-terminal row take the prescribed
isometrically normalized Bessel section of a unit vector in
the product lowest compact type of Theorem~\ref{thm:localmass156}.
Use its weighted height from Theorem~\ref{thm:idealordinary554}.
On the same fields and intervals, with \(b=\exp(x^\varepsilon)\),
one has for an absolute \(c>0\)
\[
      \mathcal M(W_\varpi,V)
        \ \ge\ (1-e^{-c b n^2})
                    \left(\frac{\sqrt T}{R b n^4}\right)^{fp}.
\]
In particular \(Q^{-fa/2}\mathcal M(W_\varpi,V)\) tends to
infinity uniformly for these sections. The right-hand side
of the diagonal error estimate of Theorem~\ref{thm:idealordinary554},
divided by the absolute identity coefficient, is at least this
quantity.
\end{theorem}
\begin{proof}
At each infinite place identify the standard marked real model
with the row by mapping the outer vectors to the specified \(e_i\),
their duals to the chosen \(f_i\), and the negative unit vector
to \(z/\varpi_\tau\). These pairings agree, and the remaining
positive complements have the same dimension and signature;
an isometry of their orthonormal bases finishes the identification.
The ordinary real section and its character transport under
this identification. The quotient measure can be normalized
by its \(L^2\) mass. A change of hyperbolic partners fixing the
outer vectors and \(z\) belongs to the decorated stabilizer
after choosing an isometry on the positive complement. Thus
the following estimate depends only on the decorated frame.

Write \(a_{s,\tau}\) for the distance from \(g_\tau^{-1}e_s\)
to \(g_\tau^{-1}X_{s-1}\) in the reference norm. For every
\(b'\in B_\tau\) the vector \(b'^{-1}e_s\) belongs to
\(e_s+X_{s-1}\). Hence
\[
     a_{s,\tau}
      \le \|g_\tau^{-1}b'^{-1}e_s\|_{\rm ref}
      \le c_s\|(b'g_\tau)^{-1}\|_{\rm op}
      =R\|b'g_\tau\|_{\rm op}.
\]
Take the infimum over \(b'\). Theorem~\ref{thm:localmass156}
shows that at least \(1-e^{-c b n^2}\) of the normalized mass
lies where, at every place simultaneously,
\(a_{s,\tau}>\sqrt T/(bn^4)\). On this set the product height
is at least \((\sqrt T/(Rbn^4))^f\). Integrating the \(p\)-th
power and applying \eqref{eq:momentcomparison555} proves the bound.

On these intervals \(\log R=O(\log x)\), since \(E\) has
polynomial conductor, and \(n\) is polynomial in \(x\).
Furthermore \(\log T=(17+o(1))m x^\varepsilon\), \(m\to\infty\),
whereas \(\log Q=(16+o(1))x^\varepsilon\), by
Theorem~\ref{thm:prime209}. Since \(p>a>0\), the logarithm of
the asserted lower bound after multiplication by \(Q^{-fa/2}\)
tends to infinity. Finally, on a diagonal row
\(\lambda_j=|c_{j,j}|\,v_j\|W_j\|_2^2/d_j\) in
Theorem~\ref{thm:flagisometry544}, with \(d_j\ge1\).
Dividing the estimate's right-hand side by \(\lambda_j\)
therefore gives exactly \(d_jQ^{-fa/2}\mathcal M(W_j,V_j)\).
\end{proof}

The lower bound concerns the absolute estimate used in the
unfolding. The integral projected period depends on signed
orbital integrals. Their phases or a sharper estimate on the
individual regions of integration are essential to an argument
using these prime-frequency rows.


\subsection{Intersecting flags at the two terminal frequencies}
\label{sec:mixedsuffix556}

Consider a term of the ordinary projected period between rows
of terminal norms \(T_0,T_1\in\{1,T\}\). Transport the second
frame by the rational isometry indexing the term. Finite
support identifies the two lattices with one \(\Lambda\);
above \(Q=g_1g_2\) the reduced outer flags agree under upper
unitriangular changes of frame. Write \(e_{b,i}\) for the
resulting outer vectors, \(X_b\) for their span, and
\(X_{b,i}\) for its first \(i\) positions, \(b=0,1\).
The line modules are \(\mathfrak a_i e_{b,i}\) as in
Section~\ref{sec:flagideals554}. A term can be nonzero only
if its two characters agree on the common unipotent group:
integration over the compact unipotent fibre gives this condition
by Lemma~\ref{lem:zero133}, with the respective terminal
vectors in the two characters.

\begin{lemma}[The determinant on the common suffix]
\leavevmode\label{lem:suffixdet556}
Suppose the characters agree, \(I=X_0\cap X_1\) has dimension
\(e>0\), and \(t=s-e\). Put \(P_b=X_{b,t}\) and \(M=\Lambda\cap I\).
Projection from \(I\) to \(X_b/P_b\) is a rational isomorphism
for both \(b\). Let \(\mathfrak c_b\) be the index ideal of the
image of \(M\) in the integral quotient
\((\Lambda\cap X_b)/(\Lambda\cap P_b)\). There is a
\(\lambda\in E^\times\) such that
\[
       \lambda\bar\lambda=T_1/T_0,\qquad
       \mathfrak c_0=(\lambda)^e\mathfrak c_1 .
\]
The index ideal of a full sublattice has valuation the length
of its cokernel at each finite prime.
\end{lemma}
\begin{proof}
Theorem~\ref{thm:success135} applies to unequal terminal norms.
It shows that the flags induce the same filtration on \(I\),
with one-dimensional successive quotients at positions
\(t+1,\ldots,s\), and that the cross pairing of \(X_0,X_1\)
has kernel \(I\) on both sides. In particular, \(I\) projects
isomorphically as asserted. Write \(S_b\) for the successor
map on the ordered outer vectors, whose last value is \(z_b\).
For \(v\in I\) with last nonzero coefficient in position \(i\),
equality of the characters gives
\[
             S_0v-S_1v\in X_{0,i}+X_{1,i}
\]
by \eqref{eq:success135}. If \(i<s\), subtracting vectors in
these spaces gives a common vector in position \(i+1\) with
the same two leading coefficients. At \(i=s\) it gives a
common vector of the two extended flags with negative norm.
Its two terminal coefficients \(a_0,a_1\) satisfy
\(a_0\bar a_0T_0=a_1\bar a_1T_1\), since addition of an
isotropic vector to its perpendicular terminal vector preserves
the norm. All common vectors at the terminal position have the
same ratio \(\lambda=a_0/a_1\): projection to either terminal
line from \((X_0+Ez_0)\cap(X_1+Ez_1)\) has kernel \(I\).
Indeed an isotropic vector in either extended flag has zero
terminal coefficient in that flag. It follows by the successive
steps that this same ratio occurs on each of the \(e\) common
graded lines.

Identify the two quotients by the rational maps from \(I\).
In their ordered suffix bases the map from side \(1\) to
side \(0\) is triangular with every diagonal entry
\(\lambda\). It sends one projected copy of \(M\) onto the
other. The ambient determinant ideal in both these bases is
\(\prod_{i>t}\mathfrak a_i\). On taking the top exterior
power, the index ideals therefore differ by
\((\lambda)^e\), which proves the claim.
\end{proof}

\begin{theorem}[A terminal-prime divisor of the prefix pairing]
\leavevmode\label{thm:terminaldivisor556}
Take the Hermitian pairing conjugate-linear in its first argument.
Let \(D_t=\det(h(e_{0,i},e_{1,j}))_{i,j\le t}\), put \(D_0=1\),
and let \(\mathfrak h_t=\prod_{i\le t}\mathfrak a_i\).
The fractional expression
\[
       \mathfrak A=(D_t/Q^t)\bar{\mathfrak h}_t\mathfrak h_t
\]
is a nonzero integral ideal. More precisely,
\(\mathfrak A=\bar{\mathfrak c}_0\mathfrak c_1\mathfrak b\)
for an integral ideal \(\mathfrak b\).
If \(T_0=T,\ T_1=1\), then for each conjugate pair of primes
over \(T\), at least one of them occurs in \(\mathfrak A\) with
exponent at least \(e\). In particular
\[
               T^{fe}\ \mid\ N_{E/\QQ}\mathfrak A .
\]
Interchanging the rows conjugates \(\mathfrak A\). At unequal
terminal norms there is no contributing term with \(X_0=X_1\).
\end{theorem}
\begin{proof}
Set \(A_b=\Lambda\cap X_b\), \(P_b^0=\Lambda\cap P_b\), and
\(U_b=A_b/M\). The pairing \(h/Q\) takes integral values on
\(A_0\times A_1\): away from \(Q\) use self-duality, and at
each prime over \(Q\) both reduced spaces are the same
isotropic space with a basis of marked vectors. It descends
to \(U_0\times U_1\) because each side is isotropic and
contains \(I\). The rank computation in
Lemma~\ref{lem:suffixdet556} makes the resulting rational
pairing nonsingular. Its determinant ideal \(\mathfrak b\)
in the two projective lattices is therefore integral and
nonzero. Here the ideal is generated locally by the determinant
of the pairing in integral local bases, using conjugate local
bases on the first side.

The injected image of \(P_b^0\) has index ideal
\(\mathfrak c_b\) in \(U_b\): the two descriptions of the
finite quotient both give \(A_b/(M+P_b^0)\). Restricting the
pairing to these prefix sublattices multiplies its determinant
ideal by \(\bar{\mathfrak c}_0\mathfrak c_1\), by determinants
of local basis changes. Computing the same determinant in the
specified \(e_{b,i}\) yields
\((D_t/Q^t)\bar{\mathfrak h}_t\mathfrak h_t\).
This also holds for the empty prefix, with the empty
determinant equal to one.

Every place of \(F\) over the split rational prime \(T\)
gives a pair \(\mathfrak P,\bar{\mathfrak P}\) in \(E\).
The valuations of \(\lambda\) at this pair sum to \(-1\),
since \(\lambda\bar\lambda=1/T\). At a prime with valuation
at most \(-1\), integrality of
\(\mathfrak c_0=(\lambda)^e\mathfrak c_1\) forces
\(v_{\mathfrak P}(\mathfrak c_1)\ge e\), with \(\mathfrak P\)
chosen to be that prime. The determinant factorization puts
the same divisor into \(\mathfrak A\). There are \(f\) such
pairs, each of residue degree one. This proves the norm
divisibility. With \(t=0\), \(\mathfrak A=\mathcal O_E\)
contradicts the divisor since \(e=s\ge1\).
\end{proof}

We next use this divisor to control part of the mixed
ordinary period in the region identified in
Theorem~\ref{thm:primemoment555}. For row \(j\) let \(a_j\) be
the real isometry acting by \(\beta_j\) on every \(e_i\), by
\(\bar\beta_j^{-1}\) on every \(f_i\), and by the identity
on its middle space, with \(\beta_j=1\) or \(\varpi\) as
prescribed. Set \(\widetilde V_j(y)=V_j(a_jy)\). Since \(a_j\)
normalizes the decorated stabilizer, this is the quotient
height obtained by multiplying all the original outer lengths
by \(\sqrt{T_j}\) on the \(e_i\) and by \(1/\sqrt{T_j}\) on
the \(f_i\). Denote by \(\widetilde U_j,\widetilde I_j\)
the quantities in Theorem~\ref{thm:idealordinary554} with
this height.

\begin{theorem}[The mixed contribution from intersections of large dimension]
\leavevmode\label{thm:mixedhalf556}
Use a pair of retained rows of unequal terminal norms. For
a fixed \(1\le t<s\), restrict the orbital sum at each marked
index \(i\) to terms with \(\dim(X_0\cap X_1)=e=s-t\).
Call the resulting signed operator \(Z_i(t)\). If \(2t\le s\), then
on the same intervals, for all sufficiently large \(x\),
\[
 \sum_i\|Z_i(t)\|_{\rm HS}
 \ \le\ |c_{0,1}|\,Q^{-fa/2}T^{-fa(e-t)/(2t)}
      \min\{\widetilde I_0\widetilde U_1,\,
                         \widetilde I_1\widetilde U_0\}.
\]
For the union of all \(1\le t\le s/2\) the same estimate holds
with the power of \(T\) omitted. Both bounds use the
unit-frequency real sections after the exact change of variables
described below.
\end{theorem}
\begin{proof}
First keep any \(0<t<s\). In a common real reference space
pair the two initial \(t\)-fold wedges after applying \(y^{-1}\)
at both ends, with the second wedge transported by \(\gamma\).
Their Hermitian pairing has absolute value \(|\tau D_t|\).
The Hermitian form on the exterior reference space is represented
by a unitary involution, so Cauchy--Schwarz bounds this by the
product of their positive lengths. Each wedge is fixed by its
own decorated stabilizer. In its row's modified majorant the
wedge has length exactly
\(\eta(\mathfrak h_t)T_j^{t/2}\).
The norm of the exterior inverse is at most the \(t\)-th
power of the operator norm. Taking the infimum independently
at each end, then multiplying over all places, gives
\[
 |N_{E/\QQ}D_t|^{1/2}
 \le \eta(\mathfrak h_t)^{2f} T^{tf/2}
      \big(\widetilde V_0(y)\widetilde V_1(\gamma^{-1}y)\big)^t .
\]
On the other hand the determinant divisor gives
\[
 |N_{E/\QQ}D_t|^{1/2}
 =Q^{tf}\eta(\mathfrak h_t)^{2f}(N\mathfrak A)^{1/2}
 \ \ge\ Q^{tf}\eta(\mathfrak h_t)^{2f}T^{ef/2}.
\]
The product of the two heights is consequently at least
\(Q^f T^{f(e-t)/(2t)}\). If \(2t\le s\) it is at least \(Q^f\).

The modified heights still admit a uniform coset count.
In the \(j\)-th exterior lattice a graded block with \(b\)
middle factors, \(l_+\) outer vectors and \(l_-\) dual outer
vectors gains a factor \(T_j^{(l_+-l_-)/2}\) in each positive
length. Here \(b+l_++l_-=j\). By Lemma~\ref{lem:twistedminimum554},
on the lattice twisted by
\(\mathfrak h_{\min(j,s)}^{-1}\) its product minimum is at least
\(\eta(\mathfrak h_{\min(j,s)})^f T^{-jf/2}\).
The initial wedge has length at most
\(\eta(\mathfrak h_{\min(j,s)})T^{j/2}\), including the terminal
step \(j=s+1\). The translated-ball count in
Theorem~\ref{thm:idealordinary554} therefore applies with box
constant
\[
  3^{(2fk)^2}\prod_{j=1}^{s+1}
                  (1+2\cdot4^j T^j)^{2f(n-j+1)}
\]
and its original exponent \(4D\). The logarithm of the new
constant is \(O(f^2k^2+fn(s+1)^2\log(2T))=o(fa\log Q)\)
by the interval estimates in Theorem~\ref{thm:bothordinary553}.
It still absorbs, together with the convergent geometric
series, into \(Q^{fa/2}\). Remove the zero integrals using
character compatibility, use the separation just proved
to exponent \(a\), and unfold the nonnegative majorant in
either direction as in Theorem~\ref{thm:mixedordinary541}.
This proves the asserted bound. For the union use the common
separation \(Q^f\) before enlarging the sum to all cosets in
the box count. This avoids any factor for the number of strata.

Finally let \(W_{1,j}\) be the normalized unit-frequency
section obtained by transporting the very same row vector;
up to a scalar of modulus one,
\(W_j(y)=J_j^{-1/2}W_{1,j}(a_jy)\), with
\(J_j=T_j^{fs(n-s)}\), by Proposition~\ref{prop:scale151}.
For \(v_j=\operatorname{vol}(N_j\backslash B_{j,\infty})\)
the exact factors are
\[
 \begin{aligned}
 \widetilde U_j
     &=J_j^{-1/2}\sup_y V_j(y)^{L_0}\|W_{1,j}(y)\|,\\
 \widetilde I_j
     &=v_jJ_j^{1/2}\int_{B_{j,\infty}\backslash G_\infty}
                        V_j(y)^a\|W_{1,j}(y)\|\,d\dot y .
 \end{aligned}
\]
They follow from the quotient measure substitution in that
proposition; in particular the change applies to the actual
real vectors in the period. The finite levels and the index
\(i\) in \(Z_i(t)\) have been used throughout the proof.
\end{proof}

For smaller intersections the determinant computation still
gives \(T^{fe}\) in the norm of the prefix ideal. Its resulting
height bound has a negative power of \(T\) when \(e<t\).
Transverse flags have no common suffix to which this
argument can be applied.


\subsection{Denominators in the common part of two outer flags}
\label{sec:denomsuffix562}

Keep the retained rows and the integral projected period of
Section~\ref{sec:mixedsuffix556} on the same field \(E=E_x\).
Put \(F=E^+\) and \(f=[F:\QQ]\). Order the rows by
\(T_0\ge T_1\), where \(T_j\in\{1,T\}\) is the terminal norm.
The prime \(T\) splits completely in \(E\). Fix an integral
orbit with compatible characters. Suppose the intersection of
its two outer spaces has dimension \(e\) with \(0<e<s\),
and put \(t=s-e\). Its common graded scalar is \(\lambda\)
in Lemma~\ref{lem:suffixdet556}. Write
\[
  \beta=\lambda^{-1},\qquad
  u=T_0/T_1,\qquad
  (\beta)=\mathfrak n\,\mathfrak d^{-1},
           \qquad (\mathfrak n,\mathfrak d)=1 .
\]
The ideals \(\mathfrak n,\mathfrak d\) are integral and
\(\beta\bar\beta=u\). Denote absolute ideal norm by \(N\).
The scalar is intrinsic to the two ordered frames modulo their
upper unipotent changes, since such changes fix each graded
vector.

\begin{lemma}[The possible scalars at a fixed denominator]
\leavevmode\label{lem:scalarden562}
One has \(\mathfrak n=\bar{\mathfrak d}\,\mathfrak p\), where
\(\mathfrak p=(1)\) for \(u=1\). For \(u=T\) the ideal
\(\mathfrak p\) is a product of one prime from each of the
\(f\) conjugate pairs over \(T\). In particular
\[
             N\mathfrak n=u^f N\mathfrak d .
\]
At each fixed denominator there are at most
\(2^{f{\bf1}_{u=T}}|\mu(E)|\) possible scalars \(\beta\),
where \(\mu(E)\) is the group of roots of unity of \(E\).
For \(X\ge1\), the number with \(N\mathfrak d\le X\) is at most
\[
             2^{f{\bf1}_{u=T}}\,8f^2\,\zeta(2)^{2f}X^2 .
\]
\end{lemma}
\begin{proof}
At a prime fixed by conjugation, twice the valuation of \(\beta\)
is the valuation of \(u\). Every prime over \(T\) is paired with
a different prime, so the valuation at a fixed prime is zero.
At a conjugate pair the two valuations are integers with sum
zero away from \((u)\). Thus the numerator exponent on one side
equals the denominator exponent on the other. At a pair over
\(u=T\) the sum is one. If the valuations are \(k,1-k\) with
\(k\ge1\), the positive exponent \(k\) exceeds the opposite
denominator exponent \(k-1\) by one, and both exponents on the
other side in this comparison are zero. Exchanging the pair
covers \(k\le0\). This proves the ideal identity.

There are at most the indicated number of choices of
\(\mathfrak p\). Two generators of the resulting fractional
ideal satisfying the norm equation differ by an algebraic
integer unit of absolute value one at every embedding. By
Kronecker's theorem their ratio is a root of unity. For the
count of denominator ideals, each ideal of norm at most \(X\)
contributes at least \(X^{-2}\) to \(\zeta_E(2)\). Comparing
Euler factors gives \(\zeta_E(2)\le\zeta(2)^{2f}\).
Finally, writing \(m=|\mu(E)|\), one has
\(\varphi(m)\le 2f\) and \(\varphi(m)^2\ge m/2\). For the
second inequality the factor \(\varphi(p^k)^2/p^k\) is at
least one except at \(2^1\), where it equals one half. This
gives \(m\le8f^2\) and proves the count.
\end{proof}

\begin{theorem}[The denominator in the archimedean separation]
\leavevmode\label{thm:denomheight562}
Let \(\mathfrak c_0,\mathfrak c_1,\mathfrak b,\mathfrak A\)
be the two suffix indices, the determinant ideal of the quotient
pairing, and the normalized prefix determinant ideal of
Theorem~\ref{thm:terminaldivisor556}. Put
\(\mathfrak g=\mathfrak c_0+\mathfrak c_1\). Then
\[
       N\mathfrak A
         =u^{fe}(N\mathfrak d)^{2e}
                         (N\mathfrak g)^2 N\mathfrak b .
\]
For the modified heights \(\widetilde V_j\) of
Section~\ref{sec:mixedsuffix556}, the two arguments of the
orbital integral satisfy
\[
 \begin{split}
 \widetilde V_0(y)\widetilde V_1(\gamma^{-1}y)
 &\ \ge\ Q^f\frac{u^{fe/(2t)}}{(T_0T_1)^{f/2}}\,
       (N\mathfrak d)^{e/t}
       \bigl(N\mathfrak g\sqrt{N\mathfrak b}\bigr)^{1/t}.
 \end{split}
\]
The inequalities apply to both equal and unequal terminal norms.
\end{theorem}
\begin{proof}
The equality \(\mathfrak c_0=(\lambda)^e\mathfrak c_1\)
compares two integral ideals. Taking their minimum exponent
at every prime gives
\[
     \mathfrak c_0=\mathfrak g\,\mathfrak d^e,\qquad
     \mathfrak c_1=\mathfrak g\,\mathfrak n^e .
\]
Use the factorization
\(\mathfrak A=\bar{\mathfrak c}_0\mathfrak c_1\mathfrak b\)
and take norms. Lemma~\ref{lem:scalarden562} yields the first
formula, including contributions at primes away from \(T\).

The Hermitian pairing of the two prefix wedges equals \(D_t\),
the cross determinant. Their lengths in the individual
modified majorants at place \(\tau\) are
\(\eta(\mathfrak h_t)T_j^{t/2}\) in the notation of
Theorem~\ref{thm:mixedhalf556}. Both wedges are invariant under
their own decorated stabilizers. Cauchy--Schwarz in the
common reference majorant, followed by the exterior operator
norm bound at each end, gives
\[
 |N_{E/\QQ}D_t|^{1/2}\ \le\
   \eta(\mathfrak h_t)^{2f}(T_0T_1)^{tf/2}
           \bigl(\widetilde V_0(y)\widetilde V_1(\gamma^{-1}y)\bigr)^t .
\]
This uses one complex absolute value for each pair of conjugate
embeddings. By the definition of \(\mathfrak A\), the left
side is \(Q^{tf}\eta(\mathfrak h_t)^{2f}(N\mathfrak A)^{1/2}\).
Substitution of the norm formula and taking \(t\)-th roots
proves the assertion.
\end{proof}

\begin{corollary}[A bound for all sufficiently large denominators]
\leavevmode\label{cor:denomperiod562}
Fix \(L\ge1\), and for \(0<t<s\) put \(e=s-t\) and
\[
   X_t(L)=L^{t/e}(T_0T_1)^{ft/(2e)}u^{-f/2}.
\]
At marked index \(i\), let \(Z_i^{\rm large}(L)\) be the
sum of the original signed orbital operators with
\(0<e<s\) and \(N\mathfrak d\ge X_t(L)\). For all sufficiently
large \(x\) on the common intervals,
\[
 \sum_i\|Z_i^{\rm large}(L)\|_{\rm HS}
     \ \le\ |c_{0,1}|\,Q^{-fa/2}L^{-a}
       \min\{\widetilde I_0\widetilde U_1,\,
                          \widetilde I_1\widetilde U_0\}.
\]
All real-section factors have the normalization in
Theorem~\ref{thm:mixedhalf556}. At a fixed remaining \(t\),
if \(X_t(L)\le1\) the complementary sum is empty. Otherwise
its possible common graded scalars number at most
\(2^{f{\bf1}_{u=T}}8f^2\zeta(2)^{2f}X_t(L)^2\).
\end{corollary}
\begin{proof}
Both \(N\mathfrak g\) and \(N\mathfrak b\) are at least one.
Theorem~\ref{thm:denomheight562} shows that every retained term
has height product at least \(Q^f L\). In the absolute
unfolding insert the two heights to the power \(a\), taking out
the common factor \(Q^{-fa}L^{-a}\). The translated-ball
count in the proof of Theorem~\ref{thm:mixedhalf556} applies
to either row separately. Its modified exterior minima and
box constant use only \(1\le T_j\le T\), so they apply also
when \(T_0=T_1=T\). The count and its convergent geometric
sum cost at most \(Q^{fa/2}\) on these intervals. Unfolding in
either direction gives the two products of section norms.
One can enlarge to all cosets after extracting the common
factor. Thus the estimate covers the union of the retained
dimensions without a factor for their number. The bound for
the remaining scalar values follows from Lemma~\ref{lem:scalarden562}.
\end{proof}

For the pair \(T,1\), the cutoff at \(L=1\) is
\(T^{f(t-e)/(2e)}\). It includes every denominator when
\(e\ge t\) and also gives a contribution estimate when \(e<t\).
For the pair \(T,T\) it is \(T^{ft/e}\). The count in the
corollary bounds the distinct scalar values in the complementary
sum. Several rational orbits can have the same scalar. Their
total signed contribution and the contribution of transverse
flags still enter the full integral period comparison.


\subsection{The two marked primes in the common flag index}
\label{sec:markedindex564}

Take a term of the integral projected period with compatible
characters on \(E=E_x\). Keep \(0<e<s\), \(t=s-e\), and the
notation of Section~\ref{sec:denomsuffix562}. In particular
\(\beta=\lambda^{-1}\), \(\beta\bar\beta=u=T_0/T_1\),
\(\mathfrak d\) is its reduced denominator, and
\(\mathfrak g=\mathfrak c_0+\mathfrak c_1\). Write \(N\) for
absolute ideal norm. The two primes in \(Q=g_1g_2\) split
completely and are different from the terminal prime \(T\).
For a fixed scalar define the squarefree ideal
\[
   \mathfrak q_\beta=
   \prod_{\substack{\mathfrak w\mid Q\\
           \beta\ \text{fails to be a unit congruent to }1
                         \bmod\mathfrak w}}\mathfrak w .
\]
The marked group uses upper triangular unipotent changes of the
reduced outer frames. In particular it fixes the specified
generators on their successive quotients. The resulting restriction
on the two suffix indices can be read over the original integer rings.

\begin{theorem}[A determinant congruence at both marked primes]
\leavevmode\label{thm:markedindex564}
For every such term,
\[
       (\beta-1)^e\mathfrak c_0\ \subseteq\ Q^e\mathcal O_E,
                  \qquad
       \mathfrak g\ \subseteq\ \mathfrak q_\beta^{\,e}.
\]
If \(\beta\ne1\), the ideal \((Q)\mathfrak q_\beta^{-1}\)
divides \((\beta-1)\mathfrak d\). In particular
\[
             N\mathfrak d\,N\mathfrak q_\beta
                   \ \ge\ \left(\frac{Q}{1+\sqrt u}\right)^{2f},
             \qquad f=[E^+:\QQ].
\]
The set of primes in \(\mathfrak q_\beta\) is stable under complex
conjugation for the prescribed terminal prime.
\end{theorem}
\begin{proof}
Use suffix coordinates \(t+1,\ldots,s\) on both outer spaces.
The ambient projective lattice for these coordinates is the same
on both sides, namely \(L=\bigoplus_{i>t}\mathfrak a_i\).
Let \(L_0,L_1\) be the images of \(\Lambda\cap(X_0\cap X_1)\)
in these coordinates. Their index ideals in \(L\) are
\(\mathfrak c_0,\mathfrak c_1\). On the common rational space,
projection to each suffix is an isomorphism by
Lemma~\ref{lem:suffixdet556}. Let \(B\) carry suffix coordinates
on side zero to those on side one. It takes \(L_0\) onto \(L_1\).
The induced filtrations on the common space are equal and the
ratio of the leading coordinates from zero to one is always
\(\beta\), by the same lemma. Thus \(B\) is upper triangular
with all diagonal entries \(\beta\).

At a prime \(\mathfrak w\mid Q\) the \(\mathfrak a_i\) are
units and the reduced outer vectors are linearly independent.
The two reduced frames differ by an upper unitriangular matrix.
For the full coordinates of a common integral vector this gives
an upper unitriangular relation modulo \(\mathfrak w\) as well,
using the inverse matrix if the frames were compared in the opposite
order. Its last \(e\) coordinates depend only on the last \(e\)
input coordinates. Lift that bottom diagonal block to an upper
unitriangular \(U_{\mathfrak w}\) over \(\mathcal O_{E,\mathfrak w}\).
We obtain
\[
       (B-U_{\mathfrak w})(L_0)_{\mathfrak w}
                            \subseteq \mathfrak w L_{\mathfrak w}.
\]
The determinant of the difference is \((\beta-1)^e\) because
it is triangular. In local bases of \(L\) and \(L_0\) the
inclusion of \(L_0\) has determinant generating the localized
index ideal \(\mathfrak c_0\). Taking top exterior powers gives
the first inclusion at \(\mathfrak w\). Away from \(Q\) we have
\((B-1)L_0\subseteq L\) since both \(L_0\) and \(B L_0=L_1\)
are integral. Taking determinants gives integrality there. These
local statements prove the global inclusion, including the
case \(\beta=1\), when the determinant is zero.

At a marked prime in \(\mathfrak q_\beta\), put
\(b=v_{\mathfrak w}(\beta)\). Then
\(v_{\mathfrak w}(\beta-1)=\min(b,0)\). This also holds when
\(b=0\), by the residue condition defining \(\mathfrak q_\beta\).
The index identity in Theorem~\ref{thm:denomheight562} reads
\(\mathfrak c_0=\mathfrak g\mathfrak d^e\), with
\(v_{\mathfrak w}(\mathfrak d)=\max(-b,0)\). The determinant
congruence gives \(v_{\mathfrak w}(\mathfrak g)\ge e\).
Multiplying over the indicated primes proves the second inclusion.

For \(\beta\ne1\) the ideal \((\beta-1)\mathfrak d\) is nonzero
and integral. It lies in each marked prime outside
\(\mathfrak q_\beta\). Taking norms of this divisibility gives
\[
     Q^{2f}\ \le\ N\mathfrak q_\beta\,N\mathfrak d\,
                                 |N_{E/\QQ}(\beta-1)|.
\]
At every embedding \(\tau\) one has \(|\tau\beta|=\sqrt u\).
The triangle inequality yields the stated bound. Finally
\(u\equiv1\bmod Q\). Indeed \(u=1\) or \(u=T\), and the
terminal generator satisfies \(\varpi\equiv1\bmod Q^2\)
by Theorem~\ref{thm:common111}. The equality
\(\bar\beta=u/\beta\) shows that being a unit with residue one
is invariant under conjugation at the marked primes.
\end{proof}

\begin{corollary}[Contribution of the common flag index to the period]
\leavevmode\label{cor:markedperiod564}
Use \(X_t(L)=L^{t/e}(T_0T_1)^{ft/(2e)}u^{-f/2}\)
for \(L\ge1\). Let \(Z_i(L)\) sum the signed terms with
\(0<e<s\) and
\[
                    N\mathfrak d\,N\mathfrak q_\beta\ \ge\ X_t(L)
\]
in the original marked operator at index \(i\). Then on the
common intervals, for sufficiently large \(x\),
\[
    \sum_i\|Z_i(L)\|_{\rm HS}
       \ \le\ |c_{0,1}|Q^{-fa/2}L^{-a}
         \min\{\widetilde I_0\widetilde U_1,\,
                         \widetilde I_1\widetilde U_0\}.
\]
If \(T_0=T_1\) and \(X_t(L)\le(Q/2)^{2f}\), every term in
that dimension with \(\beta\ne1\) is included.
\end{corollary}
\begin{proof}
The height estimate in Theorem~\ref{thm:denomheight562} and
\(N\mathfrak b\ge1\) give
\[
 \widetilde V_0(y)\widetilde V_1(\gamma^{-1}y)
  \ \ge\ Q^f (T_0T_1)^{-f/2}u^{fe/(2t)}
                    (N\mathfrak d\,N\mathfrak q_\beta)^{e/t}
\]
by the second inclusion of Theorem~\ref{thm:markedindex564}.
The assumed inequality makes this at least \(Q^fL\).
Extracting its negative \(a\)-th power once on the whole set
of terms gives the common factor \(Q^{-fa}L^{-a}\).
The two unfoldings and translated-ball count in
Corollary~\ref{cor:denomperiod562} cost at most \(Q^{fa/2}\)
times the smaller displayed product of section norms. These
estimates use all cosets after the common factor is extracted
and thus apply to the indicated subset. The final assertion
is the norm inequality in Theorem~\ref{thm:markedindex564}
at \(u=1\).
\end{proof}

For an equal-norm pair with \(\mathfrak d=\mathcal O_E\) the
scalar is a root of unity, by Kronecker's theorem. If it differs
from one then \(\mathfrak q_\beta=(Q)\) for all sufficiently
large \(x\). Indeed its order is at most \(8f^2\) by
Lemma~\ref{lem:scalarden562}. Each marked prime exceeds this
bound since \(f\) is polynomial in \(x\) and
\(g_i\ge(k\exp(x^\varepsilon))^8\). Reduction at a prime of
residue characteristic prime to the order is injective on those
roots of unity. This gives \(N\mathfrak d\,N\mathfrak q_\beta
=Q^{2f}\) for all such scalars. The scalar \(1\) fixes every
marked graded line and can occur on several rational orbits.
Those terms and the terms at transverse flags still require
control in the projected period.


\subsection{The common middle lattice and the terminal divisor}
\label{sec:commonmiddle570}

Use the original integral projected period on \(E=E_x\) and
the marked levels \(Q=g_1g_2\) of Section~\ref{sec:markedindex564}.
Keep a supported pair with compatible characters and a proper
nonzero intersection \(I=X_0\cap X_1\). Write \(e=\dim_E I\),
\(t=s-e\), \(T_0\ge T_1\), \(u=T_0/T_1\), and \(f=[E^+:\QQ]\).
The scalar \(\beta=\lambda^{-1}\) has norm \(\beta\bar\beta=u\)
and reduced factorization \((\beta)=\mathfrak n/\mathfrak d\).
Set \(A_i=\Lambda\cap X_i\), \(M=\Lambda\cap I\). On
\(I^\perp/I\) consider the lattices
\[
     \Lambda'=(\Lambda\cap I^\perp)/M,\qquad
     U_i=A_i/M,\qquad D=U_0\oplus U_1 .
\]
The \(U_i\) are disjoint and isotropic. Their cross pairing has
rank \(t\) by Lemma~\ref{lem:suffixdet556}. Put \(S=ED\),
\(R=\Lambda'\cap S\), \(R_1=R\cap Q^{-1}D\), and
\(C=\Lambda'\cap S^\perp\), taking the orthogonal complement
inside \(I^\perp/I\). The space \(S^\perp\) projects isometrically
onto each rational middle space \(X_i^\perp/X_i\): after
removing \(I\), the sum of the two \(U_i\) spans a nondegenerate
space in which the perpendicular of either \(EU_i\) is itself.
View the two integral middle lattices as lattices \(J_i\) in
this one rational space by those isometries. Thus \(C\subset J_i\).
Write \(\zeta_i\in J_i\) for the image of the original terminal
vector \(z_i\) in this identification.

\begin{lemma}[The common middle index]
\leavevmode\label{lem:middleindex570}
The two lattices \(J_i\) are self-dual, and their common
sublattice \(C\) satisfies
\[
        [R_1:D]=Q^{2ft},\qquad
        [J_0:C]=[J_1:C]=\frac{N\mathfrak b}{[R:R_1]} .
\]
Here \(\mathfrak b\) is the determinant ideal of \(h/Q\)
on \(U_0\times U_1\) from
Theorem~\ref{thm:terminaldivisor556}, and each bracket is
the order of the finite quotient as an abelian group.
\end{lemma}
\begin{proof}
For a primitive isotropic sublattice of a self-dual lattice,
the induced lattice in its perpendicular modulo itself is
self-dual. Indeed an integral functional on the quotient
pulls back to one on the perpendicular lattice, extends to
the full lattice because the perpendicular is saturated,
and is therefore represented by an integral vector.
Vanishing on the isotropic sublattice puts the representative
in its perpendicular. The ambiguity of representative is
exactly the original primitive sublattice. This proves the
claim first for \(\Lambda'\), then for \(J_i\), since \(U_i\)
is primitive in \(\Lambda'\).

At each \(\mathfrak w\mid Q\) the reductions of the two
\(A_i\) agree in \(\Lambda\) by the marked flag condition.
They also agree after first taking \(\Lambda\cap I^\perp\)
and then quotienting by \(M\): the first sublattice is
saturated, so reduction embeds it in the original
reduction, and the second step is a common quotient.
Both \(U_i\) reduce to the same \(t\)-dimensional subspace
since they are primitive. The kernel of reduction
\(D/\mathfrak wD\to\Lambda'/\mathfrak w\Lambda'\)
therefore has dimension \(t\), using local lattices here.
Locally, dividing lifts of this kernel by a uniformizer
gives precisely \(R_1/D\). There is no quotient away
from \(Q\), and both marked primes split completely.
Multiplication of the local orders gives \(Q^{2ft}\).

For a nondegenerate integral Hermitian lattice \(B\) put
\(\delta(B)=[B^\sharp:B]^{1/2}\), with the dual taken in
its rational span. Enlarging a lattice by finite index \(j\)
divides \(\delta\) by \(j\), by dualizing the inclusion.
In integral local bases the Gram matrix for \(D\) consists
of its cross pairing and its conjugate transpose. Thus the
determinant ideal is \((Q)^{2t}\mathfrak b\bar{\mathfrak b}\),
also for projective lattices by the change-of-basis rule.
Its absolute norm is \([D^\sharp:D]\). Consequently
\(\delta(D)=Q^{2ft}N\mathfrak b\) and
\(\delta(R)=N\mathfrak b/[R:R_1]\).

The projection of \(\Lambda'\) to \(S\) has image \(R^\sharp\).
This follows by extending integral functionals on the
saturated \(R\) and applying self-duality. Its kernel is
\(C\), so reduction modulo \(R\) identifies
\(\Lambda'/(R\oplus C)\) with \(R^\sharp/R\).
Projection to \(S^\perp\) gives the same identification
with \(C^\sharp/C\). Hence \(\delta(R)=\delta(C)\).
Since \(J_i\) is self-dual and contains \(C\) with full
rank, \([J_i:C]=\delta(C)\) for both \(i\).
\end{proof}

\begin{theorem}[A further divisor from the common terminal line]
\leavevmode\label{thm:middleterminal570}
One has \(\zeta_0=\beta\zeta_1\). Put
\[
   L_i=J_i\cap\zeta_i^\perp,\qquad L_c=C\cap\zeta_0^\perp,
   \qquad m_i=[L_i:L_c].
\]
There is an integral ideal \(\mathfrak k\) contained in
\(\mathfrak q_\beta\) of Theorem~\ref{thm:markedindex564}
such that
\[
  N\mathfrak b
    =[R:R_1]\,m_0\,u^f N\mathfrak d\,N\mathfrak k,
       \qquad m_1=u^f m_0 .
\]
In particular \(u^f N\mathfrak d\,N\mathfrak q_\beta\)
divides the integer \(N\mathfrak b\).
\end{theorem}
\begin{proof}
The last step in Lemma~\ref{lem:suffixdet556} supplies a
common vector in the two extended flags with coefficients
\(1,\beta\) on their terminal vectors, after scalar
normalization. Thus \(z_0-\beta z_1\in X_0+X_1\).
Projecting to \(S^\perp\) after quotient by \(I\) kills
this difference and proves the first equality. Each
\(\zeta_i\) is primitive in its own \(J_i\), by the
integral terminal condition of Section~\ref{sec:flagideals554}.
Hence the image of \(h(\zeta_i,-)\) on \(J_i\) is the unit
ideal, by self-duality. Let
\(\mathfrak h_i=h(\zeta_i,C)\). This is a nonzero integral
ideal since \(C\) has full rank in the middle space.
The kernels of the two evaluations on \(J_i,C\) are
\(L_i,L_c\). Passing to their quotients gives the exact
index equality
\[
                   [J_i:C]=m_i\, N\mathfrak h_i .
\]
Since \(h\) is conjugate-linear in its first argument,
\(\mathfrak h_0=(\bar\beta)\mathfrak h_1\). Taking the
smaller prime exponent at every prime in these two
integral ideals gives
\[
  \mathfrak k=\mathfrak h_0+\mathfrak h_1,\qquad
  \mathfrak h_0=\mathfrak k\,\bar{\mathfrak n},\qquad
  \mathfrak h_1=\mathfrak k\,\bar{\mathfrak d}.
\]
In particular \(N\mathfrak h_0=u^fN\mathfrak d\,N\mathfrak k\).

At both marks the reduced terminal vectors agree modulo the
common reduced outer space, by the terminal part of the
marking in Section~\ref{sec:bothordinary553}. The two terminal
frequencies have the same marked reduction since
\(\varpi\equiv1\bmod Q^2\). An element of \(C\) has an
integral representative perpendicular to both outer spaces:
lift it to \(\Lambda\cap I^\perp\) and use its
orthogonality to \(S\). Pairing against this representative
gives \(h(\zeta_0-\zeta_1,c)\in Q\mathcal O_E\) for every
\(c\in C\). Thus
\[
                (\bar\beta-1)\mathfrak h_1
                               \subseteq Q\mathcal O_E .
\]
At a marked prime where \(\bar\beta\) fails to be a unit
with residue one, the valuation of
\((\bar\beta-1)\bar{\mathfrak d}\) is zero. The congruence
therefore puts \(\mathfrak k\) in each such prime. These
are the primes of \(\mathfrak q_\beta\), by its conjugation
invariance in Theorem~\ref{thm:markedindex564}. This
proves the ideal inclusion, also when \(\beta=1\) and
\(\mathfrak q_\beta\) is the unit ideal. Substituting the
two ideal norms in Lemma~\ref{lem:middleindex570} proves
both exact formulas and the divisibility.
\end{proof}

\begin{corollary}[The mixed period through the first smaller intersection]
\leavevmode\label{cor:middleperiod570}
Put \(D_\beta=N\mathfrak d\,N\mathfrak q_\beta\).
Every term under consideration satisfies
\[
  \widetilde V_0(y)\widetilde V_1(\gamma^{-1}y)
      \ \ge\ Q^f\,
        \frac{u^{\,f(e+1)/(2t)}}{(T_0T_1)^{f/2}}\,
        D_\beta^{(2e+1)/(2t)} .
\]
For unequal terminal norms, write \(Z_i(t)\) for the full
signed sum in this intersection dimension at marked index \(i\).
Uniformly for sufficiently large \(x\) on the common intervals,
\[
 \sum_i\|Z_i(t)\|_{\rm HS}
   \ \le\ |c_{0,1}|\,Q^{-fa/2}T^{-fa(e+1-t)/(2t)}
       \min\{\widetilde I_0\widetilde U_1,\,
                          \widetilde I_1\widetilde U_0\}.
\]
In particular the estimate with the power of \(T\) omitted
holds on the whole union \(1\le t<s,\ 2t\le s+1\),
including \(t=e+1\).
\end{corollary}
\begin{proof}
In Theorem~\ref{thm:denomheight562} insert
\(N\mathfrak g\ge(N\mathfrak q_\beta)^e\) and
\(N\mathfrak b\ge u^fD_\beta\). The first inequality
comes from Theorem~\ref{thm:markedindex564} and the
second from Theorem~\ref{thm:middleterminal570}. Taking
the indicated \(t\)-th roots gives the first assertion.
For \(T_0=T,T_1=1\), use \(D_\beta\ge1\); the resulting
common height factor is \(Q^fT^{f(e+1-t)/(2t)}\).
Extract its negative \(a\)-th power in the absolute
unfolding of the signed sum as in
Theorem~\ref{thm:mixedhalf556}. The translated-ball sum
costs at most \(Q^{fa/2}\) and is valid with either row
unfolded first. This proves the second bound with the same
real section norms and original finite operators. On the
specified union the common height is at least \(Q^f\).
Extract that factor once before summing over cosets to
obtain the last assertion.
\end{proof}

For scalar one the integer \(m_0=m_1\) measures the
positive lattices directly. For example \(m_0=1\) means
that both are exactly \(L_c\) in the common rational
positive column, even when the terminal norm is \(T\).
Their self-dual enlargements at the primes dividing
\(\varpi\) in Lemma~\ref{lem:completedcolumn535} then
also agree. This lattice comparison leaves the signed
sum of the real orbital integrals on that column to be
estimated.


\subsection{The positive lattice in the signed period}
\label{sec:positivecomparison576}

Fix two rows on \(E=E_x\) on the interval of
Section~\ref{sec:missing}, with \(Q=g_1g_2\), \(F=E^+\),
\(f=[F:\QQ]\), and order them so that \(T_0\ge T_1\).
Write \(u=T_0/T_1\). Consider supported terms \(\gamma\)
with compatible characters and \(0<\dim(X_0\cap X_1)<s\).
The two positive spaces in an individual term are identified
by the common space \(S^\perp\cap\zeta_0^\perp\) of
Section~\ref{sec:commonmiddle570}. Indeed the two terminal
classes there are proportional, and taking their perpendicular
in \(S^\perp\) removes the same negative line. Composing the
identifications with the row coordinates defines a rational
isometry \(j_\gamma:P_0\to P_1\) between the fixed positive
spaces, of dimension \(k\). Write \(L_0,L_1\) for their original
positive lattices. The two copies of
\(L_c=C\cap\zeta_0^\perp\) in these lattices depend on the
term. Denote them by \(C_{b,\gamma}\subset L_b\), \(b=0,1\).
The isometry takes \(C_{0,\gamma}\) onto \(C_{1,\gamma}\).
Put
\[
 \mathfrak a_{b,\gamma}
   =\operatorname{Ann}_{\mathcal O_E}(L_b/C_{b,\gamma}),
 \qquad m_{b,\gamma}=[L_b:C_{b,\gamma}].
\]
Theorem~\ref{thm:middleterminal570} gives
\(m_{1,\gamma}=u^f m_{0,\gamma}\).

\begin{lemma}[Comparison on the original marked lattices]
\leavevmode\label{lem:scaledmark576}
Use the orientation \(0\to1\) for the marked index \(i_\gamma\).
Multiplication by any \(\alpha\in\mathfrak a_{0,\gamma}\)
makes \(\alpha j_\gamma:L_0\to L_1\) integral. At every prime
\(\mathfrak w\mid Q\) its reduction in the marked spaces is
\(\overline{\alpha}\,i_{\gamma,\mathfrak w}\), where the bar
here denotes residue. The inverse statement holds for
\(\mathfrak a_{1,\gamma}\) and \(i_{\gamma,\mathfrak w}^{-1}\).
Moreover \(N\mathfrak a_{b,\gamma}\le m_{b,\gamma}\), with
\(N=N_{E/\QQ}\) on ideals.
\end{lemma}
\begin{proof}
For each \(w\in L_0\) the element \(\alpha w\) belongs to the
image of \(L_c\); the image of the same element at the other
end is \(\alpha j_\gamma w\in L_1\). Locally this element of
\(L_c\) lifts to a vector of \(\Lambda\cap I^\perp\). Its
class modulo \(I\) lies in \(S^\perp\) and pairs to zero with
both terminal classes. Its lift therefore pairs to zero with
both outer spaces and both terminal vectors. In the two
marked reductions the same lift gives the same class modulo
the common reduced outer space. At the marks the terminal
norm is a unit. The positive lattice there reduces to the
perpendicular of the terminal class in the middle quotient,
since evaluation on the terminal class splits off the
nondegenerate line. Reading the two images of the lift
in the prescribed bases gives precisely the change of
coordinates \(i_{\gamma,\mathfrak w}\). This is also true
in the conjugate split factor, using the dual bases. Applied
to \(\alpha w\) the assertion proves the reduction formula,
even if \(\alpha\) has zero residue. Interchange the ends for
the inverse. Finally, in the elementary divisor decomposition
of each local quotient the annihilator exponent is the largest
elementary divisor exponent, bounded by their sum. Taking
ideal norms gives the stated bound.
\end{proof}

\begin{theorem}[Separation with the positive lattice indices]
\leavevmode\label{thm:positivecompare576}
Consider two such terms \(\gamma,\delta\) in the fixed pair of
rows, allowing different intersection dimensions and different
common graded scalars. For \(b=0,1\) set
\(\mathfrak d_b=\mathfrak a_{b,\gamma}\cap\mathfrak a_{b,\delta}\).
If their marked indices occupy a common output entry in one
row family of the projection \(H\), put \(d=1\). If the
indices are equal one may put \(d=2\). Whenever
\(j_\gamma\ne j_\delta\), one has
\[
 Q^d\ \le\ 4k\,\bigl(N(\mathfrak d_0\mathfrak d_1)\bigr)^{1/(2f)}
 \ \le\ 4k\,u\,(m_{0,\gamma}m_{0,\delta})^{1/f}.
\]
Both inequalities include the primes dividing the lattice indices.
\end{theorem}
\begin{proof}
Write \(j=j_\gamma\), \(j'=j_\delta\). On \(P_0\) form the trace
\[
 h=\operatorname{tr}\big((j^{-1}-{j'}^{-1})(j-j')\big)
    =2k-\operatorname{tr}(j^{-1}j')-\operatorname{tr}({j'}^{-1}j).
\]
The two traces on the right are complex conjugates by unitarity,
so \(h\in F\). At each embedding of \(F\) it is the squared
Hilbert--Schmidt norm of the difference of the two unitary
maps. If these maps differ, this gives \(0<h_\tau\le4k\) at
every embedding. For local elements
\(\alpha\in\mathfrak d_0,\beta\in\mathfrak d_1\) the two maps
\(\alpha(j-j')\) and \(\beta(j^{-1}-{j'}^{-1})\) are integral
in opposite directions. Thus \((h)\mathfrak d_0\mathfrak d_1\)
is an integral ideal of \(E\).

At a marked prime reduce these two integral maps by
Lemma~\ref{lem:scaledmark576}. Write \(i,i'\) for the two
residue indices. Their product has trace equal to the reduction
of \(\alpha\beta\) times
\(2k-\operatorname{tr}(i^{-1}i')-\operatorname{tr}({i'}^{-1}i)\).
In a common row entry the relative index is unipotent by
Lemma~\ref{lem:whittakermatrix526}. Its inverse and its
contragredient in the conjugate factor are unipotent as well.
The displayed expression is therefore zero at every prime of
\(E\) above both marks. If \(i=i'\), each of the two integral
maps has zero reduction, giving divisibility of their product
by the square of the local prime. Taking generators of the
two localized ideals proves in both cases
\[
           (h)\mathfrak d_0\mathfrak d_1\subseteq Q^d\mathcal O_E.
\]
All marked primes have ramification index one. Taking absolute
norms, then using \(|N_{E/\QQ}h|\le(4k)^{2f}\), proves the
first inequality. Each intersection ideal contains the product
of its two constituent ideals. Lemma~\ref{lem:scaledmark576}
and \(m_{1,\nu}=u^fm_{0,\nu}\) for \(\nu=\gamma,\delta\)
give the second inequality.
\end{proof}

\begin{corollary}[Integral recovery and the contribution at large index]
\leavevmode\label{cor:indexperiod576}
In any collection with \(m_{0,\gamma}^{1/f}\le B\) and
\(4kuB^2<Q\), the signed projected array of each rational
isometry is recovered by selecting entries from the total
array of the collection in each row family. Its squared norm
is additive over distinct isometries. In particular, at
either equal terminal norm the criterion is \(4kB^2<Q\).

For an individual absolutely convergent orbit of identical
words, allowing either frequency at each end, the arithmetic
average fixes its operator \(B_\gamma\) if
\[
                  4k\sqrt u\,m_{0,\gamma}^{1/f}<Q^2 .
\]
Here the average is the one over \(M_{i_\gamma}\) in
Theorem~\ref{thm:orbitaverage532}.

For \(L\ge1\), let \(Z_i(L)\) sum the signed terms at index
\(i\) over all proper nonzero intersections for which
\[
 m_{0,\gamma}[R:R_1]D_\beta^{\,2e+1}
      \ \ge\ L^{2t}(T_0T_1)^{ft}u^{-f(e+1)},\qquad
       D_\beta=N\mathfrak d\,N\mathfrak q_\beta .
\]
Here \(e=\dim I\), \(t=s-e\) and the ideals and lattices are
those of Section~\ref{sec:commonmiddle570}, evaluated in each
term. On the same intervals and retained rows one has
\[
 \sum_i\|Z_i(L)\|_{\rm HS}\ \le\
 |c_{0,1}|\,Q^{-fa/2}L^{-a}
       \min\{\widetilde I_0\widetilde U_1,
                       \widetilde I_1\widetilde U_0\}.
\]
\end{corollary}
\begin{proof}
Theorem~\ref{thm:positivecompare576} with \(d=1\) makes the
output supports of distinct isometries disjoint under the
first inequality. Coordinate projections onto these supports
are integral and orthogonal. This proves the array statements
for finite sums and then for absolutely convergent sums by
continuity.

For \((a,b)\in M_{i_\gamma}\) apply the comparison to
\(j_\gamma\) and \(a j_\gamma b^{-1}\), where \(b,a\) act on
the source and target lattices. These translates have the same
marked index by the definition of \(M_{i_\gamma}\).
The translated common sublattices are \(bC_{0,\gamma}\) and
\(aC_{1,\gamma}\); the two quotient annihilator ideals are
unchanged because \(b,a\) are integral lattice automorphisms.
Thus the first bound in Theorem~\ref{thm:positivecompare576}
uses \(N(\mathfrak d_0\mathfrak d_1)\le m_{0,\gamma}m_{1,\gamma}\).
The stated strict inequality forces \(aj_\gamma=j_\gamma b\).
The common positive perpendicular has full rank \(k\) in both
columns, as in the construction of \(j_\gamma\).
The intertwining identity of Theorem~\ref{thm:common133}
therefore fixes \(B_\gamma\) under every pair in the average.

For the last assertion retain the factor \(m_{0,\gamma}[R:R_1]\)
when substituting Theorem~\ref{thm:middleterminal570} into
Theorem~\ref{thm:denomheight562}. The factorization there,
\(N\mathfrak g\ge (N\mathfrak q_\beta)^e\) from
Theorem~\ref{thm:markedindex564}, and
\(N\mathfrak k\ge N\mathfrak q_\beta\) give
\[
 \big(\widetilde V_0\widetilde V_1\big)^{2t}
   \ \ge\ Q^{2ft}\frac{u^{f(e+1)}}{(T_0T_1)^{ft}}\,
                          m_{0,\gamma}[R:R_1]D_\beta^{\,2e+1}.
\]
Every term included in \(Z_i(L)\) has height product at least
\(Q^fL\). Extract its negative \(a\)-th power once before
unfolding over all the cosets. The translated-lattice count
of Theorem~\ref{thm:mixedhalf556} costs at most \(Q^{fa/2}\)
and applies in both directions and at both frequencies by
Corollary~\ref{cor:denomperiod562}. This gives the claimed
estimate without summing a separate error for each dimension.
\end{proof}

For example, at unequal frequencies and \(t=e+2\), every term
with \(m_{0,\gamma}\ge T^f\) is in the last estimate with
\(L=1\), independently of its scalar. Different terms with
one and the same rational isometry can contribute to the
same projected entries with different archimedean phases.
The separation bound controls distinct isometries. A bound
for the signed sums within the remaining isometries and for
transverse terms is still needed in the integral comparison
of Section~\ref{sec:missing}.


\subsection{The determinant comparison at transverse flags}
\label{sec:cofactor581}

Fix two marked rows over \(E=E_x\) in the interval of
Section~\ref{sec:missing}. Put \(F=E^+\), \(f=[F:\QQ]\) and
\(Q=g_1g_2\). We allow either original terminal frequency at
each end. Consider a cell \(\gamma\) with common positive
hyperplane lattice \(R=R_\gamma\) of rank \(k-1\).
Use the two self-dual column lattices \(\widehat L_s,\widehat L_t\)
and determinant map \(d_\gamma^\circ\) of
Lemma~\ref{lem:completedcolumn535}. Write \(j_\gamma:P_s\to P_t\)
for the rational isometry extending the identification on \(R\)
with that determinant. Let \(i_\gamma\) be the cell's prescribed
period index, oriented from source to target. At a prime
\(\mathfrak p\) of \(E\) define
\[
 v_{\mathfrak p}(\mathfrak t_\gamma)
   =\operatorname{length}_{\mathcal O_{E,\mathfrak p}}
        (R^\sharp/R)_{\mathfrak p},\qquad
 h_\gamma=(N_{E/\QQ}\mathfrak t_\gamma)^{1/(2f)} .
\]
Here \(h_\gamma=h_{R_\gamma}\) in the notation of
Theorem~\ref{thm:mixedtransport535}, and \(R^\sharp\) denotes the
Hermitian dual inside the rational hyperplane. The completed lattices and the original lattices
agree at the two marks.

\begin{lemma}[A cofactor identity at a marked prime]
\leavevmode\label{lem:cofactormark581}
For every \(\alpha\in\mathfrak t_\gamma\), the maps
\(\alpha j_\gamma\) and \(\alpha j_\gamma^{-1}\) are integral
between the respective completed lattices. At every prime over
\(Q\) their reductions are \(\bar\alpha i_\gamma\) and
\(\bar\alpha i_\gamma^{-1}\), respectively. In particular, if
\(n=v_{\mathfrak p}(\mathfrak t_\gamma)>0\) at such a prime,
both unscaled maps have denominator exponent at most \(n-1\)
there.
\end{lemma}
\begin{proof}
The integral assertions away from the marks follow from
Theorem~\ref{thm:mixedtransport535}. Fix one of the two split
factors over a marked place, with integer ring \(D\) and
uniformizer \(\pi\). Every calculation below takes place in this
discrete valuation ring or its fraction field. At either end choose a basis of the
completed lattice in this factor and its dual basis in the
conjugate factor. Use the same basis of the abstract lattice
\(R\) at the source and target in each factor. Its columns in
the first factor give \(k\)-by-\((k-1)\) matrices \(A_s,A_t\).
The conjugate basis paired against the first factor gives
\((k-1)\)-by-\(k\) matrices \(B_s,B_t\). Isometry gives
\[
   G=B_sA_s=B_tA_t,\qquad \Delta=\det G\ne0 .
\]
In these bases \(G\) is the matrix of the pairing on \(R\).
Thus \(\Delta\) generates \(\mathfrak t_\gamma D\) by elementary
divisors. Put \(d=\det j_\gamma\) in the two chosen column bases.
It is a unit with residue \(\det i_\gamma\) by
Lemma~\ref{lem:completedcolumn535}.

Write \(a(A)\) for the cofactor row satisfying
\(a(A)y=\det[A,y]\). For \(B\) let \(n(B)\) be the cofactor
column with \(m n(B)=\det\left[\begin{smallmatrix}B\\m\end{smallmatrix}\right]\)
for each row \(m\). Then \(B n(B)=0\) and
\(a(A)n(B)=\det(BA)\) by Cauchy--Binet. In particular
\([A_b,n(B_b)]\) has determinant \(\Delta\) at both ends
\(b=s,t\). The isometry identifies the first \(k-1\) columns,
takes the last line to the last line by orthogonality, and
takes \(n(B_s)\) to \(d\,n(B_t)\) by its determinant. Evaluating
on this basis gives
\begin{equation}\label{eq:cofactor581}
 \Delta j_\gamma
       = A_t\operatorname{adj}(G)B_s+d\,n(B_t)a(A_s).
\end{equation}
Here \(\operatorname{adj}\) is the classical adjugate, including
the value \(1\) for a one-by-one matrix. Both summands are integral.

For each vector of \(R\), the two marked reductions of its images
agree after the prescribed change of coordinates. Indeed the
vector comes from one representative in \(\Lambda\cap A\) in
the common rational perpendicular. In the two middle quotients
its reductions are the classes of that same representative
modulo the same reduced isotropic space. This argument applies
to the conjugate factor as well, regardless of the rank of
the reduced hyperplane. Writing \(i=i_\gamma\bmod\mathfrak p\)
in the chosen bases, we get
\[
    \bar A_t=i\bar A_s,\qquad \bar B_t=\bar B_s i^{-1}.
\]
For an invertible matrix \(u\), expansion of a determinant by
its last row gives \(n(Bu^{-1})=(\det u)^{-1}u\,n(B)\).
Reduce the right side of \eqref{eq:cofactor581} using these
identities and \(\bar d=\det i\). The result is \(i\) times the
same expression for the identity transport on the source,
which equals \(\bar\Delta i\). This polynomial calculation
also holds when \(\bar\Delta=0\). Interchanging endpoints gives
the inverse formula. Multiplication by any local multiple of
\(\Delta\) proves the stated congruences. When \(\pi\mid\Delta\),
all entries of \(\Delta j_\gamma^{\pm1}\) consequently belong
to \(\pi D\), giving the one-power improvement.
\end{proof}

The local exponent in the lemma is sharp for the lattice data
used there. In rank two over one split factor take \(A_s=A_t=(1,0)^{\mathsf T}\),
\(B_s=(\pi^n,1)\), \(B_t=(\pi^n,1+\pi)\), and \(i=1,d=1\).
The bases in the conjugate column spaces are the dual bases.
The common pairing is \(\pi^n\), and
\eqref{eq:cofactor581} gives
\(j=\left(\begin{smallmatrix}1&-\pi^{1-n}\\0&1\end{smallmatrix}\right)\).
For \(n=1\) it is integral with nonidentity reduction. For
\(n>1\) its denominator exponent is \(n-1\).

For each direction \(\sigma=+,-\) let \(D_\gamma^\sigma\)
be the ideal of scalars \(\alpha\in\mathcal O_E\) such that
\(\alpha j_\gamma^\sigma\) is integral on the completed lattices
and has reduction \(\bar\alpha i_\gamma^\sigma\) at both marks.
Here the minus sign denotes the inverse and the plus sign
the original map. The preceding lemma gives
\(\mathfrak t_\gamma\subseteq D_\gamma^\sigma\).
The ideal can be calculated more precisely. At an unmarked prime
use the ordinary denominator ideal of the matrix. At a marked
prime, lift the residue matrix \(i_\gamma^\sigma\) to an integral
matrix \(I_\sigma\), and put \(M_\sigma=\Delta j_\gamma^\sigma\)
using \eqref{eq:cofactor581} in the appropriate direction.
Then
\[
  v_{\mathfrak p}(D_\gamma^\sigma)=
  \max\{0,\ 1+v_{\mathfrak p}(\Delta)
                -\min_{a,b}v_{\mathfrak p}
                   ((M_\sigma-\Delta I_\sigma)_{ab})\}.
\]
The minimum for a zero matrix is \(+\infty\). In fact the single
condition \(\alpha(j_\gamma^\sigma-I_\sigma)\in\pi M_k(D)\)
is equivalent to integrality with the prescribed reduction.
It is unchanged by the choice of the lift.

\begin{theorem}[Separation of transverse marked terms]
\leavevmode\label{thm:transverseseparate581}
Take two hyperplane cells \(\gamma,\delta\) in the same pair
of rows and put \(D^\sigma=D_\gamma^\sigma\cap D_\delta^\sigma\).
Suppose their indices contribute to a common entry in one row
family of the signed projection. Set \(m=1\), or take \(m=2\)
if the two indices agree. If \(j_\gamma\ne j_\delta\), then
\[
 Q^m\le 4k\,\bigl(N_{E/\QQ}(D^+D^-)\bigr)^{1/(2f)}
                  \le 4k(h_\gamma h_\delta)^2 .
\]
\end{theorem}
\begin{proof}
Let \(j=j_\gamma,\ j'=j_\delta\). The element
\[
 z=\operatorname{tr}\bigl((j^{-1}-{j'}^{-1})(j-j')\bigr)\in F
\]
is the squared Hilbert--Schmidt distance of the two unitary maps
at each embedding. Thus \(0<z_\tau\le4k\) everywhere when
the maps differ. Multiplication by elements of \(D^+,D^-\)
makes the corresponding differences integral in opposite
directions. At a marked prime their reductions are the scaled
differences of the indices in both directions, by definition.
The trace of the product vanishes on reduction for a common
row entry, since the relative index and its inverse are
unipotent by Lemma~\ref{lem:whittakermatrix526}. For equal
indices each factor has entries divisible by the local
uniformizer. Using local generators of the two ideals gives
\((z)D^+D^-\subseteq Q^m\mathcal O_E\). The marked primes are
unramified. Take norms and use
\(\lvert N_{E/\QQ}z\rvert\le(4k)^{2f}\) to obtain the first
inequality. The product \(\mathfrak t_\gamma\mathfrak t_\delta\)
lies in each \(D^\sigma\), giving the second inequality.
\end{proof}

\begin{corollary}[Integral recovery and the transverse orbit trace]
\leavevmode\label{cor:transversetrace581}
In any absolutely convergent collection of these cells with
\(h_\gamma\le B\) and \(4kB^4<Q\), every projected array
belonging to a distinct rational transport can be selected
integrally in each row family. Its squared norm is additive
over the transports.

For an individual absolutely convergent orbit of identical
words, the signed trace of Proposition~\ref{prop:hyperplanetrace534}
survives the arithmetic average whenever
\[
 4k\,\bigl(N_{E/\QQ}(D_\gamma^+D_\gamma^-)\bigr)^{1/(2f)}
       < Q^2 .
\]
In particular \(4k h_\gamma^2<Q^2\) suffices. Both statements
hold for the two original terminal frequencies in any combination.
\end{corollary}
\begin{proof}
Under the first cutoff, Theorem~\ref{thm:transverseseparate581}
excludes a common output entry for different maps. Selecting
disjoint coordinates in an orthogonal row basis is integral
and gives the norm identity. One may also use the first bound
in that theorem separately for each pair.

For the second assertion let \((a,b)\) stabilize the index
\(i_\gamma\) under the two endpoint actions. The transport
\(aj_\gamma b^{-1}\) has the same marked index. Since \(a,b\)
are integral automorphisms of the completed lattices and their
reductions stabilize this index, its two ideals of scalars
are exactly \(D_\gamma^\pm\), by conjugating the defining
conditions and using the inverse automorphisms. The proof of
the theorem with \(m=2\) therefore forces \(aj_\gamma=j_\gamma b\).
The unitary identification induced by \(j_\gamma\) on the
two column types is fixed by every member of the averaging
group. Taking its inner product with the orbital operator
before and after averaging gives the same signed trace, which
is the sum of the branch traces in
Proposition~\ref{prop:hyperplanetrace534}. The inclusion
\(\mathfrak t_\gamma\subseteq D_\gamma^\pm\) gives the
stated sufficient cutoff.
\end{proof}

All the comparisons use the original two marked primes. Their
conclusion concerns different rational transports and the
trace in a single orbit. By Theorem~\ref{thm:everydiagonal536},
infinitely many signed transverse terms can have the same transport
and the same marked index even when \(D_\gamma^+=D_\gamma^-=\mathcal O_E\).

\section{Gram determinants and height truncation}\label{nav:chapter-transverse-determinants}

\subsection{Gram determinants and the transverse orbital sum}
\label{sec:transversesum582}

Fix retained rows \(0,1\) at one \(x\) in
Section~\ref{sec:missing}, with terminal norms \(-T_0,-T_1\).
Transport the second frame into the first row as in
Section~\ref{sec:mixedsuffix556}. Write \(E=E_x\), \(F=E^+\),
\(f=[F:\QQ]\), \(Q=g_1g_2\), and \(d=s+1\). In the ordered
frames the first \(s\) vectors are isotropic and the last
is the terminal vector. For \(1\le j\le d\) let \(D_j\) be the
determinant of the cross pairing of their first \(j\) vectors.
Write \(\mathfrak h_j=\prod_{i\le\min(j,s)}\mathfrak a_i\)
for the product of the prescribed outer ideals, and put
\(\eta(\mathfrak h_j)=(N_{E/\QQ}\mathfrak h_j)^{-1/(2f)}\).
A change of either frame by an allowed upper unitriangular
matrix fixes all the \(D_j\), by taking determinants of its
initial triangular blocks.

\begin{proposition}[The two frame heights and the transverse determinant]
\leavevmode\label{prop:transverseheight582}
At each real embedding \(\tau\) choose either complex
extension to \(E\). Define
\[
 C_j=\eta(\mathfrak h_j)^2(T_0T_1)^{j/2},\qquad
 \mathcal H_\gamma=
 \prod_\tau\max\left(\{1\}\cup
       \{(|\tau D_j|/C_j)^{1/j}:1\le j\le d\}\right).
\]
For every argument \(y\) of the orbital integral,
\[
  \widetilde V_0(y)\widetilde V_1(\gamma^{-1}y)
                      \ \ge\ \mathcal H_\gamma .
\]
When the two outer spaces are transverse and the common positive
lattice \(R\) has rank \(k-1\), let \(I_\Lambda\) and \(h_R\)
be the saturation index and discriminant height of
Corollary~\ref{cor:actualgram535}. Put
\[
    P_\gamma=\prod_\tau
       \left(1+\frac{T_0T_1|\tau D_s|^2}{\tau\mathcal G}\right).
\]
Then
\[
 \mathcal H_\gamma\ \ge\
 \frac{Q^f}{(T_0T_1)^{f/2}}
       \bigl(I_\Lambda h_R^f P_\gamma\bigr)^{1/(2d)} .
\]
\end{proposition}
\begin{proof}
Take the exterior product of the first \(j\) vectors in each
frame. The decorated stabilizer fixes this vector, including for
\(j=d\), since the terminal vector may change only by an
earlier vector. In row \(b\) the exterior vector has length
\(\eta(\mathfrak h_j)T_b^{j/2}\) in the modified majorant of
Theorem~\ref{thm:mixedhalf556}. Indeed the outer vectors there
are mutually orthogonal with lengths \(\eta(\mathfrak a_i)\sqrt{T_b}\).
The terminal vector is orthogonal to them and has length
\(\sqrt{T_b}\). In the reference majorant the Hermitian form
is represented by a unitary involution. Its exterior pairing
satisfies Cauchy--Schwarz. Applying the two inverse changes
of frame to the wedges costs at most the \(j\)-th power of
each operator norm. Each change can first be multiplied by
an arbitrary member of its own decorated stabilizer. Taking
their infima proves that the local product of the two heights
is at least \((|\tau D_j|/C_j)^{1/j}\). Each height is at
least one since the changes are isometries between compatible
majorants and their singular values come in reciprocal pairs.
Take a maximum before multiplying over the real places.

In the hyperplane case \(\mathcal G=D_d\overline{D_d}
-T_0T_1D_s\overline{D_s}\) is totally positive by
Theorem~\ref{thm:cross132}. It follows that
\[
 |N_{E/\QQ}D_d|=|N_{F/\QQ}\mathcal G|\,P_\gamma .
\]
Corollary~\ref{cor:actualgram535} evaluates the first factor as
\(Q^{2df}I_\Lambda h_R^f/(N_{E/\QQ}\mathfrak h_d)^2\).
The \(j=d\) term of the maximum gives
\(\mathcal H_\gamma^d\ge |N_{E/\QQ}D_d|^{1/2}/C_d^f\).
Substitution cancels the fractional outer ideal and proves
the assertion.
\end{proof}

\begin{theorem}[An absolute estimate for a union of transverse terms]
\leavevmode\label{thm:transversesum582}
Let \(L\ge1\). For both terminal frequencies, take all supported
transverse double cosets for which \(\mathcal H_\gamma\ge Q^f L\).
Write \(B_\gamma\) for each original signed ordinary orbital
operator with its finite weight. On the common intervals, for
all sufficiently large \(x\),
\[
 \sum_\gamma\|B_\gamma\|_{\rm HS}
 \ \le\
 |c_{0,1}|\,Q^{-fa/2}L^{-a}
       \min\{\widetilde I_0\widetilde U_1,\,
                     \widetilde I_1\widetilde U_0\}.
\]
Here \(a,\widetilde I_b,\widetilde U_b\) have precisely the
normalization of Theorem~\ref{thm:mixedhalf556}. The same bound
holds for the sum of the norms after grouping by marked index,
including within any one rational transport. The sum includes
every transverse hyperplane term with
\[
         I_\Lambda h_R^f P_\gamma\ge
                    (T_0T_1)^{df} L^{2d}.
\]
\end{theorem}
\begin{proof}
The norm of the rank-one integrand is the product of the two
section norms, and the finite phases have absolute value one.
Put \(F_b(y)=\widetilde V_b(y)^a\|W_b(y)\|\). On every one
of the stated cosets Proposition~\ref{prop:transverseheight582}
bounds the product of the two section norms by
\[
             Q^{-fa}L^{-a} F_0(y)F_1(\gamma^{-1}y).
\]
After extracting that factor, enlarge the nonnegative sum to
all double cosets in the row. Its unfolding in
Theorem~\ref{thm:mixedordinary541} is the integral over
\(N_0\backslash G_\infty\) of \(F_0\) times the translated
coset sum of \(F_1\). The box estimate in the proof of
Theorem~\ref{thm:mixedhalf556} bounds that coset sum by
\(Q^{fa/2}\widetilde U_1\) for all sufficiently large \(x\).
Its constants depend on the outer ideals only through the
packing ratio in which they cancel. The argument uses one
row at a time and covers equal frequencies as in
Corollary~\ref{cor:denomperiod562}. Integration gives
\(\widetilde I_0\), and inversion of the cosets gives the
other product of norms. Restoring the pair volume proves the
estimate before any signed summation. The triangle inequality
proves its grouped form. The hyperplane cutoff follows from the
second inequality of Proposition~\ref{prop:transverseheight582}.
\end{proof}

\begin{corollary}[The cross pairing at the terminal prime]
\leavevmode\label{cor:terminalprefix582}
Write \(b_T=\#\{b\in\{0,1\}:T_b=T\}\). For \(1\le j\le s\)
with \(D_j\ne0\), take the cross pairing on the two integral
prefix lattices at each \(\mathfrak p\mid T\), using the conjugate
factor for the first lattice. Let \(u_{\mathfrak p,j}\) be the
nullity of its reduction. The estimate of
Theorem~\ref{thm:transversesum582} with \(L=1\) includes every
transverse term satisfying, for at least one such \(j\),
\[
             \sum_{\mathfrak p\mid T}u_{\mathfrak p,j}\ge fjb_T .
\]
In particular, for unequal frequencies it suffices that the
prefix pairing vanish on reduction at one prime of \(E\)
over each place of \(F\) above \(T\). For \(T_0=T_1=T\)
vanishing at both factors suffices. The union costs the single
bound in the theorem.
\end{corollary}
\begin{proof}
The ideal
\(\mathfrak A_j=(D_j/Q^j)\overline{\mathfrak h_j}\mathfrak h_j\)
is integral. Away from \(Q\) this is integrality of the cross
pairing on integral bases. Over \(Q\) the two reduced outer
spaces coincide and are totally isotropic, giving one local
uniformizer in every entry before taking the determinant.
At \(\mathfrak p\mid T\), at least \(u_{\mathfrak p,j}\)
elementary divisors of the pairing matrix in integral bases
are divisible by the uniformizer. Its determinant has the same
valuation as \(\mathfrak A_j\) there. These primes have norm
\(T\), since \(T\) splits completely. Thus
\[
  |N_{E/\QQ}D_j|^{1/2}
   =Q^{jf}\eta(\mathfrak h_j)^{2f}(N_{E/\QQ}\mathfrak A_j)^{1/2}
   \ge Q^{jf}\eta(\mathfrak h_j)^{2f}
                    T^{\sum_{\mathfrak p\mid T}u_{\mathfrak p,j}/2}.
\]
Using the \(j\)-th wedge in the maximum gives
\(\mathcal H_\gamma\ge Q^f\) under the stated inequality.
\end{proof}

\begin{corollary}[All powers of a repeated transverse transport]
\leavevmode\label{cor:cyclicuniform582}
On any retained diagonal row of terminal norm \(-T_0\), take
the cyclic groups of Theorem~\ref{thm:everydiagonal536} at a
precision giving \((\gamma-1)\Lambda\subseteq QT_0\Lambda\).
One may do so for every rational positive hyperplane. For
\(M\ge1\), sum over nonidentity powers \(\gamma^n\) with
\(|n|\ge M\), even over the union of all these groups in the
row, counting each contributing double coset once. Their
ordinary signed operators satisfy
\[
       \sum\|B_{\gamma^n}\|_{\rm HS}
       \le |c_{0,0}|\,Q^{-fa/2}M^{-af}
                                \widetilde I_0\widetilde U_0 .
\]
In particular \(M=1\) controls the whole indicated collection
at the identity marking and identity rational transport.
\end{corollary}
\begin{proof}
The construction of Theorem~\ref{thm:everydiagonal536} allows
the stated precision on the original lattice and gives
\(h(e_1,\gamma^n e_1)=nt\) with \(t\ne0\). For \(n=1\),
pair the integral line \(\mathfrak a_1e_1\) with its image
under \(\gamma-1\). The original line is isotropic. Consequently
\((t)\overline{\mathfrak a_1}\mathfrak a_1
\subseteq QT_0\mathcal O_E\).
It follows that
\[
   |N_{E/\QQ}(nt)|^{1/2}
      \ge |n|^f(QT_0)^f\eta(\mathfrak a_1)^{2f}.
\]
On this diagonal \(C_1=\eta(\mathfrak a_1)^2T_0\).
The first wedge thus gives \(\mathcal H_{\gamma^n}\ge
Q^f|n|^f\). All the indicated cosets occur in the single sum
of Theorem~\ref{thm:transversesum582} with \(L=M^f\). The
construction gives their common marking and transport.
\end{proof}

The estimates retain the archimedean factors of the original
sections. For the full comparison one still needs the remaining
transverse terms at small frame height and a bound at the
normalization of the integral theta maps.


\subsection{Separation of rational transports by their discriminants}
\label{sec:transportpacking589}

Fix two retained rows on a period field \(E=E_x\), with
\(F=E^+\), \(f=[F:\QQ]\), and \(Q=g_1g_2\). Consider cells
with common positive hyperplane lattice of rank \(k-1\).
For a cell \(\gamma\), let \(j_\gamma\) be its rational isometry
between the two rank-\(k\) positive spaces, \(i_\gamma\) its
marked index, and \(\mathfrak t_\gamma\) its discriminant ideal
as in Lemma~\ref{lem:cofactormark581}. Put
\(h_\gamma=(N_{E/\QQ}\mathfrak t_\gamma)^{1/(2f)}\).

\begin{theorem}[A product of discriminant ideals for a common projected entry]
\leavevmode\label{thm:transportproduct589}
Suppose two cells \(\gamma,\delta\) contribute to the same
entry of one projected row family and have distinct rational
transports. At a marked prime of \(E\), call a cell singular
if its discriminant ideal is divisible by that prime. Let
\(\mathfrak q_{\gamma,\delta}\) be the product of the marked
primes at which both cells are singular, together with those
at which both are nonsingular and the two indices agree locally.
Then
\[
 Q\,(N_{E/\QQ}\mathfrak q_{\gamma,\delta})^{1/(2f)}
 \ \le\ h_\gamma h_\delta
       \left(\prod_\tau
       \|j_{\gamma,\tau}-j_{\delta,\tau}\|_{\rm HS}^2\right)^{1/f}
 \ \le\ 4k h_\gamma h_\delta .
\]
The norms are for the positive Hermitian metrics of the two
rows at the real places of \(F\). In particular a common
entry requires \(Q\le4k h_\gamma h_\delta\). For equal
marked indices and identical sets of singular marked primes
it requires \(Q^2\le4k h_\gamma h_\delta\).
\end{theorem}
\begin{proof}
Write \(j=j_\gamma,\ j'=j_\delta\). In \(F\) put
\[
 z=2k-\operatorname{tr}(j^{-1}j')
          -\operatorname{tr}({j'}^{-1}j).
\]
The two traces are conjugates because the maps are unitary.
At every real embedding \(z\) is the squared Hilbert--Schmidt
distance of the maps, positive since \(j\ne j'\) and at most
\(4k\). The two compositions in this formula use one map
from each cell. Thus \(z\mathfrak t_\gamma\mathfrak t_\delta\)
is integral by Lemma~\ref{lem:cofactormark581}, which clears
both directions between the completed column lattices.

At a marked prime choose local generators \(a,b\) of
\(\mathfrak t_\gamma,\mathfrak t_\delta\), and write
\(\pi\) for the uniformizer. The same lemma identifies the
four reductions of \(aj,aj^{-1},bj',b{j'}^{-1}\) with
the corresponding scaled marked maps and their inverses.
Consequently \(abz\) on reduction equals \(\bar a\bar b\)
times \(2k-\operatorname{tr}(i_\gamma^{-1}i_\delta)
-\operatorname{tr}(i_\delta^{-1}i_\gamma)\). A common row
entry makes both relative marked maps unipotent, by
Lemma~\ref{lem:whittakermatrix526}, so this difference is zero.
If \(\pi\) divides both \(a\) and \(b\), each of the four
scaled maps has entries in \(\pi\) and \(ab\) is in \(\pi^2\).
It follows term by term that \(abz\in\pi^2\).
If \(a,b\) are units and the marked maps agree, the two
differences \(j-j'\) and \(j^{-1}-{j'}^{-1}\) have entries
in \(\pi\). The trace of their composition equals \(z\),
giving the same extra divisibility. This proves
\[
             (z)\mathfrak t_\gamma\mathfrak t_\delta
                     \ \subseteq\ Q\,\mathfrak q_{\gamma,\delta}.
\]
Here \(Q\) splits completely, so its local exponent is one
at each marked prime. Take absolute ideal norms. The
\(2f\)-th root of \(|N_{E/\QQ}z|\) is the geometric mean
of its positive real values. The display in the statement
follows. In the last case every marked prime occurs in
\(\mathfrak q_{\gamma,\delta}\).
\end{proof}

\begin{corollary}[Integral selection and the number of transports at an entry]
\leavevmode\label{cor:transportpacking589}
Retain cells with \(h_\gamma\le H\), \(H\ge1\), in an absolutely
convergent family. If \(4kH^2<Q\), the signed arrays grouped by
rational transport have disjoint support in each row family.
Each group can be recovered by integral coordinate selection,
and its squared norm adds to those of the others. On the
prescribed intervals this holds for
\(H=\exp((8-\delta)x^\varepsilon)\) with any fixed
\(0<\delta<8\) for all sufficiently large \(x\).

For arbitrary \(H\) the number of distinct rational transports
in a single output entry is at most
\[
             \left(1+2\sqrt{k}\,H/\sqrt Q\right)^{2fk^2}.
\]
For a fixed full marked index the bound can be replaced by
\[
                   2^{4f}\left(1+2\sqrt{k}\,H/Q\right)^{2fk^2}.
\]
All assertions allow the original terminal norm at each end
to be either \(1\) or \(T\) in absolute value.
\end{corollary}
\begin{proof}
The first conclusion follows from the necessary condition for
a common entry in Theorem~\ref{thm:transportproduct589}.
Coordinate selection has matrix entries zero and one, and
disjoint support gives orthogonality in the given row basis.
Here \(\log Q=(16+o(1))x^\varepsilon\) and \(k\) is polynomial
in \(x\) by Theorem~\ref{thm:orbits204}. These estimates imply
the asserted explicit cutoff.

Choose one cell for each distinct transport at a given entry.
Use orthonormal bases at infinity and place the \(f\) matrices
of each transport in \(\operatorname{Mat}_k(\CC)^f\), with
the maximum of the Hilbert--Schmidt norms. They lie in the
ball of radius \(\sqrt{k}\). The theorem implies that distinct
points have distance at least \(\sqrt Q/H\), since the maximum
dominates the geometric mean. Open balls of half this radius
about the points are disjoint. Comparing their volumes with
the enclosing ball gives the first bound in real dimension
\(2fk^2\), initially for every finite subset and therefore
for the whole set. For a fixed full marked index group the
chosen cells by their sets of singular marked primes. There
are at most \(2^{4f}\) choices, since each of the two rational
primes has \(2f\) factors. Within a group the distance is at
least \(Q/H\) by the last assertion of the theorem. The same
volume argument proves the other bound.
\end{proof}

\begin{corollary}[The projected sum below the frame-height cutoff]
\leavevmode\label{cor:smalltransverse589}
In a transverse hyperplane row pair use the height
\(\mathcal H_\gamma\) of Proposition~\ref{prop:transverseheight582},
the two terminal parameters \(T_0,T_1\), and \(d=s+1\).
For \(L\ge1\) put
\[
        H_L=(T_0T_1)^d L^{2d/f},\qquad
        m_L=(1+2\sqrt{k}\,H_L/\sqrt Q)^{2fk^2}.
\]
Retain the cells with \(\mathcal H_\gamma<Q^fL\).
Let \(S_j\) be their signed projected array for a single
rational transport \(j\) in a fixed row family, with the
actual arithmetic coefficients. Then, for absolutely
convergent arrays,
\[
          \left\|\sum_j S_j\right\|_{\rm HS}^2
                     \le m_L\sum_j\|S_j\|_{\rm HS}^2 .
\]
If \(4kH_L^2<Q\), the two squared norms in this comparison
agree after omitting \(m_L\). When \(T_0=T_1=1\) and \(L=1\)
the retained set is empty. All the complementary hyperplane
terms satisfy Theorem~\ref{thm:transversesum582}.
\end{corollary}
\begin{proof}
The determinant inequality of Proposition~\ref{prop:transverseheight582}
gives
\[
       I_\Lambda h_\gamma^f P_\gamma
                            < (T_0T_1)^{df}L^{2d}
\]
on the retained set. The lattice index \(I_\Lambda\) is at
least one and \(P_\gamma\ge1\) by total positivity. Thus
\(h_\gamma<H_L\). At most \(m_L\) of the signed groups can
occur in a given entry by Corollary~\ref{cor:transportpacking589}.
Cauchy--Schwarz entry by entry and summation give the
inequality, including for Hilbert-valued entries. Absolute
convergence justifies passage from finite subsums. Under the
smaller cutoff the groups have disjoint supports. In the
unit-terminal case with \(L=1\), the strict inequality
would require \(h_\gamma<1\), impossible for a lattice
discriminant ideal. The complement is exactly covered by
the frame-height bound in the cited theorem.
\end{proof}

This estimate controls how many distinct rational transports
can contribute to one output entry. The size of the signed
sum within a single transport at low frame height still
enters the integral period comparison.


\subsection{Counting cross matrices below the frame-height cutoff}
\label{sec:smallframes590}

Fix a pair of retained rows with their ordered rational frames
as in Section~\ref{sec:transversesum582}. Write \(O_E=\mathcal O_E\)
and \(A=T_0T_1\). For \(L\ge1\) we count the relevant transverse
hyperplane terms in this pair with \(\mathcal H_\gamma<Q^fL\).

\begin{theorem}[Possible rational flag pairs at small frame height]
\leavevmode\label{thm:smallframes590}
Put \(B=\lfloor\log_2(Q^fL)\rfloor\). The number of distinct
rational double cosets of the decorated flag stabilizers
represented in this range is at most
\[
   N_L=2^d\binom{B+f}{f}
        \prod_{j=1}^d
          \left(1+2(2\sqrt A\,L^{1/f})^j\right)^{2f}.
\]
This includes pairs with vanishing proper leading minors.
The bound is independent of the norms of the outer fractional
ideals.
\end{theorem}
\begin{proof}
Write \(Z\) for the cross matrix and \(D_j\) for its leading
minors, with \(D_0=1\). Put
\(\mathfrak l_j=(\overline{\mathfrak h_j}\mathfrak h_j)^{-1}\).
Integrality of the pairing on the two prefix lattices gives
\(D_j\in\mathfrak l_j\), also for \(j=d\) by the primitive
terminal vector. At each mark the outer spaces coincide on
reduction and are totally isotropic. The two terminal vectors
agree modulo that reduced outer space by the marking in
Section~\ref{sec:bothordinary553}. Their norms reduce to
\(-1\) since \(\varpi\equiv1\bmod Q^2\). The cross matrix
therefore reduces to \(-E_{dd}\), with \(E_{dd}\) the last
diagonal matrix unit. Consequently
\[
  D_j\in Q^j\mathfrak l_j\quad(1\le j<d),\qquad
              D_d+D_s\in Q^d\mathfrak l_d .
\]
Indeed the determinant of \(Z+E_{dd}\) is \(D_d+D_s\).
The outer ideals are units at both marks, and
\(\mathfrak l_s=\mathfrak l_d\) away from them as well.
For a fixed earlier tuple, differences between last entries
thus belong to \(Q^d\mathfrak l_d\).

At an embedding \(\tau\) let \(m_\tau\ge1\) be the factor
of \(\mathcal H_\gamma\) there. Set
\(b_\tau=\lfloor\log_2 m_\tau\rfloor\). Then
\(\sum_\tau b_\tau\le B\), giving at most
\(\binom{B+f}{f}\) tuples of nonnegative integers.
For a fixed tuple the \(j\)-th minor lies in discs of radii
\[
    r_{j,\tau}=\eta(\mathfrak h_j)^2 A^{j/2}2^{j(b_\tau+1)} .
\]
The difference \(z\) of two distinct possibilities, conditional
on earlier entries, belongs to \(Q^j\mathfrak l_j\). Hence
\[
       \prod_\tau|\tau z|=|N_{E/\QQ}z|^{1/2}
                   \ge Q^{jf}\eta(\mathfrak h_j)^{2f}.
\]
This follows by taking the norm of the integral ideal
\((z)/(Q^j\mathfrak l_j)\). After division by \(r_{j,\tau}\)
in each coordinate, the maximum distance of the two points
is at least \((2\sqrt A\,L^{1/f})^{-j}\), since
\(2^{\sum_\tau b_\tau}\le Q^fL\). Open polydiscs of half
that radius about the rescaled points are disjoint inside
the unit polydisc enlarged by the same radius. Volume
comparison in real dimension \(2f\) gives the \(j\)-th
factor of the product in the theorem. This counts zero
whenever it is allowed. Iterating the conditional counts
bounds the number of full minor tuples.

It remains to pass from minors to rational flag pairs.
The reduction in Lemma~\ref{lem:forms203} applies with
the two fixed terminal norms. More explicitly, the
concatenated Gram determinant is nonzero by total positivity
of \(\mathcal G\). A common cross matrix determines the
pair of decorated flags up to rational isometry by Witt
extension, because the two diagonal Gram blocks are fixed,
respectively \(-T_0E_{dd}\) and \(-T_1E_{dd}\).
Both blocks are preserved by upper unitriangular changes.
In these coordinates the simple characters use the
successor map sending the last outer vector to its own
terminal vector, as in Lemma~\ref{lem:suffixdet556}.
Thus the adjacent descent calculation in
Lemma~\ref{lem:forms203} has the same coefficients on
both frames. Relevance forces reversed consecutive blocks
with alternating signs on their antidiagonals. The nonzero
minor indices specify their endpoints, and successive
ratios of endpoint minors determine the leading scalar
of a block of length \(a\) up to an \(a\)-th root of unity.
There are at most \(a\) choices. Their product over the
blocks is at most \(2^d\), proving the result.
\end{proof}

\begin{proposition}[Integral cross matrices in each rational class]
\leavevmode\label{prop:smallcross590}
Let
\[
       Y=4^f\eta(\mathfrak h_d)^{4f} A^{df}L^{2d}.
\]
In each rational class of Theorem~\ref{thm:smallframes590},
the integral cross matrices that occur have at most
\(Y^{d(d-1)}\) equivalence classes under left lower and
right upper unitriangular changes over \(O_E\).
\end{proposition}
\begin{proof}
Each outer fractional ideal contains \(O_E\), since it is a
nonpositive power of the extended different. In particular
the chosen frame vectors, and hence the entries of \(Z\),
are integral. Both unitriangular actions fix \(E_{dd}\).
The substitution \(Z\mapsto(Z+E_{dd})/Q\) therefore
respects integral equivalence in both directions and has
integral values by the marked congruence above. Its last
determinant is nonzero, since
\(|\tau D_d|>\sqrt A\,|\tau D_s|\) and \(A\ge1\).

For clarity, the prefix products in the integral reduction
of Theorem~\ref{thm:crossclasses230} can be bounded using
one endpoint maximum here. If the last block is a singleton,
adding \(E_{dd}\) changes just its scalar. Otherwise an
earlier row has its unique pivot in the final column and
a lower row operation removes the added entry in the rational
normal form. After division by \(Q\), the absolute norms
of the pivot products at proper block endpoints \(j\) are
\(|N_{E/\QQ}(D_j/Q^j)|\), and at the final endpoint the
norm is \(|N_{E/\QQ}((D_d+D_s)/Q^d)|\). The initial product
is one. For a partial block the logarithm of the pivot-product
norm interpolates linearly between its endpoint values,
because all pivots of an individual nonsingleton block have
the same absolute norm.

Before taking dyadic boxes one has
\(|\tau D_j|\le\eta(\mathfrak h_j)^2 A^{j/2}m_\tau^j\)
and \(\prod_\tau m_\tau<Q^f L\). Every proper endpoint
norm is therefore at most \(\eta(\mathfrak h_j)^{4f}
A^{jf}L^{2j}\le Y\). For the last endpoint the triangle
inequality gives
\(|\tau(D_d+D_s)|\le2\eta(\mathfrak h_d)^2 A^{d/2}m_\tau^d\),
which accounts for the factor \(4^f\). Thus every intermediate
pivot-product norm is also at most \(Y\).

For a pivot product \(b_j\) through position \(j<d\), the
integral minor calculation in Theorem~\ref{thm:crossclasses230}
puts the lower entries below that pivot and the allowed upper
entries to its right in \(b_j^{-1}O_E\). Row and column
reduction needs only their classes modulo \(O_E\), at most
\(|N_{E/\QQ}b_j|\) choices per position. This proof uses
the shifted integral matrix and its pivot positions, and
so applies to the shift by \(E_{dd}\) used here. There are
at most \(2\sum_{j<d}(d-j)=d(d-1)\) such positions. The
bound follows.
\end{proof}

\begin{corollary}[A finite sum of signed groups below the cutoff]
\leavevmode\label{cor:smallgroups590}
For an absolutely convergent array in
Corollary~\ref{cor:smalltransverse589}, choose an integral
representative for each contributing cell. Group its signed
contribution by rational transport, rational flag class, and
integral cross-matrix class of the representative. If the
resulting arrays are \(S_\beta\), then
\[
           \left\|\sum_\beta S_\beta\right\|_{\rm HS}^2
               \le m_L N_L Y^{d(d-1)}
                            \sum_\beta\|S_\beta\|_{\rm HS}^2 .
\]
Here \(m_L\) is the entry bound of
Corollary~\ref{cor:smalltransverse589}. If
\(\log L=O(x^C)\) for fixed \(C\), the logarithm of the
displayed factor is \(o(\exp(\nu x^\varepsilon))\) on the
common intervals for every fixed \(\nu>0\).
\end{corollary}
\begin{proof}
At any output entry there are at most \(m_L\) distinct
transports. For each, the theorem bounds the number of
rational classes and the proposition bounds the number of
matrix classes. Cauchy--Schwarz at each entry gives the
display upon summation. Every choice of representatives
gives this inequality with the original signs and weights
of the contributions in each group.

The field degree, \(d\), and \(k\) are polynomial in \(x\).
Also \(\log Q=O(x^\varepsilon)\) and
\(\log T=O(\sqrt{x}\,x^\varepsilon)\) by
Theorem~\ref{thm:prime209} and Theorem~\ref{thm:orbits204}.
If \(R=|\operatorname{disc}(F)|^{1/f}\), then
\(\log R=O_\psi(\log x)\) as in
Proposition~\ref{prop:archdifferent555}, and
\(\eta(\mathfrak h_d)=R^{s(s+1)/2}\).
In the binomial term one may use
\(\binom{B+f}{f}\le(e(1+B/f))^f\). Substitution in all
three factors bounds their combined logarithm by a polynomial
in \(x\). This implies the stated asymptotic.
\end{proof}

The finite count bounds the number of groups that can meet
an output entry. An individual group can contain several
integral cells. A bound for their signed orbital sum in the
normalization of the integral period remains necessary.


\subsection{Integral orbits below the frame-height cutoff}
\label{sec:smallorbits591}

Fix the same pair of rows and cutoff \(L\) as in
Theorem~\ref{thm:smallframes590}. When the row lattices are
isometric, identify them with \(\Lambda\) as in
Section~\ref{sec:mixedordinary541}. Write
\(\Gamma=\operatorname{Aut}(\Lambda,h)\) and let \(N_b\) be
the integral unipotent subgroup of row \(b\). The period
uses a subset of \(N_0\backslash\Gamma/N_1\). We count the
individual double cosets in the relevant transverse
hyperplane part below the cutoff, including those with equal
marked index.

\begin{lemma}[Integral changes of an ordered frame]
\leavevmode\label{lem:framelifts591}
Every upper unitriangular change over \(O_E\) of the ordered
frame \((e_1,\ldots,e_s,z)\) is induced on this frame by an
element of \(N_b\), for either terminal norm.
\end{lemma}
\begin{proof}
Use the integral splitting \(A\oplus J\oplus A'\) of
Lemma~\ref{lem:projectivewitt554}, with partners \(f_i\).
For \(i<j\le s\), the change \(e_j\mapsto e_j+c e_i\)
extends by \(f_i\mapsto f_i-\bar c f_j\), fixing the other
lines and \(J\). This preserves the hyperbolic pairing.
It is integral for \(c\in O_E\) since
\(\mathfrak a_j\subseteq\mathfrak a_i\) and
\(\bar{\mathfrak a}_i^{-1}\subseteq\bar{\mathfrak a}_j^{-1}\).

For \(e_i\) preceding \(z\), choose \(a\in J\) such that
\(h(a,z)=c\). This is possible since \(O_Ez\) is saturated
in the projective lattice \(J\), so it is a direct summand.
Extend its coordinate functional to \(J\) and use
self-duality. Take \(t\in O_E\) with \(t+\bar t=1\),
whose existence was proved in Lemma~\ref{lem:projectivewitt554}.
Fix \(A\) and \(f_j\) for \(j\ne i\), and put
\[
 w\longmapsto w+e_i h(a,w)\quad(w\in J),\qquad
 f_i\longmapsto f_i-a-t h(a,a)e_i .
\]
The new \(f_i\) pairs to zero with the new middle vectors.
Its norm is \(h(a,a)(1-t-\bar t)=0\). The other block
pairings are preserved. The new terminal vector is
\(z+c e_i\). The inclusion \(O_E\subseteq\mathfrak a_i\)
and its dual \(\bar{\mathfrak a}_i^{-1}\subseteq O_E\)
give integrality on all three summands. The inverse is
integral by the finite geometric series for the nilpotent
difference from the identity. Both types of maps act as the
identity on the successive outer lines and on the full
middle quotient. Hence they belong to \(N_b\). The elementary
changes generate the upper unitriangular group by eliminating
entries one superdiagonal at a time.
\end{proof}

Let \(n=2s+1+k\), \(h_0=k-1\ge1\), \(D=2fh_0\), and
\(\Delta_E=|\operatorname{disc}(E)|\). Define
\[
 G_0=\eta(\mathfrak h_d)^{4f}(T_0T_1)^{df}
                         Q^{2df}L^{2d},\qquad
 C=2D^{D/2}\Delta_E^{h_0/2}G_0 .
\]
In the next lemma the frames are fully ordered and each
individual vector is distinguished.

\begin{lemma}[Integral embeddings with a prescribed cross matrix]
\leavevmode\label{lem:smallglue591}
Fix an integral cross matrix \(Z\) occurring in the stated
range. The number of \(\Gamma\)-orbits of pairs of ordered
frames in \(\Lambda\) with that cross matrix and the given
two diagonal Gram blocks is at most
\[
 P_L=G_0^{\,8fd+8fn}
       (1+2C^2)^{2fh_0^2} C^{4fh_0D}.
\]
\end{lemma}
\begin{proof}
Let \(U\) be the free lattice on the \(2d\) ordered vectors
with their specified Gram matrix. Its determinant is
\((-1)^d\mathcal G\), where
\(\mathcal G=D_d\bar D_d-T_0T_1D_s\bar D_s\) is totally
positive by Proposition~\ref{prop:transverseheight582}.
For \(G=N_{F/\QQ}\mathcal G\) that proposition gives
\[
  1\le G\le |N_{E/\QQ}D_d|
       \le\eta(\mathfrak h_d)^{4f}(T_0T_1)^{df}
                               \mathcal H_\gamma^{2d}\le G_0 .
\]
The first inequality uses integrality of the ordered vectors.
The index of \(U\) in its Hermitian dual \(U^\sharp\) is \(G^2\),
by taking the absolute \(E/\QQ\) norm of the Gram determinant.
An embedding gives the saturation \(B=\Lambda\cap EU\).
Pulling it back to the marked rational span gives
\(U\subseteq B\subseteq U^\sharp\). There are at most
\((G_0^2)^{4fd}\) choices. Indeed \(B/U\) is generated
over \(\mathbb Z\) by at most \(4fd\) elements since \(B\) has
that rank, and a padded ordered list of generators specifies
its subgroup of \(U^\sharp/U\).

Write \(R=\Lambda\cap(EU)^\perp\). The nondegenerate
hyperplane case gives a totally positive lattice of rank
\(h_0\) here. Restricting the integral pairing maps
\(\Lambda\) onto \(B^\sharp\) with kernel \(R\), upon using
orthogonal projection to \(EB\). Surjectivity follows by
extending a functional from the saturated direct summand
\(B\) and applying self-duality of \(\Lambda\). The same
reasoning with \(R\) gives
\[
 [R^\sharp:R]=[\Lambda:B\oplus R]
          =[B^\sharp:B]=G^2/[B:U]^2\le G_0^2 .
\]
The last equality follows by dualizing \(U\subseteq B\).

We bound the number of isometry types of \(R\).
Its integral trace pairing \(\operatorname{Tr}_{E/\QQ}h\)
is positive definite, with trace dual
\(\mathfrak D_{E/\QQ}^{-1}R^\sharp\). Indeed one tests trace
integrality against all \(O_E\)-multiples of each vector and
uses the defining trace dual of \(O_E\). Its Euclidean
covolume is at most \(\Delta_E^{h_0/2}G_0\), since the
index of \(R^\sharp\) in its multiple by the inverse
different is \(\Delta_E^{h_0}\). All nonzero vectors have
trace length at least one. Minkowski's second theorem for the
unit ball in dimension \(D\), which contains the cube of
side \(2/\sqrt D\), bounds the product of the successive
minima by \(D^{D/2}\Delta_E^{h_0/2}G_0\).
There are \(D\) independent vectors of length at most \(C\).
They span the underlying rational space, hence contain an
\(E\)-basis \(x_1,\ldots,x_{h_0}\).

In this basis every Gram entry is an algebraic integer of
absolute value at most \(C^2\) at every embedding, by
Cauchy--Schwarz and the trace length bound. Two distinct
algebraic integers differ by at least one in maximum
archimedean modulus, by the norm of their difference.
Packing open polydiscs of radius \(1/2\) therefore gives at
most \((1+2C^2)^{2fh_0^2}\) possible Gram matrices. For
each, put \(M=\sum_i O_Ex_i\). Hadamard's inequality at
each embedding gives
\([M^\sharp:M]\le C^{4fh_0}\). The intermediate lattice
\(M\subseteq R\subseteq M^\sharp\) is specified by a
subgroup with at most \(D\) abelian generators. This gives
the remaining factor \(C^{4fh_0D}\) in the definite count.

For each \(B\) and each isometry type of \(R\), an ambient
integral lattice containing \(B\oplus R\) lies inside
\(B^\sharp\oplus R^\sharp\). The quotient of these endpoints
has order at most \(G_0^4\). The ambient lattice has
\(\mathbb Z\)-rank \(2fn\), so at most
\((G_0^4)^{2fn}\) subgroups can occur as its image.
Each choice specifies the entire lattice by inverse image,
with the ordered frame vectors already marked in \(U\).
Any two identifications of such a lattice with \(\Lambda\)
differ by an element of \(\Gamma\). This bounds the number
of orbits by the product claimed.
\end{proof}

\begin{theorem}[Number of individual low-height terms on an original row]
\leavevmode\label{thm:smallorbits591}
For either terminal frequency at either end, put
\[
          K_L=N_L\,Y^{d(d-1)}\,P_L\,3^{\,2(2fk)^2},
\]
where \(N_L,Y\) are given in Theorem~\ref{thm:smallframes590}
and Proposition~\ref{prop:smallcross590}. The number of
relevant supported transverse hyperplane double cosets below
the cutoff in the fixed pair of rows is at most \(K_L\).
In particular, if \(B_\gamma\) is the signed projected
coefficient of an individual such coset, then
\[
  \left\|\sum_\gamma B_\gamma\right\|_{\rm HS}^2
                     \le K_L\sum_\gamma\|B_\gamma\|_{\rm HS}^2 .
\]
For \(\log L=O(x^c)\) with fixed \(c\), one has
\(\log K_L=o(\exp(\nu x^\varepsilon))\) on the original
intervals for every fixed \(\nu>0\).
\end{theorem}
\begin{proof}
The preceding two results of Section~\ref{sec:smallframes590}
give at most \(N_LY^{d(d-1)}\) cross matrices up to the two
integral unitriangular actions. Each double coset admits a
representative with one of those exact matrices. In fact, on
replacing \(\gamma\) by \(n_0\gamma n_1\), the first frame
in the cross pairing changes by \(n_0^{-1}\) and the second
by \(n_1\). Lemma~\ref{lem:framelifts591} realizes both
specified frame changes. These are elements of the very
groups used to define the double cosets of
Section~\ref{sec:mixedordinary541}. Restricting to the
admissible classes at the original marking can only reduce
their number.

For a fixed exact matrix, classify the ordered concatenated
frames in the target lattice up to \(\Gamma\) by
Lemma~\ref{lem:smallglue591}. If \(\gamma,\delta\) give the
same orbit, there is \(u\in\Gamma\) fixing the first ordered
frame pointwise with \(\delta^{-1}u\gamma\) fixing the second
row's original ordered frame pointwise. For either row \(b\),
the stabilizer of its ordered frame modulo earlier vectors
maps to the isometry group of the positive rank-\(k\)
lattice \(L_b\), with kernel \(N_b\). Indeed it fixes the
terminal class, and the middle rational space is the
orthogonal sum of its line and \(EL_b\). Trivial action on
the positive lattice forces trivial action on the entire
middle quotient, which defines the unipotent radical inside
this stabilizer. The finite group of isometries of \(L_b\)
injects into matrices modulo three on its underlying lattice
of rank \(2fk\), as proved in
Theorem~\ref{thm:uniformheight552}. Its order is at most
\(3^{(2fk)^2}\). Thus one orbit of ordered embeddings
accounts for at most the square of this number of
\(N_0,N_1\) double cosets. Multiplying proves the count.
Cauchy--Schwarz for a finite sum in the Hilbert--Schmidt
space proves the inequality with all signs and finite
projectors in place.

The field degree, \(n,k,d\), and the logarithms of the
parameters in \(G_0\) grow at most polynomially in \(x\),
as in Corollary~\ref{cor:smallgroups590}. The extension
\(E/F\) is unramified at finite places on these fields,
so \(\Delta_E=|\operatorname{disc}(F)|^2\).
It follows directly from the formulas for \(C,P_L,K_L\)
that \(\log K_L\) also grows at most polynomially in \(x\).
This proves the last assertion.
\end{proof}

The estimate applies to the actual double cosets, even when
several of them contribute to a single entry of the finite
projection. Comparison with the integral theta map in
Section~\ref{sec:missing} still requires estimates on the
sizes of their orbital integrals with the prescribed sections.


\subsection{Height truncation of the projected ordinary period}
\label{sec:truncatedordinary592}

Keep the fields, rows, and finite projection of
Section~\ref{sec:mixedordinary541}. We integrate on the
original real quotient, using \(N_j\) for the integral
unipotent subgroup of row \(j\). Its real section \(W_j\)
has nonzero finite \(L^2(N_j\backslash G_\infty)\) norm \(w_j\).
Let \(V_j'\) be either the height \(V_j\) of
Theorem~\ref{thm:idealordinary554} or the translated height
\(\widetilde V_j\) of Theorem~\ref{thm:mixedhalf556}.
For \(q\ge0\) and \(R\ge1\) put
\[
 \mathfrak M_j(q)=w_j^{-2}\int (V_j')^q\|W_j\|^2,\qquad
 p_j(R)=w_j^{-2}\int_{V_j'\le R}\|W_j\|^2 .
\]
These integrals have the same quotient measure. Set
\[
 D=\sum_{i=1}^{s+1}i(n-i+1),\qquad
 C_*=3^{(2fk)^2}\prod_{i=1}^{s+1}
            (1+2\cdot4^iT^i)^{2f(n-i+1)} .
\]

\begin{theorem}[Bounds in terms of the actual squared section norms]
\leavevmode\label{thm:twoenergies592}
Let \(\mathcal A\) be any set of supported integral double
cosets in one fixed pair of rows. Write \(B_\gamma\) for the
signed ordinary orbital integral, including its original finite
coefficient and normalized column in the finite projection.
If \(V_0'(y)V_1'(\gamma^{-1}y)\ge H>0\) for all these
cosets and all \(y\), then for \(\sigma>4D\) and \(u\ge0\),
\[
 \frac{\sum_{\gamma\in\mathcal A}\|B_\gamma\|_{\rm HS}}
      {|c_{0,1}|\,w_0w_1}
 \le H^{-u}C_*(1-2^{-(\sigma-4D)})^{-f}
       \sqrt{\mathfrak M_0(\sigma+2u)\mathfrak M_1(\sigma+2u)}
\]
whenever the two moments are finite.
With the integration on each coset restricted by
\(V_0'(y)\le R_0\) and \(V_1'(\gamma^{-1}y)\le R_1\), write
\(B_\gamma(R_0,R_1)\) for the resulting coefficient. For
every such set of cosets one has
\[
 \frac{\sum_{\gamma\in\mathcal A}
                    \|B_\gamma(R_0,R_1)\|_{\rm HS}}
      {|c_{0,1}|\,w_0w_1}
 \le\sqrt{M(R_0)M(R_1)p_0(R_0)p_1(R_1)},\qquad
 M(R)=C_* R^{4D}\binom{\lfloor\log_2 R\rfloor+f}{f}.
\]
Here \(c_{0,1}\) is the same nonzero volume factor as in the
unfolded ordinary period. If \(R_0R_1<H\), all the truncated
integrals in this set vanish.
\end{theorem}
\begin{proof}
For either height write its local factors as \(\nu_{j,\tau}\).
For fixed \(y\), the number of cosets in \(N_j\backslash\Gamma\)
with \(2^{a_\tau}\le\nu_{j,\tau}(\delta y)<2^{a_\tau+1}\)
at every \(\tau\) is at most \(C_*2^{4D\sum a_\tau}\).
For the translated height this is the exterior lattice count
in Theorem~\ref{thm:mixedhalf556}. For the original height
the count in Theorem~\ref{thm:idealordinary554} is smaller,
since \(T^{(i+1)/2}\le T^i\). The \(a_\tau\) are nonnegative
as the positive norms are compatible with the Hermitian form.
Consequently
\[
 \sup_y\sum_{\delta\in N_j\backslash\Gamma}
             V_j'(\delta y)^{-\sigma}
 \le C_*(1-2^{-(\sigma-4D)})^{-f}=:P_\sigma .
\]
For \(V_j'(\delta y)\le R\), the sum of the \(a_\tau\) is at
most \(\lfloor\log_2 R\rfloor\). Counting their tuples and
bounding each box by \(C_*R^{4D}\) gives \(M(R)\).

On the disjoint union of the coset integration domains,
temporarily write \(v_0=V_0'(y)\), \(v_1=V_1'(\gamma^{-1}y)\)
and \(F_b\) for the norm of the corresponding section at
that argument. The norm of its rank-one integrand is
\(F_0F_1\). Its phase has modulus one and its normalized
projector column has norm at most one, as in
Theorem~\ref{thm:flagisometry544}. The separation
\(v_0v_1\ge H\) bounds \(F_0F_1\) by \(H^{-u}\) times
the product
\[
 (F_0v_0^{u+\sigma/2}v_1^{-\sigma/2})
 (F_1v_1^{u+\sigma/2}v_0^{-\sigma/2}).
\]
Apply Cauchy--Schwarz over the whole disjoint union. In the
integral of the first square, enlarge the nonnegative sum to
all double cosets. For each fixed \(y\) in
\(N_0\backslash G_\infty\) the weight from the second row
becomes \(\sum_{\delta\in N_1\backslash\Gamma}
V_1'(\delta y)^{-\sigma}\).
Indeed a fundamental domain of
\(N_0\cap\gamma N_1\gamma^{-1}\) is tiled by translates
of one for \(N_0\), and inversion of the double cosets lists
each left \(N_1\)-coset once. This is the counting-measure
unfolding in Section~\ref{sec:mixedordinary541}. Tonelli's
theorem applies to these nonnegative functions.
The first squared integral is thus at most
\(P_\sigma w_0^2\mathfrak M_0(\sigma+2u)\). Reversing the
two rows gives the corresponding bound for the second square.
Multiplication proves the first estimate.

For the truncated comparison apply the same Cauchy--Schwarz
inequality to \(F_0,F_1\) on the part of the union satisfying
both cuts. Unfolding the first square counts second-row
cosets of height at most \(R_1\). It is at most
\(M(R_1)w_0^2p_0(R_0)\). The opposite square is at most
\(M(R_0)w_1^2p_1(R_1)\). If the product of the cutoffs is
less than \(H\), the integration set is empty.
\end{proof}

For example, in the part of
Theorem~\ref{thm:transversesum582} with
\(\mathcal H_\gamma\ge Q^fL\), take \(V_j'=\widetilde V_j\),
\(\sigma=4D+1\), and \(u=a\) as in
Theorem~\ref{thm:idealordinary554}. The required moment has
order \(4D+1+2a\le L_0\) for sufficiently large \(x\).
This follows from \(2a=L_0-4fn^2\) and \(D\le n(s+1)(s+2)/2\),
with \(n/(s+2)^2\) tending to infinity on these intervals.
It is finite by the weighted supremum and integral bounds
of Theorem~\ref{thm:mixedhalf556}, since
\(\int(V_j')^{L_0}\|W_j\|^2
 \le\widetilde U_j\int\|W_j\|<\infty\).
The resulting dimensionless factor multiplying the geometric
mean of these moments is at most \(Q^{-fa/2}L^{-a}\), since
\(\log(2^fC_*)=o(fa\log Q)\) there. Moreover the
normalized moments for \(\widetilde V_j\) are exactly the
moments for \(V_j\) of the transported unit-frequency vector.
The substitution \(y\mapsto a_jy\) in
Theorem~\ref{thm:mixedhalf556} cancels its Haar factor
against the squared normalization of the section. The
volume of \(N_j\backslash B_{j,\infty}\) cancels against
the same volume in \(w_j^2\).

\begin{theorem}[The exact diagonal with both real heights restricted]
\leavevmode\label{thm:exactcore592}
Take one retained row \(j\), repeated at both ends with its
original marking, word and terminal frequency. Use the
original height \(V_j'=V_j\). Let \(\mathcal Z_j(R)\) be its
projected ordinary diagonal block with both heights restricted
to at most \(R\), and let \(\mathcal B_j\) be its unrestricted
flag-preserving block. For \(1\le R<(Q/2)^{f/2}\),
\[
       \mathcal Z_j(R)=p_j(R)\mathcal B_j,\qquad
       \mathcal Z_j(R)^*\mathcal Z_j(R)
                  =p_j(R)^2m_j\lambda_j^2\,1_{V_{\rho_j}} .
\]
Here \(V_{\rho_j}\) is the positive column space,
\(m_j\) counts the admissible flag-preserving double cosets,
and \(\lambda_j=|c_{j,j}|w_j^2/\dim V_{\rho_j}\).
At prime terminal norm \(T\), for the prescribed unit
lowest-type vectors of Theorem~\ref{thm:localmass156},
\[
           p_j(R)\le \exp(-c b n^2),\qquad b=\exp(x^\varepsilon),
\]
uniformly on the same intervals for sufficiently large \(x\)
with the constant of that local mass estimate.
\end{theorem}
\begin{proof}
For a supported isometry that first fails to transport an
outer vector modulo its predecessors,
Theorem~\ref{thm:idealordinary554} gives
\(V_j(y)V_j(\gamma^{-1}y)\ge(Q/2)^f\) throughout its
integral. Such a coset cannot contribute. If all the outer
vectors transport but the terminal class changes, the two
real unipotent radicals agree after transport and the two
characters on them differ by
Proposition~\ref{prop:terminalvanishing553}. Both height
cutoffs are invariant under the common radical, since the
heights are left invariant under their decorated stabilizers.
Integration of the nontrivial character on its compact
unipotent lattice quotient still gives zero.

For an admissible flag-preserving \(\gamma\) both heights
in its orbital integral coincide. The cutoff is left
invariant under the real decorated group. Compact averaging
on the positive column therefore gives a scalar matrix by
irreducibility, just as for the unrestricted coefficient
in Theorem~\ref{thm:flagisometry544}. Its trace before
the finite phase and unitary middle action is
\(\int_{V_j\le R}\|W_j\|^2\). Hence the original coefficient,
at its same marked index, is multiplied by precisely
\(p_j(R)\). Summing these terms gives the first identity.
Their normalized finite columns have disjoint supports and
norm one. The unrestricted Gram operator is thus
\(m_j\lambda_j^2\,1_{V_{\rho_j}}\), giving the second identity.

For the prime-frequency estimate let
\(R_F=|\operatorname{disc}(F)|^{1/f}\).
The proof of Theorem~\ref{thm:primemoment555} gives
normalized mass at least \(1-e^{-c b n^2}\) to
\(V_j\ge(\sqrt T/(R_F b n^4))^f\). This lower threshold
exceeds \((Q/2)^{f/2}\), as one sees by comparing the
\(f\)-th roots using
\(\log T=(17+o(1))m x^\varepsilon\), \(m\to\infty\),
\(\log Q=(16+o(1))x^\varepsilon\), and
\(\log(R_F n^4)=O_\psi(\log x)\).
This proves the asserted upper bound for the mass in the cut.
\end{proof}

Thus the original finite projection retains an exactly
computable part of the real diagonal. At prime terminal
frequency this truncation captures exponentially little
of its squared section norm. The comparison in
Section~\ref{sec:missing} requires the signed contributions
from the complementary integration region and the
normalization of the integral theta map.


\subsection{Two determinant bounds for the signed ordinary period}
\label{sec:markedpencil602}

Fix two retained rows on the same \(E=E_x\) and interval
\([V,V^{3/2}]\) as in Section~\ref{sec:missing}. Work with the
supported cells with compatible characters and with
\(0<\dim(X_0\cap X_1)<s\). Order the rows so that
\(T_0\ge T_1\), and put \(u=T_0/T_1\), \(F=E^+\),
\(f=[F:\QQ]\) and \(Q=g_1g_2\). Both rational primes split
completely in \(E\). The positive spaces \(P_0,P_1\) have
dimension \(k\). A cell \(\gamma\) has a rational isometry
\(j_\gamma:P_0\to P_1\), a prescribed marked index \(i_\gamma\),
and common sublattices \(C_{b,\gamma}\subset L_b\) of index
\(m_{b,\gamma}\), as in Section~\ref{sec:positivecomparison576}.
In particular \(m_{1,\gamma}=u^f m_{0,\gamma}\). All marked
indices below are those of the original period.

\begin{theorem}[A trace bound for overlap of two projected columns]
\leavevmode\label{thm:traceoverlap602}
Let \(j_\gamma\ne j_\delta\) and
\(\rho=\operatorname{rank}_E(j_\gamma-j_\delta)\).
If the two columns of the integral projection \(H\) meet in
one output coordinate, then
\[
 Q\ \le\ 2\rho\min\bigl\{
    (N(\mathfrak a_{1,\gamma}\mathfrak a_{0,\delta}))^{1/(2f)},
    (N(\mathfrak a_{1,\delta}\mathfrak a_{0,\gamma}))^{1/(2f)}
                         \bigr\}
 \ \le\ 2\rho\sqrt u\,(m_{0,\gamma}m_{0,\delta})^{1/(2f)} .
\]
Here \(\mathfrak a_{b,\gamma}\) annihilates \(L_b/C_{b,\gamma}\)
and \(N\) is the absolute ideal norm in \(E\).
\end{theorem}
\begin{proof}
Set \(M=j_\gamma^{-1}j_\delta\) and \(z=k-\operatorname{tr} M\).
At every archimedean embedding \(M\) is unitary for a positive
Hermitian form. It has \(k-\rho\) eigenvalues equal to one,
counted with multiplicity, as it is diagonalizable there.
The real part of \(z\) is a sum of \(1-\operatorname{Re}\lambda\)
over its other \(\rho\) eigenvalues. Every summand is strictly
positive and \(|z|\le2\rho\). In particular \(z\ne0\).

Localize at any finite prime and take generators \(a,b\) of
\(\mathfrak a_{1,\gamma},\mathfrak a_{0,\delta}\). The maps
\(a j_\gamma^{-1}\) and \(b j_\delta\) between the row
lattices are integral, so \(abz\) is integral. At a mark these
maps reduce to the corresponding scaled marked maps by
Lemma~\ref{lem:scaledmark576}, also when one scalar has zero
residue. A shared coordinate means that the relative marked
matrix is conjugate to an upper unitriangular matrix. This
follows directly from the formula for \(H\) in
Lemma~\ref{lem:whittakermatrix526}. The relative matrix in
the conjugate split factor is also unipotent. Its trace at
either factor is \(k\) in the residue field. Thus \(abz\)
reduces to zero. We have proved the ideal inclusion
\[
       (z)\mathfrak a_{1,\gamma}\mathfrak a_{0,\delta}
                                      \subseteq Q\mathcal O_E .
\]
Taking norms and using the archimedean bound gives the first
estimate in this orientation. Interchanging the cells gives
the other one. The annihilator norm is at most the order of
the quotient, by elementary divisors. Substituting both row
indices proves the last inequality.
\end{proof}

The ranks at individual marked places also carry information.
Write \(r_v=\operatorname{rank}(i_{\gamma,v}-i_{\delta,v})\)
for \(v\mid Q\) in \(F\), and \(q_v\) for its residue size.
For a full sublattice \(C\subset L\), write \([L:C]_{\rm id}\)
for the integral ideal whose local exponent is the length of
\(L/C\). Its absolute norm is the order \([L:C]\).
For this pair of cells define
\[
 C_b=C_{b,\gamma}\cap C_{b,\delta},\qquad
 J_b=[L_b:C_b]_{\rm id},\qquad
 D_{\gamma,\delta}=(N(J_0J_1))^{1/(2f)}.
\]

\begin{theorem}[Rational rank and rank at the two marks]
\leavevmode\label{thm:rankmarks602}
For distinct transports let \(\rho\) be their rational
difference rank. Set \(b_v=\max(0,\rho-r_v)\).
At each \(v\) where the relative marked matrix is unipotent,
one can replace \(b_v\) by \(\max(1,\rho-r_v)\). Then
\[
     \prod_{v\mid Q}q_v^{2b_v}
       \ \le\ 4^{f\rho}D_{\gamma,\delta}^{\,f},\qquad
     D_{\gamma,\delta}\le u(m_{0,\gamma}m_{0,\delta})^{1/f}.
\]
In particular, two equal full marks for distinct transports
require \((Q/2)^{2\rho}\le D_{\gamma,\delta}\).
\end{theorem}
\begin{proof}
At a finite prime fix bases of \(L_0,L_1,C_0,C_1\) over the
local integer ring. Let \(U_b\) be the matrix of
\(C_b\hookrightarrow L_b\) and \(\lambda_b=\det U_b\), a
generator of \(J_b\). Consider the two integral polynomials
\[
 \begin{split}
   F_0(Y)&=\det\bigl((j_\delta-(1+Y)j_\gamma)U_0\bigr),\\
   F_1(Y)&=\det\bigl((j_\delta^{-1}-(1+Y)j_\gamma^{-1})U_1\bigr).
 \end{split}
\]
Each map used on \(C_b\) takes values in the opposite lattice.
At a mark its reduction on \(C_b\) is its prescribed marked
map composed with \(\bar U_b\). To see this for every vector of
\(C_b\), use its representative from the common positive lattice
of the cell. Reducing that representative at the two ends gives
the same vector in the common reduced middle quotient. Reading
it in the two row bases gives the indicated formula, as in
the proof of Lemma~\ref{lem:scaledmark576}. This applies
separately to each of the two cells and to both directions.

Let \(c_b\) be the coefficient of \(Y^{k-\rho}\) in \(F_b\).
Over the fraction field, factor \(j_\gamma\) from the first
determinant and \(j_\gamma^{-1}\) from the second. The relative
matrices that remain are \(M\) and a conjugate of \(M^{-1}\).
The two factored determinants cancel in their product. Hence
\[
 c_0c_1=\lambda_0\lambda_1 h,\qquad
 h=\prod_{\lambda\ne1}(\lambda-1)(\lambda^{-1}-1),
\]
where the product runs over the eigenvalues of \(M\) different
from one, counted with multiplicity. The two factors from the
eigenvalue one have signs \((-1)^{k-\rho}\) which cancel.
The coefficient description shows that \(h\in E\).
At every complex embedding the eigenvalues lie on the unit
circle and \(0<h\le4^\rho\). The ideal \((h)J_0J_1\) is
integral at all primes by the integrality of \(c_0,c_1\).

At a marked prime the constant matrix in the first determinant
has reduction \((i_\delta-i_\gamma)\bar U_0\), of rank at most
\(r_v\). In the second determinant the inverse difference has
the same rank bound. The conjugate split factor has that bound
as well, since taking inverse transposes preserves difference
rank. In each term of the coefficient \(c_b\), exactly \(\rho\)
columns come from the constant matrix. Every minor of order
\(\rho\) of this matrix is divisible by the uniformizer to
power \(\max(0,\rho-r_v)\), by diagonal reduction over the
local integer ring. Expanding the determinant on these columns
gives that divisibility for \(c_b\).

If the relative matrix is unipotent, the whole reduced
polynomial \(F_b\) is a scalar multiple of \(Y^k\). Indeed the
two reduced maps factor on the right through \(\bar U_b\);
after taking its determinant, the remaining characteristic
polynomial is that of a unipotent matrix shifted by \(1+Y\).
This argument includes \(\det\bar U_b=0\). As \(\rho>0\),
\(c_b\) then has positive valuation. At each prime above \(v\)
the valuation of \((h)J_0J_1\) is therefore at least \(2b_v\).
There are two primes of \(E\) over each such \(v\). Taking the
square root of the absolute ideal norm and applying the
archimedean bound proves the product estimate. The injection
\(L_b/C_b\hookrightarrow (L_b/C_{b,\gamma})\oplus
(L_b/C_{b,\delta})\) proves the upper bound for
\(D_{\gamma,\delta}\). For equal marks all \(r_v=0\) and
\(\prod_{v\mid Q}q_v=Q^f\), giving the final assertion.
\end{proof}

\begin{corollary}[An integral decomposition and its archimedean error]
\leavevmode\label{cor:properperiod602}
Assume \(s\ge2\). Put \(B=Q/(2k\sqrt u)\). For the fixed pair
of rows retain the proper-intersection cells with
\(m_{0,\gamma}^{1/f}<B\). In each row family of \(H\), the
signed projected sum on the retained cells with a given
rational transport is recovered from the retained total by
selecting coordinates. This holds over the integers of the
coefficient field at every coefficient prime.

Let \(S\) denote the full contribution of the proper nonzero
intersections under consideration, in the normalization
\(d^{-1/2}H\) of Lemma~\ref{lem:whittakermatrix526}, with
\(d=[G:U]\). Let \(S_j^{<}\) be its retained part for transport
\(j\), including the signs and the ordinary real orbital
integrals. Put
\[
 H_*=\max\left\{1,\;Q^f\min_{1\le t\le s-1}
   \left(\frac{\max(1,B)u^{s-t+1}}{(T_0T_1)^t}\right)^{f/(2t)}
                    \right\}.
\]
Use the translated heights of Theorem~\ref{thm:mixedhalf556}.
With the squared section moments and constants \(D,C_*\) of
Theorem~\ref{thm:twoenergies592}, write
\(P_\sigma=C_*(1-2^{-(\sigma-4D)})^{-f}\).
For \(\sigma>4D,\ b\ge0\) with finite indicated moments,
\[
 \frac{\left|\|S\|_{\rm HS}
              -(\sum_j\|S_j^{<}\|_{\rm HS}^2)^{1/2}\right|}
      {|c_{0,1}|w_0w_1}
 \ \le\ P_\sigma H_*^{-b}
       \sqrt{\mathfrak M_0(\sigma+2b)\mathfrak M_1(\sigma+2b)} .
\]
At either equal terminal norm, the integral selection includes
all \(m_{0,\gamma}^{1/f}\le\exp((16-\delta)x^\varepsilon)\)
for any fixed \(0<\delta<16\) and sufficiently large \(x\)
on the prescribed intervals. When \(T_0=T_1=1\) and \(B>1\),
\(H_*=Q^f B^{f/(2(s-1))}\).
\end{corollary}
\begin{proof}
If two retained cells have different transports, an overlap
would contradict Theorem~\ref{thm:traceoverlap602}, since
\(\rho\le k\) and both index inequalities are strict. Thus
in the retained projection each coordinate belongs to at most
one transport. Selecting all its coordinates gives precisely
the signed array for that transport. The selection matrix
has entries zero and one. Orthogonality gives
\(\|\sum_j S_j^{<}\|_{\rm HS}^2=\sum_j\|S_j^{<}\|_{\rm HS}^2\);
the common real normalization does not enter the integral
selection. The prescribed estimate for \(Q\) is
\(\log Q=(16+o(1))x^\varepsilon\), and \(\log(2k)=
o(x^\varepsilon)\), by Theorem~\ref{thm:orbits204}.

For each omitted cell set \(e=\dim(X_0\cap X_1)\) and
\(t=s-e\). In the height inequality in the proof of
Corollary~\ref{cor:indexperiod576}, the integer factors
\([R:R_1]\) and \(D_\beta^{2e+1}\) are at least one, and
\(m_{0,\gamma}\ge\max(1,B)^f\). Taking the \(2t\)-th root
shows that the product of the two translated heights is at
least the \(t\)-term defining \(H_*\). Each height is also at
least one. Apply Theorem~\ref{thm:twoenergies592} to this whole
set of omitted cells. It bounds the sum of their operator
norms by the right side times \(|c_{0,1}|w_0w_1\); all
finite projection columns have already been normalized there.
The triangle inequality and then the reverse triangle
inequality give the stated comparison. The formula with both
terminal norms one follows by minimizing over \(t\).
\end{proof}

The integral selection isolates a rational isometry. Several
orbital integrals can still occur with that isometry and the
same marked index. Estimating their signed sum and the size
of the section moments at the normalization of
Section~\ref{sec:missing} remains necessary.

\section{Integral lifts and the signed transverse period}\label{nav:chapter-signed-transverse-integrals}

\subsection{Rational isometries at a fixed transverse marking}
\label{sec:hyperplanerank603}

Consider the transverse cells of a pair of original marked rows
over \(E_x\) for which the common positive lattice \(R_\gamma\)
is a hyperplane. Use \(F=E_x^+\), \(f=[F:\QQ]\),
\(Q=g_1g_2\) and the two completed column lattices
\(\widehat L_b\) of Lemma~\ref{lem:completedcolumn535}.
At the marks these are the original lattices. The rational
isometry \(j_\gamma\) on the positive columns extends the
identification on \(R_\gamma\) and has the prescribed determinant.
Let \(\mathfrak t_\gamma\) be the discriminant ideal of this
hyperplane lattice and \(h_\gamma=(N\mathfrak t_\gamma)^{1/(2f)}\).
Both terminal frequencies are allowed independently.

\begin{lemma}[Matched sublattices of prescribed determinant]
\leavevmode\label{lem:hyperplanelattice603}
There are full sublattices \(C_{b,\gamma}\subset\widehat L_b\)
whose local index ideals at both ends are \(\mathfrak t_\gamma\).
The map \(j_\gamma\) carries \(C_{0,\gamma}\) onto \(C_{1,\gamma}\).
At every marked prime the reduction of this map into the
ambient target lattice is the prescribed \(i_\gamma\) composed
with the reduction of \(C_{0,\gamma}\hookrightarrow\widehat L_0\).
The analogous assertion holds in the other direction.
\end{lemma}
\begin{proof}
Work over the integer ring at a finite prime of \(E_x\). Take
the column bases and dual conjugate bases in the construction
of Lemma~\ref{lem:cofactormark581}. Let \(A_b\) have as columns
the common hyperplane basis at end \(b\), and let \(B_b\) have
as rows the pairings with a common conjugate hyperplane basis.
These bases can be chosen independently in the two variables
of the pairing, including at a self-conjugate prime.
The square matrix \(G=B_bA_b\) is the same at both ends and
\(\Delta=\det G\) generates the discriminant ideal locally.
If \(B\) has \(k-1\) rows, define the column \(n(B)\) by
\(v n(B)=\det\left[\begin{smallmatrix}B\\v\end{smallmatrix}\right]\)
for every row \(v\). Use the \(k\) columns
\[
                         [\,A_b,\ n(B_b)\,]
\]
for the sublattice at end \(b\). Expansion in cofactors gives
determinant \(\Delta\). The last column lies on the normal
line since \(B_b n(B_b)=0\). The map \(j_\gamma\) matches
the first \(k-1\) columns and takes the last column to
\(d_\gamma n(B_1)\), where \(d_\gamma\) is its determinant
in the ambient bases. This follows by taking determinants of
the two displayed bases. By Lemma~\ref{lem:completedcolumn535}
\(d_\gamma\) is a unit at every finite prime, so the indicated
sublattices match in both directions.

At a mark, write \(i=i_\gamma\). The common representatives
give \(\bar A_1=i\bar A_0\) and \(\bar B_1=\bar B_0 i^{-1}\);
bars here denote reduction. Expansion along the last row gives
\(n(B i^{-1})=(\det i)^{-1}i n(B)\) over the residue field.
Since \(\bar d_\gamma=\det i\), the last columns also have
matching reductions under the prescribed map. This includes
vanishing reductions of columns. The inverse conclusion follows
on the same bases. Outside the finite support of
\(\mathfrak t_\gamma\) the sublattice is the ambient one.
Taking the intersection of these local prescriptions inside
the rational space gives the required global lattices. Indeed
their quotients are supported at finitely many primes, so
the Chinese remainder theorem recovers each prescription on
localization.
\end{proof}

\begin{theorem}[Rank and distance at one marked index]
\leavevmode\label{thm:hyperplanerank603}
Let \(j_\gamma\ne j_\delta\), put
\(\rho=\operatorname{rank}_{E_x}(j_\gamma-j_\delta)\), and
let \(r_v=\operatorname{rank}(i_{\gamma,v}-i_{\delta,v})\).
At \(v\mid Q\) set \(b_v=\max(0,\rho-r_v)\), replacing zero
by one when the relative marked matrix is unipotent. Write
\(\|\ \|_{\tau,\mathrm{op}}\) for the operator norm in the
positive Hermitian metrics at a real place of \(F\). Then
\[
 \left(\prod_{v\mid Q}q_v^{b_v}\right)^{1/f}
 \le h_\gamma h_\delta
       \left(\prod_{\tau}\|j_\gamma-j_\delta\|_{\tau,\mathrm{op}}
                          \right)^{\rho/f}
 \le 2^\rho h_\gamma h_\delta .
\]
If the two indices have the same determinant at the marks,
then \(\rho\ge2\). In particular, at one fixed full marked
index one has \((Q/2)^\rho\le h_\gamma h_\delta\). All cells
with \(h_\gamma<Q/2\) at that index have the same rational
transport, including when the discriminant ideals meet
different subsets of the marks.
\end{theorem}
\begin{proof}
The determinant of \(M=j_\gamma^{-1}j_\delta\) is a scalar
\(\xi\in E_x^\times\) preserving the integral determinant
line at every prime. At infinity \(|\tau\xi|=1\). Equal
marked determinants give \(\xi-1\in Q\mathcal O_{E_x}\),
including at the conjugate split factors by duality. A
nonzero such element has absolute norm at least \(Q^{2f}\);
the archimedean bound is \(2^{2f}\). Hence \(\xi=1\).
A unitary matrix is diagonalizable. If its difference from
the identity had rank one, its determinant would be its
unique eigenvalue different from one. This excludes rank one.

For the inequality work on
\(C=C_{0,\gamma}\cap C_{0,\delta}\subset\widehat L_0\).
Let \(J\) be its index ideal. The injection into the sum of
the two quotients shows locally that
\(v_{\mathfrak p}(J)\le
 v_{\mathfrak p}(\mathfrak t_\gamma\mathfrak t_\delta)\).
In the polynomial \(\det(M-(1+Y)I)\) let \(c\) be the
coefficient of \(Y^{k-\rho}\). Up to sign this is
\(\prod_{\lambda\ne1}(\lambda-1)\) on the eigenvalues of
\(M\), with multiplicities. It is nonzero. In a basis of
\(C\) the determinant of \(j_\delta-(1+Y)j_\gamma\) equals
this polynomial times the determinant of \(j_\gamma\) and the
determinant of \(C\hookrightarrow\widehat L_0\). The first
factor outside the polynomial is a local unit and the second
generates \(J\). The determinant on \(C\) is integral by
Lemma~\ref{lem:hyperplanelattice603}.

At a mark the constant term matrix on \(C\) has rank at most
\(r_v\) on reduction. The coefficient in question uses \(\rho\)
columns of that matrix. Its minors contribute the power
\(\max(0,\rho-r_v)\), as in Theorem~\ref{thm:rankmarks602}.
If the relative matrix is unipotent, the determinant on \(C\)
reduces to a scalar multiple of \(Y^k\). Thus
\((c)J\) has valuation at least \(b_v\) in each split factor
and is integral elsewhere. Taking absolute norms, and then
\(2f\)-th roots, gives the first inequality because
\[
     |\tau c|=\prod_{\lambda\ne1}|1-\lambda|
              \le\|j_\gamma-j_\delta\|_{\tau,\mathrm{op}}^\rho .
\]
The equality of absolute values at conjugate embeddings
follows from unitarity. Each operator difference has norm at
most two. For one full marked index all \(b_v=\rho\) and the
two rational primes split completely. This gives the remaining
assertions.
\end{proof}

\begin{corollary}[Multiplicity of transports at a marked index]
\leavevmode\label{cor:rankpacking603}
For \(H\ge1\), let
\[
   r_H=\min\{k,\lfloor 2\log H/\log(Q/2)\rfloor\},\qquad
   n(H)=
   \begin{cases}
      1,&r_H<2,\\
      (1+10H/Q)^{2fr_H(k+1)},&r_H\ge2 .
   \end{cases}
\]
At any fixed full marked index the number of distinct rational
transports among the cells with \(h_\gamma\le H\) is at most
\(n(H)\).
\end{corollary}
\begin{proof}
Choose one cell for each distinct transport. Every difference
of two chosen transports has rank between \(2\) and \(r_H\),
by Theorem~\ref{thm:hyperplanerank603}. If a difference has
rank \(\rho\), its largest archimedean operator norm is at
least \(Q/H^{2/\rho}\ge Q/H\). Fix one chosen transport as
origin and compose with its inverse. The resulting matrices
are unitary and their distances are unchanged. At each
embedding each such matrix can be written
\[
                  I+W\operatorname{diag}(\lambda_a-1)W^*
                     \qquad(1\le a\le r_H).
\]
Use orthonormal eigenvectors for the nontrivial eigenvalues
and pad with eigenvectors of eigenvalue one. The \(k\)-by-\(r_H\)
matrix \(W\) has operator norm one when \(r_H\ge2\), and the
list of eigenvalues has maximum absolute value one. For two
choices of parameters the distance of the resulting matrices
is at most \(4\|W-W'\|_{\mathrm{op}}+
\max_a|\lambda_a-\lambda'_a|\), by adding and subtracting
the products with one factor changed. Give the parameter
space the maximum norm in the two components and across
embeddings. The tuples for different transports are separated
by at least \(Q/(5H)\) and lie in its unit ball. This space
has real dimension \(2fr_H(k+1)\). Disjoint open balls of
radius \(Q/(10H)\) about those tuples fit inside the ball of
radius \(1+Q/(10H)\). The ratio of volumes proves the bound
for every finite subset and hence for the whole set.
\end{proof}

\begin{corollary}[The transverse period on both sides of a frame cutoff]
\leavevmode\label{cor:oneindexperiod603}
Write \(\mathcal P\) for the normalized signed ordinary
contribution of the transverse hyperplane cells in this pair
of rows. Fix \(L\ge1\), put \(d_0=s+1\), and write
\[
       H_L=(T_0T_1)^{d_0}L^{2d_0/f},\qquad
       \alpha=\prod_{v\mid Q}\prod_{a=1}^{k-1}(1-q_v^{-a})^{-1}.
\]
For cells of frame height \(\mathcal H_\gamma<Q^fL\) from
Proposition~\ref{prop:transverseheight582}, let \(Z_{j,i}\)
denote the sum of their signed ordinary operators at the
original index \(i\) and transport \(j\), before applying
the normalized marked projection. Then
\[
 \begin{split}
 \|\mathcal P\|_{\rm HS}
 &\le \sqrt{\alpha\, n(H_L)}
             \left(\sum_{i,j}\|Z_{j,i}\|_{\rm HS}^2\right)^{1/2}\\
 &\quad + |c_{0,1}|w_0w_1 P_\sigma (Q^fL)^{-b}
             \sqrt{\mathfrak M_0(\sigma+2b)\mathfrak M_1(\sigma+2b)} .
 \end{split}
\]
Here the translated section moments and the allowed
\(\sigma,b\) are as in Corollary~\ref{cor:properperiod602}.
For \(H_L\le Q/2\) one may use \(1\) in place of \(n(H_L)\).
In particular when both terminal norms are one this holds at
\(L=(Q/2)^{f/(2(s+1))}\).
\end{corollary}
\begin{proof}
The frame-height estimate gives \(h_\gamma<H_L\) on the
retained set, since \(I_\Lambda\) and \(P_\gamma\) in
Proposition~\ref{prop:transverseheight582} are at least one.
At a fixed full index at most \(n(H_L)\) distinct transports
therefore occur. If \(H_L\le Q/2\), the strict inequality
and Theorem~\ref{thm:hyperplanerank603} give at most one.
Apply Cauchy--Schwarz at each index to the sum of the
\(Z_{j,i}\). The squared norm of its image under
\([G:U]^{-1/2}H\) is at most \(\alpha\) times its squared
norm on raw indices, by Proposition~\ref{prop:amplitudes126}.
For the omitted cells the product of the two heights is at
least \(Q^f L\) at every integration point. Theorem~\ref{thm:twoenergies592}
bounds their absolute sum by the second term. Adding these
two bounds proves the assertion on the original field and
interval.
\end{proof}

The estimate leaves the signed operators for a single pair
\((j,i)\) in the retained sum. Proposition~\ref{prop:repeatedtransport534}
exhibits distinct rational partners even for trivial marked
index and discriminant height one. Their real orbital
coefficients enter this sum separately.


\subsection{Summing integral lifts of a transverse rational class}
\label{sec:signedslice604}

Fix a source row and a target row in the same integral class
of Section~\ref{sec:mixedordinary541}. They lie in a common
Hermitian space \(V\) over \(E=E_x\), with \(F=E^+\) and
\(G=U(V)\) over \(F\). Write \(b=s,t\) for the source and
target subscripts, respectively, and use \(l\) for the length
of each isotropic flag. Put \(Y_b=X_b+E z_b\), the span of
its outer space and terminal vector. Choose a rational Witt
decomposition at each end with middle space \(P_b\perp E z_b\).
Such a decomposition can be chosen inside \(z_b^\perp\) for the
outer hyperbolic pairs, since \(z_b\) is nonisotropic. It
identifies the positive quotient of each row with \(P_b\).
Write \(N_b\) for its algebraic unipotent group and
\(\psi_b\) for its real character on the prescribed section.

For a transverse hyperplane cell \(\gamma\), the subspace
\[
                    C_\gamma=(Y_t+\gamma Y_s)^\perp
\]
is positive of dimension \(k-1\), by Theorem~\ref{thm:cross132}.
Its projections to \(P_t\) and, via \(\gamma^{-1}\), to \(P_s\)
are hyperplanes \(H_t,H_s\). Fix these projected hyperplanes
and fix the rational transport \(j:P_s\to P_t\), including
its determinant in Lemma~\ref{lem:completedcolumn535}.
Thus the two identifications on the common hyperplane induce
\(j|_{H_s}\). Fix also the full marked index \(i\) when
summing the integral terms. In this section the groups that
act on the rational cells are written independently of the
finite support and the arithmetic subgroups used to integrate
the real sections.

\begin{theorem}[Rational classes in the orthogonal complement]
\leavevmode\label{thm:smallcomplement604}
Assume the specified hyperplanes and transport occur.
Put \(V_b'=H_b^\perp\) in \(V\), and
\(N_b'=N_b\cap U(V_b')\), embedding the latter unitary
group by the identity on \(H_b\). Each \(V_b'\) has dimension
\(2l+2\) and signature \((l+1,l+1)\) at every real place.
There is a rational isometry \(\eta_*:V\to V\) taking
\(H_s\) onto \(H_t\) by \(j|_{H_s}\) for which the rational
\(N_t(F),N_s(F)\) double classes of transverse hyperplane
pairs with these data are in bijection with
\[
 N_t'(F)\backslash
  \{a\in SU(V_t')(F):Y_t+a\eta_*Y_s=V_t'\}
       /\,\eta_*N_s'(F)\eta_*^{-1}.
\]
The cells at the original finite level map into this set.
\end{theorem}
\begin{proof}
At the target \(C_\gamma\) lies in \(X_t^\perp\cap z_t^\perp\).
Its projection to \(P_t\) is injective since the kernel lies in
the isotropic \(X_t\), whereas \(C_\gamma\) is positive.
It is therefore the graph of a rational linear map
\(A_t:H_t\to X_t\). The identical argument applied to
\(\gamma^{-1}C_\gamma\) gives \(A_s:H_s\to X_s\).
Extend \(A_b\) by zero on the orthogonal complement of \(H_b\)
inside the middle space. In coordinates consisting of \(X_b\),
the middle space, and its outer dual summand \(X_b^\vee\),
respectively, set
\[
 n(A_b)=
 \begin{pmatrix}
 I&A_b&-\tfrac12 A_b A_b^\dagger\\
 0&I&-A_b^\dagger\\
 0&0&I
 \end{pmatrix}.
\]
The adjoint uses the middle Hermitian form and the pairing of
the two outer summands. Multiplication with the Witt form
shows this matrix is unitary. The two terms linear in \(A_b\)
cancel the cross terms from the middle space, and the two
half terms cancel its quadratic term. It lies in \(N_b(F)\),
takes \(H_b\) onto the required graph, fixes \(X_b\) and
\(z_b\) pointwise, and has determinant one. Its outer simple
coordinates and its terminal character coordinate vanish,
so \(\psi_b(n(A_b))=1\).

Consequently \(\eta=n(A_t)^{-1}\gamma n(A_s)\) takes the
embedded \(H_s\) onto \(H_t\) by the prescribed isometry.
Take one such \(\eta\) as \(\eta_*\). The determinant of a
rational cell map in \(V\) is determined by \(j\) and the
two fixed decorated frames. Indeed the determinant-line
construction of Lemmas~\ref{lem:flagdeterminant534}
and~\ref{lem:completedcolumn535} composes the ambient
determinant with fixed isomorphisms at its two ends, using
the normalized terminal vector at each end. Hence
\(\eta\eta_*^{-1}\) fixes \(H_t\) pointwise and has determinant
one. It is exactly an element of \(SU(V_t')(F)\).
The perpendicular of \(Y_t+\eta Y_s\) is \(H_t\), so their
sum is \(V_t'\). Conversely the displayed spanning condition
forces the perpendicular to be precisely \(H_t\), and forces
transversality since the two summands both have dimension
\(l+1\). The restriction to the hyperplane and the determinant
then recover \(j\) by its uniqueness. In the ambient signature
\((l+k,l+1)\) the removal of \(k-1\) positive lines gives the
asserted signature.

Finally suppose two such representatives satisfy
\(\eta'=u_t\eta u_s^{-1}\) with \(u_b\in N_b(F)\).
Each \(u_b\) preserves \(Y_b\). Taking perpendiculars of the
two frame sums gives \(u_t H_t=H_t\). As \(N_t\) acts
trivially on the middle quotient, \(u_t\) fixes \(H_t\)
pointwise. Applying the argument at the source gives the
same conclusion for \(u_s\) on \(H_s\). Thus both lie in the
indicated smaller unipotent groups. The reverse implication
is immediate by inclusion. This proves the bijection.
\end{proof}

For the next theorem use the full column-valued ordinary
sections \(W_b\). An integral cell has its arithmetic
stabilizer \(\Gamma_\gamma\) and finite scalar
\(c_\gamma=c_{t,s}\epsilon_\gamma\), where \(c_{t,s}\) is
the volume of the row pair and \(\epsilon_\gamma\) its finite
character value, as in Proposition~\ref{prop:partnerweights531}.
All sums below are on the fixed rows and index. Write
\(G_\infty=G(F\otimes_\QQ\RR)\), with the analogous convention
for subgroups.

\begin{theorem}[The signed sum over the integral lifts of one class]
\leavevmode\label{thm:liftphase604}
For a class \([a]\) of Theorem~\ref{thm:smallcomplement604}
choose \(\eta=a\eta_*\), and put
\[
 D_\eta=N_t\cap\eta N_s\eta^{-1}
        =N_t'\cap\eta N_s'\eta^{-1}.
\]
Discard a class if the two real characters differ on \(D_\eta\);
its contribution is zero. For each integral lift at \(i\) write
\(\gamma=u_t\eta u_s^{-1}\) with \(u_b\in N_b(F)\), and
put \(\Lambda_{\gamma,\eta}=u_t^{-1}\Gamma_\gamma u_t\).
On choosing Haar measure on \(D_{\eta,\infty}\), the entire
contribution at \(i\) from this class before the marked
projection is the product \(m_i(\eta) I_\eta\), where
\[
 \begin{split}
 m_i(\eta)&=\sum_{\gamma\mapsto[a]}
   c_\gamma\,\psi_t(u_t)\overline{\psi_s(u_s)}
     \operatorname{vol}(\Lambda_{\gamma,\eta}\backslash D_{\eta,\infty}),\\
 I_\eta&=\int_{D_{\eta,\infty}\backslash G_\infty}
                   W_t(g)W_s(\eta^{-1}g)^*\,d\dot g .
 \end{split}
\]
The first sum is finite. It admits the following exact
cancellation test. For a nonempty set of lifts put
\(\Delta=\bigcap_\gamma\Lambda_{\gamma,\eta}\),
\(d_\gamma=[\Lambda_{\gamma,\eta}:\Delta]\), and
\(M=\operatorname{lcm}_\gamma d_\gamma\).
Write the total phase of a lift as \(\zeta^{e_\gamma}\)
for a common primitive \(n\)-th root of unity. Then \(m_i(\eta)=0\)
if and only if the cyclotomic polynomial \(\Phi_n\) divides
\[
        \sum_{\gamma\mapsto[a]}(M/d_\gamma)X^{e_\gamma}
                                 \qquad\text{in }\ZZ[X],
\]
using exponents in \(\{0,\ldots,n-1\}\). All factors in
this polynomial come from the original lifts and the fixed
class representative.
\end{theorem}
\begin{proof}
If \(d\in D_\eta\) and \(v\in H_t\), then \(d v-v\in X_t\).
Reading \(d\) in the transported radical gives
\(d v-v\in\eta X_s\) as well. Transversality gives zero.
Thus \(d\) fixes \(H_t\) pointwise, which proves the equality
of groups. In the original integral \eqref{eq:orb133} change
variable by left multiplication by \(u_t\). Its stabilizer
becomes \(\Lambda_{\gamma,\eta}\), and the two sections supply
the two indicated character factors. In particular the finite
phase has been retained even if a unipotent change of variables
in this calculation is nonintegral at a finite place.
Lemma~\ref{lem:zero133} identifies the transformed stabilizer
as a uniform lattice in \(D_{\eta,\infty}\).
Integrating along its compact quotient gives zero for different
characters and gives exactly its covolume for matching ones.
This proves the factorization, with compatible quotient measures,
for each lift. Changing \(u_t,u_s\) for the same lift changes
them by corresponding elements of \(D_\eta(F)\); equality of
the characters cancels the change of phase. The resulting
conjugation in the unipotent group preserves its Haar measure.

Every prefix exterior vector, including the last one, is fixed
by the respective \(N_b(F)\), since the changes of the ordered
frame are unitriangular. Hence the frame height
\(\mathcal H_\gamma\) of Proposition~\ref{prop:transverseheight582}
is fixed throughout the rational class. Take \(L\) large enough
to put that height strictly below \(Q^fL\). The finite count
in Theorem~\ref{thm:smallorbits591} proves finiteness of its
supported relevant integral lifts on the retained rows.

The lattices in the group \(D_\eta(F)\) just obtained are
commensurable. Indeed each imposes preservation of finitely
many full row lattices after rational conjugation and, if a
level is present, a finite congruence. For two full lattices
\(L_1,L_2\), select a positive integer \(a\) with
\(aL_1\subset L_2\subset a^{-1}L_1\). An automorphism of
\(L_1\) congruent to one modulo \(a^2\) preserves \(L_2\):
its difference from the identity takes \(a^{-1}L_1\) into
\(aL_1\), and so does the difference for its inverse. Increasing
the modulus imposes each required finite level. The reduction
of a lattice modulo an integer is finite. Applying this
argument in both directions, then intersecting finitely many
groups, proves \(d_\gamma<\infty\).

The real character of a rational unipotent element has finite
order here. The original character is adelic and is trivial on
rational elements. Its finite additive factors have finite
order at any given argument and are one almost everywhere.
The real factor is their inverse. The other factor
\(\epsilon_\gamma\) is already a finite character value.
Now divide \(m_i(\eta)\) by
\(c_{t,s}\operatorname{vol}(\Delta\backslash D_{\eta,\infty})/M\).
The result is exactly the displayed integer polynomial at
\(\zeta\). Its vanishing is equivalent to divisibility by the
minimal polynomial \(\Phi_n\), including \(n=1\). This gives
the test.
\end{proof}

\begin{corollary}[The squared norm at fixed hyperplanes and transport]
\leavevmode\label{cor:signedslice604}
Let \(Z\) be this full signed sum at \(i\), for the fixed
\(H_s,H_t,j\). Use \(K=U(H_t)_\infty\), acting on the source
by its identification with \(U(H_s)_\infty\). For each common
irreducible summand \(\mu\) in the two compact restrictions,
let \(V_{b,\mu}\) denote its subspace and choose a unitary
intertwiner \(J_\mu:V_{s,\mu}\to V_{t,\mu}\), extended by
zero on its orthogonal complement. Write \(d_\mu=\dim V_{s,\mu}\)
and \(L_\infty=U(V_t')_\infty\). Define
\[
 b_\mu(\eta)=
 \int_{L_\infty\backslash G_\infty}
 \int_{D_{\eta,\infty}\backslash L_\infty}
 \operatorname{tr}\!\left(J_\mu^* W_t(hg)
                              W_s(\eta^{-1}hg)^*\right)
                                  \,d\dot h\,d\dot g .
\]
Then, summing over the supported relevant classes with a lift at \(i\),
\[
                \|Z\|_{\rm HS}^2=
           \sum_\mu \frac1{d_\mu}
            \left|\sum_{[a]} m_i(a\eta_*)\,b_\mu(a\eta_*)\right|^2 .
\]
These are identities for the original sections, with either
terminal frequency. They apply to the full absolutely
convergent expansion and to a cutoff by frame height. The
normalized marked column at \(i\) has norm one, so adjoining
that column preserves the norm in this formula.
\end{corollary}
\begin{proof}
The inclusions \(D_{\eta,\infty}\subset L_\infty\subset G_\infty\)
are inclusions of unimodular groups. Decompose the Haar
quotient on \(D_{\eta,\infty}\backslash G_\infty\) first over
\(L_\infty\backslash G_\infty\). This gives the displayed
integral for \(\operatorname{tr}(J_\mu^* I_\eta)\) and explains
why its inner integral runs over a unitary group of dimension
\(2l+2\) in its defining representation. Absolute integrability
justifies disintegration; for the ordinary sections it also
follows from the unfolded nonnegative majorant in the proof of
Theorem~\ref{thm:twoenergies592} with \(b=0\) and
\(\sigma=4D+1\) in the notation of Corollary~\ref{cor:properperiod602}.
The finiteness of the moments at this order on the prescribed
rows follows from Theorem~\ref{thm:mixedhalf556}, as explained
after Theorem~\ref{thm:twoenergies592}. Thus the tensor integrals
can be grouped in the indicated order. A height cutoff keeps
whole rational classes by the prefix invariance proved above.

All the \(I_\eta\) intertwine the same actions of \(K\).
Indeed an element of \(K\) fixes both frames, commutes with
the action on the orthogonal complement, and acts on the
source compact column via \(j|_{H_s}\). Change variables by
that compact element in the quotient integral. Its action
preserves measure, giving the intertwining identity as in
Theorem~\ref{thm:common133}. Each restriction has multiplicity
one by compact unitary branching. Schur's lemma gives
\(I_\eta=\sum_\mu b_\mu(\eta)J_\mu/d_\mu\).
The operators \(J_\mu\) are orthogonal in Hilbert--Schmidt
inner product, with squared norms \(d_\mu\). Summing with
the signed coefficients from Theorem~\ref{thm:liftphase604}
and taking the squared norm proves the result.
\end{proof}

This formula sums the phases of every integral lift before
the norm is taken. A uniform estimate for the sums over
the remaining rational classes, with the original real
sections and the different hyperplanes combined, is still
needed in Section~\ref{sec:missing}.


\subsection{Finite local factors of the signed sum of integral lifts}
\label{sec:localphase605}

Retain a pair of rows, the two projected hyperplanes, the rational
transport, and the full index \(i\) on \(E_x\) from
Section~\ref{sec:signedslice604}. Fix a relevant rational class
with representative \(\eta\) and at least one supported integral
lift at \(i\). Its scalar \(m_i(\eta)\) in
Theorem~\ref{thm:liftphase604} multiplies one real tensor. We
calculate this scalar with the original outer characters. In
particular the finite tests below are taken before any sum
over marked indices or terminal projections.

Set \(U=N_t\times N_s\), and embed \(D=D_\eta\) in \(U\) by
\(d\mapsto(d,\eta^{-1}d\eta)\). For a finite place \(v\) of
\(F=E_x^+\) let \(U_v=U(F_v)\), \(D_v=D(F_v)\), and let
\(K_v\) be the product of the two local integral unipotent
subgroups of the rows. Thus \(U(F)\cap\prod_v K_v\) is the
product of the two integral unipotent groups in the original
double quotient. For \(w=(w_t,w_s)\in U_v\) write
\(\gamma(w)=w_t\eta w_s^{-1}\). Let \(\Omega_v\) be the set
of double classes \(K_v w D_v\) for which \(\gamma(w)\) is
an integral isometry with the prescribed local support and,
where imposed, marking \(i_v\). Put
\[
   K^D_{v,w}=D_v\cap w^{-1}K_v w,\qquad
   z_v(w)=\epsilon_v(\gamma(w))\,
         \overline{\psi_{t,v}(w_t)}\,\psi_{s,v}(w_s).
\]
Here \(\epsilon_v\) is the local factor of the ordinary outer
finite phase in Proposition~\ref{prop:partnerweights531};
\(\psi_{b,v}\) is its component at \(v\) for the adelic character
with archimedean component \(\psi_b\). At residual characteristic \(p\)
each of these local additive character values is a root of
unity of \(p\)-power order.

\begin{theorem}[Normalization by an actual integral lift]
\leavevmode\label{thm:localphase605}
The integral lifts of the fixed class at \(i\), up to the
original integral unipotents, are in bijection with
\(\prod_{v<\infty}\Omega_v\). Each set in the product is finite,
and almost all are singletons. At \(v\) choose a class \(\omega_v^*\)
maximizing the Haar volume of \(K^D_{v,w}\). The local choices
are realized together by an integral lift \(\gamma_*\) on the
specified rows and index. For \(\omega\in\Omega_v\) put
\[
       n_v(\omega)=
       \frac{\operatorname{vol}(K^D_{v,\omega_v^*})}
            {\operatorname{vol}(K^D_{v,\omega})},\qquad
       A_v=\sum_{\omega\in\Omega_v}n_v(\omega)z_v(\omega).
\]
The ratio \(n_v(\omega)\) is a nonnegative power of the
residual characteristic. For any integral lift \(\gamma\)
with local classes \((\omega_v)\), one has
\[
 \frac{\operatorname{vol}(\Lambda_{\gamma,\eta}\backslash D_\infty)}
      {\operatorname{vol}(\Lambda_{\gamma_*,\eta}\backslash D_\infty)}
       =\prod_v n_v(\omega_v),\qquad
 \frac{m_i(\eta)}
      {c_{t,s}\operatorname{vol}(\Lambda_{\gamma_*,\eta}\backslash D_\infty)}
       =\prod_v A_v .
\]
Almost all \(A_v=1\). The normalized scalar in the second
identity is an algebraic integer.
\end{theorem}
\begin{proof}
First check the local functions. In a fixed pair of rows the
finite support conditions are membership in the local lattice
stabilizers and the specified congruences on the local frames
and columns. Their conjunction is the original global finite
support condition. Integral unipotent changes preserve that
condition and the full marked index. Under such a left change
in \(U_v\) the two finite seeds multiply their phase by
\(\psi_{t,v}(k_t)\overline{\psi_{s,v}(k_s)}\), by their
character transformation in \eqref{eq:seed127}, equivalently
by the change of representatives in the proof of
Proposition~\ref{prop:partnerweights531}. These two factors
cancel those from the explicit characters in \(z_v\).
A right change in \(D_v\) fixes \(\gamma(w)\) and also fixes
the combined character. Indeed the ratio of the two finite
adelic characters on \(D\) is one on \(D(F)\), since each
global character is trivial on rational points and the two
real characters agree there by relevance. Rational points
are dense in \(D(\mathbb A_{F,f})\): logarithm and exponential
are inverse polynomial maps for this unipotent group in
characteristic zero, and additive approximation holds in its
rational Lie algebra. The continuous ratio is therefore one
at each finite place. This proves that \(z_v(\omega)\) is
well defined. The stabilizer volumes are also well defined,
since a right change conjugates the stabilizer inside the
unimodular group \(D_v\). A continuous local additive character
at any specified argument has \(p\)-power order, since
multiplication of that argument by a large power of \(p\)
puts it in the character kernel.

Write \(K=\prod_vK_v\). Rational points of \(U\) are dense
in its finite adeles by the same additive approximation.
Thus in \(K\backslash U(\mathbb A_{F,f})/D(\mathbb A_{F,f})\)
each double class has a rational representative. If \(w,w'\)
are two rational representatives of the same double class,
there is a finite adelic \(d\in D\) with \(w'd^{-1}w^{-1}\in K\).
This is an open condition on \(d\). By density it holds for
some \(d\in D(F)\), and the resulting left multiplier belongs
to \(U(F)\cap K\). This proves injectivity in the rational
description as well. The stabilizer of \(\eta\) in \(U(F)\)
is exactly the displayed copy of \(D(F)\). A rational
\(\gamma(w)\) integral in both directions at every finite
place is an isometry of the given global lattices, by
intersection of their localizations. This proves the
description by the restricted product of the local supported
double classes, retaining the original finite conditions.

Start with any one supported global lift. Any single local
choice can replace its class at that place and gives an
adelic element, hence a global lift by the preceding paragraph.
Infinitely many local choices at one place, or alternatives
at infinitely many different places, would therefore give
infinitely many distinct integral lifts of the same rational
class. Their finiteness was proved in
Theorem~\ref{thm:liftphase604}. Consequently the restricted
product is the entire stated product, and one can choose all
the maximizing stabilizers at once. Almost everywhere the
initial rational representative is integral in \(U\), the
support is the basic class, and all the finite phases are one.

For a rational \(w\) representing a lift, its lattice in \(D\)
is \(D(F)\cap\prod_v K^D_{v,w}\), under the fixed embedding
in \(U\). This follows by conjugating the two original integral
unipotent stabilizers, just as in Theorem~\ref{thm:liftphase604}.
Choose compatible real and finite Haar measures and write
\(V_D=\operatorname{vol}(D(F)\backslash D(\mathbb A_F))\).
This volume is finite: a rational central series with vector
group quotients reduces compactness to that of \(\mathbb A_F/F\).
By density at the finite places a fundamental domain modulo
\(\prod_v K^D_{v,w}\) is given by the real quotient by the
displayed lattice. Equivalently, a real fundamental domain
times \(\prod_v K^D_{v,w}\) is an adelic fundamental domain
up to null sets. Its volume gives
\[
  \operatorname{vol}(\Lambda_{\gamma,\eta}\backslash D_\infty)
                    \prod_v\operatorname{vol}(K^D_{v,w})=V_D .
\]
Use volume one for the usual integral subgroups at almost all
places. Division by the same formula for \(\gamma_*\) proves
the first identity.

Every compact subgroup of \(D_v\) is pro-\(p\). In a faithful
unipotent matrix representation it preserves a lattice:
the span over local integers of its orbit on a fixed lattice
is bounded, full, and stable. Its finite reduction consists
of unipotent matrices, and thus has \(p\)-power order by
Cauchy's theorem. The successive congruence kernels have
elementary abelian \(p\)-group subquotients. This proves the
assertion. Two compact open subgroups of \(D_v\) have open
intersection, of \(p\)-power index in each. Their volume ratio
is a power of \(p\); choosing the largest volume for the
numerator makes that exponent nonnegative.

For rational \(w\), the product of \(z_v(w)\) is the full phase
\(\epsilon_\gamma\psi_t(w_t)\overline{\psi_s(w_s)}\),
by triviality of the adelic characters on rational elements.
Insert this identity and the covolume ratios into the finite
sum in Theorem~\ref{thm:liftphase604}. Distributing a product
of finite sums gives the second identity. The local weights
are integers and the local phases roots of unity, proving
integrality.
\end{proof}

\begin{corollary}[Local cancellation and the terminal valuation]
\leavevmode\label{cor:phasevaluation605}
At \(v\) of residual characteristic \(p\), write all the
phases as \(\zeta^{e(\omega)}\) for a primitive \(p^h\)-th
root of unity, with exponents between \(0\) and \(p^h-1\).
Put \(a_e=\sum_{e(\omega)=e}n_v(\omega)\) and
\(W_v=\sum_e a_e\). For \(h=0\) one has \(A_v=W_v>0\).
For \(h>0\) the factor vanishes exactly when, for every
\(0\le r<p^{h-1}\),
\[
        a_r=a_{r+p^{h-1}}=\cdots=a_{r+(p-1)p^{h-1}} .
\]
The whole signed scalar vanishes exactly when some local
factor does. If the number of classes with maximal stabilizer
volume at \(v\) is prime to \(p\), then \(A_v\) is a unit
at every coefficient prime above \(p\).

Let \(T\) be the terminal rational prime of the original
period. When all factors are nonzero, the valuation at any
prime above \(T\), normalized by \(\operatorname{val}_T(T)=1\),
obeys
\[
 0\ \le\
 \operatorname{val}_T\!\left(
   \frac{m_i(\eta)}
        {c_{t,s}\operatorname{vol}(\Lambda_{\gamma_*,\eta}\backslash D_\infty)}
                         \right)
 \ \le\ \sum_v
   \frac{[\QQ(\zeta_{p^h}):\QQ]}{e_{v,T}f_{v,T}}\,
                       \frac{\log W_v}{\log T}.
\]
Here \(p,h\) are local to \(v\), and \(e_{v,T},f_{v,T}\)
are the ramification and residue degrees at \(T\) of that
cyclotomic field. Only finitely many summands can be nonzero.
\end{corollary}
\begin{proof}
For \(h>0\) the minimal polynomial of \(\zeta\) is
\(\sum_{j=0}^{p-1} X^{jp^{h-1}}\). A polynomial of degree
less than \(p^h\) is a multiple of it precisely when its
quotient has degree less than \(p^{h-1}\) and multiplication
repeats each coefficient of that quotient in the indicated
positions. This proves the test. Reducing \(A_v\) at \(p\)
sends every phase to one. Exactly the maximal-volume classes
have weight one; every other weight is a positive power of
\(p\). Hence the residue is the asserted count modulo \(p\).

Every conjugate of \(A_v\) has absolute value at most \(W_v\).
For a nonzero \(A_v\) its absolute field norm is a positive
integer at most \(W_v^{[\QQ(\zeta_{p^h}):\QQ]}\). The norm
of the chosen prime ideal over \(T\) is \(T^{f_{v,T}}\).
Taking logarithms bounds its exponent in \((A_v)\); division
by \(e_{v,T}\) gives the stated normalization. The product
formula of Theorem~\ref{thm:localphase605} adds these valuations.
At a singleton place \(W_v=1\).
\end{proof}

This gives the integral arithmetic coefficient for one rational
class relative to an actual lift at its prescribed index. The
sizes of these factors, the real tensors in
Corollary~\ref{cor:signedslice604}, and their sum over classes
and hyperplanes still enter the full period comparison.


\subsection{Finite residue sums for the signed local factors}
\label{sec:orbitcoordinates606}

In Theorem~\ref{thm:localphase605}, fix the two original rows
and a rational representative \(\eta\). The integral matrices
in its rational unipotent orbit may have large denominators
in the variables \(w_t,w_s\) that produce them. We give a finite
range of variables at every finite place, independent of the
full marked index to be tested. The terminal norms, row lattices
and field \(E_x\) are those of Section~\ref{sec:missing}.

View the space \(M\) of \(E_x\)-linear maps between the two
ambient row spaces as an \(F=E_x^+\) vector space. The group
\(U=N_t\times N_s\) acts linearly on \(M\) by
\((w_t,w_s)z=w_t z w_s^{-1}\). Write \(D\) for the stabilizer
of \(\eta\) in this product. At a finite place \(\mathcal O_v\)
has uniformizer \(\pi_v\) and residue size \(q_v\).

\begin{lemma}[Polynomial coordinates from the two rational flags]
\leavevmode\label{lem:orbitcoordinates606}
Set \(r=\dim_F U-\dim_F D\). There are explicitly determined
rational polynomials \(P_1,\ldots,P_r\) in the matrix entries
over \(F\) and elements \(X_1,\ldots,X_r\) in the rational
Lie algebra of \(U\) such that, with
\[
 s(t)=\exp(t_1X_1)\cdots\exp(t_rX_r),\qquad z(t)=s(t)\eta ,
\]
one has \(P(z(t))=t\) and \(z(P(z))=z\) on the orbit.
Here \(s(t)\eta\) denotes the indicated action. These
identities hold over \(F\) and every \(F_v\).
In these coordinates any Haar quotient measure on \(U_v/D_v\)
is a constant multiple of product additive Haar measure on
\(F_v^r\).
\end{lemma}
\begin{proof}
Take bases in which each row radical is upper unitriangular,
by ordering its two outer flags and the middle block in Witt
coordinates. On the source dual reverse the order, so that
the dual action is upper unitriangular too. In a tensor basis
of target times source dual the product group acts by upper
unitriangular matrices. On passing from \(E_x\) to \(F\),
refine each diagonal identity block with an \(F\)-basis. All
changes can therefore be made in fixed rational coordinates.

Start with \(U_1=U\). If \(U_j\) moves \(\eta\), take the
largest coordinate \(a\) on which \(u\eta-\eta\) is nonzero
for some \(u\in U_j\). The function
\(\lambda(u)=(u\eta-\eta)_a\) is a nonzero homomorphism to
the additive group. Indeed all larger coordinates of that
difference vanish and an upper unitriangular transformation
fixes its \(a\)-th coordinate. Applying this to
\(uv\eta-\eta=u(v\eta-\eta)+(u\eta-\eta)\) proves additivity.
On a one-parameter exponential it is a linear function of
the parameter, since it is an additive polynomial in
characteristic zero. Exponentials exhaust each group here:
initially this follows from finite logarithm and exponential
for the row radicals, and subsequently the kernel of a
character is the exponential of the kernel of its differential.
Thus \(a\) is the largest index on which some \(X\eta\)
is nonzero in a basis of the current rational Lie algebra.
Here \(X\eta\) is the derivative of the action. Linear
algebra gives \(X_j\) with \((X_j\eta)_a=1\).

Put \(U_{j+1}=\ker\lambda\). For a current point \(y=u\eta\)
read \(t_j=y_a-\eta_a\), and replace \(y\) by
\(\exp(-t_jX_j)y\). This lies in the orbit of \(U_{j+1}\),
and is polynomial in the initial coordinates of \(y\), since
each exponential is a finite polynomial. Iterating produces
the polynomials \(P_j\) and reverses the product defining
\(s(t)\). Each kernel has codimension one and still contains
\(D\); the process stops precisely at \(D\). This proves
existence, uniqueness of the ordered parameters and all
claimed identities. At each step \(U_j\) is the product of
the one-parameter section and its normal kernel. Conjugating
the kernel by this section preserves Haar measure: the
adjoint acts unipotently on its Lie algebra and has determinant
one. Haar measure is thus a constant times the product of
additive measure in \(t_j\) and Haar measure on the kernel.
Iteration and division by the measure on \(D_v\) prove the
measure assertion.
\end{proof}

Let \(\mathcal L_v\) be the full \(\mathcal O_v\)-lattice of
integral row maps in \(M(F_v)\); it suffices here to impose
integrality in the forward direction. In a basis of this
lattice express each \(P_j\) as a polynomial. Take \(a_{j,v}\)
to be the larger of zero and the maximum of the negatives
of its coefficient valuations, omitting zero coefficients.
The value zero is used for an empty maximum. Put
\[
                     B_v=\prod_{j=1}^r\pi_v^{-a_{j,v}}\mathcal O_v .
\]
Represent \(U_v\) as a pair of matrices on the original row
lattices, using bases over \(\mathcal O_v\) after restriction
of scalars. Choose \(c_v\ge0\) so that every element of \(U_v\)
congruent to the identity modulo \(\pi_v^{c_v}\) in these
lattices lies in \(K_v\). For \(c_v=0\) the condition means
that its difference from the identity is integral. Such an
exponent is given by the two finite levels; \(c_v=0\) is
available almost everywhere. Integral unipotent matrices have
integral inverses by the finite geometric series.

In the chosen row lattice bases form the pair of matrix
polynomials
\[
              s(t+\delta)s(t)^{-1}-1 .
\]
Replace \(t_j\) by \(\pi_v^{-a_{j,v}}T_j\) and leave
\(\delta_j\) as indeterminates. Let \(b_v=c_v\) plus the
larger of zero and the maximum of the negatives of the
coefficient valuations of these polynomials, with the same
zero convention. All these quantities use the given rows
and rational matrices and require no choice of integral
lifts of the class.

\begin{theorem}[Local cancellation on finite residue classes]
\leavevmode\label{thm:residueorbits606}
For a fixed original full index \(i\), test one representative
of each residue class in
\[
       B_v/(\pi_v^{b_v}\mathcal O_v)^r,\qquad
       M_v=q_v^{\,r b_v+\sum_j a_{j,v}} .
\]
At \(t\), test \(z(t)\) against the original local integrality
and marking conditions. If supported, attach the phase
\(z_v(s(t))\) of Theorem~\ref{thm:localphase605}; otherwise
attach zero. These values are independent of the representatives.
The sum of the values is zero if and only if the signed local
factor \(A_v\) is zero, for each globally supported index.
If \(m_v\) is the least number of these residue classes in
any one supported \(K_v\)-orbit, the sum is exactly \(m_v A_v\).
In particular,
\[
                       |\Omega_v|\le W_v\le M_v .
\]
All \(a_{j,v}\) and \(b_v\) vanish outside a finite set determined
by the rational input, row lattices and original levels. The
conclusions include \(r=0\), with a single test.
\end{theorem}
\begin{proof}
If \(z\) is supported, it is in \(\mathcal L_v\). Polynomial
evaluation gives \(P(z)\in B_v\). Conversely \(z(t)\) and the
source and target matrices of \(s(t)\) give the actual local
orbit point and its two character arguments. For \(t\in B_v\)
and \(\delta\in(\pi_v^{b_v}\mathcal O_v)^r\), the definition
of \(b_v\) gives
\[
                        s(t+\delta)s(t)^{-1}\in K_v .
\]
Indeed all monomials in the matrix difference contain at
least one \(\delta\), since the difference vanishes identically
at \(\delta=0\). Substituting \(\pi_v^{b_v}\) times integers
in these variables makes all its entries divisible by
\(\pi_v^{c_v}\) on the whole indicated box. Thus each residue
class belongs to a single orbit; the support and phase are
constant there. Counting all residue classes gives \(M_v\).

Give the additive parameter box its product Haar measure.
By Lemma~\ref{lem:orbitcoordinates606} each small residue
class has the same quotient measure. The measure of
\(K_v wD_v/D_v\) is
\(\operatorname{vol}(K_v)/\operatorname{vol}(K^D_{v,w})\):
integrate the indicator of \(K_vw\) first in the right
\(D_v\) variable; over every point of this orbit its
integral in that variable is \(\operatorname{vol}(K^D_{v,w})\)
using representatives in \(K_vw\). Thus the orbit with
largest stabilizer volume contains the least number \(m_v\)
of small residue classes. All other orbit counts divided by
\(m_v\) are precisely the integers \(n_v(\omega)\) of
Theorem~\ref{thm:localphase605}. This proves the equality of
signed sums. Dividing the number of supported residue classes
by \(m_v\ge1\) gives \(W_v\le M_v\).

Only finitely many coefficients of the rational polynomials
are involved. The rational bases used to construct them give
integral lattice bases away from finitely many primes, and
the levels are full there. Clearing the denominators of the
bases in both directions, of the polynomials \(P_j,s,s^{-1}\),
and including the finite level primes, therefore leaves
\(a_{j,v}=b_v=0\) at every other place. At the exceptional
places the same polynomial construction takes place over
the fraction field and the residue classes use the resulting
denominator exponents.
\end{proof}

\begin{corollary}[An ideal norm bound for the signed coefficient]
\leavevmode\label{cor:orbitideal606}
Form the integral ideal of \(F\)
\[
           \mathfrak b_\eta=
              \prod_{v<\infty}\mathfrak p_v^{\,r b_v+\sum_j a_{j,v}} .
\]
For each supported full index, divide \(m_i(\eta)\) by the
row volume and by the covolume of its lift \(\gamma_*\) from
Theorem~\ref{thm:localphase605}. Every complex conjugate of
the resulting algebraic integer has absolute value at most
\(N_{F/\QQ}\mathfrak b_\eta\). The number of integral lifts
at that index is at most this same ideal norm. In the upper
bound for its normalized terminal valuation in
Corollary~\ref{cor:phasevaluation605}, one can use
\((r b_v+\sum_j a_{j,v})\log q_v\) in place of \(\log W_v\).
\end{corollary}
\begin{proof}
The local factors are weighted sums of roots of unity. Every
conjugate is bounded by \(W_v\). Multiply the local bounds of
Theorem~\ref{thm:residueorbits606} over the finite places and
use the two product statements of Theorem~\ref{thm:localphase605}.
The valuation bound follows by the same substitution in the
norm estimate. Outside the indicated finite set there is one
residue class and the supported factor is a unit of absolute
value one at every embedding.
\end{proof}

This bound is effective for given rational frames and uses their
actual finite levels. A bound for these ideals across the
rational classes that contribute on \([V,V^{3/2}]\), together
with control of the real tensors and their signed sum, is
still required for the unrestricted period estimate.


\subsection{The signed coefficient in the measures of the two rows}
\label{sec:rowmeasures607}

Use the original transverse hyperplane rows over \(E=E_x\),
with \(F=E^+\) totally real and \(\Delta_F=|\operatorname{disc}F|\).
For the class representative \(\eta\), let \(U=N_t\times N_s\),
\(D=D_\eta\), and \(r=\dim_F U-\dim_F D\).
Write \(\Gamma_b=N_b(F)\cap\prod_{v<\infty}K_{b,v}\) for the
original integral row radical. Fix rational invariant volume
forms on the two row radicals and give \(U\) their product
form \(\omega_U\). Use Lebesgue measure on \(F_\infty\) and
give \(\mathcal O_v\) additive volume one at each finite
place. These choices specify a real measure \(du_\infty\)
and a product of finite measures on \(U\). Put
\[
              v_U=\operatorname{vol}_{du_\infty}
                  ((\Gamma_t\times\Gamma_s)\backslash U_\infty).
\]
It is the product of the two row unipotent covolumes.

Take the rational directions \(X_j\) constructed in
Lemma~\ref{lem:orbitcoordinates606}. Give \(D\) the invariant
form whose value on a tangent basis \(Y\) at the identity is
\(\omega_U(X_1,\ldots,X_r,Y)\), using its embedding in \(U\).
This specifies the Haar measure on \(D_\infty\) in both
\(m_i(\eta)\) and \(I_\eta\) of Theorem~\ref{thm:liftphase604}.
The measure on the original real group \(G_\infty\) is the one
used for the ordinary period; on \(D_\infty\backslash G_\infty\)
use its quotient by the specified \(D_\infty\) measure.

\begin{lemma}[Covolume with additive integers of volume one]
\leavevmode\label{lem:unipotentmeasure607}
For an \(F\)-split unipotent group \(B\) of dimension \(d\),
give its local groups the measures of any one rational
invariant top form with the additive conventions above.
Then \(\operatorname{vol}(B(F)\backslash B(\mathbb A_F))
=\Delta_F^{d/2}\).
\end{lemma}
\begin{proof}
For one additive coordinate, finite additive approximation
reduces every finite coordinate to the product of the rings
of integers. The residual real lattice is \(\mathcal O_F\)
in its real embeddings. The matrix of embeddings of an
integral basis has determinant of absolute value
\(\sqrt{\Delta_F}\), by the trace definition of the
discriminant. This proves the assertion in dimension one.
A rational scalar multiplying a form has product of its
normalized absolute values equal to one, by the number field
product formula.

For the induction choose a rational codimension-one ideal in
the nilpotent Lie algebra containing its commutator and use
exponential coordinates. The corresponding normal unipotent
subgroup has additive quotient and a rational one-parameter
section. In particular the group on rational points maps
onto the rational points of the quotient. Disintegrate first
over its additive adelic quotient. Each fibre has the
covolume of the normal subgroup modulo its rational points:
the section identifies the fibre with that quotient up to
conjugation, and conjugation by the section has Jacobian one
on the normal subgroup. Choose compatible rational forms at
the identity. The Haar measures disintegrate as their product
at all places, and changing a rational form by a scalar costs
one globally. Multiplying the induction hypothesis by
\(\sqrt{\Delta_F}\) proves the assertion. Dimension zero
uses the trivial group and the same product formula.
\end{proof}

For each \(v\), use \(a_{j,v},b_v\) and the residue classes of
Theorem~\ref{thm:residueorbits606}. Write \(S_{v,i}(\eta)\)
for the sum of their original phases with zero on unsupported
classes. In particular this is the sum over the explicit
residue classes themselves, including their multiplicities.
Define two ideals of \(F\) by
\[
        \mathfrak a_\eta=\prod_v\mathfrak p_v^{\sum_j a_{j,v}},
          \qquad \mathfrak c_\eta=\prod_v\mathfrak p_v^{b_v}.
\]
Both products use finitely many nonzero exponents.

\begin{theorem}[The exact finite factor with row covolumes]
\leavevmode\label{thm:rowfactor607}
For a supported full index of the given relevant rational
class set \(C_i(\eta)=\Delta_F^{r/2}m_i(\eta)/(c_{t,s}v_U)\).
With the original row volume \(c_{t,s}\) and the measures above,
\[
              C_i(\eta)=\prod_{v<\infty}q_v^{-r b_v}S_{v,i}(\eta).
\]
Almost all factors are one. For every complex embedding of
the coefficient field,
\[
      |\sigma C_i(\eta)|\le N_{F/\QQ}\mathfrak a_\eta,\qquad
      (N_{F/\QQ}\mathfrak c_\eta)^r C_i(\eta)
                       \ \text{is an algebraic integer}.
\]
Consequently, when \(C_i(\eta)\ne0\) and \(T\) is the terminal
rational prime on the interval, every valuation normalized by
\(\operatorname{val}_T(T)=1\) satisfies
\[
             \operatorname{val}_T C_i(\eta)
                       \ge -r\sum_{v\mid T} b_v .
\]
The last expression uses complete splitting of \(T\) in \(F\).
\end{theorem}
\begin{proof}
Write \(du_v,dd_v\) for the local measures on \(U_v,D_v\).
In the coordinates \(s(t)d\) of
Lemma~\ref{lem:orbitcoordinates606}, the quotient \(du_v/dd_v\)
is exactly additive product measure on \(t\). The lemma gives
proportionality. At the origin the derivative in the parameter
directions is \(X_1,\ldots,X_r\); evaluation of the two forms
on those vectors followed by a basis of the stabilizer gives
ratio one. Hence each residue class in that lemma's coordinates
has quotient volume \(q_v^{-rb_v}\) at the precision of
Theorem~\ref{thm:residueorbits606}. The integral \(L_{v,i}\) of
the signed finite test on \(U_v/D_v\) is therefore
\(q_v^{-r b_v}S_{v,i}(\eta)\).

Alternatively integrate on each \(K_v\)-orbit. The volume
calculation in Theorem~\ref{thm:residueorbits606} gives
\[
 \frac{L_{v,i}}{\operatorname{vol}_{du_v}(K_v)}
       =\sum_{\omega\in\Omega_v}
          \frac{z_v(\omega)}{\operatorname{vol}_{dd_v}(K^D_{v,\omega})}.
\]
The product of the sums enumerates exactly the integral lifts
at \(i\) by Theorem~\ref{thm:localphase605}. For each lift the
product of the phases is its full phase, and its real
stabilizer covolume is the reciprocal product of the displayed
stabilizer volumes times \(\Delta_F^{\dim_F D/2}\), by
Lemma~\ref{lem:unipotentmeasure607} and finite approximation
for \(D\). Thus
\[
       m_i(\eta)=c_{t,s}\Delta_F^{\dim_F D/2}
                          \frac{\prod_v L_{v,i}}{\prod_v\operatorname{vol}_{du_v}(K_v)} .
\]
Finite approximation for \(U\) gives
\(v_U\prod_v\operatorname{vol}_{du_v}(K_v)
=\Delta_F^{\dim_F U/2}\) by the same lemma. Substitution
proves the asserted formula. All real stabilizer measures
introduced here cancel against the associated real quotient
measures in \(m_i(\eta)I_\eta\).

Each conjugate of the finite sum \(S_{v,i}(\eta)\) has
absolute value at most the total number of residue classes,
\(q_v^{r b_v+\sum_j a_{j,v}}\), and each such sum is an
algebraic integer. Its denominator in the formula therefore
cancels the factor \(q_v^{r b_v}\) in the absolute bound.
Multiplication gives the two assertions about
\(\mathfrak a_\eta,\mathfrak c_\eta\). Only places over \(T\)
contribute to the \(T\)-valuation of the denominator. They
have \(q_v=T\), which proves the valuation inequality.
\end{proof}

If the original real Haar measures on the two row radicals
are given in a different convention, let their product be
\(\kappa\,du_\infty\). Use \(\kappa\,dd_\infty\) on the
stabilizer. The row covolume and \(m_i(\eta)\) both acquire
the factor \(\kappa\), so the same \(C_i(\eta)\) applies
with those given row measures. The quotient tensor \(I_\eta\)
acquires \(1/\kappa\), as the measure on \(G_\infty\) is fixed.

\begin{proposition}[The value ideals of the orbit coordinates]
\leavevmode\label{prop:valueideals607}
For each \(P_j\) let \(J_j\) be the fractional ideal of \(F\)
whose localization is generated by the values
\(P_j(\mathcal L_v)\), with \(\mathcal L_v\) the lattice of
integral row maps. Set \(J=\prod_{j=1}^r J_j\), taking the
unit ideal for \(r=0\). Then
\[
                 |\sigma C_i(\eta)|\le (N_{F/\QQ}J)^{-1}
\]
for every coefficient embedding. The local ideals can be
calculated on finite grids using the row lattices. For each
common compact summand \(\mu\), the ratio
\(|b_\mu(\eta)|/(N_{F/\QQ}J)\) is unchanged when the
parameters \(t_j\) are individually multiplied by nonzero
elements of \(F\) and the stabilizer form is adjusted as above.
\end{proposition}
\begin{proof}
Each coordinate polynomial is nonzero. On a full lattice a
nonzero polynomial has some nonzero value: multiples of a
rational basis in that lattice give a Cartesian product of
infinite sets, on which a nonzero polynomial cannot vanish
identically. Choose such a global integral input for each
polynomial. Its value is a unit away from finitely many
places. The coefficients on the row lattice are themselves
integral almost everywhere, so the stated local ideals
define nonzero fractional ideals.

For an explicit local calculation take a \(\ZZ_p\)-basis of
\(\mathcal L_v\) at residual characteristic \(p\). In the
resulting scalar variables \(P_j\) is a polynomial over \(F_v\);
write \(d_a\) for its degree in variable \(a\). Evaluate it
at each integer tuple with coordinates \(0,\ldots,d_a\).
The fractional \(\mathcal O_v\)-ideal of these finitely many
values is the whole value ideal. Indeed the Newton expansion
in products of \(\binom{x_a}{h_a}\) has as coefficients the
iterated forward differences of the polynomial at zero.
Repeated one-variable finite differences give this expansion
on integer inputs, hence as a polynomial identity. Each
coefficient is an integer combination of the sampled values.
For \(x_a\in\ZZ_p\) the binomial polynomials are integral by
continuity and density of nonnegative integers. This proves
one containment, and the sample belongs to the lattice for
the reverse containment.

By the coordinate lemma the supported parameters lie in
\(\prod_j (J_j)_v\). Its additive volume is the product of
the local inverse ideal norms, with the convention that
\(\pi_v^a\mathcal O_v\) has volume \(q_v^{-a}\) for any
integer \(a\). Every conjugate phase has modulus one, so
integrating the absolute value on this box bounds the local
factor \(L_{v,i}\) in the proof of Theorem~\ref{thm:rowfactor607}.
Taking the product proves the inequality.

For \(t'_j=\ell_j t_j\), the new directions are
\(X'_j=X_j/\ell_j\) and the ideals are \(J'_j=(\ell_j)J_j\).
Put \(L=|N_{F/\QQ}\prod_j\ell_j|\). The real measure on
\(D_\infty\) prescribed by the volume form scales by \(1/L\);
the quotient measure in each real tensor therefore scales by
\(L\). Thus \(b_\mu\) and \(N_{F/\QQ}J\) acquire precisely
the same positive factor, proving the invariance.
\end{proof}

\begin{corollary}[A signed formula in the original row normalization]
\leavevmode\label{cor:rowperiod607}
Fix the projected hyperplanes, the transport, and the marked
index as in Corollary~\ref{cor:signedslice604}. In the sum over
supported relevant rational classes let \(r_\eta=\dim_F U-\dim_F
D_\eta\), and integrate \(b_\mu(\eta)\) with the corresponding
stabilizer measure just specified. The complete contribution \(Z\)
of these cells then satisfies
\[
 \frac{\|Z\|_{\rm HS}^2}{|c_{t,s}|^2v_U^2}
   =\sum_\mu\frac1{d_\mu}
       \left|\sum_{[\eta]}\Delta_F^{-r_\eta/2}
                  b_\mu(\eta)\prod_{v<\infty}
                          q_v^{-r_\eta b_v}S_{v,i}(\eta)\right|^2 .
\]
It holds with any cutoff by frame height as well.
\end{corollary}
\begin{proof}
Insert Theorem~\ref{thm:rowfactor607} into
Corollary~\ref{cor:signedslice604} before taking the squared
norm. The row covolumes are fixed throughout the sum.
Each scalar times its real tensor is the original contribution
of its class; the absolute integrability established in that
corollary justifies summation with the indicated measures.
\end{proof}

At a fixed rational class the absolute value of its summand
inside the last square is consequently at most
\(\Delta_F^{-r_\eta/2}|b_\mu(\eta)|/(N_{F/\QQ}J)\).
Summing the signed contributions over the rational classes
and the different hyperplanes at the required theta scale,
and proving the cohomology estimate of Section~\ref{sec:missing},
remain open on the common intervals.


\subsection{Summing the signed period over different hyperplanes}
\label{sec:hyperplanesum610}

Fix \(x\) on an interval of Section~\ref{sec:missing}. Keep the
field \(E=E_x\), its real subfield \(F\) of degree \(f\), the two
original rows, the full marked index \(i\), and the rational
transport \(j:P_s\to P_t\) of Corollary~\ref{cor:signedslice604}.
Either terminal frequency is allowed at either end. Identify the two
positive spaces by \(j\) and write \(P=P_t\). A transverse cell
of common rank \(k-1\) now gives one rational hyperplane \(H\subset P\).
Choose one complex embedding of \(E\) above each real place \(\tau\)
of \(F\). Put \(K=\prod_\tau U(P_\tau)\) and write \(V_s,V_t\)
for the two irreducible column spaces, with their given unitary
metrics. Inner products are linear in the second argument.

Let \(Z_H:V_s\to V_t\) be the signed contribution at this \(H\),
before the marked projection. In the notation of
Corollaries~\ref{cor:signedslice604} and~\ref{cor:rowperiod607} it is
\begin{equation}
  Z_H=\sum_\mu \frac{J_{H,\mu}}{d_\mu}
                 \sum_{[\eta]\ \mathrm{at}\ H}m_i(\eta)b_{H,\mu}(\eta).
  \label{eq:hyperplaneslice610}
\end{equation}
Here \(J_{H,\mu}\) is the partial isometry between the two
copies of the common compact summand. The factor \(m_i(\eta)\)
includes the finite phases and the covolume of the real stabilizer
in the measures of Theorem~\ref{thm:rowfactor607}. We allow the
frame cutoff of Corollary~\ref{cor:oneindexperiod603} throughout.
It keeps whole rational classes, by Theorem~\ref{thm:liftphase604}.
We first compute the angles between the spaces containing the
different \(Z_H\).

\begin{lemma}[The matrix coefficient for two unit normals]
\leavevmode\label{lem:normalkernel610}
Let \(k\ge3\), let \(e\) be a unit vector in \(\CC^k\), and let
\(K_e=U(k-1)\) fix \(e\). An irreducible representation of \(U(k)\)
with a \(K_e\)-fixed vector has highest weight
\((a,0,\ldots,0,-b)\), where \(a,b\ge0\). Its fixed space is a line.
There are unit vectors \(\xi_{a,b}(v)\) on the corresponding fixed
lines for all unit \(v\), equivariant under \(U(k)\). If
\(z=\langle v,w\rangle\), \(d^2=1-|z|^2\), and \(q=\min(a,b)\), then
\begin{equation}
 \begin{split}
 \langle \xi_{a,b}(v),\xi_{a,b}(w)\rangle
   &=\Phi_{a,b}(z)=\sum_{m=0}^q c_m z^{a-m}\bar z^{b-m}d^{2m},\\
 c_0&=1,\qquad
 c_{m+1}=-\frac{(a-m)(b-m)}{(m+1)(k-1+m)}c_m .
 \end{split}
 \label{eq:normalkernel610}
\end{equation}
Put \(p=k-2\) and \(h=|a-b|\). At \(d>0\),
\begin{equation}
 \begin{aligned}
 |\Phi_{a,b}(z)|&\le \min\{1,\,2^{h/2}p!(q d^2)^{-p}\}
                      &&(q\ge1),\\
 |\Phi_{a,b}(z)|&=(1-d^2)^{h/2}
                      &&(q=0,\ h>0).
 \end{aligned}
 \label{eq:normaldecay610}
\end{equation}
For \(a=b=0\) the coefficient is one everywhere.
\end{lemma}
\begin{proof}
The unitary branching rule used in Theorem~\ref{thm:common133}
gives multiplicity one for each interlacing weight. Interlacing
with the zero weight of \(U(k-1)\) forces the asserted highest
weight and permits every pair \(a,b\).

To fix the phases take the trace-free part of
\(\operatorname{Sym}^a\CC^k\otimes\operatorname{Sym}^b\overline{\CC^k}\).
Here the trace contracts one vector with one conjugate vector.
When both degrees are positive the contraction is surjective.
Indeed, on identifying each tensor \(s\) with the polynomial
\(\langle s,T(v)\rangle\), where
\(T(v)=v^{\otimes a}\otimes\bar v^{\otimes b}\), its adjoint is,
up to a nonzero positive factor, multiplication by \(\|v\|^2\),
which is injective. This polynomial identification follows by
expanding in monomials in \(v,\bar v\), which are independent as
polynomials in the real and imaginary coordinates. The trace-free
space contains the highest vector \(e_1^{\otimes a}\otimes
\bar e_k^{\otimes b}\). Its dimension, obtained by subtracting
the dimension in degrees \(a-1,b-1\) when both are positive, is
\[
 \frac{a+b+k-1}{k-1}
       \binom{a+k-2}{k-2}\binom{b+k-2}{k-2}.
\]
This is the Weyl dimension of the asserted highest weight, so
the trace-free space is exactly that irreducible representation.
The same dimension computation with zero subtracted applies when
one of the degrees is zero.

Let \(\Pi\) be the orthogonal projection onto this space. The
\(T(v)\) for unit \(v\) span the tensor space by the polynomial
identification and homogeneity. Thus \(\Pi T(e)\ne0\), since
the unitary group is transitive on unit vectors. Set
\(\xi_{a,b}(v)=\Pi T(v)/\|\Pi T(e)\|\) on the unit sphere.
For \(v=e\), the coefficient against \(T(w)\), normalized to
take value one at \(e\), is a harmonic polynomial in \(w\) of
bidegree \((a,b)\). Applying the complex Laplacian
\(\sum_j\partial_{w_j}\partial_{\bar w_j}\) contracts one
index of each kind, so trace-freeness gives harmonicity. Put
\(t=\langle e,w\rangle\) and \(u=\|w-te\|^2\). Invariance under
\(U(e^\perp)\) forces the polynomial to be a sum of the monomials
\(t^{a-m}\bar t^{b-m}u^m\). Indeed the scalar circle on
\(e^\perp\) forces equal perpendicular degrees and transitivity on
its unit sphere gives a multiple of \(u^m\) in those degrees.
The perpendicular Laplacian sends \(u^{m+1}\) to
\((m+1)(k-1+m)u^m\). Equating coefficients after applying the
full Laplacian gives the recurrence, and evaluation at \(e\)
gives \(c_0=1\). Restriction to the sphere and equivariance give
\eqref{eq:normalkernel610} for every pair. These are the disk
polynomials from the unitary sphere construction of
Koornwinder~\cite{XXII-XVII-koornwinder610}.

For the estimate let \(\zeta\) be the first coordinate of a
uniform unit vector in \(\CC^{k-1}\). With \(p=k-2\), the variable
\(Y=1-|\zeta|^2\) has density \(p y^{p-1}\) on \([0,1]\).
To see this, normalize \(k-1\) independent standard complex
Gaussians to unit length. Their squared lengths are independent
unit exponentials. Conditional on the sum \(S\) of the last \(p\)
squared lengths, the probability that the first squared length
exceeds \(tS/(1-t)\) is \(\exp(-tS/(1-t))\). Integration gives
\(\Pr(|\zeta|^2\ge t)=(1-t)^p\) for \(0<t<1\).
It follows that \(\mathbb E|\zeta|^{2m}=m!/(k-1)_m\), where
\((s)_m=s(s+1)\cdots(s+m-1)\) and \((s)_0=1\). Expansion and
rotation invariance of \(\zeta\) now give
\[
  \Phi_{a,b}(z)=\mathbb E\big[(z+\mathrm i d\zeta)^a
                          (\bar z+\mathrm i d\bar\zeta)^b\big].
\]
Indeed only equal powers of \(\zeta\) and \(\bar\zeta\) survive,
and their coefficients are those of \eqref{eq:normalkernel610}.
The product of the absolute values of the two bases is
\(\lvert z+\mathrm i d\zeta\rvert\lvert z-\mathrm i d\zeta\rvert
\le |z|^2+d^2|\zeta|^2=1-d^2Y\), by the arithmetic--geometric
mean inequality on the squared moduli. Each base has modulus at
most \(|z|+d\le\sqrt 2\). For \(q\ge1\) the expectation is
therefore bounded in absolute value by
\[
  2^{h/2}p\int_0^1 e^{-q d^2y}y^{p-1}\,dy
       \le 2^{h/2}p!(q d^2)^{-p}.
\]
The coefficient of two unit vectors also has modulus at most one.
For \(q=0\) the sum in \eqref{eq:normalkernel610} has just one term.
\end{proof}

\begin{theorem}[The signed sum over all projected hyperplanes]
\leavevmode\label{thm:hyperplanesum610}
Give \(\mathcal V=\operatorname{Hom}(V_s,V_t)\) its
Hilbert--Schmidt norm and the action
\(A\mapsto\rho_t(g)A\rho_s(g)^{-1}\) of \(K\), using \(j\) on
the source. In an orthogonal isotypic decomposition
\(\mathcal V=\bigoplus_\lambda \mathcal H_\lambda\otimes M_\lambda\),
write \(\Pi_\lambda\) for the projection. Only the types
\(\lambda=((a_\tau,0,\ldots,0,-b_\tau))_\tau\) can occur in
any \(Z_H\). Choose a unit normal \(n_{H,\tau}\) at every place.
There is a unique \(b_{H,\lambda}\in M_\lambda\) such that
\[
  \Pi_\lambda Z_H=\xi_\lambda(n_H)\otimes b_{H,\lambda},\qquad
  \xi_\lambda(n_H)=\bigotimes_\tau
                       \xi_{a_\tau,b_\tau}(n_{H,\tau}).
\]
In particular \(b_{H,\lambda}\) is given by contracting
\(\Pi_\lambda Z_H\) from \eqref{eq:hyperplaneslice610} with the
indicated unit vector in its first factor. Then
\begin{equation}
 \begin{split}
  \left\|\sum_H Z_H\right\|_{\rm HS}^2
   &=\sum_{\lambda}\sum_{H,H'}K_\lambda(H,H')
                      \langle b_{H,\lambda},b_{H',\lambda}\rangle,\\
  K_\lambda(H,H')&=\prod_\tau
       \Phi_{a_\tau,b_\tau}(\langle n_{H,\tau},n_{H',\tau}\rangle).
 \end{split}
 \label{eq:sumallplanes610}
\end{equation}
This holds for every subset of hyperplanes and for the full
absolutely convergent transverse sum at \(i,j\).
\end{theorem}
\begin{proof}
The proof of Corollary~\ref{cor:signedslice604} gives
equivariance of each class integral under \(\prod_\tau U(H_\tau)\),
which fixes the normal vectors pointwise. After transporting the
source by \(j\), this says precisely that \(Z_H\) is fixed by this
subgroup in \(\mathcal V\). An isotypic projection commutes with
the subgroup. Its fixed subspace in each irreducible tensor factor
is the product of the fixed lines of Lemma~\ref{lem:normalkernel610},
or is zero. The multiplicity space has trivial action. This proves
the asserted factorization, even for an irreducible type of
multiplicity greater than one. Expand the squared norm, using
orthogonality of the isotypic spaces and the inner product of
pure tensors. Formula~\eqref{eq:normalkernel610} gives every entry
of \eqref{eq:sumallplanes610}. Replacing \(n_{H,\tau}\) by
\(u n_{H,\tau}\), \(|u|=1\), multiplies its fixed vector by
\(u^{a_\tau-b_\tau}\) and the multiplicity vector by the inverse
phase, leaving each summand unchanged. Absolute summability of
the original tensors was proved in Corollary~\ref{cor:signedslice604}.
It implies \(\sum_H\|Z_H\|_{\rm HS}<\infty\), justifying the
double sum at fixed \(x\).

In the trivial type the contribution is
\(\|\sum_H b_{H,0}\|^2\). For irreducible columns this type occurs
exactly when the transported representations are equivalent.
Indeed its vectors are \(K\)-intertwiners. If \(U:V_s\to V_t\)
is a unitary intertwiner, Schur's lemma gives the unit basis
\(U/\sqrt{\dim V_s}\) for that line. Its contribution is then
\(\big|\sum_H\operatorname{tr}(U^*Z_H)\big|^2/\dim V_s\).
\end{proof}

For a rational normal line \(\ell\) in \(P\), let
\(M_\ell=\widehat L_t\cap\ell\), where \(\widehat L_t\) is the
self-dual target lattice of Lemma~\ref{lem:completedcolumn535}.
Set \(B_\ell=[M_\ell^\sharp:M_\ell]^{1/(2f)}\), taking the dual
inside \(\ell\). For two lines let
\(d_\tau(\ell,\ell')=(1-|\langle n_\tau,n'_\tau\rangle|^2)^{1/2}\)
with unit vectors in their respective fibres.

\begin{proposition}[Separation of rational normals in the row lattice]
\leavevmode\label{prop:normalseparation610}
For a cell with common lattice \(R\) projecting to \(H=\ell^\perp\)
one has \(1\le B_\ell\le h_R\). For distinct lines,
\[
     \prod_\tau d_\tau(\ell,\ell')^2\ \ge\ (B_\ell B_{\ell'})^{-f}.
\]
Put \(D=2(k-1)\). The number of lines with \(B_\ell\le B\), for
\(B\ge1\), is at most \((2B)^{Df}\). In a product of closed
projective balls of radii \(0<R_\tau\le1\) with
\(R_*=\prod_\tau R_\tau\ge B^{-f}\), their number is at most
\[
                   3^{Df} B^{Df} R_*^D .
\]
\end{proposition}
\begin{proof}
Put \(L=\widehat L_t\), \(S=L\cap H\), and \(M=L\cap\ell\).
The orthogonal projections of \(L\) onto \(H,\ell\) have
images \(S^\sharp,M^\sharp\). In fact \(S,M\) are saturated,
so an integral functional extends locally to \(L\) and is represented
there by self-duality. The same argument uses the paired factors
at a split place. The projections both induce isomorphisms from
\(L/(S+M)\) onto the respective discriminant modules. If the
projection lies in \(S\) or \(M\), one subtracts it to land in
the other lattice, giving injectivity. The image of \(R\) lies
in \(S\) with finite index by Theorem~\ref{thm:mixedtransport535}.
Changing lattice by this index multiplies the discriminant order
by \([S:R]^2\), as follows by taking the index in the lattice and
in its conjugate dual. Consequently
\([R^\sharp:R]=[S:R]^2[M^\sharp:M]\), which proves the height claim.

Write \(M_\ell=Iv\) and \(M_{\ell'}=Jw\) for fractional ideals
of \(E\). For the Hermitian form put \(s=\langle v,v\rangle\),
\(t=\langle w,w\rangle\), \(u=\langle v,w\rangle\).
The ideals \(I\bar I(s)\) and \(J\bar J(t)\) have absolute norms
\(B_\ell^{2f}\) and \(B_{\ell'}^{2f}\), respectively. For
example, the dual of \(Iv\) is \(s^{-1}\bar I^{-1}v\).
Integrality of \(L\) gives \(\bar I J(u)\subseteq\mathcal O_E\)
and its conjugate inclusion. It follows that the ideal
\[
                I\bar I J\bar J(st-u\bar u)
\]
is integral. The generator in parentheses is totally positive,
since two distinct rational lines remain independent under every
embedding. The absolute norm of this ideal is at least one.
Dividing it by the norms of the two one-line ideals and taking
a square root gives the product bound, since at \(\tau\) the
ratio of Gram determinants is \(d_\tau^2\).

The projective distance \(d_\tau\) is the operator norm of the
difference of the two orthogonal line projections. This follows
by diagonalizing their difference on the span of the lines,
where the two eigenvalues for distinct lines are \(\pm d_\tau\).
In particular it satisfies the triangle inequality. The invariant
probability of a projective ball of radius \(r\le1\) is \(r^D\).
Indeed the Gaussian argument of Lemma~\ref{lem:normalkernel610}
in complex dimension \(k\) gives
\(\Pr(1-|\langle e,n\rangle|^2\le r^2)=r^{2(k-1)}\).
Distinct lines of heights at most \(B\) satisfy
\(\prod_\tau d_\tau\ge B^{-f}\). The open products of balls of
common radius \(1/(2B)\) about these lines are disjoint, proving
the first count by product probability. For the second count take
the smaller radii \(r_\tau=R_\tau/(2B R_*^{1/f})\) about each
line in the indicated product. They satisfy \(r_\tau\le R_\tau/2\)
and \(\prod_\tau(2r_\tau)=B^{-f}\), so the products of open balls
are again disjoint. Their union lies in the product with radii
\(R_\tau+r_\tau\) about the original centre. Dividing its volume
bound by the volume of one small product gives
\((1+2B R_*^{1/f})^{Df}\le 3^{Df}B^{Df}R_*^D\).
The volume bound for an enlarged radius greater than one follows
from total mass one. These arguments apply to every finite set
and hence to all the lines.
\end{proof}

\begin{corollary}[Exact cross terms at height one]
\leavevmode\label{cor:heightoneplanes610}
At \(i,j\) there are at most \(k\) distinct hyperplanes with
\(B_{H^\perp}=1\). In \eqref{eq:sumallplanes610}, on any set
of these hyperplanes the squared norm in type \(\lambda\) equals
\[
   (1-\kappa_\lambda)\sum_H\|b_{H,\lambda}\|^2+
                    \kappa_\lambda\left\|\sum_H b_{H,\lambda}\right\|^2 .
\]
Here \(\kappa_\lambda=0\) if some \(a_\tau\ne b_\tau\). Otherwise
\[
   \kappa_\lambda=\prod_\tau
          \frac{(-1)^{a_\tau}}{\binom{k-2+a_\tau}{a_\tau}} .
\]
In particular, if the two columns have different scalar central
characters at any real place, the squared norm of this entire
height-one sum is \(\sum_H\|Z_H\|_{\rm HS}^2\).
\end{corollary}
\begin{proof}
Proposition~\ref{prop:normalseparation610} forces the product
of the squared distances of two distinct height-one lines to be
at least one. Each squared distance is at most one, so the normals
are orthogonal at every place. Any one place gives the cardinality
bound. At \(z=0\) the polynomial in \eqref{eq:normalkernel610}
vanishes for unequal degrees. For \(a=b\) only \(m=a\) survives,
giving \((-1)^a a!/(k-1)_a\). The diagonal of the kernel is one
and every off-diagonal entry is the same real \(\kappa_\lambda\).
Expanding the two squares proves the formula, including the
constant type. Finally, at each place the scalar \(uI\) acts on
\(\mathcal V\) by the target central character divided by the source
central character. On a type \((a_\tau,0,\ldots,-b_\tau)\) its
value is \(u^{a_\tau-b_\tau}\). Different central characters
therefore force unequal degrees in every occurring type at that
place. Summing the resulting equalities proves the last assertion.
\end{proof}

\begin{theorem}[An error bound at arbitrary bounded height]
\leavevmode\label{thm:angularerror610}
For a set of \(N\ge1\) hyperplanes of normal heights at most \(B\),
put \(p=k-2\), \(D=2(k-1)\), and define, for \(q=\min(a,b)\)
and \(h=|a-b|\),
\[
 \omega_{a,b}=
 \begin{cases}
  \min\{1,\,4^p2^{h/2}p!/q^p\},&q\ge1,\\
  \min\{1,\,(8p/(\mathrm e h))^p\},&q=0,\ h>0,\\
  1,&a=b=0 .
 \end{cases}
\]
For every occurring type set
\[
   D_\lambda(B,N)=\min\left\{N-1,\,
     (4/3)^f(3B)^{Df}\prod_\tau\omega_{a_\tau,b_\tau}\right\}.
\]
With \(S_\lambda=\sum_H\|b_{H,\lambda}\|^2\) one has
\[
 \left|\left\|\sum_H\Pi_\lambda Z_H\right\|_{\rm HS}^2-S_\lambda\right|
              \le D_\lambda(B,N)S_\lambda .
\]
\end{theorem}
\begin{proof}
Fix a line from the set as centre. For each other line all its
distances from the centre are positive. Assign the unique
integers \(n_\tau\ge0\) with
\(2^{-n_\tau-1}<d_\tau\le2^{-n_\tau}\).
The product of the upper radii is at least \(B^{-f}\) for an
occupied box by Proposition~\ref{prop:normalseparation610}.
That proposition bounds the number in the box by
\((3B)^{Df}2^{-D\sum_\tau n_\tau}\).
At each place the modulus of the kernel there is bounded by
\(\omega_{a,b}2^{2p n_\tau}\). For \(q\ge1\) this follows
from \eqref{eq:normaldecay610} and the bound one. For \(q=0\),
\(h>0\), the modulus is at most \(\exp(-h d_\tau^2/2)\).
Use \(y^p e^{-y}\le(p/\mathrm e)^p\) at \(y=h d_\tau^2/2\)
and the same lower bound on \(d_\tau\). The case \(a=b=0\)
is immediate. Thus the sum of \(|K_\lambda|\) off the diagonal
in this box is at most
\((3B)^{Df}(\prod_\tau\omega_{a_\tau,b_\tau})4^{-\sum_\tau n_\tau}\),
since \(D-2p=2\). Summing over all nonnegative indices gives
the factor \((4/3)^f\). The row sum is also at most \(N-1\).
These bounds hold at every centre. Bound the off-diagonal terms
in \eqref{eq:sumallplanes610} using Cauchy--Schwarz in
\(M_\lambda\) and \(2uv\le u^2+v^2\). Symmetry of the entry
moduli yields the maximum off-diagonal row sum times \(S_\lambda\),
as asserted.
\end{proof}

To apply the theorem to the projected period in
Corollary~\ref{cor:oneindexperiod603}, retain its frame cutoff and
take \(B=H_L\) from that corollary. Its proof gives \(h_R<H_L\)
for every retained cell, so Proposition~\ref{prop:normalseparation610}
applies at both terminal frequencies. For each occupied \(i,j\)
write \(N_{i,j}\) for its number of retained hyperplanes and
\(S_{i,j,\lambda}\) for the sum in Theorem~\ref{thm:angularerror610}.
If \(E_L\) denotes the omitted-cell bound (the second term of
Corollary~\ref{cor:oneindexperiod603}), the original normalized
marked contribution satisfies
\[
 \|\mathcal P\|_{\rm HS}\le E_L+
  \left(\alpha\,n(H_L)
    \sum_{i,j,\lambda}\bigl(1+D_\lambda(H_L,N_{i,j})\bigr)
                              S_{i,j,\lambda}\right)^{1/2}.
\]
Indeed the preceding theorem bounds the squared norm at each
\(i,j\) before the marked projection. The operator bound for
that projection is precisely the one in the cited corollary,
which also permits \(n(H_L)=1\) when \(H_L\le Q/2\).
All norms here use the actual signed class coefficients in
\eqref{eq:hyperplaneslice610}. The constant type can still
contain the full squared sum of their traces. Bounding those
traces, the remaining class integrals and the omitted cells at
the scale required in Section~\ref{sec:missing} is open, as is
the required uniform estimate for the residual cohomology on
the same intervals.


\subsection{The determinant of a pair of frames and its signed period}
\label{sec:frameorientation611}

Keep the two rows, the rational transport \(j\), the marked index
\(i\) and the field \(E=E_x\) of Section~\ref{sec:hyperplanesum610}.
Fix one of the hyperplanes in that section and write \(H_b\subset
P_b\) for it at end \(b=s,t\). Put \(F=E^+\), \(f=[F:\QQ]\),
\(\ell_b=P_b\cap H_b^\perp\), and \(d=l+1\), where \(l\) is the
number of outer isotropic vectors. Use the integral Witt splitting
of Lemma~\ref{lem:projectivewitt554} at each end. Denote the
ordered frame by \(y_b=(e_{b,1},\ldots,e_{b,l},z_b)\), its span by
\(Y_b\), and its exterior product by \(\Phi_b\). In the smaller
spaces \(V'_b=H_b^\perp\) of Theorem~\ref{thm:smallcomplement604},
a rational cell has a representative \(\eta\) with
\(V'_t=Y_t\oplus\eta Y_s\). We extend its map \(V'_s\to V'_t\)
to the whole space using \(j|_{H_s}\). Its determinant is fixed by
the determinant of \(j\).

Write \(R_b(\alpha)\) for multiplication by \(\alpha\) on \(\ell_b\)
and the identity on its perpendicular, for
\(\alpha\in E^1=\{\alpha\in E^\times:\alpha\bar\alpha=1\}\).
Let \(C(\eta)\) be the \(d\)-by-\(d\) matrix with entries
\(\langle y_{t,u},\eta y_{s,v}\rangle\), with the form linear
in its second variable. Its class under lower unitriangular row
operations and upper unitriangular column operations over \(E\)
will be denoted by \([C]\).

\begin{theorem}[The rational classes above one cross matrix]
\leavevmode\label{thm:frameorientation611}
Fix an occurring \([C]\), and choose one rational representative
\(\eta_0\) above it. Every rational unipotent double class above
\([C]\) at the given hyperplanes and transport has exactly one
representative of the form
\[
       \eta_\alpha=R_t(\alpha)\eta_0R_s(\alpha)^{-1},
                  \qquad \alpha\in E^1 .
\]
More precisely, the indicated representatives form a system of
representatives without repetition. The parameter of any class
is read from its top exterior vector
\[
  \Omega_\eta=\Phi_t\wedge\eta\Phi_s\in\det V'_t,\qquad
             \Omega_\eta=\alpha\Omega_{\eta_0}.
\]
Here the representative used to form \(\Omega_\eta\) lies in the
smaller spaces as above.
\end{theorem}
\begin{proof}
On \(Y_b\) the reduced unipotent radical \(N'_b\) gives all upper
unitriangular changes. An upper unitriangular change on the outer
space extends by its inverse adjoint on the opposite space.
To send \(z_b\) to \(z_b+v\), \(v\in X_b\), use the unipotent
matrix \(n(A)\) in the proof of Theorem~\ref{thm:smallcomplement604}
with \(A z_b=v\), vanishing on the positive middle space. Both
maps fix \(H_b\) pointwise. Together they realize every required
change. Acting by \(N'_t,N'_s\) thus gives exactly the indicated
unitriangular equivalence on cross matrices. It preserves
\(\Omega_\eta\), since the frame wedges are fixed at their
respective ends and the target unipotent has determinant one
on \(V'_t\).

Take two classes over \([C]\). By the preceding frame changes
choose representatives \(\eta,\eta'\) with identical cross
matrices. The ordered bases \((y_t,\eta y_s)\) and
\((y_t,\eta' y_s)\) of \(V'_t\) have identical Gram matrices.
Their diagonal blocks are the two fixed frame Gram matrices.
The rational linear map \(g_t\) taking the first basis to the
second is therefore unitary and fixes \(y_t\) pointwise.
Then \(g_s=\eta'^{-1}g_t\eta\) fixes \(y_s\) pointwise.
The determinants of \(g_t,g_s\) agree because the two class
maps have the same determinant. At either end a unitary map
stabilizing the ordered frame up to an upper unitriangular
change acts on the middle quotient by the identity on \(z_b\)
and by a norm-one scalar on \(\ell_b\). That scalar is its
determinant, since the two outer diagonal blocks have determinant
one. After division by \(R_b\) of that scalar the map lies in
\(N'_b\), by its trivial action on the middle quotient and its
unitriangular action on the outer flag. It follows that the two
classes are related by the asserted simultaneous rotations.

Conversely the rotations normalize the two radicals, fix each
frame pointwise, preserve the transport on \(H_s\) and leave the
determinant of the class map fixed. They preserve the spanning
condition. The top exterior vector changes by
\(\det(R_t(\alpha)|_{V'_t})=\alpha\). Since it is nonzero, distinct
parameters cannot belong to the same class. The rational bijection
of Theorem~\ref{thm:smallcomplement604} now gives the assertion
for the original rational double classes.
\end{proof}

Let \(Q=g_1g_2\) be the product of the two marked primes and
\(\mathfrak h_d\) the product of the prescribed outer ideals.
Thus the integral frame span at either end has determinant lattice
\(\mathfrak h_d\Phi_b\). Let \(\mathcal G>0\) at all real places
be the transverse Gram invariant of
Proposition~\ref{prop:transverseheight582}. Define
\[
       X_C=\frac{(N_{E/\QQ}\mathfrak h_d)^2
                    |N_{F/\QQ}\mathcal G|}{Q^{2df}} .
\]
It depends only on \([C]\) and the given rows. In what follows
choose \(\eta_0\) from a class having a supported integral lift.
Let \(L'=\operatorname{pr}_{V'_t}\Lambda_t\) be the orthogonal
projection of the target row lattice. There is a fractional ideal
\(\mathfrak a\) of \(E\) specified by the equality of lattices
in the determinant line
\[
   Q^{-d}\mathfrak h_d^2\,\Omega_{\eta_0}
                            =\mathfrak a\,\det L'.
\]

\begin{proposition}[An integral ideal bounding the possible rotations]
\leavevmode\label{prop:rotationideal611}
The ideal \(\mathfrak a\) is integral. For every \(\alpha\) with
a supported integral lift at \(i\), the ideal
\((\alpha)\mathfrak a\) is also integral. With
\(A=\Lambda_t\cap H_t\) and \(\delta(A)=[A^\sharp:A]^{1/2}\)
one has
\[
             N_{E/\QQ}\mathfrak a=X_C\,\delta(A)\le X_C^2 .
\]
The number of such \(\alpha\) is at most
\begin{equation}
   |\mu(E)|\prod_{\{w,\bar w\},\ w\ne\bar w}
       \bigl(1+v_w(\mathfrak a)+v_{\bar w}(\mathfrak a)\bigr)
       \ \le\ |\mu(E)|\,C_\epsilon^{\,f}X_C^{2\epsilon}
       \qquad(\epsilon>0),
 \label{eq:rotationcount611}
\end{equation}
where the product counts each split pair once and \(C_\epsilon\)
is independent of the field, rows and interval. At the frame
cutoff \(\mathcal H_\gamma<Q^f L\) one may use
\(X_C\le(T_0T_1)^{df}L^{2d}\), where \(T_0,T_1\) are the
absolute terminal norms. If \(X_C=1\), only roots of unity can
occur.
\end{proposition}
\begin{proof}
Take an actual integral lift \(\gamma=u_t\eta_\alpha u_s^{-1}\),
with \(u_b\in N_b(F)\) in the full row radicals. Its frame wedge
is \((\bigwedge^{2d}u_t)\Omega_{\eta_\alpha}\), since each radical
fixes its own frame wedge. Its projection to \(\det V'_t\) is
exactly \(\Omega_{\eta_\alpha}\). To verify the last assertion,
in the Witt decomposition order the outer, middle and opposite
blocks consecutively. The map \(\operatorname{pr}_{V'_t}u_t|_{V'_t}\)
is block upper triangular with diagonal the outer unitriangular
matrix, the identity on \(E z_t+\ell_t\), and the inverse adjoint
on the opposite block. Projection removes only the \(H_t\)
component of the middle block. Its determinant is one.

Multiplying the original integral pair of frame wedges by
\(\mathfrak h_d^2\) puts it in \(\bigwedge^{2d}\Lambda_t\).
At each factor above \(Q\), the two reduced frames agree up to an
upper unitriangular change. Their \(d\) corresponding differences
are divisible by the uniformizer after an integral lift of that
change. In the concatenated exterior product replace the second
\(d\) columns by those differences. This proves divisibility by
the \(d\)-th power at each marked prime, including both halves
of each split place. The outer ideals are units there. Thus
the same original wedge, multiplied by \(Q^{-d}\mathfrak h_d^2\),
is still integral everywhere. Projection gives the inclusion in
\(\det L'\) and proves integrality of \((\alpha)\mathfrak a\),
including \(\mathfrak a\) at \(\alpha=1\).

For a lattice in a nondegenerate Hermitian space use its
determinant covolume \(\delta\). On the determinant line a
lattice \(Iv\) has covolume \(N_{E/\QQ}I\) times
\(|N_{F/\QQ}\langle v,v\rangle|\). Local Gram determinants
give \([M^\sharp:M]^{1/2}\) for an integral lattice \(M\)
and reciprocal covolumes for dual lattices. The two primitive
perpendicular lattices \(A\) and \(B=\Lambda_t\cap V'_t\)
have equal discriminant orders. Indeed the projections of the
self-dual row identify \(\Lambda_t/(A+B)\) with each discriminant
module, by the same extension of functionals as in
Proposition~\ref{prop:normalseparation610}. In particular
\(L'=B^\sharp\) and \(\delta(L')=1/\delta(A)\).
The Gram determinant of the basis \((y_t,\eta_0y_s)\) is
\((-1)^d\mathcal G\), by Theorem~\ref{thm:cross132}.
The covolume of \(Q^{-d}\mathfrak h_d^2\Omega_{\eta_0}\)
is therefore \(X_C\). Division by the covolume of \(\det L'\)
gives \(N_{E/\QQ}\mathfrak a=X_C\delta(A)\).

For any supported lift its common integral lattice \(R\)
projects isometrically onto a full sublattice of \(A\).
Indeed it lies in \(X_t^\perp\cap z_t^\perp\), and its projection
along \(X_t\) has values in the integral middle lattice of our
Witt splitting and in \(H_t\). Hence
\(\delta(A)\le\delta(R)=h_R^f\). Corollary~\ref{cor:actualgram535}
gives \(X_C=I_\Lambda h_R^f\), with saturation index
\(I_\Lambda\ge1\). This proves the upper bound and \(X_C\ge1\).
The frame estimate of Proposition~\ref{prop:transverseheight582},
whose additional factor \(P_\gamma\) is at least one, gives the
stated cutoff bound.

Since \(\alpha\bar\alpha=1\), at a self-conjugate prime its
valuation is zero. At two conjugate primes it is determined by
one integer between \(-v_w(\mathfrak a)\) and
\(v_{\bar w}(\mathfrak a)\). Two parameters with the same ideal
differ by a unit all of whose complex conjugates have absolute
value one. Such a unit is a root of unity. Indeed the powers
of the unit have bounded integral characteristic-polynomial
coefficients in degree \(2f\), so only finitely many values occur.
This proves the first bound. For a split pair of residue size
\(q_v\) put \(m=v_w(\mathfrak a)+v_{\bar w}(\mathfrak a)\).
If \(q_v\ge 2^{1/\epsilon}\), then \(1+m\le q_v^{\epsilon m}\).
At any place over a rational prime \(p<2^{1/\epsilon}\) use
\((1+m)q_v^{-\epsilon m}\le(1-p^{-\epsilon})^{-2}\), by summing
the corresponding geometric derivative series. There are at most
\(f\) such places over each \(p\). Taking
\(C_\epsilon=\prod_{p<2^{1/\epsilon}}(1-p^{-\epsilon})^{-2}\)
proves the second bound. If \(X_C=1\) the integral ideal has norm
one and every supported parameter is a norm-one unit.
\end{proof}

\begin{theorem}[One real tensor for an entire cross-matrix class]
\leavevmode\label{thm:rotationperiod611}
Keep a relevant \(\eta_0\) and choose a rational invariant Haar
form on \(D_0=N_t\cap\eta_0N_s\eta_0^{-1}\). At \(\eta_\alpha\)
transport it by conjugation by \(R_t(\alpha)\). Use the
corresponding quotient measures in the real tensors
\(I_{\eta_\alpha}\) of Theorem~\ref{thm:liftphase604}. If
\(\omega_s,\omega_t\) are the scalar central characters of the
two compact column representations at infinity, put
\(\chi_\infty(\alpha)=\omega_t(\alpha_\infty)
                         \overline{\omega_s(\alpha_\infty)}\).
The full signed contribution of the classes over \([C]\) at
the indicated \(H,j,i\) is
\begin{equation}
      I_{\eta_0}\sum_{\alpha\in E^1}
                    \chi_\infty(\alpha)m_i(\eta_\alpha).
 \label{eq:rotationperiod611}
\end{equation}
The coefficients use the transported stabilizer measures and are
zero for a class without a supported lift. The sum is finite by
Proposition~\ref{prop:rotationideal611}. In particular for
equivalent columns and a unitary \(K\)-intertwiner \(U_j\), the
trace against \(U_j\) of this contribution is
\[
       \operatorname{tr}(U_j^*I_{\eta_0})\sum_{\alpha\in E^1}
                                                 m_i(\eta_\alpha).
\]
\end{theorem}
\begin{proof}
Each rotation fixes the ordered frame and normalizes its radical.
It preserves the Bessel character, which uses the outer simple
coordinates and the pairing against the terminal vector. It
therefore carries the stabilizer to that for the rotated pair
and preserves relevance. In the real quotient integral substitute
\(g=R_t(\alpha)g'\). The Haar form on the stabilizer was transported,
and the Haar measure on the ambient group is left invariant.
Equivariance of the two original sections under the compact
column groups gives
\[
 I_{\eta_\alpha}
       =\rho_t(R_t(\alpha)) I_{\eta_0}\rho_s(R_s(\alpha))^{-1}.
\]
At every place \(R_b(\alpha)|_{P_b}\) is the scalar \(\alpha I\)
times the element acting by \(\alpha^{-1}\) on \(H_b\) and by one
on its normal line. The latter elements correspond under \(j\).
The map \(I_{\eta_0}\) intertwines these actions of \(U(H_b)\)
by Corollary~\ref{cor:signedslice604}. Thus the displayed tensor
is exactly \(\chi_\infty(\alpha)I_{\eta_0}\). Multiplication by
its own signed lift coefficient and summation give
\eqref{eq:rotationperiod611}. For equivalent representations
the central characters agree, giving the trace formula.
\end{proof}

The finite factors in \eqref{eq:rotationperiod611} can all be
computed with one rational parameter space. Use the slice
\(s(t)\) and directions \(X_j\) of Lemma~\ref{lem:orbitcoordinates606}
at \(\eta_0\). At \(\eta_\alpha\) conjugate the target and source
radicals separately by \(R_t(\alpha),R_s(\alpha)\), respectively,
and transport \(s(t)\) and \(X_j\) by this automorphism of
\(N_t\times N_s\). The determinant of its adjoint action on each
row radical is one. Indeed it is an \(F\)-rational character of
the norm-one torus, whose characters over \(E\) are \(u\mapsto u^n\).
Conjugation sends \(n\) to \(-n\), so the only \(F\)-rational
one has \(n=0\). Thus the stabilizer form constructed from the
row forms and these directions in Theorem~\ref{thm:rowfactor607}
is precisely the transported form if that construction is used
at \(\eta_0\). In the same additive parameters one tests the
original levels and phases on the rotated orbit point
\(R_t(\alpha)(s(t)\eta_0)R_s(\alpha)^{-1}\) and on the
correspondingly transported pair of unipotents. The notation
\(s(t)\eta_0\) here denotes the two-sided action on a row map.
Consequently \(m_i(\eta_\alpha)\) is given by the residue formula
of that theorem with the same factor
\(c_{t,s}v_U\Delta_F^{-r/2}\) throughout the sum, where
\(r=\dim_F(N_t\times N_s)-\dim_FD_0\).

Every frame cutoff from Corollary~\ref{cor:oneindexperiod603}
keeps or drops an entire \([C]\), since its leading minors are
unchanged by the two triangular operations. Summing
\eqref{eq:rotationperiod611} over the retained cross matrices
therefore gives the actual \(Z_H\) of \eqref{eq:hyperplaneslice610}.
In particular the constant component in \eqref{eq:sumallplanes610}
uses one real trace per cross-matrix class and its signed finite
coefficient. The size of their total, including the variation of
the real traces and the remaining period terms, still enters
the uniform comparison in Section~\ref{sec:missing}.


\subsection{Finite characters in the sum of frame determinants}
\label{sec:rotationcharacters612}

Fix an occurring \([C]\), a base class \(\eta_0\) with an
integral lift, and the same \(H,j,i\) as in
Theorem~\ref{thm:rotationperiod611}. Give its stabilizers the
transported measures constructed there. Let \(\mathcal T\) denote
the norm-one torus of \(E/F\). For a finite place \(v\) write
\(\mathcal T_v=\mathcal T(F_v)\) and let \(\mathcal T_v^0\)
consist of its integral units, using both factors at a split place.
For \(a\in\mathcal T_v\) let \(L_v(a)\) be the signed local
integral from Theorem~\ref{thm:rowfactor607}, using the transported
parameters at \(R_t(a)\eta_0 R_s(a)^{-1}\). It includes the original
support and full marking, and has value zero on empty support.
The normalization of the coefficient of \(I_{\eta_0}\) in
\eqref{eq:rotationperiod611} is then
\begin{equation}
  A_i=\frac{\Delta_F^{r/2}}{c_{t,s}v_U}
             \sum_{\alpha\in E^1}\chi_\infty(\alpha)m_i(\eta_\alpha)
       =\sum_{\alpha\in E^1}\chi_\infty(\alpha)
                                      \prod_{v<\infty}L_v(\alpha).
  \label{eq:rotationnormalized612}
\end{equation}

For each end \(b\) let \(P_b^\ell\) be the orthogonal projection
of the ambient row onto its fixed normal line \(\ell_b\), and
let \(\mathfrak d_b\) be the ideal of integral scalars that clear
the denominators of \(P_b^\ell\) on \(\Lambda_b\). Put
\(\mathfrak d=\mathfrak d_t\cap\mathfrak d_s\).

\begin{proposition}[Local unit subgroups for the original signed tests]
\leavevmode\label{prop:rotationunits612}
At an unmarked place let
\[
 K_{i,v}=\{u\in\mathcal T_v^0:u-1\in\mathfrak d\mathcal O_{E_v}\},
                  \qquad E_v=E\otimes_F F_v .
\]
At \(v\mid Q\) also require
\(\overline{R_t(u)}\,i_v=i_v\,\overline{R_s(u)}\) on the marked
positive column. Then \(L_v(au)=L_v(a)\) for every \(u\in K_{i,v}\).
The support of \(L_v\) is contained in
\(\{a:a\mathfrak a\mathcal O_{E_v}\subseteq\mathcal O_{E_v}\}\)
with \(\mathfrak a\) from Proposition~\ref{prop:rotationideal611}.
It consists of finitely many cosets of \(K_{i,v}\).

There is the following explicit test at a mark. In its chosen
split factor let \(\pi\) be a uniformizer, set
\(D=\max\{v_\pi(\mathfrak d_s),v_\pi(\mathfrak d_t)\}\), and
lift \(i_v\) to an integral isomorphism \(J_v\) between the
positive column lattices. On these columns write \(P_b=P_b^\ell\).
For a unit \(u\) in the chosen factor the condition is
\[
 v_\pi(u-1)\ \ge\ D+
  {\bf1}_{\,\pi^D(P_tJ_v-J_vP_s)\not\equiv0\bmod\pi}.
\]
In particular \(K_i=\prod_vK_{i,v}\) contains all integral
norm-one units congruent to one modulo the ideal \(Q\mathfrak d\),
and
\[
     N_{E/\QQ}(Q\mathfrak d)
              \le Q^{2f}(T_0T_1)^{2f} X_C^4 .
\]
\end{proposition}
\begin{proof}
Write \(M_b=\Lambda_b\cap\ell_b\). Self-duality and extension
of functionals from this primitive line give \(P_b^\ell\Lambda_b
=M_b^\sharp\). Thus \(\mathfrak d_b\) is the annihilator of
\(M_b^\sharp/M_b\). If \(M_b=Iv\), this is
\(I\bar I(\langle v,v\rangle)\), a conjugation-invariant integral
ideal of norm \([M_b^\sharp:M_b]\).
Now \(R_b(u)=1+(u-1)P_b^\ell\). For integral \(u\) of norm one,
both it and its inverse preserve the row lattice exactly when
\(u-1\in\mathfrak d_b\mathcal O_{E_v}\), since \(u^{-1}-1\)
is a unit multiple of \(u-1\).

These integral rotations fix the two frames pointwise. On an
integral row map their action on both sides preserves the frame
conditions, including the primitive terminal class in the given
row lattice. At a mark the induced middle matrices multiply the
marked index on the left and on the inverse right. This gives
the extra condition in the statement. The outer finite phase is
unchanged by these operations. On the right all successive source
vectors are fixed. On the left their images at a mark agree
modulo \(Q\) with combinations of the fixed target vectors.
Thus the simple-coordinate pairings that give the phase agree
modulo \(Q\) there. Away from \(Q\) their integral changes are
in the kernel of the normalized additive character. This is the
local computation in Proposition~\ref{prop:partnerweights531}
and Theorem~\ref{thm:orbitaverage532}, applied to the two
pointwise stabilizers. The same argument for the inverses
preserves the zero test. The unipotent characters in the
transported coordinates are unchanged because rotations fix the
terminal vector and the outer flag coordinates. This proves
invariance of the whole integrand, hence of \(L_v\).

The projection of the concatenated wedge in the proof of
Proposition~\ref{prop:rotationideal611} works over \(F_v\)
with this same local support. It gives
\(a\mathfrak a\mathcal O_{E_v}\subseteq\mathcal O_{E_v}\) whenever
an integral lift occurs. At a split place this confines the
valuation to a finite interval. At a nonsplit place the local
torus is already compact. The subgroup contains an open unit
congruence subgroup, giving finiteness modulo \(K_{i,v}\).

At a mark the terminal norm is a unit. The integral Witt
decomposition and the terminal line both split, so the projector
denominators on the positive column are exactly the ones computed
on the ambient lattice. The two conjugate exponents are equal.
For \(v_\pi(u-1)\ge D\) the congruence on \(i_v\) is equivalent
to \((u-1)(P_tJ_v-J_vP_s)\equiv0\bmod\pi\). Its coefficient
matrix has valuations at least \(-D\). Reduction after multiplying
by \(\pi^D\) gives exactly the stated extra digit and is
independent of the lift \(J_v\), since its change is divisible
by \(\pi\).

In the self-dual middle row lattice the primitive perpendicular of
\(M_b\) contains \((\Lambda_b\cap H_b)\oplus\mathcal O_E z_b\)
with finite index. Their discriminant covolumes give
\(\delta(M_b)\le T_b^f\delta(\Lambda_b\cap H_b)\le T_b^f X_C\),
using the two-end version of the projection argument in
Proposition~\ref{prop:rotationideal611}. Here \(T_b\) denotes
the absolute terminal norm at end \(b\). Consequently
\(N_{E/\QQ}\mathfrak d_b\le T_b^{2f}X_C^2\).
The intersection of the two integral ideals contains their
product. Multiplying the two bounds and inserting the factor
\(N_{E/\QQ}(Q)=Q^{2f}\) proves the last assertion.
\end{proof}

Put \(\mathcal A=\mathcal T(\mathbb A_{F,f})/K_i\),
\(\mathcal B=\operatorname{im}(E^1\to\mathcal A)\),
\(W_i=E^1\cap K_i\), and \(h_i=|\mathcal A/\mathcal B|\).
Let \(h_E\) denote the ideal class number of \(E\).

\begin{theorem}[A finite character formula for the complete coefficient]
\leavevmode\label{thm:rotationcharacters612}
One has \(W_i\subseteq\mu(E)\) and
\(h_i\le h_E N_{E/\QQ}(Q\mathfrak d)\).
If \(\chi_\infty\) is nontrivial on \(W_i\), then \(A_i=0\).
Otherwise let \(\mathcal X\) be the set of unitary characters
\(\theta\) on \(\mathcal A\) with
\(\theta([\alpha])=\chi_\infty(\alpha)\) for every \(\alpha\in E^1\).
It has \(h_i\) elements. With \(\theta_v\) its restriction at \(v\),
\begin{equation}
        A_i=\frac{|W_i|}{h_i}
          \sum_{\theta\in\mathcal X}\prod_{v<\infty}
              \ \sum_{a\in\mathcal T_v/K_{i,v}} L_v(a)\theta_v(a).
 \label{eq:rotcharacters612}
\end{equation}
Only finitely many local factors differ from one. For identical
central characters \(\mathcal X\) is the full character group of
the finite group \(\mathcal A/\mathcal B\).
\end{theorem}
\begin{proof}
A global norm-one element integral and invertible at every finite
place is a root of unity by the unit argument in
Proposition~\ref{prop:rotationideal611}. Modulo all local compact
units a finite torus idele is specified by one integer at each
split pair, zero almost everywhere. These are exactly the ideals
\(J/\bar J\), as one can choose an exponent on a single factor
of each pair to realize any tuple. For a principal \(J=(t)\)
the tuple is realized by \(t/\bar t\in E^1\). The quotient by the
global points and all compact units thus has order at most \(h_E\).
Passing to \(K_i\) costs at most the number of units modulo
\(Q\mathfrak d\) by Proposition~\ref{prop:rotationunits612}.
This is bounded by the absolute ideal norm. This proves finiteness
and the claimed bound.

Write \(L(a)=\prod_v L_v(a_v)\). Almost everywhere the base class
has unit local factor by Theorem~\ref{thm:rowfactor607}, the ideal
\(\mathfrak a\) is the unit ideal, and \(K_{i,v}=\mathcal T_v^0\).
Thus the only possible local coset is the identity coset, with
value one, by Proposition~\ref{prop:rotationunits612}. Hence \(L\)
has finite support on \(\mathcal A\). In \eqref{eq:rotationnormalized612}
the elements above a point of \(\mathcal B\) form a coset of \(W_i\).
The character sum on such a coset vanishes if the restriction is
nontrivial. Otherwise it is \(|W_i|\) times the character of the
image. This gives a well-defined character on \(\mathcal B\).
It extends to \(\mathcal A\), since the quotient is finite abelian.
Explicitly adjoin one generator at a time, taking an \(n\)-th root
of the value of its \(n\)-th power when its order modulo the
previous subgroup is \(n\). The extensions differ by exactly the
\(h_i\) quotient characters. Their average at a point restricts
to the prescribed value on \(\mathcal B\) and is zero off
\(\mathcal B\) by character orthogonality. Insert this average
in the finite sum over \(\mathcal A\). Each character on its
restricted product of local factors factors on the finitely
supported tuples, so expansion of the product gives
\eqref{eq:rotcharacters612}.
\end{proof}

There are two estimates in the original arithmetic normalization.
Write \(D(\mathfrak a)\) for the product of the split-place factors
in \eqref{eq:rotationcount611}, omitting the prefactor \(|\mu(E)|\).
Then
\[
 |A_i|^2\le |W_i|\,|\mu(E)|D(\mathfrak a)
              \prod_{v<\infty}\sum_{a\in\mathcal T_v/K_{i,v}}
                                                  |L_v(a)|^2.
\]
Indeed at most \(|\mu(E)|D(\mathfrak a)/|W_i|\) points of
\(\mathcal B\) with nonzero coefficient can occur. Above each
one all \(|W_i|\) global representatives have the same finite
coefficient. Apply Cauchy--Schwarz on these image points and
extend the squared sum to \(\mathcal A\), where it factors.
For \(\chi_\infty=1\) there is also an integral bound that does
not involve \(h_i\). On each locally supported coset use one
representative and a precision \(b_v(a)\) from
Theorem~\ref{thm:residueorbits606} in the transported parameters.
Let \(B_v\) be the maximum of these finitely many precisions.
A place with empty support makes the whole coefficient zero.
For nonzero \(A_i\) and a valuation above the split terminal
prime normalized by \(v_T(T)=1\), one has
\[
                        v_T(A_i)\ge -r\sum_{v\mid T}B_v .
\]
In fact \(q_v^{rB_v}L_v(a)\) is an algebraic integer for every
coset, since it is an integral multiple of its sum of roots of
unity. Almost everywhere \(B_v\) can be chosen zero. The rational
sum in \eqref{eq:rotationnormalized612} is finite, all its character
weights are one here, and \(q_v=T\) at \(v\mid T\). Multiplication
by the product of the local denominators proves the inequality.

\begin{corollary}[The root-of-unity sum on the original unit-terminal rows]
\leavevmode\label{cor:rotationroots612}
Suppose both terminal vectors have norm \(-1\) and one supported
cell in the cross-matrix class has \(h_R=1\). Then
\(K_i=\prod_v\mathcal T_v^0\) and \(W_i=\mu(E)\). In particular
the entire contribution \eqref{eq:rotationperiod611} vanishes
when \(\chi_\infty|_{\mu(E)}\ne1\), at every permitted frame
cutoff that includes this class. When \(X_C=1\), the normalized
coefficient for every \(\chi_\infty\) is exactly
\[
 A_i=
 \begin{cases}
 |\mu(E)|\,\displaystyle\prod_{v<\infty}L_v(1),
                     &\chi_\infty|_{\mu(E)}=1,\\
 0,                  &\chi_\infty|_{\mu(E)}\ne1 .
 \end{cases}
\]
For identical columns this multiplies the single base trace by
\(|\mu(E)|\) in the full signed sum. The multiplier is a unit at
the terminal prime.
\end{corollary}
\begin{proof}
At either end the image of \(R\) is self-dual in \(H_b\) and
lies in the positive integral column. Thus it splits off
orthogonally by integral projection. The unit terminal line
also splits off from the self-dual middle row. The remaining
normal lattice is therefore self-dual and \(\mathfrak d_b=(1)\).
Theorem~\ref{thm:mixedtransport535} makes \(j\) integral with
integral inverse, inducing \(i_v\) at every mark. It carries
normal line to normal line, so \(P_tj=jP_s\). Taking \(J_v=j\)
in Proposition~\ref{prop:rotationunits612} gives zero extra digits
at all marks. This proves the subgroup assertions. If \(X_C=1\),
the equality \(X_C=I_\Lambda h_R^f\) gives \(h_R=1\) at every
supported lift. Proposition~\ref{prop:rotationideal611} permits
only roots of unity, and the local contribution for every root
equals that for one by unit invariance. Summing the character on
\(\mu(E)\) gives the displayed values. Since the odd terminal prime
\(T\) splits completely in \(E\), the group \(\mu(E)\) injects
into \(\QQ_T^\times\), whose torsion has order \(T-1\). Indeed
reduction is injective on torsion in \(\ZZ_T^\times\) because the
\(T\)-adic logarithm is injective on \(1+T\ZZ_T\). Hence
\(T\nmid|\mu(E)|\).
\end{proof}

For sufficiently distant intervals a retained row of terminal
norm \(-T\) forces \(h_R>b_x^5\) by
Proposition~\ref{prop:retainedtest535}. Thus the evaluation at
\(X_C=1\) covers every retained pair for which this value occurs.
The finite character formula applies at all other values as well.
The needed bound for the sum of its coefficients against the real
traces, and the residual cohomology bound on these fields, remain
open in Section~\ref{sec:missing}.



\subsection{Moments of the integral frame determinants}
\label{sec:orientationmoment613}

Fix the two marked rows, the positive hyperplanes, their rational
transport and the residue index as in
Theorem~\ref{thm:rotationperiod611}. The terminal norms \(T_s,T_t\)
can independently be \(1\) or \(T\). At a norm-\(T\) end impose
the hyperplane test of Theorem~\ref{thm:columns235}, with a fixed
selection parameter \(A_0>0\). Let \([C]\) range over the relevant
transverse cross matrices of positive Gram invariant that have a
supported integral lift at these data. Define \(n_i(C)\) to be the
number of rational unipotent double classes above \([C]\) with
such a lift. The quantities to be summed are
\[
 w_C=\prod_{j:D_j\ne0}\prod_{\tau}(1+|\tau D_j|^2)^2,\qquad
 X_C=Q^{-2df}(N_{E/\QQ}\mathfrak h_d)^2|N_{F/\QQ}\mathcal G|.
\]
Here \(D_j\) is the leading cross minor of order \(j\), \(d=l+1\),
and \(\tau\) runs over one complex embedding above each real place.
Write \(Y_0=1\) on two unit-terminal rows. If at least one terminal
norm is \(T\), let \(Y_0=T^f(k|\operatorname{disc}E|^{1/(2f)}
/(2\pi\exp(A_0+4)))^f\), the hyperplane threshold in the original
positive lattice at that end.

\begin{theorem}[Moments over all supported frame determinants]
\leavevmode\label{thm:orientationmoment613}
Fix \(0\le \sigma<1\). On the common intervals of
Theorem~\ref{thm:columns235}, one has, uniformly in the indicated
data,
\[
 \sum_{[C]}\frac{n_i(C)\,X_C^\sigma}{w_C}
 \ \le\
 \begin{cases}
  \exp\!\bigl(-(64-o(1))fd\,x^\varepsilon\bigr),
                         &T_s=T_t=1,\\
  \exp\!\bigl(-(132-34\sigma-o(1))fd\,x^\varepsilon\bigr),
                         &\max(T_s,T_t)=T .
 \end{cases}
\]
The sum includes arbitrarily large values of \(X_C\). For a fixed
\([C]\) the multiplicity counts each rational double class once,
even when it has several integral lifts.
\end{theorem}

We first give the tail estimate used in the proof. It retains the
length of the final reversed block and also applies to tails beyond
any prescribed cutoff. The positive parameter \(\lambda\) below is
independent of the interval exponent \(\varepsilon\).

\begin{lemma}[The last block with a lower bound on the Gram determinant]
\leavevmode\label{lem:orientationtail613}
Let \(h\) be the final block length, \(j=d-h\), \(Z\ge1\), and
\(Z_*=\max(Z,Y_0)\). Put \(u=\sigma+2\lambda\) for \(\lambda>0\).
Choose \(0<\delta<1\) with \(4h-2-2hu-\delta\ge0\). Set
\[
 B_\delta=9\int_0^\infty
       \frac{4r^{5-\delta}}{(1+r^2)^3}\,dr+2,\qquad
 C_\lambda=\prod_{p<2^{1/\lambda}}(1-p^{-\lambda})^{-2},
\]
where the product is over rational primes. Then
\begin{align*}
 \sum_{\substack{[C]\text{ with last length }h\\ X_C\ge Z}}
       \frac{n_i(C)X_C^\sigma}{w_C}
 &\le |\mu(E)|\, C_\lambda^{\,f} B_\delta^{fj}
                  \bigl(9(2+4h/\delta)\bigr)^f\\[-2pt]
 &\quad{}\cdot Q^{-(4-\delta)jf-(4-\delta/h)df}
          Z_*^{-(2-1/h-u-\delta/(2h))}.
\end{align*}
\end{lemma}
\begin{proof}
The prescribed outer ideals are nonpositive powers of the extended
different by Section~\ref{sec:flagideals554}. Thus
\(N_{E/\QQ}\mathfrak h_d\le1\), and every rational frame vector is
integral. Put \(B=\prod_\tau|\tau D_d|\). At each embedding the
identity
\(\mathcal G=D_d\bar D_d-T_sT_tD_{d-1}\bar D_{d-1}>0\) gives
\[
           X_C\le B^2/Q^{2df},\qquad B\ge Q^{df}\sqrt{Z_*}
           \quad\hbox{on the indicated tail}.
\]
For the second inequality use \(X_C\ge Z\) and
\(X_C=I_\Lambda\delta(R)\), by Corollary~\ref{cor:actualgram535}.
At a retained norm-\(T\) end the common lattice \(R\) projects
isometrically into a primitive hyperplane of the original positive
lattice, so \(\delta(R)>Y_0\) by
Proposition~\ref{prop:retainedtest535}. At either other end
\(X_C\ge1\) by Proposition~\ref{prop:rotationideal611}.
That proposition also bounds \(n_i(C)\) by
\(|\mu(E)|C_\lambda^f X_C^{2\lambda}\). We may henceforth sum
\(X_C^u/w_C\) over cross-matrix classes and restore the indicated
constant at the end.

The proof of Theorem~\ref{thm:smallframes590} gives reversed
consecutive blocks, \(D_a\in Q^a\mathcal O_E\) at every nonterminal
endpoint, and \(D_d+D_{d-1}\in Q^d\mathcal O_E\). These congruences
use the actual two marks, where all the terminal vectors have the
same unit reduction. For \(j=0\) take \(D_0=1\).
Fix all prefix block scalars through \(j\). If \(h\ge2\), then
\(D_{d-1}=0\) and \(D_d=\pm D_j c^h\) for the last scalar \(c\).
Write \(D_j=Q^j a\). The nonzero integer \(a\) satisfies
\(a(c/Q)^h\in\mathcal O_E\). Hence \(c/Q\) belongs to the
fractional ideal with exponents
\(-\lfloor v_{\mathfrak p}(a)/h\rfloor\). Its norm is at least
\(|N_{E/\QQ}a|^{-1/h}\). In coordinates
\(v_\tau=|\tau D_j|^{1/h}\tau c\), distinct possibilities are
therefore separated by
\[
 \prod_\tau|v_\tau-v'_\tau|\ \ge\
 |N_{E/\QQ}D_j|^{1/(2h)}Q^f|N_{E/\QQ}a|^{-1/(2h)}
      =Q^{df/h}.
\]
Here \(|v_\tau|^h=|\tau D_d|\). If \(h=1\), use
\(v_\tau=\tau D_d\) itself. For a fixed prefix the shifted
congruence makes every difference a multiple of \(Q^d\), and
the final scalar is recovered uniquely from \(D_d/D_j\) up to
its fixed sign. Thus the same separation holds also in this case.

Take a box \(|v_\tau|\le r_\tau\), set
\(R_0=(\prod_\tau r_\tau)^{1/f}\), and put \(P=Q^{d/h}\).
Nonempty boxes require \(R_0\ge P Z_*^{1/(2hf)}\). After rescaling
to the unit polydisc their centers have sup distance at least
\(P/R_0\). Disjoint open polydiscs of radius \(P/(2R_0)\) fit
inside the polydisc of radius \(1+P/(2R_0)\), giving at most
\(9^f(R_0/P)^{2f}\) choices. For \(b=4h-2-2hu-\delta\) bound
the support indicator by \((R_0/P)^{bf} Z_*^{-b/(2h)}\).
On this box \(X_C^u\le Q^{-2udf}\prod_\tau r_\tau^{2hu}\).
Integrate the count against the positive density
\(4h r^{2h-1}(1+r^{2h})^{-3}\,dr\) in every coordinate. Its
tail from \(|v_\tau|\) is the original last-endpoint weight.
Tonelli's theorem applies first to finite subsums and then to
the whole nonnegative sum. The one-variable integral left after
the indicated powers are inserted is
\[
   \int_0^\infty \frac{4h r^{2h+1+b+2hu}}{(1+r^{2h})^3}\,dr
          \le 2+4h/\delta .
\]
Indeed its exponent in the numerator is \(6h-1-\delta\). Split
at one and use the denominator bounds \(1\) and \(r^{6h}\).
The equality \(Q^{-2udf}P^{-(2+b)f}
=Q^{-(4-\delta/h)df}\) gives the stated last-endpoint factor.
Counting \(c\) itself includes all of its possible roots.

For completeness fix a composition of the prefix of length \(j>0\)
with endpoint set \(J\). At \(a\in J\) a nonzero value lies in
\(Q^a\mathcal O_E\), and distinct values are separated in the
product of complex distances by \(Q^{af}\). The same polydisc
argument with \(h=1\), \(u=0\), no additional cutoff, and
indicator exponent \(2-\delta\) gives a factor
\((B_\delta-2)^f Q^{-(4-\delta)af}\).
The endpoint \(j\) accounts for the claimed power of \(Q\).
Each preceding endpoint ratio has at most its block length many
roots. Their product over the prefix is at most \(2^{j-|J|}\).
Summing over subsets of \(\{1,\ldots,j-1\}\) gives at most
\(B_\delta^{fj} Q^{-(4-\delta)jf}\), since
\((B_\delta-2)^f+2\le B_\delta^f\). For \(j=0\) the prefix
factor is one. Multiplying the factors proves the lemma.
\end{proof}

\begin{proof}[Proof of Theorem~\ref{thm:orientationmoment613}]
Apply the lemma with \(Z=1\), \(\lambda=1/\sqrt{\log(2x)}\),
and \(\delta=1/\log(2Q)\). The condition on \(b\) holds for all
\(h\) once \(x\) is large, uniformly for \(\sigma\) in any fixed
compact subinterval of \([0,1)\). The prefix integral gives
\(B_\delta=O(1/\delta)\) by splitting it at one. Also
\(\log C_\lambda\le
2^{1+1/\lambda}\log(1/(1-2^{-\lambda}))=x^{o(1)}\).
By Lemma~\ref{lem:scalarden562}, \(|\mu(E)|\le8f^2\), and
\(f\) is polynomial in \(x\). All positive logarithms of the
prefactors and the cost \(\log d\) of summing the lengths are
\(o(fd\,x^\varepsilon)\). This uses \(d\to\infty\) polynomially,
\(\log Q=(16+o(1))x^\varepsilon\), and fixed \(\varepsilon>0\).

On two unit-terminal rows the power of \(Q\) from \(d\) alone
gives the first assertion. On any pair with a retained norm-\(T\)
row, \(\log Y_0=(34+o(1))fd\,x^\varepsilon\), by
Theorem~\ref{thm:prime209} with \(m=2d\) and the polynomial
bounds for \(k\) and the root discriminant. After division by
\(fd\,x^\varepsilon\) the saving for a given \(h\) is at least
\[
       64(1+j/d)+34(2-1/h-\sigma)-o(1).
\]
If \(h\ge\sqrt d\), the term \(34/h\) tends uniformly to zero.
If \(h<\sqrt d\), then \(j/d=1-o(1)\) and \(64j/d\) absorbs
\(34/h\). The second assertion follows.
\end{proof}

In particular at \(\sigma=0\) the norm-\(T\) estimate of
Theorem~\ref{thm:powers236} retains its leading exponent after
counting the rational frame determinants over each label, also
on a pair with one unit-terminal end. The weights in this theorem
are explicit functions of the cross matrices. The signed
coefficient in the projected period still contains the finite
integrals in \eqref{eq:rotationnormalized612} multiplied by the
real tensor in \eqref{eq:rotationperiod611}. Bounding their total
over the hyperplanes and transports is required in
Section~\ref{sec:missing}.


\subsection{Evaluation of the signed Bessel integral on the open cell}
\label{sec:besselsampling615}

Retain the fields, interval, marked rows, positive hyperplanes,
transport and residue index of Theorem~\ref{thm:orientationmoment613}.
Write \(d=s+1\), let \(k\ge2\) be the dimension of the positive
terminal space, and set \(A=T_0T_1\). Here \(T_0,T_1\) are the
positive absolute norms of the two negative terminal vectors.
Consider the rational cross matrices for which every leading
minor \(D_j\), \(1\le j\le d\), is nonzero and the transverse
Gram invariant \(\mathcal G\) is totally positive. We give an
estimate on this open cell that remains uniform as a local
Gram invariant tends to zero. Its analytic right side consists
of explicitly specified derivatives.

At a real place, take \(s\) hyperbolic pairs and a middle space of
signature \((k,1)\). A representative on the open cell exchanges
the \(i\)-th outer vector with its dual, with multipliers
\(c_i,\bar c_i^{-1}\), and acts by an element of \(U(k,1)\) on
the middle space. Indeed left and right unitriangular elimination
diagonalizes a cross matrix with nonzero leading minors, with
successive diagonal entries equal to the ratios of these minors.
The perpendicular middle is then identified by an isometry.
In bases with terminal vectors \(\sqrt{T_0}e_-\) and
\(\sqrt{T_1}e_-\), let \(\zeta\) be their cross pairing. Write
\[
 z=(c_1,\ldots,c_s,\zeta),\qquad
 U_j(z)=\prod_{i=1}^j z_i,\qquad
 \Omega_A=(\mathbb C^\times)^s\times
                       \{\zeta:|\zeta|^2>A\}.
\]
The product \(U_j\) is the leading minor in these coordinates.
One can choose a smooth middle representative throughout
\(\Omega_A\). In the plane of one positive vector and \(e_-\),
use the real boost with diagonal \(\cosh t\) and off-diagonal
\(\sinh t\), followed by a phase on the negative line chosen
to give the prescribed \(\zeta\), where
\(|\zeta|=\sqrt A\cosh t\). Extend by the identity on the
perpendicular space. The compact groups at the ends account for
the choice of positive direction.

For a smooth group coefficient \(F\), denote its signed transform
on this section by \(J(z)\) whenever the orbital integral is
defined. Its integrand is
\[
 \chi_0(n_0)\overline{\chi_1(n_1)}
  \rho_0(h_0)F(h_0^{-1}n_0^{-1}g(z)n_1h_1)\rho_1(h_1)^{-1}
\]
over the two compact groups and the two flag unipotents, with
the given Haar measures. The unipotent stabilizer on this cell
is trivial. In the block order consisting of the outer lines,
the middle space and the opposite lines in reverse order, the
first unipotent is strictly block upper triangular after
subtraction of the identity. Transport of the other unipotent
by this representative makes it strictly block lower triangular
after the same subtraction. Their intersection is the identity.
Thus the displayed integral uses the original characters on
both sides. The estimate below also applies to Hilbert-space
valued profiles, in particular to a matrix with the
Hilbert--Schmidt norm.

Set \(a_i=2(d-i)+k-1\) for \(1\le i\le s\), and use ordinary
complex Lebesgue measure \(dz\) in the chart. The weight is
\[
  \rho_A(z)=\prod_{i=1}^s |c_i|^{2a_i}
                        (|\zeta|^2-A)^{k-1}.
\]
The powers can be read directly from the Haar measures. In the
decomposition \(n_0 w m n_1\), the density in \(m\) includes
the determinant of conjugation on the right unipotent. Right
translation by \(m_0\) replaces \(n_1\) by \(m_0^{-1}n_1m_0\);
its Jacobian is the inverse of this determinant. For a pair of
outer indices \(i<j\) the two complex coordinates have weights
\(c_i/c_j\) and \(c_i\bar c_j\). Their two real Jacobians
multiply to \(|c_i|^4\). The \(k+1\) complex coordinates to the
middle contribute \(|c_i|^{2(k+1)}\), and the imaginary diagonal
coordinate contributes \(|c_i|^2\). Multiplication by the
multiplicative measure \(d^2c_i/|c_i|^2\) gives precisely
\(|c_i|^{2a_i}d^2c_i\). For the middle quotient take a negative
unit vector \((v,b)\in\mathbb C^k\times\mathbb C\), with
\(\|v\|^2-|b|^2=-1\). Euclidean measure disintegrated at this
level of the Hermitian form is unitary invariant, since unitary
matrices preserve both the form and real volume. In polar
coordinates for \(v\), fixing \(\|v\|^2=|b|^2-1\) gives
\(\pi^k(|b|^2-1)^{k-1}/(k-1)!\) times probability measure in
its direction. Passing from \(b\) to the terminal pairing of
modulus \(\sqrt A\,|b|\) contributes the constant \(A^{-k}\)
in this density. We use \(\rho_A\,dz\) explicitly on both sides
of the trace estimate. The middle factor converting it to the
stated invariant shell measure is \(\pi^k/(A^k(k-1)!)\).

Use radial variables \(u_i=\log|c_i|\) for \(i\le s\) and
\(u_d=t>0\), and the ordinary angles \(\theta_j=\arg z_j\).
For subsets \(I,B\) of \(\{1,\ldots,d\}\), let
\(\partial_u^I\partial_\theta^B\) differentiate once in each
of the indicated variables. Define
\[
 \mathcal S_A(H)=\sum_{I,B\subseteq\{1,\ldots,d\}}
       \int_{\Omega_A}\rho_A(z)
       \|\partial_u^I\partial_\theta^B H(z)\|^2\,dz,
 \qquad
 K_{k,s}=2(k-1)\prod_{i=1}^s 2(a_i+1).
\]

\begin{theorem}[A trace estimate at every positive terminal Gram value]
\leavevmode\label{thm:hardycell615}
For a smooth Hilbert-space valued \(H\) on \(\Omega_A\) with
\(\mathcal S_A(H)<\infty\), every \(z\in\Omega_A\) satisfies
\[
 \rho_A(z)\|H(z)\|^2
      \le \frac{\mathcal S_A(H)}{K_{k,s}|U_d(z)|^2}.
\]
For compactly supported \(H\) only the terms with
\(I=\{1,\ldots,d\}\) are needed on the right.
\end{theorem}
\begin{proof}
We first assume compact support. Fix the angles. In the
\(i\)-th outer variable \(u\), the radial energy measure is
\(e^{(2a_i+2)u}\,du\). Integrating a derivative from \(u\) to
infinity and applying Cauchy--Schwarz costs the reciprocal
weight integral \(e^{-(2a_i+2)u}/(2a_i+2)\). After multiplication
by the point weight \(e^{2a_i u}\), the factor is
\(1/((2a_i+2)|c_i|^2)\).

In the terminal variable the radial measure is
\(A^k\sinh^{2k-1}t\cosh t\,dt\). The corresponding inverse
weight integral satisfies
\[
 \int_t^\infty\frac{ds'}{\sinh^{2k-1}s'\cosh s'}
 \le\frac1{\cosh^2t}\int_t^\infty
                   \frac{\cosh s'}{\sinh^{2k-1}s'}\,ds'
 =\frac{\sinh^{-2(k-1)}t}{2(k-1)\cosh^2t}.
\]
Multiplying by \(A^{-k}(A\sinh^2t)^{k-1}\) leaves
\(1/(2(k-1)|\zeta|^2)\). This cancellation takes place before
any product over embeddings. It applies for arbitrarily small
positive \(t\).

Repeated integration of the full mixed radial derivative from
the evaluation point to infinity recovers \(H\) up to sign.
Cauchy--Schwarz on the product of these half-lines multiplies
the inverse weight integrals just computed. To evaluate at the
prescribed angles, use the Fourier series on the \(d\)-torus.
For a smooth Hilbert-space valued function \(g\), Cauchy--Schwarz
with Fourier weights \(\prod_j(1+n_j^2)\), followed by Parseval,
gives
\[
  \|g(\theta)\|^2\le
  \left(\frac{1}{2\pi}\sum_{n\in\mathbb Z}(1+n^2)^{-1}\right)^d
     \sum_B\int_{[0,2\pi]^d}\|\partial_\theta^B g\|^2\,d\theta .
\]
The constant in parentheses is at most \(5/(2\pi)<1\), since
\(\sum_{n\ge1}n^{-2}\le1+\int_1^\infty y^{-2}dy=2\). Apply this
inequality to \(H\) at the chosen radii, then the radial
inequality to each angular derivative. Enlarge the resulting
outward domains to \(\Omega_A\). This proves the bound for
compact support.

For the general assertion fix the evaluation point. In each
radial variable insert a smooth cutoff equal to one below \(R\)
and zero above \(R+1\), with first derivative bounded independently
of \(R\). Below the evaluation point insert another cutoff equal
to one on its entire outward half-line. The second cutoff may
be supported away from \(t=0\) or \(u_i=-\infty\), as appropriate.
Use the proof with the integrals still restricted to the outward
domain. Terms differentiating an upper cutoff are supported
where some variable exceeds \(R\) and tend to zero in the
weighted \(L^2\) norm by finiteness of \(\mathcal S_A(H)\).
Dominated convergence treats the term in which all derivatives
hit \(H\). This proves the stated inequality.
\end{proof}

We next sum these evaluations in the archimedean frames belonging
to the same integral flags as the prescribed period. Write
\(\eta_j=(N_{E/\mathbb Q}\mathfrak h_j)^{-1/(2f)}\) as in
Proposition~\ref{prop:transverseheight582}, with \(\eta_0=1\).
For a supported rational class \(C\) set
\[
 V_{j,\tau}(C)=\eta_j^{-2}\tau D_j,\qquad
 z_{j,\tau}(C)=V_{j,\tau}(C)/V_{j-1,\tau}(C),\qquad V_{0,\tau}=1.
\]
These coordinates divide the \(i\)-th outer vector at each
end by its norm \(\eta_i/\eta_{i-1}\). The terminal vector keeps
its given norm because \(\eta_d=\eta_s\). In particular
\(\mathcal G>0\) at \(\tau\) gives
\(|z_{d,\tau}|^2>A\). Constant phase changes in any of the two
ordered frames can be included in the arguments of the
profiles. Put
\[
 W(z)=\prod_{j=1}^d(1+|U_j(z)|^2)^2,\qquad
 W_C=\prod_\tau W(z_\tau(C)).
\]
Let \(X_C,n_i(C)\) be the normalized integral determinant and
number of supported frame determinants of
Theorem~\ref{thm:orientationmoment613}.

\begin{proposition}[Packing all normalized leading minors]
\leavevmode\label{prop:openpacking615}
For \(0<\delta<1\), put
\[
 B_\delta=9\int_0^\infty\frac{4r^{5-\delta}}{(1+r^2)^3}\,dr .
\]
On the open cell with the fixed data above, at either choice of
each terminal norm,
\[
 \sum_{[C]}\frac1{W_C}
      \le B_\delta^{fd}\,Q^{-(4-\delta)fd(d+1)/2}.
\]
\end{proposition}
\begin{proof}
The congruences of Theorem~\ref{thm:smallframes590} put \(D_j\)
in \(Q^j(\bar{\mathfrak h}_j\mathfrak h_j)^{-1}\) for \(j<d\).
For \(j=d\), fixing \(D_s\) puts \(D_d\) in a translate of the
same ideal with \(j=d\). If \(D_j,D'_j\) are distinct possible
values at a step with fixed prefix, taking the ideal norm of
their difference gives
\[
  \prod_\tau|V_{j,\tau}-V'_{j,\tau}|\ge Q^{jf}.
\]
For \(j<d\), comparison with zero gives
\(\prod_\tau |V_{j,\tau}|\ge Q^{jf}\). At \(j=d\) the
determinant formula and total positivity give
\[
 Q^{2df} X_C=(N_{E/\mathbb Q}\mathfrak h_d)^2
                    |N_{F/\mathbb Q}\mathcal G|
       \le\prod_\tau |V_{d,\tau}|^2 .
\]
Every supported class has \(X_C\ge1\) by
Proposition~\ref{prop:rotationideal611}. This supplies the same
lower bound also in the translated last lattice.

At a fixed step write \(P=Q^j\). In the box
\(|V_{j,\tau}|\le r_\tau\), put
\(R=(\prod_\tau r_\tau)^{1/f}\), taking the radii positive.
A nonempty box has \(R\ge P\). After dividing the coordinates by
their radii, distinct centers have sup distance at least \(P/R\).
Disjoint open polydiscs of radius \(P/(2R)\) lie in the
polydisc of radius \(1+P/(2R)\). Hence the count is bounded by
\((1+2R/P)^{2f}\le9^f(R/P)^{(4-\delta)f}\).
Integrate this bound against the product of the densities
\(4r_\tau(1+r_\tau^2)^{-3}dr_\tau\). The one-coordinate tail
of this density is \((1+|V_{j,\tau}|^2)^{-2}\).
Tonelli's theorem gives the bound \(B_\delta^f Q^{-(4-\delta)jf}\)
at step \(j\), uniformly in the earlier minors.
On this cell all the blocks have length one, so the ordered
minor tuple determines its rational cross-matrix class by
successive division. Summing from \(j=d\) down to \(1\)
therefore gives the claimed exponent.
\end{proof}

\begin{corollary}[The sampled real profiles on the two marked rows]
\leavevmode\label{cor:openprofiles615}
Let \(0\le\sigma<1\), choose \(\lambda>0\) with
\(\sigma+2\lambda\le1\), and let \(C_\lambda\) be the finite
constant in Lemma~\ref{lem:orientationtail613}. Take one smooth
profile \(J_\tau\) per real place with
\(\mathcal S_A(WJ_\tau)<\infty\). Then
\begin{align*}
 &\sum_{[C]}n_i(C)X_C^\sigma W_C
       \prod_\tau \rho_A(z_\tau(C))\|J_\tau(z_\tau(C))\|^2\\
 &\qquad\le
 \frac{|\mu(E)|\,C_\lambda^f B_\delta^{fd}}
      {K_{k,s}^f Q^{\,2df+(4-\delta)fd(d+1)/2}}
             \prod_\tau\mathcal S_A(WJ_\tau).
\end{align*}
For fixed \(\sigma<1\) the coefficient of the product of
energies is at most
\(\exp(-(32-o(1))fd(d+1)x^\varepsilon)\) on the period
intervals.
\end{corollary}
\begin{proof}
Apply Theorem~\ref{thm:hardycell615} to \(WJ_\tau\).
Multiplying over the places leaves \(W_C^{-1}\) after inserting
the \(W_C\) on the left. The remaining determinant denominator
is \(\prod_\tau |V_{d,\tau}|^2\ge Q^{2df}X_C\).
Proposition~\ref{prop:rotationideal611} gives
\(n_i(C)\le|\mu(E)|C_\lambda^f X_C^{2\lambda}\).
Since \(X_C\ge1\), the residual power
\(X_C^{\sigma+2\lambda-1}\) is at most one. Now apply
Proposition~\ref{prop:openpacking615}. For the asymptotic
keep \(\lambda>0\) fixed and take \(\delta=1/\log(2Q)\).
Splitting the defining integral at one gives
\(B_\delta=O(1+\delta^{-1})\). Use
\(\log Q=(16+o(1))x^\varepsilon\),
\(|\mu(E)|\le8f^2\) from Lemma~\ref{lem:scalarden562}, and the
polynomial field degree on these intervals. These factors
give the asserted estimate.
\end{proof}

The corollary applies directly to smooth compactly supported
pieces of a signed orbital profile. To use it for the full
projected period one needs a bound for the displayed mixed
derivative energies relative to the energy of the prescribed
uncut coefficient, including its Haar constants. The cells
with a vanishing proper leading minor and the finite signed
orbital factors in \eqref{eq:rotationnormalized612} also enter
the comparison of Section~\ref{sec:missing}.


\subsection{Flag minors on every transverse block of the signed period}
\label{sec:pivotperiod619}

Fix the retained rows at \(x\), with \(E=E_x\), \(F=E^+\),
\(f=[F:\QQ]\), \(Q=g_1g_2\), \(d=s+1\), and terminal
norms \(-T_0,-T_1\), as in Section~\ref{sec:transversesum582}.
Use the original marked level at both primes. Consider relevant
transverse terms for which
\(\mathcal G=D_d\bar D_d-T_0T_1D_s\bar D_s\) is totally positive.
For one such term write \(Z\) for the cross matrix and \(w\) for its
pivot permutation in Lemma~\ref{lem:forms203}. This permutation
reverses consecutive blocks. The leading square of a proper prefix
inside a reversed block is singular. A different square minor of
the same row prefix retains the missing pivot.

For \(I\subseteq\{1,\ldots,d\}\) put
\(\mathfrak a_I=\prod_{i\in I}\mathfrak a_i\), where
\(\mathfrak a_i=\mathfrak d^{-(d-i)}\), also for \(i=d\).
Here \(\mathfrak d\) is the extension of the different of \(F\)
from Section~\ref{sec:flagideals554}. Put
\(\eta_I=(N_{E/\QQ}\mathfrak a_I)^{-1/(2f)}\). Write
\([j]=\{1,\ldots,j\}\), \(K_j=w([j])\), and take all subsets in
increasing order when forming a determinant. Set
\[
 M_j=\det Z_{[j],K_j},\qquad
 V_{j,\tau}=\frac{\tau M_j}{\eta_{[j]}\eta_{K_j}},\qquad
 B_j=\prod_{\tau}|V_{j,\tau}|.
\]
At \(j=0\) all three quantities equal one. The product takes one
complex embedding over each real place.

\begin{proposition}[An integral minor for each row prefix]
\leavevmode\label{prop:pivotflags619}
The \(M_j\) are nonzero invariants of the two decorated flags.
For a supported integral lift and \(j<d\),
\[
   M_j\in Q^j(\overline{\mathfrak a_{[j]}}\mathfrak a_{K_j})^{-1},
   \qquad B_j\ge Q^{jf}.
\]
At the final index \(M_d=D_d\) and
\[
   B_d=Q^{df}(X_C P_\gamma)^{1/2},
   \qquad P_\gamma=\prod_\tau
       \left(1+\frac{T_0T_1|\tau D_s|^2}{\tau\mathcal G}\right).
\]
Inside any reversed block with endpoints \(a,a+h\), for
\(0\le m\le h\) one has
\[
 |V_{a+m,\tau}|=
 |V_{a,\tau}|^{1-m/h}|V_{a+h,\tau}|^{m/h}.
\]
\end{proposition}
\begin{proof}
Left multiplication by a lower unitriangular matrix acts on the
first \(j\) rows by a determinant-one matrix. On the right, apply
the minors formula for multiplication by an upper unitriangular
matrix. Starting with the monomial pivot representative, precisely
one column subset \(K_j\) has nonzero \(j\)-minor on these rows.
A minor in columns \(L\) after the change can receive a contribution
from \(K_j\) only if the ordered indices of \(K_j\) are componentwise
at most those of \(L\), by upper triangularity. For \(L=K_j\) the
multiplier is one. Every changed matrix has the same pivot permutation,
as seen from the ranks of its northwest rectangles. This proves
invariance and identifies \(M_j\) as the value in the monomial
representative. This argument allows arbitrary unitriangular changes
over \(E\).

Compute on an actual integral lift. The row lattices contain
\(\mathfrak a_iY_{b,i}\) in each of the two frames \(Y_b\), including
the integral terminal vector \(Y_{b,d}\). Pairing their exterior
products gives
\(M_j\overline{\mathfrak a_{[j]}}\mathfrak a_{K_j}
\subseteq\mathcal O_E\). At the two split marks those fractional
ideals are units and \(Z\equiv-E_{dd}\bmod Q\), by the marking in
Theorem~\ref{thm:smallframes590}. For \(j<d\) the row set omits \(d\),
so every entry of the selected square is zero modulo both primes,
even if \(d\in K_j\). Its determinant gains \(Q^j\). Taking absolute
ideal norms and square roots gives the bound for \(B_j\).
At \(d\), the determinant identity in
Proposition~\ref{prop:transverseheight582} gives
\(|N_{E/\QQ}D_d|=|N_{F/\QQ}\mathcal G|P_\gamma\).
The definition of \(X_C\) in Theorem~\ref{thm:orientationmoment613}
gives the asserted equality.

Within the block the selected columns are
\([a]\cup\{a+h-m+1,\ldots,a+h\}\). The alternating pivot signs cancel
the sign of reversal in the selected \(m\)-square, giving
\(M_{a+m}=D_a c^m\) for the first pivot \(c\).
The new denominator factor at step \(m\) is
\(\eta_{\{a+m\}}\eta_{\{a+h+1-m\}}\). It is independent of \(m\),
since the two exponents of the different have constant sum.
This proves the interpolation formula, including the last block
when one selected column is the terminal column.
\end{proof}

\begin{theorem}[Summation over all supported block types]
\leavevmode\label{thm:pivotpacking619}
Fix the hyperplanes, transport and full marked index of
Theorem~\ref{thm:orientationmoment613}. Write
\(\mathcal W_C=\prod_{\tau,j}(1+|V_{j,\tau}|^2)^2\),
\(L=d(d+1)/2\), and let \(h\) be the length of the final block.
Put \(Z_0=\max(1,Z,Y_0)\), with \(Y_0\) the retained hyperplane
threshold of that theorem. For \(0<\delta<1\), \(\lambda>0\),
\(\sigma\ge0\), and
\(u=\sigma+2\lambda<h-\delta/2\), one has
\[
 \sum_{\substack{[C]\ \mathrm{with\ final\ length}\ h\\X_C\ge Z}}
       \frac{n_i(C)X_C^\sigma}{\mathcal W_C}
 \le
 2^{2d}|\mu(E)|\,C_\lambda^f B_\delta^{fd}\,
 Q^{-(4-\delta)fL}\,Z_0^{-(h-u-\delta/2)} .
\]
Here \(B_\delta\) is the integral in
Proposition~\ref{prop:openpacking615} and \(C_\lambda\) is as in
Lemma~\ref{lem:orientationtail613}.
\end{theorem}
\begin{proof}
Fix a composition with that final length. At each proper index
inside a block use
\(\prod_\tau(1+|V_{j,\tau}|^2)^2\ge B_j^4\).
If \(j=d-h+m\) is inside the last block, the proposition,
\(P_\gamma\ge1\), and interpolation give the stronger bound
\[
                  B_j^4\ge Q^{4jf} X_C^{2m/h}.
\]
All these interior indices together thus contribute their full
powers \(Q^{4jf}\) to the denominator, and the last block contributes
the additional factor \(X_C^{h-1}\).
Apply \(n_i(C)\le|\mu(E)|C_\lambda^f X_C^{2\lambda}\) from
Proposition~\ref{prop:rotationideal611}. For the endpoint sum we
are left with \(X_C^{v}\), \(v=u-h+1\). If \(v<0\), use
\(X_C^{v}\le Z_0^{v}\) and take \(v'=0\). Otherwise take \(v'=v\).
Thus the remaining numerator is \(Z_0^{v-v'}X_C^{v'}\), and
\(0\le v'<1-\delta/2\).

At a nonterminal block endpoint \(j\) use the normalized complex
coordinates of the proposition. Distinct determinant values are
separated in the product of the \(f\) complex distances by
\(Q^{jf}\), by the same ideal containment applied to their
difference. Fixing the earlier endpoint values gives the same
separation at \(j=d\), since \(D_d+D_s\) belongs to
\(Q^d(\overline{\mathfrak a_{[d]}}\mathfrak a_{[d]})^{-1}\)
and \(D_s\) is then fixed or zero. Each endpoint tuple accounts
for at most the product of its block lengths many classes by
Lemma~\ref{lem:forms203}, also for the two frequencies by
Theorem~\ref{thm:smallframes590}.

For clarity put \(P=Q^j\), and give the \(f\) discs radii
\(r_\tau>0\) with geometric mean \(R\). A nonempty box has
\(R\ge P\). After rescaling the radii to one, disjoint polydiscs of
radius \(P/(2R)\) give a count at most \(9^f(R/P)^{2f}\).
Integration against \(\prod_\tau4r_\tau(1+r_\tau^2)^{-3}dr_\tau\)
inserts the reciprocal weights at this endpoint. At every
nonterminal endpoint use \(R/P\ge1\) to replace \(2f\) by
\((4-\delta)f\). The resulting bound is
\(B_\delta^f Q^{-(4-\delta)jf}\).

At the final endpoint \(R/P\ge Z_0^{1/(2f)}\), and
\(X_C^{v'}\le P^{-2fv'}\prod_\tau r_\tau^{2v'}\) in that box.
Multiply the point count by this bound. In the count increase
the exponent \(2f\) to \((4-2v'-\delta)f\), extracting
\(Z_0^{-(1-v'-\delta/2)}\). After insertion of the numerator,
the one-dimensional integral has the same integrand
\(4r^{5-\delta}(1+r^2)^{-3}\) as before. Thus the final endpoint
costs \(B_\delta^f Q^{-(4-\delta)df}\) times that power of \(Z_0\).
Multiplying by \(Z_0^{v-v'}\) gives the claimed exponent.
All integrations have nonnegative integrands. Sum from the last
endpoint backwards using the conditional bounds. Interior indices
supply an exponent \(4\) in place of \(4-\delta\).
There are at most \(2^{d-1}\) compositions and the product of the
lengths of their blocks is at most \(2^d\). Since \(B_\delta\ge1\)
(integrate over \(r\ge1\) using \((1+r^2)^3\le8r^6\)),
the displayed bound follows.
\end{proof}

We next give a direct estimate on the original integral terms.
For \(b=0,1\) and an argument \(y\) on its real row define
\[
 T_{b,j}(y)=\prod_\tau
 \frac{\|\bigwedge_{i=1}^j y_\tau^{-1}Y_{b,i,\tau}\|}
      {\eta_{[j]}},\qquad T_{b,0}=1 .
\]
Use the common compatible positive reference norm on exterior
powers. Each prefix wedge is fixed by its decorated stabilizer,
so these are well-defined functions on the row quotient. Choose
\(p_j\ge0\), and let \(H_b(j;y)\) be the value at \(j\) of the
least concave majorant on \([0,d]\) of the data
\(\log T_{b,e}(y)\), \(0\le e\le d\). Put
\[
  A_b(y)=\exp\left(\sum_{j=1}^d p_j H_b(j;y)\right),\qquad
  \mathcal E_b=\int_{N_b\backslash G_\infty}
          (V'_b(y))^\nu A_b(y)^2\|W_b(y)\|^2\,dy .
\]
Here \(N_b\) is the integral row unipotent in
Theorem~\ref{thm:twoenergies592}. Use its height \(V'_b\), constants
\(D,C_*\), parameter \(\nu>4D\), and original quotient measure.
For a last block of length \(h\) define
\[
                 b_h=\sum_{m=1}^h (m/h)p_{d-h+m}.
\]

\begin{theorem}[A bound in the two original row measures]
\leavevmode\label{thm:pivotintegral619}
On any fixed pair of the retained rows, let \(\gamma\) range over
any subset of its relevant supported transverse integral double
cosets of positive Gram invariant. Let \(B_\gamma\) be its original
signed orbital operator including its finite coefficient and
normalized marked column as in Theorem~\ref{thm:twoenergies592}.
If the two displayed energies are finite, then
\begin{align*}
 \sum_\gamma (X_C P_\gamma)^{b_{h(\gamma)}/2}
                   \|B_\gamma\|_{\rm HS}
 &\le |c_{0,1}|\,C_*(1-2^{-(\nu-4D)})^{-f}\\
 &\qquad{}\cdot Q^{-f\sum_{j=1}^d j p_j}
                 (\mathcal E_0\mathcal E_1)^{1/2}.
\end{align*}
In particular \(p_j=2\) for every \(j\) gives the factor
\(Q^{-fd(d+1)}\) with moment \((X_C P_\gamma)^{(h(\gamma)+1)/2}\).
The same right side bounds the norm of the sum over any of these
terms after the corresponding weights have been inserted.
\end{theorem}
\begin{proof}
For a block endpoint \(e\) the pairing of the two prefix wedges,
evaluated at \(y\) and \(\gamma^{-1}y\), is \(D_e\) at each real
place. The Hermitian pairing on the exterior reference space has
operator norm one for its compatible positive norm. Therefore
\[
              B_e\le T_{0,e}(y)T_{1,e}(\gamma^{-1}y).
\]
This holds also for \(e=0\). At any \(j\), interpolate the logarithms
of the endpoint bounds according to its block. The linear
interpolant in each row is at most the corresponding concave
majorant. In the other direction, at every endpoint below \(d\)
we have \(\log B_e\ge fe\log Q\). At \(d\) we have the additional
term \(\tfrac12\log(X_C P_\gamma)\). Linear interpolation puts
the fraction \(m/h\) of that term at index \(d-h+m\).
Multiplication with the nonnegative weights \(p_j\) now gives,
pointwise on the actual integration domain,
\[
 Q^{f\sum_j j p_j}(X_C P_\gamma)^{b_h/2}
      \le A_0(y)A_1(\gamma^{-1}y).
\]

The Hilbert--Schmidt norm of the rank-one integrand is bounded by
the product of the section norms. The finite phase has absolute
value one, and each normalized marked column has norm at most one.
After extracting the displayed factor, enlarge the resulting
positive sum to all the integral double cosets of the pair using
\(F_b=A_b\|W_b\|\). For completeness put \(v_0=V'_0(y)\) and
\(v_1=V'_1(\gamma^{-1}y)\). Split \(F_0F_1\) as the product of
\(F_0(v_0/v_1)^{\nu/2}\) and \(F_1(v_1/v_0)^{\nu/2}\).
Cauchy--Schwarz on the disjoint integration domains bounds the
integral sum by the two square norms. Unfolding the first one
gives the original row integral of \(F_0^2v_0^\nu\) times the
translated sum of inverse heights \( (V'_1)^{-\nu}\).
That sum is bounded by \(C_*(1-2^{-(\nu-4D)})^{-f}\) by the box
count in Theorem~\ref{thm:twoenergies592}. Unfold the second
square with the rows interchanged. These operations use the same
integral stabilizers and Haar measures as the original signed
operators. They hold for the nonnegative measurable weights by
Tonelli's theorem. Taking square roots and restoring the pair
volume proves the inequality.
\end{proof}

The last estimate includes transverse terms with vanishing proper
leading minors at both terminal frequencies. To use it in
Section~\ref{sec:missing} one still needs quantitative values for the
two energies of the prescribed sections relative to their identity
coefficient, including \(C_*\), and the remaining period terms.
For the underlying leading-minor convention see
\srcloc{Formalism/SharpRestriction/CuspSpanNormBounds/BruhatCellsKronecker/CellGeo}{436}.


\subsection{Exterior moments at the two terminal frequencies}
\label{sec:exteriormoments620}

We evaluate the change of frequency in the two integrals of
Theorem~\ref{thm:pivotintegral619}. Throughout, use the same
retained fields, matrices, marked columns and Haar measures as
in that theorem. Write \(d=s+1\), \(n=2s+k+1\),
\(f=[F:\QQ]\), and \(R=|\operatorname{disc}F|^{1/f}\).
For the nonnegative parameters \(p_1,\ldots,p_d\), put
\[
                P=\sum_{j=1}^d j p_j,\qquad
                G(e)=\sum_{i=1}^{e}(d-i).
\]
For a row of terminal norm \(-T_b\), choose at each real place
the Witt identification sending its terminal vector to
\(\beta_\tau z\), where \(z\) is the middle unit negative vector
and \(|\beta_\tau|=\sqrt{T_b}\). Here \(T_b=1\) or \(T\).
The same identification sends its outer vectors to
\(e_1,\ldots,e_s\). The adapted positive norm gives \(e_i\)
length \(R^{d-i}\) and \(z\) length one. Thus the exterior
normalization of a prefix in Proposition~\ref{prop:pivotflags619}
is \(\eta_{[e]}=R^{G(e)}\).

\begin{proposition}[Transport of the two exterior moments]
\leavevmode\label{prop:exteriortransport620}
Use the translated height \(V'_b(y)=V_b(a_\beta y)\) in
Theorem~\ref{thm:pivotintegral619}, where \(a_\beta\) is the
terminal change of variables in Proposition~\ref{prop:scale151}.
Write \(A_1\) for the exterior weight of that row after replacing
its terminal vector by \(z\), and \(\mu_1\) for the probability
measure of its transported unit-frequency section. If
\(w_b^2=\int_{N_b\backslash G_\infty}\|W_b\|^2\), then
\[
 \frac{\mathcal E_b}{w_b^2}
       =T_b^{\,fP}\int_{B_b\backslash G_\infty}
                       V_b(g)^\nu A_1(g)^2\,d\mu_1(g).
\]
Here \(B_b\) is the full real decorated stabilizer and \(N_b\)
is the original integral unipotent. The identity also holds for
infinite integrals. Both moments in Theorem~\ref{thm:pivotintegral619}
are finite if \(4D<\nu\) and \(\nu+2P\le L_0=8(f+1)(n+1)^3\).
In particular \(p_j=2\) for all \(j\) and \(\nu=4D+1\) satisfy
these conditions on the retained intervals.
\end{proposition}
\begin{proof}
At a real place \(a_\beta\) multiplies all outer vectors by
\(\beta_\tau\), multiplies their duals by
\(\bar\beta_\tau^{-1}\), and fixes the middle space. With
\(y=a_\beta^{-1}g\), the vectors \(y^{-1}e_i\) and
\(y^{-1}(\beta_\tau z)\) are consequently \(\beta_\tau\) times
the corresponding unit-model vectors at \(g\). For every prefix
of order \(e\) the positive exterior norm gains
\(T_b^{e/2}\) at that place. Its ideal normalization is the
same for the two models. The logarithms of the global prefix
norms therefore gain the affine function
\(e\mapsto fe\log(T_b)/2\). Adding an affine function commutes
with taking the least concave majorant, by subtracting it from
each candidate majorant. The exterior weight at \(y\) is
\(T_b^{fP/2}A_1(g)\).

The section norm and all the weights are invariant under \(B_b\).
Passing from the integral unipotent quotient to the real
stabilizer quotient multiplies the numerator and \(w_b^2\) by
the same finite fibre volume. By Proposition~\ref{prop:scale151},
the squared section normalization cancels the Haar Jacobian
under \(y=a_\beta^{-1}g\). The two normalized measures are
thus related by pushforward, and \(V'_b(y)=V_b(g)\).
Integration proves the formula.

For finiteness write \(v_\tau(g_\tau)\) for the local factor of
\(V_b(g)\). A unit-model prefix wedge is fixed by the
decorated stabilizer and has adapted norm \(\eta_{[e]}\).
The norm of its image under the inverse of a representative is
at most \(\eta_{[e]}\) times the \(e\)-th power of the operator
norm of that inverse. Compatibility of the positive norms
with the Hermitian form identifies this operator norm with
that of the forward isometry. Take the infimum over the left
stabilizer and multiply over the places. The resulting prefix
bound is \(V_b(g)^e\). Its logarithm is affine in \(e\), hence
\(A_1(g)\le V_b(g)^P\). The unit section has finite moment of
order \(L_0\), since its weighted supremum times its unweighted
\(L^1\)-norm bounds that square integral by
Proposition~\ref{prop:weight163}. At this fixed row a change of adapted norm costs a
finite factor. Since \(V_b\ge1\), this proves finiteness in
the stated range. Finally \(D\le nd(d+1)/2\) and \(d\le n\);
thus \(4D+1+2d(d+1)\le L_0\).
\end{proof}

Let \(\mathcal P\) be a subdivision of the integer interval
\([0,d]\) with its two endpoints included. For each \(j\) let
\(\lambda_{j,e}\) be the coefficients of linear interpolation at
\(j\) between its two adjacent subdivision points \(e\).
At a subdivision point the sole coefficient is one. Define
\(\omega_e=\sum_{j=1}^d p_j\lambda_{j,e}\). In particular
\(\sum_e e\omega_e=P\). At one place in the unit-frequency
model put
\[
 t_{e,\tau}(g)=
 \frac{\|g^{-1}(Y_1\wedge\cdots\wedge Y_e)\|}{R^{G(e)}},
 \qquad Y_i=e_i\ (i<d),\quad Y_d=z,\quad t_{0,\tau}=1 .
\]
Use the local isometric embedding \(J_{1,\tau}\) for this
word and its actual positive compact column. For any unit vector
\(u\) of the local lowest compact type set
\[
 I_{b,\tau}(\mathcal P)=
 \int_{B_{b,\tau}\backslash G_\tau}
       v_\tau(g)^\nu\prod_{e\in\mathcal P}
       t_{e,\tau}(g)^{2\omega_e}\,
       \|(J_{1,\tau}u)(g)\|^2\,d\dot g .
\]

\begin{theorem}[A product of single-place integrals for each block type]
\leavevmode\label{thm:blockproduct620}
In the finite range of Proposition~\ref{prop:exteriortransport620}
these integrals are independent of \(u\). Between the maximum and
the sum of \(\prod_\tau I_{b,\tau}(\mathcal P)\) over the \(2^{d-1}\)
subdivisions lies the normalized value
\(\mathcal E_b/(w_b^2 T_b^{fP})\). More precisely, fix a subdivision
whose final block has length \(h\), and let \(\mathcal A\) be any
set of the supported positive-Gram transverse cosets with those
block endpoints. Then
\begin{align*}
 \frac{\sum_{\gamma\in\mathcal A}(X_C P_\gamma)^{b_h/2}
                 \|B_\gamma\|_{\rm HS}}
      {|c_{0,1}|w_0w_1}
 &\le C_*(1-2^{-(\nu-4D)})^{-f}
              \left(\frac{\sqrt{T_0T_1}}{Q}\right)^{fP}\\
 &\qquad{}\cdot\prod_\tau
          \sqrt{I_{0,\tau}(\mathcal P)I_{1,\tau}(\mathcal P)} .
\end{align*}
Here \(B_\gamma\), \(X_C\), \(P_\gamma\) and \(b_h\) are the
operators and weights of Theorem~\ref{thm:pivotintegral619}.
\end{theorem}
\begin{proof}
For fixed unit-model data at \(g\), join the logarithms of
the global prefix norms by line segments along \(\mathcal P\).
Every such polygon is bounded above by the least concave
majorant. The upper boundary of the convex hull of the
finitely many prefix points is itself one of these polygons.
It attains the majorant simultaneously at every \(j\).
Consequently \(A_1^2\) is the maximum over \(\mathcal P\) of
\(\prod_{\tau,e}t_{e,\tau}^{2\omega_e}\), since \(p_j\ge0\).

For one polygon the integrand weight factors over the real
places. Each local weight is right invariant under the
reference maximal compact group. Compress multiplication by
its truncation at a positive constant to the local isometric
image of the lowest compact type. This positive operator
commutes with that irreducible compact representation and
therefore is scalar. For a product of truncations the global
compression is the product of those scalars, since the real
embedding is a tensor product by Corollary~\ref{cor:rel155}.
This holds on all unit vectors of the product type. Let the
truncation increase and apply monotone convergence. Finiteness
also follows locally from the proof of the proposition.
For any subdivision the resulting integral is
\(\prod_\tau I_{b,\tau}(\mathcal P)\). Integrating the maximum
gives the claimed lower and upper bounds.

On cosets of the specified type, use the polygon on the actual
block endpoints on each of the two rows. At its endpoints the
cross determinant is the pairing of the corresponding prefix
wedges. Interpolation of Proposition~\ref{prop:pivotflags619}
gives the pointwise lower bound
\(Q^{fP}(X_C P_\gamma)^{b_h/2}\) for the product of the two
polygon weights, exactly as in Theorem~\ref{thm:pivotintegral619}.
The square-norm unfolding there applies to these weights on
their original row quotients. The translated-height count is
the same \(C_*(1-2^{-(\nu-4D)})^{-f}\). The affine frequency
shift in the proposition also holds for the polygon, because
\(\sum_e e\omega_e=P\). Substituting its product integrals gives
the displayed inequality, with the original finite columns and
their phases in each \(B_\gamma\).
\end{proof}

\begin{theorem}[Size of the absolute majorant at prime frequency]
\leavevmode\label{thm:exteriorcost620}
Put \(K_x=bn^4\) with \(b=\exp(x^\varepsilon)\), and
\(\delta_x=\exp(-c b n^2)\) for the constant of
Theorem~\ref{thm:localmass156}, decreased if necessary.
For each prescribed lowest-type section one has
\[
 \frac{\mathcal E_b}{w_b^2}\ \ge\
 (1-\delta_x)
     \left(\frac{T_b}{(RK_x)^{d-1}}\right)^{fP},\qquad
 \prod_\tau I_{b,\tau}(\mathcal P)\ \ge\
 (1-\delta_x)(RK_x)^{-2f\sum_e\omega_e G(e)} .
\]
Suppose that at least one row has terminal norm \(-T\).
After division by \(Z^{b_h/2}\), for any
\(1\le Z\le\exp(Mfd x^\varepsilon)\) with fixed \(M>0\), the
right side of Theorem~\ref{thm:blockproduct620} is at least
\[
 (1-\delta_x) C_*
        \exp\bigl((15-o(1)) fP d\,x^\varepsilon\bigr).
\]
For the right side of Theorem~\ref{thm:pivotintegral619},
normalized by \(|c_{0,1}|w_0w_1\) and divided by the same
factor, the coefficient \(15\) can be replaced by \(16\).
Both estimates are uniform in the nonnegative \(p_j\) in the
finite range. They concern the analytic upper bounds for the
period.
\end{theorem}
\begin{proof}
In the localized mass estimate of Theorem~\ref{thm:localmass156}
all successive Gram--Schmidt lengths at the prime frequency
satisfy \(a_{i,\tau}>\sqrt T K_x^{-(d-i)}\) for \(i<d\)
on the same set of measure at least \(1-\delta_x\) for the
normalized row. The last length is at least \(\sqrt T\),
since subtracting any earlier isotropic vector from the terminal
one preserves its negative Hermitian norm. Each length scales by
\(\sqrt T\) under the frequency transport in the proposition.
This gives a set of the same mass in the unit model on which the
lengths are at least \(K_x^{-(d-i)}\) at every place. The
unit model for an initially unit row also admits the prime
translate, by Corollary~\ref{cor:rel155}. Multiplying successive
lengths computes the positive prefix wedge exactly, by the
Gram determinant identity. Thus \(t_{e,\tau}\ge(RK_x)^{-G(e)}\)
simultaneously at all indices and places. Integrating the
polygon weight over this region with \(v_\tau\ge1\) gives
the second lower bound.

The concave majorant at \(j\) also dominates the chord from
zero to \(d\). Since \(G(d)=d(d-1)/2\), this gives
\(\log A_1\ge-fP(d-1)\log(RK_x)/2\) on the same region.
Integrate its square and apply the exact frequency change.
This proves the first lower bound.

For the two fixed polygons one has
\(\sum_e\omega_e G(e)\le(d-1)P\), since \(G(e)\le(d-1)e\).
Their product in the bound therefore costs at most
\(2(d-1)\log(RK_x)\) in the lower estimate for the logarithm
per \(fP\). For the concave majorants the cost is
\((d-1)\log(RK_x)\). Further, \(b_h\le P/d\) because at
each index \(j=d-h+m\) one has \(m/h\le j/d\).
The division by \(Z^{b_h/2}\) thus costs at most
\(M x^\varepsilon/2\) per \(fP\). For \(P=0\) the asserted
bounds follow directly from the total mass. For \(P>0\)
use \(\log(RK_x)=(1+o(1))x^\varepsilon\) since the conductor
and \(n\) are polynomially bounded, and
\(\log Q=(16+o(1))x^\varepsilon\).
Theorem~\ref{thm:prime209} with \(m=2d\) gives
\(\log\sqrt{T_0T_1}\ge(17-o(1))d x^\varepsilon\).
Now \(d\to\infty\) along the intervals. Substituting these
estimates gives the two stated constants.
\end{proof}

In particular the known lower cutoff \(Y_0\) of
Theorem~\ref{thm:orientationmoment613} lies in this range of
\(Z\) at prime frequency. At \(p_j=2\) both absolute bounds
grow even after extracting that available Gram weight. On a
repeated row its identity coefficient is
\(|c_{b,b}|w_b^2/\dim V_{\rho_b}\), by
Theorem~\ref{thm:flagisometry544}. Consequently normalization by
that coefficient multiplies the bound by \(\dim V_{\rho_b}\).
The signed orbital operators may cancel. Their inner products
before the passage to absolute values remain available in the
comparison of Section~\ref{sec:missing}.


\subsection{A compact integral for the full signed orbital sum}
\label{sec:signedkernel621}

Fix one retained pair of rows on the fields and intervals of
Section~\ref{sec:missing}. We keep both finite markings, both
terminal frequencies and the original normalized columns. Write
\(b=\exp(x^\varepsilon)\), and use \(l=0,1\) for the two rows.
As before \(n=2s+k+1\), \(d=s+1\), \(f=[F:\QQ]\) and
\(R=|\operatorname{disc}F|^{1/f}\). In each row make the real
Witt identification of Proposition~\ref{prop:exteriortransport620}.
In its standard signed coordinates put
\[
 w(g)=\|g\|_{\rm op},\qquad
 v^0_{l,\tau}(g)=\inf_{h\in B_{l,\tau}}w(hg),\qquad
 V_l^0=\prod_\tau v^0_{l,\tau}.
\]
We write \(g_l(y)=a_l y\) in these row coordinates for the change
to unit terminal frequency. Take \(A\ge1\). The cut on row \(l\)
retains \(y\) if \(v^0_{l,\tau}(g_l(y)_\tau)\le A\) at every
place. For an integral double coset \(\gamma\) let \(B_\gamma(A)\)
be its ordinary signed orbital operator with the two cuts imposed
at \(y,\gamma^{-1}y\), and let \(B_\gamma\) have the full domain.

\begin{proposition}[A uniform bound for the discarded part]
\leavevmode\label{prop:fulltail621}
Put \(q=4n^2\), \(D=\sum_{i=1}^d i(n-i+1)\),
\(\nu=4D+1\) and \(P_\nu=C_*(1-2^{-(\nu-4D)})^{-f}\) as in
Theorem~\ref{thm:twoenergies592}. For all sufficiently large
\(x\) on the common intervals one has \(\nu\le n^2\). There
is a constant \(C\) uniform in these data such that for every
set \(\mathcal A\) of supported integral cosets of the pair,
\begin{align*}
 \frac{\sum_{\gamma\in\mathcal A}
             \|B_\gamma-B_\gamma(A)\|_{\rm HS}}
      {|c_{0,1}|w_0w_1}
 &\le 2P_\nu \exp(Cfsq\log b) A^{-(q-\nu)/2},\\
 \frac{\sum_{\gamma\in\mathcal A}\|B_\gamma\|_{\rm HS}}
      {|c_{0,1}|w_0w_1}
 &\le P_\nu \exp(Cfs\nu\log b).
\end{align*}
The same second bound holds for the cut integrals. For the two
exterior moments of Theorem~\ref{thm:pivotintegral619} one also
has, whenever \(P=\sum_j j p_j\) and \(\nu+2P\le q\),
\[
       \mathcal E_l/w_l^2
       \ \le\ T_l^{fP}\exp\bigl(Cfs(\nu+2P)\log b\bigr).
\]
The constants here may depend on the fixed character and
\(\varepsilon\).
\end{proposition}
\begin{proof}
Let \(d\mu_l\) be the probability measure of the transported
unit-frequency row. Theorem~\ref{thm:moment286} gives
\(\int(V_l^0)^t\,d\mu_l\le\exp(Cfst\log b)\) for
\(0\le t\le q\), using H\"older's inequality for nonintegral
\(t\). This theorem works on the product lowest compact type.
It bounds the first positive coordinate moment and all the
successive ratios needed for the standard quotient height.
The first-coordinate statement in the source is at
\srcloc{Formalism/SharpRestriction/CuspSpanNormBounds/BesselShellMomentBounds}{1634}.
For the adapted height one has
\(V_l(g)\le R^{fs}V_l^0(g)\). Indeed the change of metric in
Proposition~\ref{prop:archdifferent555} normalizes the
stabilizer and has operator norm \(R^s\) at each place.
As \(\log R=O(\log x)=o(\log b)\), its cost fits the same
moment bound after increasing \(C\). Proposition~\ref{prop:exteriortransport620}
and \(A_1\le V_l^P\) from its proof give the exterior estimate.

Use the translated heights \(V'_l(y)=V_l(g_l(y))\). Under the
normalized row measure the event discarded on row \(l\) implies
\(V_l^0(g_l(y))> A\), since all its local factors are at least
one. Its weighted squared norm satisfies
\[
 w_l^{-2}\int_{\rm discarded}(V'_l)^\nu\|W_l\|^2
       \le R^{fs\nu}A^{-(q-\nu)}
                          \int(V_l^0)^q\,d\mu_l .
\]
Here \(\nu\le n^2\) eventually since \(D\le nd(d+1)/2\) and
\(d^2/n\to0\) on the retained intervals. In the square-norm
unfolding of Theorem~\ref{thm:twoenergies592}, insert the
discard indicator in the norm of row zero, leaving row one
unrestricted. The inverse-height count is still \(P_\nu\).
Cauchy--Schwarz now takes the square root of the last bound
times the unrestricted moment of order \(\nu\) on the other
row. Interchanging the rows bounds the other discarded region.
Their union accounts for the difference in the first inequality.
Omitting both indicators bounds the absolute integrals over all
cosets. This proves the other assertions, including the estimate
on every restricted sum.
\end{proof}

We next express the integrands without choosing the scale of a
Bessel functional. At a single real place put \(G=U(p_+,p_-)\),
\(p_++p_-=n\), \(B=N\rtimes C_0\) with \(C_0\) the definite
compact group, and give \(C_0\) probability Haar measure as in
Section~\ref{sec:relative155}. Write \(\chi\) for the unit-frequency
character on \(N\), \(\rho\) for the compact column representation,
\(\sigma\) for the row discrete series and \(d_\sigma\) for its
formal degree. On its lowest compact type use a unit vector \(u\)
and put \(c_u(t)=\langle\sigma(t)u,u\rangle\), with the inner
product linear in the first variable. If \(a\in B\), write
\(a=n_a c_a\) in this decomposition. Define the matrix
\[
 K(y,z)=d_\sigma\int_B c_u(y^{-1}a z)\,
                \overline{\chi(n_a)}\,\rho(c_a)^*\,da .
\]
For \(S\ge1\) write \(K_S\) for the same integral restricted
by \(w(y^{-1}a z)\le S\). Both matrices take values in
\(\operatorname{End}(E_\rho)\).

\begin{lemma}[The relative kernel and its uniform tail]
\leavevmode\label{lem:relkernel621}
For the isometric unit-frequency embedding \(J\),
\[
                  K(y,z)=(Ju)(z)(Ju)(y)^* .
\]
Put \(H=10^4 n^4(B_x+1)\) where \(B_x\ge1\) bounds the absolute
Harish--Chandra coordinates of the row. Put
\(\kappa=4n^2+2n+9\), and let \(v_B\) be the Haar measure in
\(B\) of \(\|a-1\|_{\rm op}\le1/2\).
If \(v^0(y),v^0(z)\le A\), for an absolute \(C\) the integral
of the operator norm of the kernel integrand, including its
factor \(d_\sigma\), is at most
\[
       D_B(A)=d_\sigma v_B H^{Cn^3} A^{8n^2}.
\]
On the part defining \(K-K_S\) the bound is
\(D_B(A) S^{-\kappa}\). The same estimates apply to the
absolute scalar coefficient before the unitary column factor.
Moreover
\[
    D_B(A)\operatorname{vol}_{B\backslash G}\{v^0\le A\}
                        \le H^{Cn^3}A^{12n^2}.
\]
\end{lemma}
\begin{proof}
On \(\sigma\boxtimes\rho^\vee\) the normalized relative form
of Theorem~\ref{thm:rel155} equals
\(\ell(v)\overline{\ell(v')}\), with
\(\ell(w\otimes\lambda)=\lambda(\Lambda w)\) and
\((Ju)(g)=\sqrt{d_\sigma}\Lambda(\sigma(g)u)\).
The hypotheses are the retained discrete series, compact column,
unitary generic character and the one-dimensional Bessel space
in Corollary~\ref{cor:rel155}. The uniqueness theorem for
these real Bessel data is
\href{https://arxiv.org/abs/0908.1728}{Jiang, Sun and Zhu,
Theorem A}. The relative form used here has the normalization
deduced in Section~\ref{sec:relative155} from
\href{https://arxiv.org/abs/1506.01452}{Beuzart-Plessis,
Sections 7.1--7.2}. It is integration
of its matrix coefficient over \((n_ac_a,c_a)\) against
\(\overline{\chi(n_a)}\). Apply it to
\(\sigma(z)u\otimes e_i^*\) and \(\sigma(y)u\otimes e_j^*\)
for an orthonormal dual basis. The row coefficient is
\(c_u(y^{-1}az)\), and the dual compact coefficient is the
\((i,j)\)-entry of \(\rho(c_a)^*\). Reading all entries and
multiplying by \(d_\sigma\) gives the first identity. It holds
on these smooth translates by the continuity of the relative
form. Absolute convergence also follows from the estimate below.

The gap choice in Lemma~\ref{lem:translate163} exceeds
\(24n^2+2\) for the present rows at sufficiently large \(x\).
Thus Theorem~\ref{thm:tensor249} bounds the matrix coefficient
by \(H^{Cn^3} w(t)^{-L}\) with \(L=8n^2+2n+9\). This bound
also holds for the operator norm of the compression
of \(\sigma(t)\) to its lowest compact type, since that theorem
allows any two unit vectors in the type.

For \(U\ge1\), a maximal \(1/(2U)\)-separated set in
\(\{a\in B:w(a)\le U\}\), in ambient matrix operator norm, has
at most \((1+4U^2)^{2n^2}\) elements. Compare the volumes of
disjoint norm balls of half this separation in the real space
of all complex \(n\)-by-\(n\) matrices. Each \(a\) in the
original ball lies in a left translate of the set measured by
\(v_B\), since \(w(a^{-1})=w(a)\) on \(G\). The Haar measure
of the ball is consequently at most
\(5^{2n^2} U^{4n^2}v_B\).

Use representatives with \(w(y),w(z)\le2A\).
Independent changes on the left by \(B\) are changes of variable
in this integral and supply only unitary factors on the two
columns. The cut on \(y^{-1}az\) is preserved in this substitution,
and the Haar measure of \(B\) is bi-invariant. On the shell
\(2^j S\le w(y^{-1}az)\le2^{j+1}S\) one has
\(w(a)\le8 A^2 2^j S\). The Haar bound times the coefficient
bound is at most \(D_B(A)(2^j S)^{4n^2-L}\) after enlarging its
absolute constant. Sum the geometric series. Taking \(S=1\)
also gives the bound on the entire integral, including its
boundary. Finally any coset with \(v^0\le A\) has a representative
of norm at most \(2A\). Multiplication by any element in the
set measured by \(v_B\) gives norm at most \(3A\).
Quotient integration yields
\[
 v_B\operatorname{vol}_{B\backslash G}\{v^0\le A\}
          \le\operatorname{vol}_G\{w\le3A\}.
\]
Use the formal-degree volume estimate of
Theorem~\ref{thm:tensor249} and absorb its dimension constant
in \(H^{Cn^3}\). This proves the last assertion.
\end{proof}

For row \(l\) at all real places use \(K_l\) for the product-group
version of this integral. Its coefficient is
\(\langle\sigma_\infty(t)u_l,u_l\rangle\) for the unit row vector
in the product lowest type. The Haar measures and formal degrees
are multiplied, as are the characters, and the compact column
is a tensor product. In \(K_{l,S}\) impose the cutoff at every
place. The local identity gives the global one for pure tensors
and then for arbitrary \(u_l\) by polarization in each local
coefficient and finite expansion in tensor bases.
Put \(v_l=\operatorname{vol}(N_l\backslash B_l)\), with \(N_l\)
the integral unipotent, and
\(J_l=T_l^{fs(n-s)}\). For physical real arguments \(y,z\) define
\[
 k_{l,S}(y,z)=\frac{w_l^2}{v_l J_l}
       \operatorname{tr}K_{l,S}\bigl(g_l(y),g_l(z)\bigr).
\]
Write \(k_l\) for the unrestricted value. From the change of
variables in Proposition~\ref{prop:scale151} and the lemma one
gets \(k_l(y,z)=W_l(y)^*W_l(z)\). Indeed \(w_l^2/v_l\) is the
squared norm of the physical section on its real stabilizer
quotient. Its unit model is isometric, and its two point values
in the physical model supply the additional factor \(J_l^{-1}\).

Let \(\Gamma_\gamma=N_0\cap\gamma N_1\gamma^{-1}\) and let
\(Y_\gamma(A)\) denote the part of
\(\Gamma_\gamma\backslash G_\infty\) satisfying the two cuts.
Write \(\theta_\gamma=\epsilon_\gamma h_{i(\gamma)}\) for the
original finite phase times its normalized marked column, so
\(\|\theta_\gamma\|\le1\). Define the scalar
\begin{align*}
 \mathcal T(A,S)
 &=\frac{1}{w_0^2w_1^2}
      \sum_{\gamma,\eta\in\mathcal A}
      \theta_\gamma^*\theta_\eta
      \int_{Y_\gamma(A)}\int_{Y_\eta(A)}
            k_{0,S}(y,z)\\[-2pt]
 &\hspace{36mm}{}\cdot
            k_{1,S}(\eta^{-1}z,\gamma^{-1}y)\,dz\,dy .
\end{align*}
With original lifts, the product of the two local covariances
is the orbital covariance of the signed term. In particular
these integrals use the same arithmetic stabilizers as the full
period. Truncation respects that covariance by the change of
variable in the proof of the lemma.

\begin{theorem}[Approximation of the full signed square on the common fields]
\leavevmode\label{thm:signedcompact621}
For each fixed \(\theta>0\) take \(A=\exp(b^\theta)\) and
\(S=\exp(b^{2\theta})\). At every sufficiently large \(x\)
on the original common intervals, for either terminal norm
on either end and every set \(\mathcal A\) as above,
\[
 \left|
 \frac{\|\sum_{\gamma\in\mathcal A} B_\gamma\|_{\rm HS}^2}
      {|c_{0,1}|^2w_0^2w_1^2}
             -\mathcal T(A,S)\right|\le \exp(-b^\theta).
\]
The sum defining \(\mathcal T(A,S)\) uses finitely many integral
cosets and compact integration ranges modulo their indicated
stabilizers. The finite Gram factors \(\theta_\gamma^*\theta_\eta\)
are retained, including those for distinct cosets at the same
marked index.
\end{theorem}
\begin{proof}
The original operator is the row-pair volume \(c_{0,1}\) times
the finite column tensored with the integral of
\(W_0(y)W_1(\gamma^{-1}y)^*\), by
Theorem~\ref{thm:mixedordinary541}. Taking the Hilbert--Schmidt
inner product of two such integrands gives exactly the indicated
product of the two scalar kernels and the finite inner product.
Thus the formula with the two row cuts but with unrestricted
kernels gives the normalized cut square. For the unrestricted
row domains the resulting double integrals are absolutely
summable by Cauchy--Schwarz on the point values and the
absolute-integral estimate in Proposition~\ref{prop:fulltail621}.
The stabilizer integrations in the kernels are performed first.

For fixed \(A\) each row's cut has representatives in a compact
set modulo \(N_l\). Indeed standard representatives in each
real quotient have norm at most \(2A\), and
\(N_l\backslash B_l\) is compact. Reversing the frequency
change preserves compactness of the chosen sets. Choose compact
representative sets \(C_0,C_1\) covering the two cuts. A coset
contributing to the cut has a representative in
\(\Gamma\cap C_0 C_1^{-1}\), which is finite by discreteness.
For each of these representatives the allowed \(y\)'s lie,
modulo \(\Gamma_\gamma\), in finitely many translates of \(C_0\).
In fact two expressions \(n_0^{-1}\gamma n_1\) giving the
same representative have their target multipliers in the
same left \(\Gamma_\gamma\)-coset. The closed cuts are therefore
compact on the quotient. The bound on \(y^{-1}az\) similarly
gives a compact subset of the real stabilizer for fixed lifts.

To track the total measure, set \(A'=R^{fs}A^f\).
On either cut the translated height of that row is at most
\(A'\). The counting bound \(M(A')\) of
Theorem~\ref{thm:twoenergies592} and the unfolding there give,
for either \(l\),
\[
 \sum_{\gamma\in\mathcal A}\operatorname{vol}Y_\gamma(A)
 \le M(A')\,v_l J_l\prod_\tau
           \operatorname{vol}_{B_{l,\tau}\backslash G_\tau}
                         \{v^0_{l,\tau}\le A\}.
\]
Indeed integrate the cut on one row and count all admissible
translates of the other. Substitution on the first real row
supplies \(J_l\) and integrating along its lattice fibre
supplies \(v_l\).

Take a common \(H\) for all places on the two rows. On a pair
of row arguments satisfying the cut the absolute full integral
or cut integral for the trace of \(K_l\) is at most
\(d_{\rho_l}\prod_\tau D_{B_{l,\tau}}(A)\), where
\(d_{\rho_l}\) is the product column dimension. Removing at
least one local tail costs at most \(f S^{-\kappa}\) times
this bound. To see this also for a vector which is not a pure
tensor, bound its coefficient by the operator norm of the
compression to the product lowest type. This operator is the
tensor product of the local compressions. Their individual
bounds in the lemma then dominate the integrand pointwise.
The trace of a unitary column factor has modulus at most its
dimension.

In the product of the two kernels, subtraction of the cut
version costs at most \(2f S^{-\kappa}\) times the product of
these majorants, with the two factors \(w_l^2/(v_lJ_l)\)
restored. Use \(|\theta_\gamma^*\theta_\eta|\le1\) for this
error. The square of the total domain volume is at most the
product of its bounds from the two rows. Inserting the last
inequality of Lemma~\ref{lem:relkernel621} at every place
cancels the Haar normalizations. The error in the normalized
cut square is at most
\[
     2 f\,d_{\rho_0}d_{\rho_1} M(A')^2
                    H^{C f n^3} A^{24 f n^2} S^{-\kappa}.
\]
The parameters \(f,n,s,\log H,\log C_*\) grow at most
polynomially in \(x\) on these intervals. The same applies
to \(\log d_{\rho_l}\) by the Weyl dimension formula and the
retained weight bounds. From the expression for \(M\) in
Theorem~\ref{thm:twoenergies592}, the logarithm of the entire
prefactor of \(S^{-\kappa}\) is bounded by a polynomial in
\(x\) times \(1+\log A\). With the chosen \(S,A\) this error
is at most \(\frac12\exp(-b^\theta)\) for sufficiently large
\(x\).

Finally the normalized squared norms before and after cutting
the rows differ by at most the product of the normalized
difference norm and the sum of the two normalized norms.
Apply both bounds of Proposition~\ref{prop:fulltail621}.
Their prefactors have logarithm polynomially bounded in \(x\),
whereas \((q-\nu)\log A/2\ge 3n^2 b^\theta/2\).
This error is also at most \(\frac12\exp(-b^\theta)\).
Addition proves the result.
\end{proof}

The formula controls the error of truncation even for the
uncut signed square with both terminal norms equal to \(T\).
To derive \eqref{eq:gap} one must still estimate the signed
compact integral itself after the required projection and
compare the resulting period with its prescribed integral
lattice. The independent bound \eqref{eq:length12} on the same
intervals remains necessary in the present reduction.


\subsection{The second moment of the signed Gram blocks}
\label{sec:fourkernels622}

Section~\ref{sec:geom11} uses the Hilbert--Schmidt squares of
off-diagonal Gram blocks as well as the deviation of a diagonal
block from the identity. We derive compact formulas for both
statistics of the ordinary marked operator. Each row-coordinate
cut has radius at most a fixed power of \(T\) in the transported
unit model. We give a separate cutoff for the stabilizer kernel.

\begin{proposition}[A simultaneous cut at all real places]
\leavevmode\label{prop:powercut622}
Keep the family of Proposition~\ref{prop:fulltail621}. For any
fixed \(L>0\) there is \(a>0\) such that, with \(A=b^{as}\),
\[
 \frac{\sum_{\gamma}\|B_\gamma-B_\gamma(A)\|_{\rm HS}}
      {|c_{0,1}|w_0w_1}\le b^{-Lsn^2}.
\]
The sum of the uncut absolute integrals and the sum of the cut
absolute integrals, in this same normalization, are each at
most \(b^{o(sn^2)}\). The estimates apply simultaneously to
all retained row pairs and their marked indices on the common
intervals. The constant \(a\) is independent of \(f,n,s,x\).
\end{proposition}
\begin{proof}
Improve the stabilizer factor in the counting proof of
Theorem~\ref{thm:twoenergies592}. An integral stabilizer of a
decorated exterior tuple modulo its integral unipotent injects
into the isometry group of the underlying positive lattice,
of rational dimension \(m=2fk\). That finite group has order
at most
\[
         M_m=\prod_{p\le m+1}
          p^{\,\sum_{j\ge0}\lfloor m/(p^j(p-1))\rfloor}
\]
by Minkowski's theorem on finite subgroups of
\(\mathrm{GL}_m(\QQ)\)
\cite[Lecture I, Theorem 1]{XXII-XVII-IV-serre211}. Its \(p\)-exponent is
at most \(mp/(p-1)^2\). Summing over integers up to \(m+1\)
gives \(\log M_m\ll m\log^2(2m)\). In the proof of
Theorem~\ref{thm:mixedhalf556} each tuple is counted with at
most this stabilizer order. The exterior-lattice box count and
its height exponent are unchanged. Hence we may use
\[
 C' = M_m\prod_{i=1}^{d}(1+2\cdot4^i T^i)^{2f(n-i+1)}
\]
in both \(P_\nu\) and \(M(U)\) of
Theorem~\ref{thm:twoenergies592}. Denote the resulting bounds
by \(P'_\nu\) and \(M'(U)\).

Here \(d=s+1=m/2\asymp x^{1/2-\varepsilon/16}\) by
Corollary~\ref{cor:actual50},
\(f\ll_\psi x^{2+8\varepsilon}\) from the field construction
in Proposition~\ref{prop:field12}, and
\(n\ge k\asymp x^{9/2+10\varepsilon}\) as in
Theorem~\ref{thm:bothordinary553}. In particular
\(fd^2/n\to0\). Since \(\log T=O(d\log b)\),
\[
 \log C'\ll fk\log^2(2fk)+fnd^3\log b=o(sn^2\log b).
\]
With \(\nu=4D+1=O(nd^2)\) also \(fs\nu=o(sn^2)\).
The unrestricted square-norm estimate of
Proposition~\ref{prop:fulltail621} with \(P'_\nu\) proves the
asserted absolute bound.

For the tail fix one real place \(\tau\), and use \(q=4n^2\).
On the event \(v_\tau^0>A\), a single factor of
\((V_l^0)^\nu\) can be bounded by
\(A^{-(q-\nu)}(v_\tau^0)^q\). At the other places use the
moment of order \(\nu\). The local estimate in the proof of
Theorem~\ref{thm:moment286} is \(b^{Cst}\) at every place
for \(1\le t\le q\). Compression of these right compact
invariant weights gives their product on the whole product
lowest type, as in Theorem~\ref{thm:blockproduct620}.
After summing over \(\tau\), the normalized adapted moment
on the discarded set is at most
\[
       f R^{fs\nu} A^{-(q-\nu)} b^{Cs(q+(f-1)\nu)}.
\]
On the other row the unrestricted moment is at most
\(R^{fs\nu}b^{Cfs\nu}\).
Take the geometric mean, multiply by \(P'_\nu\), and add
the estimate with the rows exchanged. Since
\(\nu/n^2\to0\) and \(\log R=O(\log x)\), the logarithm of
this bound divided by \(s n^2\log b\) is at most
\(-2a+O(1)+o(1)\) for fixed \(a\). Choose \(a\) large enough
for the required \(L\). Finally \(\log T=(34+o(1))d\log b\);
thus \(A\le T^{O_L(1)}\) for all sufficiently large \(x\).
\end{proof}

Consider a finite collection of source and target rows as in
Theorem~\ref{thm:mixedordinary541}, with the same \(E,n,s,k\)
at each \(x\). Let \(Z_{ri}\) be the marked block from source
\(i\) to target \(r\), and \(Z_i\) the map to the direct sum
of the targets. One may keep any specified subset of the
supported double cosets in each block. Write
\(\lambda_{ri}=|c_{ri}|w_rw_i\), with zero for pairs in
different lattice isometry classes. For a compatible pair use
the actual coefficient \(c_{ri}\), including its phase if any.

Define a scalar \(\mathcal U_{rtij}(A,S)\) for the four ordered
blocks \((r,i),(r,j),(t,j),(t,i)\) as follows. Take a coset
\(\gamma_h\) of each block, in this order, and use
\(y_h\in Y_{\gamma_h}(A)\), \(z_h=\gamma_h^{-1}y_h\).
The fiber coordinates are carried back to their original rows
before evaluating a kernel. Let \(\theta_h\) be the original
normalized finite column with its phase for this coset. On the
product of the four cut domains integrate
\[
 (\theta_1^*\theta_2)(\theta_3^*\theta_4)\,
 k_{r,S}(y_1,y_2)k_{j,S}(z_2,z_3)
 k_{t,S}(y_3,y_4)k_{i,S}(z_4,z_1).
\]
Sum these integrals and multiply by
\(\overline{c_{ri}}c_{rj}\overline{c_{tj}}c_{ti}\).
If one of the blocks is structurally zero set
\(\mathcal U_{rtij}=0\). Also write \(\mathcal U_{ri}^{(2)}\)
for \(\lambda_{ri}^2\mathcal T(A,S)\) of
Theorem~\ref{thm:signedcompact621} applied to that single
block. These formulas use only the finite data, local
discrete-series coefficients and the Haar normalizations
specified in Section~\ref{sec:signedkernel621}.

\begin{theorem}[Mixed Gram moments on the original signed rows]
\leavevmode\label{thm:fourkernels622}
For each fixed \(L>0\) one can choose constants \(a,c>0\)
so that \(A=b^{as}\), \(S=b^{cfn}\) give, with
\(\delta=b^{-Lsn^2}\),
\begin{align*}
 \left|\|Z_{ri}\|_{\rm HS}^2-\mathcal U_{ri}^{(2)}\right|
   &\le \delta\lambda_{ri}^2,\\
 \big|\operatorname{tr}(Z_{ri}^*Z_{rj}Z_{tj}^*Z_{ti})
           -\mathcal U_{rtij}\big|
   &\le\delta\lambda_{ri}\lambda_{rj}\lambda_{tj}\lambda_{ti}.
\end{align*}
Consequently the original off-diagonal or diagonal Gram block
\(K_{ij}=Z_i^*Z_j\) satisfies
\[
   \left|\|K_{ij}\|_{\rm HS}^2-\sum_{r,t}\mathcal U_{rtij}\right|
       \le\delta\left(\sum_r\lambda_{ri}\lambda_{rj}\right)^2 .
\]
All estimates are uniform on the same intervals and in all
retained subsets. They apply also when indices in the four
ordered blocks coincide.
\end{theorem}
\begin{proof}
For uncut kernels and cut row arguments, take tensor products
of the original marked column with the rank-one operator
\(W_{\rm target}(y_h)W_{\rm source}(z_h)^*\).
The adjoints on factors one and three give precisely the
two finite inner products. Multiplying the four matrices and
taking the trace contracts each of the four repeated real
rows with itself. The contractions in cyclic order are
\(W_r(y_1)^*W_r(y_2)\), \(W_j(z_2)^*W_j(z_3)\),
\(W_t(y_3)^*W_t(y_4)\) and \(W_i(z_4)^*W_i(z_1)\).
They are the exact kernels \(k_l\) of
Section~\ref{sec:signedkernel621}. This proves the cut-domain
formula with full kernels and retains the phases of distinct
cosets. The absolute integrability follows from the product
of the four absolute-integral bounds in
Proposition~\ref{prop:powercut622}.

First changing the four operators to their cut-domain
versions costs at most \(b^{-(L+1)sn^2}\) times the four
pair factors for large \(x\), on taking \(a\) from that
proposition with \(L+3\). Indeed subtract one factor at a
time and use
\(|\operatorname{tr}(D_1^*D_2D_3^*D_4)|
       \le\prod_h\|D_h\|_{\rm HS}\).
This follows from the Hilbert--Schmidt Cauchy--Schwarz
inequality and \(\|D_h\|_{\rm op}\le\|D_h\|_{\rm HS}\).
The other three norms are at most their pair factors times
\(b^{o(sn^2)}\).

Next truncate the four relative kernels at \(S\).
There are at most \(4f\) choices of a local tail when the
product is subtracted one row at a time. By
Lemma~\ref{lem:relkernel621} its contribution in that place
costs \(S^{-\kappa}\) relative to the product of the absolute
kernel majorants. For each edge \(u,v\), the sum of its cut
domain volumes has upper bound
\[
 M'(A')\left(
   v_uJ_u\, v_vJ_v
      \prod_{\tau}\operatorname{vol}X_{u,\tau}
                      \operatorname{vol}X_{v,\tau}\right)^{1/2},
 \quad X_{u,\tau}=\{v^0_{u,\tau}\le A\},\quad A'=R^{fs}A^f .
\]
This is the geometric mean of the two unfolding bounds in the
proof of Theorem~\ref{thm:signedcompact621}, with the improved
stabilizer factor. Multiply these four edge bounds. Each
occurrence of a row in the cycle is incident to two edges,
so the volume of its local coset ball occurs once, and its
\(v_uJ_u\) cancels the denominator from its kernel.
The four \(w_u^2\)'s and four absolute pair coefficients
give precisely the four \(\lambda\)'s. Thus after division by
their product the truncation error is at most
\[
  4 f\, d_{\rho_r}d_{\rho_t}d_{\rho_i}d_{\rho_j}\,
       M'(A')^4 H^{Cfn^3} A^{48 f n^2}S^{-\kappa}.
\]
The exponent of \(A\) is the sum of the four exponents
\(12fn^2\) from the integrated kernel bound. Repeated rows
give repeated factors in these products.

For the present choice of \(A\), \(\log A'=O_a(fs\log b)\).
In \(M'(A')\) the box exponent costs \(O_a(Dfs\log b)\).
The binomial factor costs \(O_a(f\log(2+s\log b))\),
from \(\binom{N+f}{f}\le(e(1+N/f))^f\).
Together with the bound on \(C'\) this gives
\(\log M'(A')=o(sn^2\log b)\).
The retained weight bound gives \(\log H=O(\log b)\) and
\(\log d_{\rho_u}=O(fn^2\log b)\) by the compact Weyl
dimension formula. With \(\kappa\ge4n^2\), taking \(c\)
sufficiently large therefore bounds the last error by
\(b^{-(L+1)sn^2}\). The same argument for two kernels has
two edge factors and exponent \(24fn^2\) in \(A\), as in
Theorem~\ref{thm:signedcompact621}. This proves both
individual estimates after adding the row-cut errors.
Finally
\(\|K_{ij}\|_{\rm HS}^2
 =\sum_{r,t}\operatorname{tr}(Z_{ri}^*Z_{rj}Z_{tj}^*Z_{ti})\).
Summing the error bounds proves the assertion.
\end{proof}

\begin{corollary}[Error in the diagonal rank statistic]
\leavevmode\label{cor:gramerror622}
Let the dimension of source \(i\) be \(r_i\).
Set \(F_i=\sum_r\lambda_{ri}^2\). With the same cuts one has
\[
 \left|\|K_{ii}-I\|_{\rm HS}^2-
  \left(\sum_{r,t}\mathcal U_{rtii}
           -2\sum_r\mathcal U_{ri}^{(2)}+r_i\right)\right|
     \le\delta(F_i^2+2F_i).
\]
\end{corollary}
\begin{proof}
The block \(K_{ii}\) is positive self-adjoint. Expand the square
as \(\|K_{ii}\|_{\rm HS}^2-2\operatorname{tr}K_{ii}+r_i\)
and use \(\operatorname{tr}K_{ii}=\sum_r\|Z_{ri}\|_{\rm HS}^2\).
The two estimates in the theorem give the stated errors.
\end{proof}

For the ordinary projected blocks, both signed statistics used
in Section~\ref{sec:geom11} can therefore be tested with finite
coset sets and compact real domains, at the stated row-pair
scales. The inequalities for the values of these integrals,
and the comparison with the integral theta maps on those same
scales, are still required in Section~\ref{sec:missing}.


\subsection{Comparison ranks in the full ordinary Gram moments}
\label{sec:markedranks623}

We estimate a portion of the signed integrals of
Section~\ref{sec:fourkernels622} on the original fields. The estimate
uses the ranks of the differences of the two actual finite markings.
It also applies to a choice of integral coordinates in the same finite
projection. Both terminal frequencies are allowed.

Write \(f=[E_x^+:\mathbb Q]\), \(m=2d=2s+2\), \(n=2s+k+1\), and
\(b=\exp(x^\varepsilon)\). The parameters in
Corollary~\ref{cor:actual50} give
\(d\asymp x^{1/2-\varepsilon/16}\), \(k\asymp x^{9/2+10\varepsilon}\),
and \(f\ll_\psi x^{2+8\varepsilon}\). In particular \(n/k\to1\).
Theorem~\ref{thm:prime209} gives \(\log T=(34+o(1))d\log b\).
In Section~\ref{sec:fourkernels622} the cut of fixed power in \(T\)
is the cut on each transported row coordinate. The separate
integration over its stabilizer uses the radius \(b^{cfn}\).

For each ordinary block write its signed orbital expansion as
\[
             Z_{ri}=\sum_{\gamma}h_{a_\gamma}\otimes P_\gamma^{ri}.
\]
Here \(h_a\) is the unit column of the prescribed finite matrix
at the index \(a\), and \(P_\gamma^{ri}\) includes the finite phase,
the real orbital integral, and its pair coefficient \(c_{ri}\).
Let \(D_i\) be the complex dimension of row \(i\)'s positive column
at infinity. With \(w_i\) and \(\lambda_{ri}=|c_{ri}|w_rw_i\) as in
Section~\ref{sec:fourkernels622}, put
\[
  \sigma_{ri}=\frac{\lambda_{ri}}{\sqrt{D_rD_i}},\qquad
  U_x=fk\{d^3\log b+k\log(2x)\}.
\]
On a repeated row \(\sigma_{ii}\) is the modulus of the actual
identity coefficient, by Theorem~\ref{thm:flagisometry544}.
For other pairs this definition only uses their own coefficient
and dimensions. Structurally zero pairs have \(\sigma_{ri}=0\).

\begin{lemma}[Absolute mass at the identity normalization]
\leavevmode\label{lem:massranks623}
There is a constant \(C\), independent of the retained rows and of
\(x\) on the common intervals, such that
\[
       \sum_\gamma \|P_\gamma^{ri}\|_{\rm HS}
                    \le \sigma_{ri}\exp(CU_x).
\]
This bound also holds if each norm of an integral is replaced by
the integral of the Hilbert--Schmidt norm of its integrand,
including the absolute pair coefficient.
\end{lemma}
\begin{proof}
The absolute unfolding in Proposition~\ref{prop:fulltail621},
with the improved counting factor of Proposition~\ref{prop:powercut622},
bounds the latter sum by \(\lambda_{ri}P'_\nu b^{Cfs\nu}\).
Here \(\nu=4\sum_{j=1}^d j(n-j+1)+1=O(nd^2)\) and
\[
  \log P'_\nu\ll fk\log^2(2fk)+fnd^3\log b+f .
\]
Indeed the finite stabilizer provides the first term, the \(d\)
lattice factors have logarithm \(O(fn d^2\log T)\), and the
dyadic series contributes \(f\log2\). As \(f,k\) are polynomial
in \(x\), the sum of these costs and \(Cfs\nu\log b\) is
\(O(fk d^3\log b)\).

For the dimensions use the actual column parameters in the proof
of Lemma~\ref{lem:weyl209}. At one real place the Weyl product
is \(\prod_{a<b'}(u_a-u_{b'})/(b'-a)\). Each inner entry
\(u_2,\ldots,u_{k-1}\) lies between successive \(v_j\), all
polynomial in \(x\) with a common bound. A factor involving
an endpoint is at most \(3\lceil xb\rceil\) by the span bound
there. There are at most \(2k-3\) such factors and at most
\(k(k-1)/2\) inner factors. Multiplying over the \(f\) real
places gives
\[
          \log D_i\ll f\{k^2\log(2x)+k\log b\}.
\]
This estimate absorbs both dimensions in
\(\lambda_{ri}=\sigma_{ri}\sqrt{D_rD_i}\) at the asserted cost.
The unfolding bounds integrals of products of the two section
norms, so it gives the asserted stronger form before integration.
\end{proof}

Let \(G=\prod_{v\mid g_1g_2}\mathrm{GL}_k(\mathbb F_{q_v})\),
using the \(f\) split factors over each marked rational prime.
For \(a,a'\in G\) define
\[
       R(a,a')=\sum_v\operatorname{rank}_{\mathbb F_{q_v}}
                         (a_v^{-1}a'_v-1).
\]
For a squared block norm write \(R=R(a_1,a_2)\) for the two
orbital terms. In the fourfold trace of
Theorem~\ref{thm:fourkernels622} use
\(R=R(a_1,a_2)+R(a_3,a_4)\), with that theorem's ordering.

\begin{theorem}[Large comparison ranks in the signed moments]
\leavevmode\label{thm:markedranks623}
For each fixed \(L>0\) one can choose \(C_L>0\) so that, with
\[
    J_x=\left\lceil C_L f
               \left(d^3+\frac{k\log(2x)}{\log b}\right)\right\rceil ,
                 \qquad  J_x/(fk)\longrightarrow0 ,
\]
the following bounds hold on every retained row pair or quadruple.
In the orbital expansion of \(\|Z_{ri}\|_{\rm HS}^2\), the sum
of absolute values of the terms with \(R\ge J_x\) is at most
\(\exp(-LU_x)\sigma_{ri}^2\). In the expansion of
\(\operatorname{tr}(Z_{ri}^*Z_{rj}Z_{tj}^*Z_{ti})\), the
corresponding sum is at most
\[
       \exp(-LU_x)\,\sigma_{ri}\sigma_{rj}\sigma_{tj}\sigma_{ti}.
\]
At any rank, a term is zero if a paired relative marking is
nonunipotent at one place. The estimates hold for any chosen
subsets of the original orbits, including the full operators.
\end{theorem}
\begin{proof}
Write \(U_0\subset G\) for the product of the upper
unitriangular subgroups. The finite matrix is the one of
Lemma~\ref{lem:whittakermatrix526}. Within each family
\(gU_0\), a column occupies one coordinate and its unnormalized
entry has modulus one. Two columns occupy the same coordinate
precisely when \(a'a^{-1}\in gU_0g^{-1}\). Consequently
\(|h_a^* h_{a'}|\) is bounded by the proportion of such families.
The test is independent at the different places for the uniform
measure on \(G/U_0\). Locally a nonunipotent matrix fixes no
decorated flag in this sense. For a unipotent one the proportion
is at most \(q_v^{-k R_v/3}\), where
\(R_v=\operatorname{rank}(a_v^{-1}a_v'-1)\), by
Lemma~\ref{lem:fixedflags302}. The two orders of the relative
matrix are conjugate, so they give the same rank and probability.
Taking the product yields the bound
\(\exp(-k R(a,a')\log q_*/3)\), \(q_*=\min(g_1,g_2)\).

The two-term expansion has one such factor. The four-term
expansion has the product of the two factors from the adjoints
on the first and third blocks. Its remaining trace has absolute
value at most the product of the four Hilbert--Schmidt norms
by two applications of Cauchy--Schwarz, as in
Theorem~\ref{thm:fourkernels622}. The two-term case uses
Hilbert--Schmidt Cauchy--Schwarz directly. Absolute convergence
and Lemma~\ref{lem:massranks623} therefore bound either
restricted sum, with \(p=2\) or \(4\) factors respectively, by
\[
       \left(\prod_{h=1}^p\sigma_h\right)
           \exp\!\left(pCU_x-\tfrac13 kJ_x\log q_*\right),
\]
where \(\sigma_h\) belongs to the corresponding block.
The same inequality holds after taking absolute values of the
real integrands before integrating. The two marked primes
satisfy \(q_*\ge(kb)^8\) by Theorem~\ref{thm:identitycoefficient551}.
Hence \(kJ_x\log q_*\ge8C_LU_x\), which proves both estimates
on increasing \(C_L\). Finally \(d^3/k\to0\) and
\(\log(2x)/\log b\to0\), giving the claimed limit.
\end{proof}

\begin{corollary}[One integral selection of marked coordinates]
\leavevmode\label{cor:integralranks623}
Choose a reference index \(a_r\) in each target row, for instance
the index of its identity isometry. Fix nonnegative weights
\(\omega_{ri}\) on a finite collection of blocks with the same
finite marked group at a given \(x\). There exists one family
\(gU_0\) for the entire collection such that the following holds.
In each target row select the coordinate of \(a_r\) in that
family from the unnormalized table and divide by its root of
unity. Let \(X_{ri}\) be the contribution to this selected
coordinate from orbits with \(R(a_r,a_\gamma)\ge J_x\).
Then, increasing \(C_L\) if necessary,
\[
       \sum_{r,i}\omega_{ri}\|X_{ri}\|_{\rm HS}^2
          \le \exp(-LU_x)\sum_{r,i}\omega_{ri}\sigma_{ri}^2.
\]
The selection has an integral right inverse on its finite
coordinates. The same assertion holds if the collection includes
several lists of orbital coefficients on these marks.
\end{corollary}
\begin{proof}
Choose the family uniformly. Set \(M_{ri}=\sum_\gamma
\|P_\gamma^{ri}\|_{\rm HS}\). For a fixed choice, the triangle
inequality bounds \(\|X_{ri}\|_{\rm HS}\) by the sum of the
norms over the high-rank terms which occupy the reference
coordinate. This sum is at most \(M_{ri}\), so its square is at
most \(M_{ri}\) times the sum itself. For each such orbit the
probability of occupying the coordinate is at most
\(\exp(-kJ_x\log q_*/3)\) by the decorated-flag count in the
theorem. Thus the expected squared norm is at most
\(M_{ri}^2\exp(-kJ_x\log q_*/3)\). Absolute summability justifies
the expectation of the series. Insert Lemma~\ref{lem:massranks623},
sum with the chosen weights, and take a family whose sum is no
greater than the mean. This argument permits arbitrary dependence
among the collision events for different rows. Inserting a
coordinate back in its selected position and multiplying by the
original root of unity gives the asserted integral right inverse.
\end{proof}

In particular the normalized diagonal tests may discard the
high-rank terms even after division by the squared or fourth
power of the identity coefficient. For the mixed Gram square
the discarded absolute contribution is at most
\(\exp(-LU_x)(\sum_r\sigma_{ri}\sigma_{rj})^2\), by summing over
its target indices. The surviving relative markings are
unipotent at every marked place with total rank below \(J_x\).
Their terms include different orbits at the very same marked
index. Their signed sum, its comparison with the prescribed
integral theta maps, and the full strict cohomology sum remain
to be bounded on the intervals in Section~\ref{sec:missing}.


\subsection{Rational isometries in the mixed ordinary Gram moments}
\label{sec:rationalgram629}

A mixed Gram term has four positive column spaces. Its four rational
transports compose to an endomorphism of the first one, even when the
two input rows are different. We show that up to an exponentially
small error this endomorphism fixes a subspace of codimension \(o(k)\)
over \(E_x\). The error is measured relative to the identity
normalizations of Section~\ref{sec:markedranks623}. We use the same
fields, auxiliary primes, translates and two terminal frequencies.

\begin{lemma}[A product cutoff at the identity normalization]
\leavevmode\label{lem:productcut629}
Use the unit-frequency row coordinates and the product height \(V_l^0\)
of Section~\ref{sec:signedkernel621}. On each side of each ordinary
orbital integral impose \(V_l^0\le A\), where \(A=b^{Cfs}\).
For \(L>0\) one can take \(C\) independent of \(x\) and of the rows
so that the sum over cosets of the integrals of the norms on the
discarded domains, including the absolute pair coefficient, is
at most \(\sigma_{ri}\exp(-LU_x)\).
At each retained point the translated height \(\widetilde V_l\)
is at most \(b^{C'fs}\) for another uniform constant.
\end{lemma}
\begin{proof}
Write \(q_0=4n^2\) and
\(\nu=4\sum_{j=1}^d j(n-j+1)+1\le n^2\). In the row probability
measure of Proposition~\ref{prop:fulltail621},
\(\int(V_l^0)^u\le b^{C_0fsu}\) for \(0\le u\le q_0\).
The adapted translated height satisfies
\(\widetilde V_l\le R^{fs}V_l^0\), where \(\log R=O(\log(2x))\).
On the deleted set insert
\(\mathbf1_{V_l^0>A}(V_l^0)^\nu
 \le A^{-(q_0-\nu)}(V_l^0)^{q_0}\).
Apply Cauchy--Schwarz in the unfolded coset integral as in
Theorem~\ref{thm:twoenergies592}, first with the deletion in one row
and then in the other. The inverse-height count is \(P'_\nu\)
of Proposition~\ref{prop:powercut622}. With a uniform \(C_1\) the
sum of the resulting integrals of norms is at most
\[
       2\lambda_{ri} P'_\nu\, b^{C_1fsq_0}
                              A^{-(q_0-\nu)/2}.
\]
This argument puts the cutoff on the product of the local heights.
It permits a large coordinate at one real place when the product
over all real places is small.

The count and the column-dimension estimates in the proof of
Lemma~\ref{lem:massranks623} give
\[
 \log P'_\nu\ll fk\log^2(2fk)+fnd^3\log b,\qquad
 \log(\lambda_{ri}/\sigma_{ri})
       \ll f\{k^2\log(2x)+k\log b\}
\]
on a nonzero block. Each right side, and \(U_x\), is
\(o(fs n^2\log b)\). Indeed \(n/k\to1\), \(d=s+1\),
\(d^2/k\to0\), and \(\log(2x)/(s\log b)\to0\) on the common
fields. Since \(q_0-\nu\ge3n^2\), the negative term in the
exponent gives the asserted bound on taking \(C\) sufficiently
large. The bound for \(\widetilde V_l\) follows from the same
change of metric. The zero block has vanishing pair coefficient.
\end{proof}

\begin{proposition}[A common lattice for each bounded transport]
\leavevmode\label{prop:matchedtransports629}
At every vertex use the completed positive lattice
\(\widehat L\) of Lemma~\ref{lem:completedcolumn535}.
Let a compatible orbital term meet the two product cutoffs of the
preceding lemma. Its rational transport \(j:P_0\to P_1\) has the
following property in both directions. In its source lattice there
is a full sublattice \(C_j\) mapped into the target lattice. On
\(C_j\) the reduction in the ambient lattices is the original
period index at every mark. If \(J_j\) is the index ideal of
\(C_j\) in its source, then
\[
               H_j=(N_{E/\QQ}J_j)^{1/(2f)}
                            \le b^{C_2d^2}.
\]
The same constant works for all the terms, including both terminal
frequencies.
\end{proposition}
\begin{proof}
Put \(e=\dim(X_0\cap X_1)\), \(t=s-e\), and \(Q=g_1g_2\).
By Theorem~\ref{thm:cross132}, the common positive quotient has
rank \(k\) or \(k-1\). If \(e>0\), compatibility forces a common
totally negative vector in the two extended flags. Indeed their
successor maps carry a common vector at position \(i\) to vectors
differing by the sum of the two initial spans through \(i\).
Subtract vectors in those spans to get a common vector at the next
position with nonzero leading coefficients. At the final position
its norm is a nonzero terminal coefficient squared in absolute value
times the negative terminal norm. This is the argument of
Lemma~\ref{lem:suffixdet556}, also for \(e=s\).
Thus the rank-\((k-1)\) case has \(e=0\). Its transport is the
extension with prescribed determinant. Lemma~\ref{lem:hyperplanelattice603}
gives the required sublattices with \(H_j=h_\gamma\).
Proposition~\ref{prop:transverseheight582} gives at a retained point
\[
     h_\gamma\le (T_0T_1)^d
          \left(\frac{\widetilde V_0\widetilde V_1}{Q^f}\right)^{2d/f}.
\]
Here the saturation index and its factor \(P_\gamma\) are at least
one. The cutoff bound and \(\log T_l=O(d\log b)\) imply the claim.

In rank \(k\), a common negative vector \(z_*\) satisfies
\(Y_b=X_b+Ez_*\). Use the isometry supplied by the shared
positive quotient. We give the sublattice in the source at end \(1\)
after putting both frames in the common ambient lattice \(\Lambda\).
Write \(I=X_0\cap X_1\),
\(U_b=(\Lambda\cap X_b)/(\Lambda\cap I)\).
The pairing \(h/Q\) gives an integral map
\(B:U_1\to U_0^\sharp\) of rank \(t\), nonsingular over \(E\), as
in Lemma~\ref{lem:markedtransport533}. Let \(\mathfrak b\) be
its determinant ideal, equal to the unit ideal when \(t=0\).
The order of its cokernel is \(N\mathfrak b\) by elementary divisors.

A vector of the original positive lattice \(L_1\) has a representative
\(w\in\Lambda\cap Y_1^\perp\). The functional \(h(-,w)/Q\) on
\(U_0\) is integral. Away from \(Q\) this follows from integrality of
the ambient pairing. At a mark \(w\) pairs to zero modulo that mark
with \(X_0\), since the two reduced isotropic spaces agree. The
functional kills \(I\). Changing \(w\) by a vector in
\(\Lambda\cap X_1\) changes the functional by an element of
\(B(U_1)\). Its class therefore defines a homomorphism
\(L_1\to\operatorname{coker}B\). Take its kernel for \(C_j\).
For a vector of that kernel choose \(u\in\Lambda\cap X_1\) with
\(h(-,w-u)=0\) on \(X_0\). The integral vector \(w-u\) still
pairs to zero with \(Y_1\) and with \(z_*\). Thus it is a
representative in the common perpendicular. It maps into the original
lattice \(L_0\), and gives the prescribed change of basis in the two
marked positive quotients because reduction of \(u\) belongs to
the common reduced outer space. This is valid also where the
divided cross pairing is singular modulo the mark. Reversing the
orientation uses the conjugate transpose pairing with cokernel of
the same order. By Lemma~\ref{lem:completedcolumn535},
\([\widehat L_b:L_b]\le T_b^f\), since an enlargement has length
one at each of \(f\) primes over the terminal prime. These
enlargements are away from both marks. Consequently
\(H_j\le\sqrt{T_1}(N\mathfrak b)^{1/(2f)}\) in this orientation.

For \(e>0\), Theorem~\ref{thm:terminaldivisor556} shows that
\(\mathfrak b\) divides
\((D_t/Q^t)\overline{\mathfrak h_t}\mathfrak h_t\), where
\(D_t\) is the determinant on the first \(t\) frame vectors. For
\(e=0\) there is equality by taking determinants in the outer
line lattices. For \(t=0\) the determinant and the ideals are one.
The exterior Cauchy--Schwarz estimate in
Proposition~\ref{prop:transverseheight582} applies to the first
\(t\) vectors without any transversality requirement. Taking
norms of the determinant ideal therefore gives
\[
 (N\mathfrak b)^{1/(2f)}
    \le Q^{-t}(T_0T_1)^{t/2}
                        (\widetilde V_0\widetilde V_1)^{t/f}.
\]
The factor \(\eta(\mathfrak h_t)^2\) in that estimate cancels
\((N\mathfrak h_t)^{1/f}\). The product bound in the lemma and
\(0\le t\le s\) give \(\log H_j=O(d^2\log b)\), also in the
reverse orientation.
\end{proof}

\begin{theorem}[Rational rank in the second and fourth moments]
\leavevmode\label{thm:rationalgram629}
In the second moment of one ordinary block put \(M=j_1^{-1}j_2\).
In the mixed trace use the block order
\((r,i),(r,j),(t,j),(t,i)\) of Theorem~\ref{thm:fourkernels622}
and put
\[
                        M=j_1^{-1}j_2j_3^{-1}j_4 .
\]
This is an automorphism of the positive space of the first source
row over \(E_x\). For each \(L>0\) there is \(C_L\) such that the
sum of the absolute values of the orbital terms with
\[
       \operatorname{rank}_{E_x}(M-1)>
         C_L\left(d^3+\frac{k\log(2x)}{\log b}\right)
\]
is at most \(e^{-LU_x}\sigma_{ri}^2\) in the second moment, or
\(e^{-LU_x}\sigma_{ri}\sigma_{rj}\sigma_{tj}\sigma_{ti}\)
in the mixed trace. The rank cutoff is \(o(k)\). In the mixed
trace the incoming positive spaces need only be linked by the
four transports appearing in the term.
\end{theorem}
\begin{proof}
First impose the product cutoffs. For any subset of tuples the sum
of the absolute differences between the uncut and cut orbital terms
is at most \(e^{-(L+1)U_x}\) times the indicated product of
normalizations for sufficiently large \(x\). To see this subtract
one factor at a time. Its sum of discarded absolute integrals is
bounded by Lemma~\ref{lem:productcut629} with a larger exponent.
The other factors have absolute mass at most
\(\sigma\exp(CU_x)\) by Lemma~\ref{lem:massranks623}, also after
cutting. Finite columns have norm one, and an absolute trace of
two or four matrices is at most the product of their
Hilbert--Schmidt norms. These bounds justify the summations.

For a retained tuple write its successive maps in composition
order, starting at the first source, as \(g_1',\ldots,g_h'\),
where \(h=2\) or \(4\). For example the first map in the fourfold
case is \(j_4\). Let \(C\) be the sublattice of the starting
\(\widehat L\) on which every partial composite lands in the
prescribed domain of the next map from
Proposition~\ref{prop:matchedtransports629}. It is constructed
by imposing these conditions in turn. At step \(l\) the quotient
of the previous lattice by the new one injects into the local
ambient lattice modulo \(C_{g_l'}\) via the preceding partial
composite. Thus for its index ideal \(J=[\widehat L:C]_{\rm id}\)
one has \(v_\mathfrak p(J)\le\sum_l v_\mathfrak p(J_{g_l'})\)
at every finite prime. The composite \(M\) carries \(C\) into
\(\widehat L\), and its action on \(C\) in the marked quotients is
the product \(M_v\) of the prescribed residue matrices in the
same order.

Put \(\rho=\operatorname{rank}_E(M-1)\) and
\(r_v=\operatorname{rank}_{\mathbb F_{q_v}}(M_v-1)\). Since the
space is totally positive, \(M\) is unitary at all complex embeddings
and diagonalizable there. The coefficient \(c\) of \(Y^{k-\rho}\) in
\(\det(M-(1+Y)I)\) is nonzero, and
\(|\tau c|\le 2^\rho\) at every embedding, by multiplying the
nonzero eigenvalue differences. In a local basis of \(C\) and of
\(\widehat L\), the polynomial for
\((M-1)|_C-Y\,\operatorname{incl}_C\) has integral entries in its
defining matrix. Its indicated coefficient is \(c\) times a generator
of \(J\). At a mark the constant matrix has residual rank at most
\(r_v\). The coefficient takes \(\rho\) columns from that matrix.
Every minor of order \(\rho\) of the constant matrix is divisible by
the uniformizer to power \((\rho-r_v)_+\), by elementary divisors.
The conjugate split factor has the same rank bound since the
marked maps there are the contragredients of the given maps, also
after composition. Taking the norm of the integral ideal \((c)J\)
and using both factors over each \(v\) proves
\[
 \prod_{v\mid Q} q_v^{(\rho-r_v)_+}
       \le 2^{f\rho}\prod_{l=1}^{h} H_{g_l'}^{\,f}.
\]
In particular
\[
 \rho\log(Q/2)
     \le f^{-1}\sum_{v\mid Q} r_v\log q_v+
                                   \sum_{l=1}^{h}\log H_{g_l'} .
\]

With all marked bases in their prescribed coordinates, in the
fourfold case \(M_v=a_1^{-1}a_2a_3^{-1}a_4\). The rank of its
difference from one is at most the sum of the ranks of
\(a_1^{-1}a_2-1\) and \(a_3^{-1}a_4-1\), by writing a product
minus identity as \(A(B-1)+(A-1)\). In the twofold case the
first rank alone suffices. Thus \(\sum_v r_v\le R\) for the
comparison rank in Section~\ref{sec:markedranks623}. For each
\(v\), \(\log q_v\le\log Q\), while
\(\log Q=(16+o(1))\log b\). The preceding inequality gives
\[
                 \rho\le (1+o(1))R/f+O(d^2)
\]
on all retained tuples. If \(\rho\) exceeds the cutoff in the
statement with \(C_L\) sufficiently large, then \(R\ge J_x\)
of Theorem~\ref{thm:markedranks623} with exponent \(L+1\).
The proof of that theorem multiplies the finite column inner
products by absolute real integrals. Those real majorants can
only decrease on the cut domains, so the same high-rank estimate
holds there. Add the discarded-domain bound. Since \(U_x\to\infty\),
the two errors fit the stated allowance. Finally \(d^3/k\to0\)
and \(\log(2x)/\log b\to0\), proving the assertion about codimension.
\end{proof}

\begin{corollary}[The fixed space on an integral lattice]
\leavevmode\label{cor:rationalfixed629}
In any retained tuple let \(D=\ker(M-1)\) and
\(D_0=D\cap\widehat L\). The lattice \(D_0\) is saturated of
rank \(k-\rho\). All partial transports can be applied integrally
on \(D_0\cap C\), and the last one returns every such vector to
itself. Over the selected split factors at the marks, the total
codimension lost on reducing \(D_0\cap C\) inside the reduction
of \(D_0\) is \(O(fd^2)\). If the finite indices coincide within
both pairs of the fourfold contraction, or within the single pair
of a twofold contraction, then \(\rho=O(d^2)\). Here equality
of indices means equality of the full marked tuples.
\end{corollary}
\begin{proof}
Saturation follows by intersecting the lattice with a rational
subspace. The quotient \(D_0/(D_0\cap C)\) injects into
\(\widehat L/C\), so the dimension lost in its reduction at a
chosen prime is at most \(v_\mathfrak p(J)\). Summing over the
selected primes costs at most
\(\log N J/\log\min(g_1,g_2)=O(fd^2)\), by the index bound in
the preceding proof. Local saturation embeds the reduction of
\(D_0\) in that of the starting ambient lattice. When paired
indices coincide \(M_v=1\) at every mark. The same proof then
gives \(\rho\log(Q/2)\le\sum_l\log H_{g_l'}=O(d^2\log b)\).
\end{proof}

The remaining ordinary terms may have \(M=1\), including contributions
from different rational partners. Their signed weights still matter
in the comparison with the prescribed integral theta period
in Section~\ref{sec:missing}.


\subsection{Operator bounds on the prescribed compact column types}
\label{sec:compactranks630}

In the ordinary Gram moments the dimension of an archimedean column
can be large. The common positive unitary group permits a bound for
an individual orbital operator in which that dimension is divided
out before taking absolute values over the rational orbits. We apply
this observation to the same signed blocks on the period fields.
Write
\[
        V_x=fk d^3\log b,\qquad
        s_{ri}=\left(\frac{\min(D_r,D_i)}{\max(D_r,D_i)}\right)^{1/2},
\]
and keep \(U_x,\sigma_{ri}\) of Section~\ref{sec:markedranks623}.

\begin{lemma}[Dimensions on restriction to a positive hyperplane]
\leavevmode\label{lem:branchratio630}
Let a highest weight \(\lambda\) of \(U(k)\), \(k\ge2\), restrict
to a branch of highest weight \(\mu\) for \(U(k-1)\).
If \(u_i=\lambda_i+k-i\) and \(B=u_1-u_k\), then
\[
          \frac{\dim V_\lambda}{\dim V_\mu}
                    \le\frac{B^{2k-3}}{(k-1)!}.
\]
For the pairs of actual column types in a nonzero ordinary cell,
every common compact branch has dimension at least
\(\max(D_r,D_i)\exp(-Cfk\log b)\), uniformly on the period fields.
\end{lemma}
\begin{proof}
Put \(v_i=\mu_i+k-1-i\). Interlacing gives
\(u_{i+1}\le v_i\le u_i-1\). Weyl's dimension formula expresses
the dimension ratio as the Vandermonde product of the \(u_i\)
divided by that of the \(v_i\) and by \((k-1)!\).
For each \(2\le a<c\le k-1\), the distinct denominator factor
\(v_{a-1}-v_c\) is at least \(u_a-u_c\).
Every other denominator difference is at least one. The numerator
has \(2k-3\) factors with an endpoint \(1\) or \(k\), each at most
\(B\). This proves the first bound, including \(k=2\).
For the original column parameters the shifted span at each
place is at most \(3\lceil xb\rceil\) by the weight bound in
Lemma~\ref{lem:weyl209}. A scalar shift has no effect on the
span. Tensor over the \(f\) places and apply the bound separately
to each of the two restrictions containing the same branch.
Since \(\log(2x)=o(\log b)\), their losses have the asserted
size. In full common dimension \(k\) any surviving intertwiner
already identifies equivalent irreducible column types and
the branch has their common dimension.
\end{proof}

\begin{theorem}[Operator mass at the scalar of the identity cell]
\leavevmode\label{thm:compactmass630}
For the ordinary block \(Z_{ri}=\sum_\gamma h_{a_\gamma}
\otimes P_\gamma^{ri}\) of Section~\ref{sec:markedranks623},
the original, uncut operators satisfy
\[
             \sum_\gamma\|P_\gamma^{ri}\|_{\rm op}
                     \le \sigma_{ri}\,s_{ri}\exp(CV_x).
\]
The constant is uniform in the rows, their words, the translate
and both terminal frequencies. The estimate holds for arbitrary
subsets of the cosets.
\end{theorem}
\begin{proof}
Let \(M_\gamma\) be the absolute pair coefficient times the
integral of the product of the two real section norms for this
orbit. The scalar counting and moment estimate in the proof of
Lemma~\ref{lem:massranks623}, before converting column dimensions,
gives
\[
      \sum_\gamma M_\gamma
        \le\lambda_{ri} P'_\nu b^{Cfs\nu}
        \le\lambda_{ri}\exp(C'V_x).
\]
Indeed \(\nu=O(nd^2)\), \(n/k\to1\), and
\(\log P'_\nu\ll fk\log^2(2fk)+fnd^3\log b\).
All fixed finite phases have absolute value one.

After integrating along the common unipotent fibre, an incompatible
cell gives zero. In a compatible cell the integral intertwines the
compact group of the common positive subspace at the real places,
as in Theorem~\ref{thm:common133}. To recall the norm calculation,
average a rank-one operator \(vw^*\) over that group with probability
Haar measure. On a common branch of dimension \(d_\mu\), its block
is a scalar times an isometry, with scalar of absolute value at most
\(\|v_\mu\|\|w_\mu\|/d_\mu\) by taking the trace. The restriction
is multiplicity-free, so the operator norm is the largest of these
scalar bounds. In full common dimension this is the same calculation
on a single irreducible. The averaged operator equals the original
after integration by the intertwining identity. Absolute convergence
allows averaging under the integral; equivalently one can first
apply the compact averaging operator to the convergent integral.
Thus
\[
       \|P_\gamma^{ri}\|_{\rm op}
             \le M_\gamma\,\frac{\exp(Cfk\log b)}{\max(D_r,D_i)} .
\]
If there is no common branch the operator is zero.
Sum this inequality, absorb \(fk\log b\) in \(V_x\), and use
\(\lambda_{ri}=\sigma_{ri}\sqrt{D_rD_i}\).
\end{proof}

\begin{theorem}[The two and four factors as operators on the incoming row]
\leavevmode\label{thm:rankoperator630}
Consider the positive-column operator \(Z_{ri}^*Z_{ri}\) or
\(Z_{ri}^*Z_{rj}Z_{tj}^*Z_{ti}\). Let \(h=2\) or \(4\), respectively,
and let \(\prod_{\rm edges}\sigma\) denote the product with these
multiplicities. For each \(L>0\) the following two bounds hold,
with a uniform \(C_L\).
\begin{enumerate}
\item The sum of the operator norms of the individual terms with
finite comparison rank \(R>C_L f d^3\) is at most
\(e^{-L V_x}\prod_{\rm edges}\sigma\). Terms with a nonunipotent
relative marking in either contracted pair are zero.
\item The sum of the operator norms of the terms with
\(\operatorname{rank}_{E_x}(M-1)>C_Ld^3\), for the rational
composite \(M\) of Theorem~\ref{thm:rationalgram629}, satisfies
the same bound. These estimates hold for the full integrals.
\end{enumerate}
In particular the error from either set of terms in the trace
against any orthogonal projection on the incoming positive column,
divided by the rank of that projection when positive, is at most
\(e^{-LV_x}\prod_{\rm edges}\sigma\).
\end{theorem}
\begin{proof}
The finite contractions in the two products are respectively
\(\langle h_{a_1},h_{a_2}\rangle\) and the product of this with
\(\langle h_{a_3},h_{a_4}\rangle\).
The proof of Theorem~\ref{thm:markedranks623} bounds their absolute
value by \(\exp(-kR\log q_*/3)\), where \(q_*=\min(g_1,g_2)\),
and proves the stated vanishing. Each remaining matrix product has
operator norm at most the product of its individual operator norms.
On any set of tuples with \(R>R_0\), Theorem~\ref{thm:compactmass630}
thus bounds the sum of norms by
\[
  \exp(ChV_x-k R_0\log q_*/3)
                    \prod_{\rm edges}(\sigma s_{ri}),
\]
where each dimension factor is evaluated on its own edge. They all
lie in \((0,1]\). Since \(\log q_*\ge 8\log b\), taking \(R_0\)
to be a sufficiently large multiple of \(fd^3\) gives (i).

For (ii) call a coset retained if its full domain meets both
product cutoffs in Lemma~\ref{lem:productcut629} at the same
integration point. On a coset not retained, the entire integral
is over the union of the two deleted domains. Therefore the sum
of the norms of these full operators is at most
\(\sigma_{ri}\exp(-A U_x)\) for any prescribed \(A\) after making
the cutoff choice in that lemma. Use this in one edge at a time,
and the full operator mass bound on the other edges. Since
\(U_x\ge V_x\), all terms involving a deleted coset together
cost at most \(\tfrac12 e^{-LV_x}\prod_{\rm edges}\sigma\)
for sufficiently large \(x\) with \(A\) large enough.

For a tuple of retained cosets the arithmetic conclusion of
Theorem~\ref{thm:rationalgram629} applies to the transports
themselves. Its proof gives
\(\operatorname{rank}_E(M-1)\le (1+o(1))R/f+O(d^2)\).
Consequently the indicated rational rank threshold implies the
finite threshold from (i), with exponent \(L+1\), on increasing
\(C_L\). Apply (i) to the full integrals in these retained
cosets. This absorbs the other half of the allowance.
Lastly the absolute trace of a compression to a subspace of
dimension \(u\) is at most \(u\) times the operator norm, by
summing its \(u\) diagonal entries in an orthonormal basis.
\end{proof}

For a repeated incoming row this calculation measures each error
against the actual identity coefficient. It also controls arbitrary
unit input vectors before a trace is taken. The low-rank rational
composites, including the identity, can depend on the relative signs
of several orbital integrals and remain in the projected period.


\subsection{Integral coordinates for bounded rational transports}
\label{sec:boundedcoordinates631}

Fix an ordinary block and use the completed source and target lattices
\(L_s,L_t\) of Lemma~\ref{lem:completedcolumn535}. In this subsection
we retain the whole integral of each compatible coset meeting both
product cutoffs of Lemma~\ref{lem:productcut629} at a common point.
The cutoff constant is fixed throughout. Put \(Q=g_1g_2\),
\(q_*=\min(g_1,g_2)\), and \(W_x=fk d^2\log b\).
We continue to use the actual index \(a\) of the finite column,
with the phases in the orbital coefficient. The local finite
matrix is described in Lemma~\ref{lem:whittakermatrix526}.
Its induced character is the one defined in
\srcloc{Formalism/LocalToGlobal/SiegelLatticeTheory/ClosedLocalOrbits/GGMirabolic}{28}.

At a place \(v\mid Q\) of \(E^+\), write \(\mathfrak p,\bar{\mathfrak p}\)
for the ordered pair of split factors of \(E\). The perfect pairing
identifies the reduction at the second with the dual of the reduction
at the first. Thus \(a_v\) specifies the two maps \(a_v\) and
\(a_v^{-*}\), where the latter denotes inverse transpose in dual bases.
For a retained rational isometry \(j\) and its marking \(a\), define
\[
 C(j,a)=\{w\in L_s:jw\in L_t,\quad
       \overline{jw}=a\,\bar w\ \text{at both factors of every mark}\}.
\]
In this formula and below a bar on a vector at a given factor means
its image in the ambient lattice modulo that prime. By definition
\(jC(j,a)=C(j^{-1},a^{-1})\) with the lattices exchanged.
In particular these lattices are intrinsic to the pair.

\begin{proposition}[The number of possible transports and markings]
\leavevmode\label{prop:transportcount631}
There are at most \(\exp(CW_x)\) distinct pairs \((j,a)\) in a retained
block, uniformly on the common fields. Each pair lies in a class
specified by \(C=C(j,a)\), \(C'=jC\), and the isometry \(j:C\to C'\).
Write \(t_{\mathfrak p}\) for the codimension of the image of \(C\)
in \(L_s/\mathfrak p L_s\). The same codimension occurs for
\(C'\) at that factor in the target. With a uniform \(H=b^{C_0d^2}\),
\[
 [L_s:C], [L_t:C']\le H^{2f},\qquad
 \sum_{v\mid Q}(t_{\mathfrak p}+t_{\bar{\mathfrak p}})\log q_v
                      \le 2f\log H .
\]
\end{proposition}
\begin{proof}
Proposition~\ref{prop:matchedtransports629}, applied in each direction,
supplies a sublattice of the stated index contained in the corresponding
canonical lattice. The two indices here are ordinary additive indices.
At each finite prime the dimension lost on reduction is the number
of positive elementary-divisor exponents. It is at most their sum.
Taking logarithms of the order of \(L_s/C\) proves the displayed
sum bound. Equality of the source and target codimensions follows
because the marked isomorphism maps the image of \(C\) onto that of
\(C'=jC\).

For a rank \(k\) lattice over \(\mathcal O_E\) the number of sublattices
of additive index at most \(Y\ge1\) is at most
\[
       Y^{k+1}\prod_{r=2}^{k+1}\zeta_E(r)\le Y^{k+1}\exp(Cf).
\]
Indeed, at a finite prime of norm \(q\) choose a basis over the local
discrete valuation ring. Project a full sublattice to the last
coordinate. If the image is \((\pi^h)\) and the intersection with
the first \(k-1\) coordinates has index \(q^l\), the lift of
\(\pi^h\) has exactly \(q^l\) possibilities modulo that intersection.
The total index is \(q^{h+l}\). Thus the local series with weight
\([L:C]^{-z}\) is \((1-q^{-z})^{-1}\) times the rank \(k-1\)
series at \(z-1\). Sublattices correspond bijectively to such
local choices equal to the original lattice almost everywhere:
intersect the local conditions in the ambient lattice, or impose
them modulo a common annihilating ideal using the Chinese remainder
theorem. The product of the resulting series at \(z=k+1\) gives
the first inequality. The Euler factors give
\(\zeta_E(r)\le\zeta(r)^{2f}\), and \(\sum_{r\ge2}\log\zeta(r)<\infty\)
by comparison with \(\sum_{n\ge2}1/(n(n-1))\).

Once \(C,C'\) are specified, any two such isometries differ by an
automorphism of \(C\). Its isometry group is finite: in an underlying
integral basis the positive trace form bounds each column of an
isometry. Minkowski's finite group bound in the proof of
Proposition~\ref{prop:powercut622} bounds its order by
\(\exp(Cfk\log^2(2fk))\) in rational dimension \(2fk\).
The restriction to the full lattice \(C\) determines the rational
isometry. For a nonempty class fix one marking \(a_0\). At a selected
factor let \(S\) be the image of \(C\), and let \(T\) be the
annihilator of the image of \(C\) at the conjugate factor, both
in the source space \(V=L_s/\mathfrak pL_s\). Each
\(u=a_{0,v}^{-1}a_v\) fixes \(S\) pointwise. Its inverse transpose
fixes the specified dual subspace pointwise, so \((u-1)V\subset T\).
Consequently \(u-1\) factors through a map \(V/S\to T\).
The number of choices is at most \(q_v^{t_{\mathfrak p}t_{\bar{\mathfrak p}}}\).
Since \(2t_{\mathfrak p}t_{\bar{\mathfrak p}}\le
k(t_{\mathfrak p}+t_{\bar{\mathfrak p}})\), their product is
at most \(H^{fk}\). Apply the sublattice count twice at \(Y=H^{2f}\).
The resulting bound is
\(H^{f(5k+4)}\exp(Cfk\log^2(2fk)+Cf)\), which has the asserted
size since \(f,k\) grow at most polynomially in \(x\).
\end{proof}

\begin{theorem}[Integral recovery on a fixed pair of matched lattices]
\leavevmode\label{thm:matchedcoordinates631}
For any one class \((C,C',j)\) in the proposition, a proportion at
least \(1-\delta_x\) of the families \(gU_0\) of the original
integral marked matrix distinguish \emph{all} distinct indices of
the class at once. Here one may take
\(\delta_x=8f q_*^{-k/4}\) for sufficiently large \(x\).
In any such family each coefficient is recovered by selecting one
coordinate and multiplying by the inverse of its root of unity.
For a matrix of columns supported on this class all determinantal
ideals of the unnormalized marked image agree with those before
applying the marked matrix. For arbitrary complex Hilbert space
coefficients \(z_a\) on its distinct indices, its unit columns satisfy
\[
 (1-\delta_x)\sum_a\|z_a\|^2
 \ \le\ \left\|\sum_a h_a\otimes z_a\right\|^2
 \ \le\ (1+\delta_x)\sum_a\|z_a\|^2 .
\]
\end{theorem}
\begin{proof}
At a selected factor use the subgroup
\[
 B_v=\{u\in{\rm GL}(V):u|_S=1,\ (u-1)V\subset T\}.
\]
It is a subgroup because its elements fix \(S\) and induce the
identity on \(V/T\), hence preserve \(T\) also under inversion.
The indices of the class are contained in \(a_0\prod_v B_v\).
There are at most \(q_v^{(t_{\mathfrak p}+t_{\bar{\mathfrak p}})r}\)
elements of \(B_v\) of difference rank \(r\): every rank \(r\)
map \(V/S\to T\) factors through a space of dimension \(r\), and
counting both factors with unrestricted entries is an upper bound.
The sum bound in the proposition gives
\(t_{\mathfrak p}+t_{\bar{\mathfrak p}}\ll fd^2=o(k)\) uniformly,
since \(\log q_*\ge8\log b\).

For uniform \(gU_0\) two indices \(a,a'\) collide at a given place
only if \(a'_va_v^{-1}\in g_v U_v g_v^{-1}\).
An element of difference rank \(r\ge1\) has probability zero if
nonunipotent and at most \(q_v^{-kr/3}\) otherwise, by
Lemma~\ref{lem:fixedflags302} as used in
Theorem~\ref{thm:markedranks623}. All relative elements for this
class are in \(a_{0,v}B_va_{0,v}^{-1}\). The sum of the probabilities
over its nonidentity elements is at most
\(\sum_{r\ge1}q_v^{-kr/4}\le2q_v^{-k/4}\) for large \(x\).
There are \(2f\) selected places. Excluding their union leaves the
claimed proportion of families in which distinct tuples occupy
distinct product coordinates. The raw nonzero entries in a family are roots of
unity. Reading them as stated is an integral left inverse on this
support over the ring of integers containing the character values.
Cauchy--Binet applied to the map and then to its left inverse gives
both inclusions of determinantal ideals, with the same argument
for fractional coefficient matrices.

For the complex upper and lower bounds fix \(a\). Each relative
element occurs for at most one \(a'\) for this \(a\). The modulus
of the normalized Gram entry is bounded by the probability of a
collision, and the probabilities for a fixed tuple multiply over
the places. The sum of the absolute off-diagonal entries in each
row is at most
\[
           \prod_{v\mid Q}(1+2q_v^{-k/4})-1
                 \le 8f q_*^{-k/4}
\]
for all sufficiently large \(x\). The diagonal entries are one.
The estimate for the quadratic form follows by pairing the
off-diagonal terms and using
\(2\|z_a\|\|z_{a'}\|\le\|z_a\|^2+\|z_{a'}\|^2\).
\end{proof}

\begin{theorem}[A common list of integral coordinates on a retained block]
\leavevmode\label{thm:commonfamilies631}
For every sufficiently large \(x\) there is a list of at most
\(Cfd^2\) families of the actual integral matrix for the block
with both of the following properties.
\begin{enumerate}
\item In every family on the list, any two retained terms in the
same coordinate have
\[
 R(a,a')\le Cfd^2,\qquad
           \operatorname{rank}_{E}(j^{-1}j'-1)\le Cd^2 .
\]
\item For each class \((C,C',j)\) at least one family on the list
separates all its distinct indices. In that family, reading the
coordinate of an index \(a\) and removing its unit phase gives its
whole signed coefficient from the class plus contributions only
from other classes satisfying the two rank bounds in (i).
\end{enumerate}
The list works for every choice of coefficients on the same retained
terms, in particular for their original signed integrals.
\end{theorem}
\begin{proof}
The number of ordered pairs of distinct indices in the block is at
most \(\exp(2C_1 W_x)\) by Proposition~\ref{prop:transportcount631}.
For a fixed pair of comparison rank \(R\) the full collision
probability is at most \(\exp(-kR\log q_*/3)\), by the independence
over places in the preceding proof. For any fixed \(A>0\), take
\(J=\lceil C_A f d^2\rceil\) with sufficiently large \(C_A\).
By the union bound a proportion at most \(\exp(-A W_x)\) of
families have even one collision with \(R>J\).
Call the complement \(\mathcal G\). On every collision in a family
in \(\mathcal G\), the two transports in the block have the same
source and target. The coefficient-divisibility argument of
Theorem~\ref{thm:rationalgram629} for these two retained cosets
gives \(\operatorname{rank}_E(j^{-1}j'-1)
\le(1+o(1))R/f+O(d^2)\). This includes collisions of equal
indices from different transports.

Choose \(m\) independent uniform families \emph{from} \(\mathcal G\).
For a fixed class the probability that a draw fails to separate it
is at most \(\delta_x/(1-e^{-A W_x})\le2\delta_x\). The number
of nonempty classes is at most \(\exp(C_1W_x)\), since each has
at least one pair and each pair specifies a unique class.
The probability that some class fails on every draw is at most
\(\exp(C_1W_x)(2\delta_x)^m\). For large \(x\),
\(-\log(2\delta_x)\ge k\log q_*/5\). Thus
\(m=\lceil C_2fd^2\rceil\) with sufficiently large \(C_2\)
makes this probability less than one. Use a successful list.
The formula for a coordinate follows directly from the integral
matrix entries, by first summing every term with the same class
and the same index with its original phase.
\end{proof}

For a given block these selected coordinates are coordinates of the
original finite projection on the retained cosets. Contributions from
other classes of small relative rank can still enter the same
coordinate. Different cosets in a single class can also contribute
with different signs at an identical index. The resulting sums must
still be compared with the integral theta period in
Section~\ref{sec:missing}. Lemma~\ref{lem:productcut629} bounds the
discarded whole cosets in the archimedean norm. At a finite place a
bound for their absolute values at infinity does not imply any
increase in their valuation.


\subsection{Recovery of a signed coefficient in the presence of other classes}
\label{sec:localrecovery632}

Fix the same block and product cutoffs as in
Section~\ref{sec:boundedcoordinates631}. Let \(A\) be any subset of
the distinct retained marked indices. In particular one can use the
entire set. At an index \(a\) the coefficient \(z_a\) denotes the
sum of all chosen contributions there, with their original phases.
We first work with arbitrary coefficients. A class below is specified
by its two canonical lattices and rational isometry as in
Proposition~\ref{prop:transportcount631}. Its marking set is the
set of indices occurring in the class, and different such sets can
overlap.

\begin{lemma}[Collisions between two classes of rational transports]
\leavevmode\label{lem:twoclass632}
Fix any two nonempty classes in the same block. In a
proportion at most \(\beta_x=4f q_*^{-k/4}\) of all families of the
integral marked matrix can a marking from the first and a
\emph{different} marking from the second occupy the same coordinate.
The classes may have different lattices and isometries. The bound
also holds when they coincide.
\end{lemma}
\begin{proof}
At a selected split factor put \(V=L_s/\mathfrak pL_s\).
By the proof of Theorem~\ref{thm:matchedcoordinates631} the two sets
are contained locally in \(a_0B_1\) and \(b_0B_2\), with
\[
 B_i=\{u\in{\rm GL}(V):u|_{S_i}=1,\ (u-1)V\subset T_i\}.
\]
All differences \(u-1\) for \(u\in B_2B_1\) kill \(S_1\cap S_2\)
and have image in \(T_1+T_2\). This follows by expanding a product
\((1+X_2)(1+X_1)\). Therefore every local relative matrix \(a'a^{-1}\)
belongs to one affine space with translation space
\[
 b_0\,\operatorname{Hom}\big(V/(S_1\cap S_2),\,T_1+T_2\big)\,a_0^{-1}.
\]
Write \(l=\operatorname{codim}(S_1\cap S_2)\) and
\(u=\dim(T_1+T_2)\). The index bounds on the two classes give
\(l+u=O(fd^2)=o(k)\), uniformly at every factor.

Consider all relative matrices \(D\) with \(\operatorname{rank}(D-I)=r\)
for a fixed \(r\ge1\). If there is one, fix \(D_0\) of the same rank
of difference from \(I\). The differences \(D-D_0\) have rank
at most \(2r\) and belong to the displayed translation space.
Their number is at most \(q_v^{2r(l+u)}\), since every map of rank at
most \(2r\) factors through a \(2r\)-dimensional space. This estimate
does not require the identity to belong to the affine space.
For each such \(D\) the probability of \(D\in g_vU_vg_v^{-1}\)
is at most \(q_v^{-kr/3}\) by Lemma~\ref{lem:fixedflags302}.
Thus the probability that the family meets any of these nonidentity
relative matrices at this place is at most
\(\sum_{r\ge1}q_v^{-kr/4}\le2q_v^{-k/4}\) for large \(x\).
A collision of different tuples requires this to happen at some
place. Sum the bounds over the \(2f\) places.
\end{proof}

Put \(J=\lceil C_0 f d^2\rceil\) with \(C_0\) large enough in
the far-pair estimate in the proof of Theorem~\ref{thm:commonfamilies631}
to give exceptional proportion at most \(e^{-W_x}\).
For \(a\in A\), let \(n_J(a;A)\) be the least number of retained
classes whose marking sets cover all
\[
             a'\in A,\qquad 0<R(a,a')\le J.
\]
It is zero for the empty set. Define
\[
 A_0=\{a\in A:n_J(a;A)\le q_*^{\,k/10}\},\qquad
       \eta_x=e^{-W_x}+4f q_*^{-3k/20}.
\]
In particular these definitions use the neighboring indices from the
whole support, including those with a different rational transport.

\begin{theorem}[Reading coefficients from the projected block]
\leavevmode\label{thm:localrecovery632}
For every \(a\in A\), the proportion of families in which another
index of \(A\) occupies its coordinate is at most
\[
                      e^{-W_x}+\beta_x n_J(a;A).
\]
There is a list of at most \(Cfd^2\) families such that the raw
integral projection of any coefficients supported on \(A\) recovers
every entry \(z_a\) for \(a\in A_0\) by reading one coordinate and
removing its unit phase. The family used can depend on \(a\). For
arbitrary Hilbert space coefficients one has the uniform lower bound
\[
    \left\|\sum_{a\in A}h_a\otimes z_a\right\|^2
                    \ge (1-\eta_x)\sum_{a\in A_0}\|z_a\|^2 .
\]
If \(Z\) is a matrix with rows supported on \(A\), \(Z_0\) consists
of its rows on \(A_0\), and \(H\) is the original unnormalized
marked matrix, then its fractional determinantal ideals satisfy
\[
                 I_u(Z_0)\subseteq I_u(HZ)\subseteq I_u(Z)
                 \qquad (u\ge1).
\]
The ideal assertion holds at every finite place over a coefficient
field containing the finite character values, whenever the entries of
\(Z\) belong to that field.
\end{theorem}
\begin{proof}
Except on a set of families of measure at most \(e^{-W_x}\), no pair
in the block of comparison rank greater than \(J\) collides.
This follows from the number of possible pairs and the individual
decorated-flag probability as in Theorem~\ref{thm:commonfamilies631}.
Fix one class containing \(a\) and a minimizing cover of the indices
within the indicated rank. For any covering class the probability
of a collision of \(a\) with a different index from it is at most
\(\beta_x\), by the lemma. A union bound proves the first assertion.
For \(a\in A_0\) the upper bound is at most \(\eta_x\).

Draw \(m\) independent uniform families. By
Proposition~\ref{prop:transportcount631}, with a fixed \(C_1\),
the probability that there is even one \(a\in A_0\) which none of
them isolates is at most \(\exp(C_1W_x)\eta_x^m\). For large \(x\),
\(\log(1/\eta_x)\ge c k\log q_*\) for a fixed \(c>0\):
\(W_x/(k\log q_*)\to\infty\) since \(\log q_*=(8+o(1))\log b\),
and \(\log f=o(k\log q_*)\). Taking \(m=\lceil C_2fd^2\rceil\)
with \(C_2\) large enough gives the required list. Choose one
isolating family for each \(a\in A_0\). Its unique nonzero
matrix entry at the chosen index is a unit root, and all other
indices of \(A\) have zero entry in that coordinate.
Multiplying by the inverse root recovers \(z_a\) with all the
original contributions at that index. The resulting coordinate
selections give an integral map carrying \(HZ\) to \(Z_0\).
Cauchy--Binet for this map and for \(H\) gives the two inclusions.

For the norm, fix a single family. Its squared norm on these
coefficients is at least the sum of \(\|z_a\|^2\) over the indices
in \(A_0\) that it isolates: their occupied coordinates are distinct,
and all remaining coordinate norms are nonnegative. The square
for the unit columns \(h_a\) is the mean of these family squares,
since the unnormalized column has exactly one unit entry in each
family. Averaging the lower bound and using the isolation probability
of each index proves the assertion.
\end{proof}

\begin{corollary}[A lower bound for the original ordinary period on an input]
\leavevmode\label{cor:fullinput632}
Fix \(L>0\) and choose the product cutoff of
Lemma~\ref{lem:productcut629} with this exponent. In a block let
\(P_a\) be the signed sum of the complete operators over the retained
cosets with index \(a\). Take \(A\) to be their set of indices.
For every vector \(v\) in the incoming positive column,
\[
 \|Z_{ri}v\|+\sigma_{ri}e^{-LU_x}\|v\|
       \ \ge\ \left((1-\eta_x)
                           \sum_{a\in A_0}\|P_a v\|^2\right)^{1/2}.
\]
More generally, if an array \(z\) on retained indices is killed
by the raw finite matrix, then every index \(a\) of its nonzero
support \(A\) satisfies
\[
              n_J(a;A)\ge (1-e^{-W_x})q_*^{k/4}/(4f).
\]
The implication about a kernel holds over any coefficient module
on which the character roots act by units.
\end{corollary}
\begin{proof}
Apply the theorem to \(z_a=P_a v\). The contribution of whole
cosets missing the joint cutoff has operator norm at most
\(\sigma_{ri}e^{-LU_x}\), by Lemma~\ref{lem:productcut629} and
the triangle inequality, as the marked columns have norm one.
Adding it back gives the first assertion. For the second, if any
family isolated \(a\) from the rest of the nonzero support, its
coordinate would be a unit times the nonzero entry \(z_a\).
Every family must therefore have a collision at \(a\). Compare
this with the first bound of Theorem~\ref{thm:localrecovery632}.
\end{proof}

There is an exact finite example illustrating the role of the number
of classes. Take the same completed positive lattice at both ends
and set \(j=1\). At one selected split place take all \(a\) fixing a
fixed codimension-two subspace \(S\) pointwise. Use the prescribed
inverse transposes at the conjugate factor and identities at all
other marks. Every pair has comparison rank at most two, and for
each \(a\) the canonical lattice has additive index
\(q_v^{2\operatorname{rank}(a-I)}\le q_v^4\). The rational
difference ranks are all zero. Yet the array with coefficient one
at each of these distinct indices is killed by \(H\).

Indeed, in a fixed family at the chosen place take a basis adapted
to its full flag in which \(S\) is spanned by a subset of the basis.
Such a basis is obtained by extending the intersections of \(S\)
with the successive flag spaces. Two indices are omitted from the
subset, so one omitted index \(i\) is greater than one. The simple
root subgroup \(1+tE_{i-1,i}\) fixes \(S\) pointwise and is in
the flag's unitriangular group. The character of that group is
nontrivial on this root: changing an adapted basis by an upper
triangular matrix scales each adjacent entry of its abelianization
by a nonzero scalar. In every occupied coordinate of the family,
left multiplication by this root subgroup preserves both the
support and the coordinate. It acts freely and multiplies the
entry by its character. Summing over each orbit gives zero by
orthogonality, hence every coordinate vanishes.
This example satisfies the two lattice index bounds on the same
finite groups. Determining the coefficients that actually occur
in the theta period requires the signed orbital integrals. In
particular, a bound for \(\sum_{a\in A_0}\|P_av\|^2\) on enough
incoming vectors and a comparison with the specified integral
map are still needed for Section~\ref{sec:missing}.


\subsection{Character traces of the signed sums at each rational isometry}
\label{sec:signedtraces635}

Consider the terms of an ordinary block for which the common
positive space has dimension \(k\). On these terms the real
operator is a scalar times the map induced by a rational isometry
of the positive spaces. Its normalized trace in a product of
blocks can be small even when the finite markings coincide.
We give estimates in terms of the norms of the actual signed
marked vectors belonging to each projective isometry.

\begin{lemma}[The lattice index in the Weyl character estimate]
\leavevmode\label{lem:indexweyl635}
Let \(M\) be a lattice in a definite Hermitian space of dimension
\(k\ge2\) over \(E=E_x\), with \(f=[E^+:\mathbb Q]\).
Let \(g\) be a rational unitary automorphism. Suppose \(gC\subset M\)
for a sublattice \(C\subset M\) of finite additive index
\([M:C]=H^{2f}\). If the distinct eigenvalues of \(g\) have
multiplicities \(m_i\), put \(q(g)=\sum_{i<j}m_i m_j\).
For a product of irreducible column representations \(\rho\) let
\(L\) be the smallest adjacent gap of its shifted highest weights
at all places, as in Proposition~\ref{prop:weyl235}. Then
\[
 \frac{|\operatorname{tr}\rho(g)|}{\dim\rho}
     \le \min\{1,\ (2k/L)^{f q(g)}H^{f(k-1)}\}.
\]
\end{lemma}
\begin{proof}
At a finite prime of \(E\) choose integral bases of \(C\) and \(M\).
If \(T_0,T_1\) are the matrices in these bases of inclusion and of
\(g|_C\), then \(\det(XT_0-T_1)=\det(T_0)\det(X-g)\) is integral.
Write \(e_v=v(\det T_0)\) with uniformizer of valuation one.
For the roots \(z_1,\ldots,z_k\) of \(\det(X-g)\), repeated by
multiplicity in a local splitting field with extended valuation,
this implies
\[
              \sum_{i=1}^k \max(0,-v(z_i))\le e_v .
\]
Indeed the maximum of the absolute values of the coefficients of
a product of polynomials is the product of the two maxima. After
dividing each factor by a coefficient of largest absolute value
this follows by multiplying two nonzero residue polynomials.
Apply this identity to the linear factors of the monic polynomial.

Let \(D(g)\) be the product of \(z_i-z_j\) over ordered pairs
whose eigenvalues differ. It lies in \(E^\times\), being invariant
under permutations induced by the Galois group. Put
\(t_i=\max(0,-v(z_i))\). For each unordered pair of distinct
eigenvalues the possible denominator valuation is at most
\(\max(t_i,t_j)\le t_i+t_j\). Each index belongs to at most
\(k-1\) unordered pairs. Thus \(v(D(g))\ge-2(k-1)e_v\).
If \(\mathfrak J\) is the index ideal of \(C\subset M\), this
shows that \(\mathfrak J^{2(k-1)}(D(g))\) is integral. Its
absolute norm is at least one, and \(N\mathfrak J=[M:C]\).

At each real place of \(E^+\) write
\(D_\tau=\prod_{i<j}|z_i-z_j|\), omitting equal eigenvalues
and retaining their multiplicities. Unitarity and complex
conjugation give \(\prod_\tau D_\tau=|N_{E/\mathbb Q}D(g)|^{1/4}\).
Hence \(\prod_\tau D_\tau\ge H^{-f(k-1)}\).
The confluent Weyl determinant computation in
Proposition~\ref{prop:weyl235} bounds the normalized character at
each place by \((2k/L)^{q(g)}/D_\tau\). Multiplication over the
places proves the assertion, together with the unitary bound one.
\end{proof}

Use the completed source and target lattices \(M_s,M_t\) of
Section~\ref{sec:boundedcoordinates631}. For a rational isometry
\(u\) between them over \(E\), define
\[
 H(u)^{2f}=
 \max\{[M_s:M_s\cap u^{-1}M_t],\
                     [M_t:M_t\cap u M_s]\}.
\]
This definition puts no conditions on the marking at either prime.
The same height applies to the inverse. For a product of \(h\) such
maps or inverses of height at most \(H\), there is a source sublattice
of index at most \(H^{2fh}\) carried into the last lattice. Indeed,
impose each successive domain condition on the original vectors.
At each step the new quotient injects into the quotient by the next
integral domain, of order at most \(H^{2f}\).

In a full common positive space the two words must have equivalent
compact representations for the orbital intertwiner to be nonzero.
Fix the induced isometries \(U_u\) between their unitary representation
spaces, using one identification with the word model at each endpoint.
Schur's lemma in Theorem~\ref{thm:common133} gives
\(P_\gamma=c_\gamma U_{u_\gamma}\), where \(c_\gamma\) includes
the scalar of the real integral, the finite phase and the pair
coefficient. This is the same transport appearing in the trace of
Corollary~\ref{cor:trace235}. Composition of the induced maps
represents composition of the rational isometries in these models.

Group isometries up to multiplication by a scalar \(\xi\in E^\times\)
with \(\xi\bar\xi=1\). Choose one occurring representative \(u\)
per class. For \(u_\gamma=\xi u\) write
\(U_{u_\gamma}=\omega_\rho(\xi) U_u\), where
\(|\omega_\rho(\xi)|=1\) by the central character. Form the signed
vector and its total squared norm in block \(e\) by
\[
       z_{e,u}=\sum_{\gamma:\,[u_\gamma]=[u]}
            c_\gamma\,\omega_\rho(u_\gamma/u)\,h_{a_\gamma},
       \qquad
       B_e^2=\sum_{[u]}\|z_{e,u}\|^2 .
\]
The notation \(u_\gamma/u\) means their scalar ratio. Here \(h_a\)
is the original unit marked column. The sums use complete orbital
integrals and converge absolutely by Lemma~\ref{lem:massranks623}
and \(P_\gamma=c_\gamma U_{u_\gamma}\). One may take any subset
of the full-dimensional terms throughout, with \(z_{e,u}\)
formed from precisely that subset.

\begin{theorem}[The second and fourth traces at polynomial lattice index]
\leavevmode\label{thm:centraltraces635}
Fix \(A,B>0\). There is a choice of the polynomial word spacing
in Theorem~\ref{thm:common111}, uniform on the original intervals,
with the following properties. In each block take all full-dimensional
terms with \(H(u_\gamma)\le x^B\), or any subset of them. Write
\(Z_e=\sum_{[u]}z_{e,u}\otimes U_u\) for this portion of the
ordinary marked operator and \(D=\dim\rho\) for its column dimension.
Then
\[
       \left|D^{-1}\|Z_e\|_{\rm HS}^2-B_e^2\right|
                    \le x^{-Afk}B_e^2 .
\]
For a fourfold trace in the block order \((r,i),(r,j),(t,j),(t,i)\)
of Theorem~\ref{thm:fourkernels622}, expand
\(D^{-1}\operatorname{tr}(Z_1^*Z_2 Z_3^*Z_4)\) by the same
projective classes. The sum of the absolute values of the terms with
\[
                     u_1^{-1}u_2u_3^{-1}u_4
                         \text{ noncentral}
\]
is at most \(x^{-Afk} B_1 B_2 B_3 B_4\). Thus the surviving sum
up to this error may be taken over exact equalities in the
projective unitary group, with the original finite inner products
and with the central character on each scalar composite.
\end{theorem}
\begin{proof}
First count distinct isometries between two given row lattices
of height at most \(H\ge1\). Their integral domain and image have
indices at most \(H^{2f}\). The sublattice count proved in
Proposition~\ref{prop:transportcount631} bounds the choices for
these two lattices together by
\(H^{4f(k+1)}\exp(Cf)\). With both fixed, their isometries, if
any, form a torsor for the integral isometry group of the domain.
This finite group embeds in \({\rm GL}_k(E)\); positivity bounds
the integral matrices of its elements in the trace metric.
Lemma~\ref{lem:conductorstabilizer532} bounds its order by
\(\exp(C_\psi k\log^2(2x))\) on the polynomial-conductor intervals,
with the constant also allowed to depend on \(\varepsilon\).
The same counts bound the number of projective representatives.
For \(H=x^B\) this number is at most \(x^{C_Bfk}\) uniformly.
Here \(f\) grows at least as a fixed positive power of \(x\) by
Theorem~\ref{thm:conductor232}.

In either trace let \(h=2\) or \(4\) be the number of maps and
write \(g\) for the composite, using \(u_1^{-1}u_2\) when \(h=2\).
The normalized real trace is \(\operatorname{tr}\rho(g)/D\), by
composition and unitarity. Every nonzero fourfold term uses
equivalent words along its cycle and therefore a common dimension.
The finite contractions are the inner products of the first two
and, for \(h=4\), of the last two vectors. Their product in
absolute value is bounded by the product of the \(h\) vector norms.
Changing representatives multiplies the real and finite factors
by cancelling unit phases.

If \(g\) is noncentral, unitarity gives at least two distinct
eigenvalues and \(q(g)\ge k-1\). Apply Lemma~\ref{lem:indexweyl635}
with lattice height at most \(H^h\), by the domain construction above.
Choose \(L/(2k)\ge x^K\) with \(K>4B+2A+4C_B+2\).
The words of Theorem~\ref{thm:common111} admit this choice:
their shifted gaps are at least \(3\Delta\), and the exponent
of the polynomial spacing \(\Delta\) may be fixed before selecting
the terminal prime, as in Corollary~\ref{cor:trace235}.
The field, interval and marked construction are those of that
theorem. The count used here depends on the field and lattices
only through \(f,k,H\) and the field conductor, so its constant
does not increase with the chosen gap.

The real factor of each noncentral term is consequently at most
\(x^{-(K-hB)f(k-1)}\). On each edge the sum of the vector norms
is at most \(x^{C_Bfk/2}B_e\), by Cauchy--Schwarz. Multiply these
bounds and use \(k-1\ge k/2\), \(h\le4\). The absolute sum over
all noncentral composites is at most \(x^{-Afk}\prod_e B_e\),
counting a repeated edge each time. For the second trace a central
composite means that the two representatives are the same. Its
terms sum exactly to \(B_e^2\). In the fourth trace a central
composite acts by its central character. This proves both formulas.
\end{proof}

\begin{corollary}[Rational scalars in the traces at the full retained index]
\leavevmode\label{cor:retainedtraces635}
Choose word spacing so that \(L/(2k)\ge x\). On each edge use
the full-dimensional cosets meeting both original product cutoffs
at one point, as in Proposition~\ref{prop:transportcount631};
retain their complete integrals. Put \(W_x=fk d^2\log b\).
For every \(A>0\) there is a uniform \(C_A\) such that, in the
second or fourth normalized trace above, the absolute sum of terms
whose projective composite has no representative satisfying
\[
            \operatorname{rank}_{E}(g-\xi I)
                 \le C_A d^2\frac{\log b}{\log x},
            \qquad \xi\in E^\times,\quad \xi\bar\xi=1 ,
\]
is at most \(e^{-A W_x}\prod_e B_e\).
The rank bound is \(o(d^3)\) on the common intervals.
\end{corollary}
\begin{proof}
The matched lattices in Proposition~\ref{prop:transportcount631}
are contained in the respective integral domains. Hence
\(\log H(u)\le C d^2\log b\) for each retained isometry and its
inverse. The same proposition counts at most \(\exp(C_1W_x)\)
representatives on each edge. Let
\(r_*=k-\max_i m_i\) for the composite \(g\). Since
\(\sum_i m_i^2\le k\max_i m_i\),
\[
             q(g)=\tfrac12(k^2-\textstyle\sum_i m_i^2)
                         \ge kr_*/2 .
\]
Lemma~\ref{lem:indexweyl635} on the composite therefore bounds the
normalized trace by \(\exp(C_2W_x-fkr_*\log(x)/2)\). The constant
allows for at most four successive index losses. Summing absolute
values as in the theorem costs at most
\(\exp(2C_1W_x)\prod_e B_e\). Thus all terms with
\(r_*>C_A d^2\log(b)/\log x\) satisfy the claimed total bound
when \(C_A\) is sufficiently large.

Now \(\log b=x^\varepsilon\) and
\(d\asymp x^{1/2-\varepsilon/16}\) in the range of
Section~\ref{sec:markedranks623}. It follows that the threshold
is \(o(d^3)=o(k)\). Below it, for large \(x\) the eigenvalue of
largest multiplicity has multiplicity greater than \(k/2\) and is
unique. The characteristic polynomial has coefficients in \(E\),
so every Galois conjugate over \(E\) of this eigenvalue has the
same multiplicity. Uniqueness forces it to be an element \(\xi\)
of \(E\). Its complex absolute value is one by unitarity, hence
\(\xi\bar\xi=1\). Diagonalizability shows that
\(\operatorname{rank}_E(g-\xi I)=r_*\). The argument is unchanged
by multiplying a representative by a rational norm-one scalar.
\end{proof}

These estimates use the norms of the marked vectors after all
phases for a given projective isometry have been included.
The central products can contribute to the fourth moment, and a
single vector \(z_{e,u}\) can vanish. The magnitudes of these
vectors and the cells of common dimension \(k-1\) still enter
the integral comparison of Section~\ref{sec:missing}.


\subsection{Scalar transports in the integral marked period}
\label{sec:scalartransports636}

The central character in Section~\ref{sec:signedtraces635} groups
rational transports which differ by a scalar. The two finite marks
can separate them, even though their real operators are proportional.
We use the same fields, lattices and complete retained orbital
integrals as in Section~\ref{sec:boundedcoordinates631}.

\begin{lemma}[High-multiplicity eigenvalues on a matched lattice]
\leavevmode\label{lem:torsioneigen636}
Let \(g\) be a composite of \(h=2\) or \(4\) successive retained
transports and inverses returning to its original positive space.
Let \(C\subset M\) be the successive matched domain in the original
completed lattice \(M\), and set \(Y=[M:C]\). In particular \(gC\subset M\)
and, on \(C\), reduction of \(g\) agrees with the composite \(g_v\)
of the prescribed markings at every chosen split factor. Write
\(t_v\) for the codimension of the image of \(C\) there.
Then \(\log Y\ll f d^2\log b\) and
\(\sum_{v\mid g_1g_2}t_v\ll f d^2\), using the same one of the
two conjugate factors for the marking at each \(v\).

An eigenvalue \(z\) of \(g\) of multiplicity \(m\) with
\(m\log2>\log Y\) is a root of unity. For every such eigenvalue
different from one, and at each chosen factor,
\[
                    \operatorname{rank}(g_v-I)\ge m-t_v .
\]
When \(m>t_v\), the marked composite has an eigenvalue different
from one in the algebraic closure of the residue field.
\end{lemma}
\begin{proof}
Construct \(C\) by imposing the next domain condition before each
map in the product, starting with its canonical matched lattice.
The successive quotients inject into the original domain quotients
on the respective edges, as in Theorem~\ref{thm:rationalgram629}.
Proposition~\ref{prop:transportcount631} bounds their indices,
giving \(\log Y\ll f d^2\log b\). Reduction on this domain is
compatible with every intermediate map and hence with composition.
At a selected prime the elementary-divisor length of \(M/C\) is
at least \(t_v\). Thus \(\sum_v t_v\log q_v\le\log Y\).
The estimate for \(\sum_v t_v\) follows from \(\log q_*\ge8\log b\),
where \(q_*=\min(g_1,g_2)\).

Let \(P\) be the monic irreducible polynomial of \(z\) over \(E\).
Then \(P^m\) divides \(\det(X-g)\) as a monic factor, since all
its roots have the same multiplicity in characteristic zero.
If a coefficient of \(P\) had negative valuation at a prime
\(\mathfrak p\) of \(E\), the maximum coefficient norm of \(P\)
there would be at least \(N\mathfrak p\), with
\(|\pi|_{\mathfrak p}=(N\mathfrak p)^{-1}\). Multiplicativity
of this norm, used in Lemma~\ref{lem:indexweyl635}, and the monic
remaining factor would force the norm of \(\det(X-g)\) there
to be at least \((N\mathfrak p)^m\). The same lemma bounds it by
the \(\mathfrak p\)-part of \(Y\). This would give \(Y\ge2^m\).
Thus \(P\) is integral and \(z\) is an algebraic integer.
All its complex conjugates have modulus one, since each comes from
an eigenvalue at an embedding of \(E\) and the space is definite
at every such embedding. To see that \(z\) is torsion, the monic
integer polynomials of its positive powers have bounded degree and
bounded coefficients, by the elementary symmetric function formula.
Only finitely many roots occur, so two powers agree.

Let \(N\) be the order of \(z\). If a marked residue characteristic
\(\ell\) divided \(N\), then \(\mathbb Q(z)\) would contain the
\(\ell\)-th roots of unity. This would give
\(\ell-1\le[\mathbb Q(z):\mathbb Q]\le2fk\).
But \(g_1,g_2\) grow faster than \(2fk+1\) on the common intervals.
Hence \(N\) is prime to \(\ell\), and a nontrivial such root has
nontrivial reduction. Indeed, if its reduction were one,
\(1+z+\cdots+z^{N-1}=0\) would give \(N=0\) in the residue field.

At a chosen factor extend the local integer ring to contain \(z\).
The number \(t_v\) of nonunit elementary divisors of \(C\subset M\)
is unchanged by this extension. Let \(S\) be the intersection of
the \(z\)-eigenspace with the extended \(M\). It is a saturated
direct summand of rank \(m\), since \(g\) is semisimple by
archimedean unitarity. The finite module \(S/(S\cap C)\) embeds
in \(M/C\) after extension and has at most \(t_v\) generators:
over a discrete valuation ring their minimal number is the
residue dimension of the subgroup killed by the uniformizer.
At least \(m-t_v\) independent reductions of \(S\cap C\) thus
survive in \(S\) and in \(M\). The marked map on these reductions
is multiplication by \(\bar z\). If \(z\ne1\), this gives the rank
bound and the asserted residual eigenvalue. Matrix ranks are
preserved by residue field extension.
\end{proof}

\begin{theorem}[Coordinate separation of scalar rotations]
\leavevmode\label{thm:scalarseparate636}
In any retained ordinary block, two distinct rational transports
in the same projective class have disjoint sets of occupied
coordinates in every family of the raw integral marked matrix.
More generally this holds for two transports whose relative
automorphism has a nonidentity eigenvalue of multiplicity \(m\)
with \(m\log2>\log Y\). Here \(Y\) is the matched index of
\(g=u_1^{-1}u_2\) in the lemma; each \(t_v\) is then strictly
smaller than \(m\) since \(t_v\log q_v\le\log Y\).

Consequently, on the full common dimension \(k\) in a retained block,
the quantity \(B_e^2\) of Theorem~\ref{thm:centraltraces635} is
exactly
\[
        B_e^2=\sum_{u}
             \left\|\sum_{\gamma:\,u_\gamma=u}
                                  c_\gamma h_{a_\gamma}\right\|^2 ,
\]
where \(u\) now ranges over actual rational isometries. The sums
at different scalars in a given projective class can be recovered
from that class vector by coordinate selection and central unit
phases, also for the unnormalized finite matrix.
\end{theorem}
\begin{proof}
If two columns share a coordinate of a family, their relative
marked matrix at every place is conjugate to an element of the
upper unitriangular group, by Lemma~\ref{lem:whittakermatrix526}.
Its eigenvalues all equal one. A nonidentity eigenvalue in the
last conclusion of Lemma~\ref{lem:torsioneigen636} therefore
precludes a common coordinate, including after any change of
family. The source and target versions of the relative matrix
are conjugate by one of the markings and give the same condition.

For distinct scalar rotations the relative map equals \(\xi I\)
with \(\xi\ne1\) and \(m=k\). On the actual fields
\[
             f d^2\log b=o(k).
\]
Indeed \(f\ll x^{2+8\varepsilon}\),
\(d\asymp x^{1/2-\varepsilon/16}\), \(\log b=x^\varepsilon\),
and \(k\asymp x^{9/2+10\varepsilon}\). Hence the hypotheses on
\(m\) and \(t_v\) hold for every pair of retained cosets for large
\(x\). The lemma also shows that \(\xi\) is a root of unity.
All cosets at one rotation thus have support disjoint from those
at another in every finite family. Orthogonality of these coordinate
subspaces proves the equality. Restriction to the corresponding
coordinates recovers each rotated vector. The required scalar
correction is a root of unity: a central character of a unitary
highest-weight representation on such a scalar is a product of
integer powers of its complex conjugates. It is therefore a unit
in the coefficient integers containing the phases. The same argument
works before the common normalization of the finite columns.
\end{proof}

\begin{theorem}[Nonidentity central products in the fourth trace]
\leavevmode\label{thm:centralfour636}
In the fourfold ordinary trace of Theorem~\ref{thm:fourkernels622}
use the original indices and full integrals of retained cosets.
Write \(g=u_1^{-1}u_2u_3^{-1}u_4\). The sum of the absolute values
of the trace terms, normalized by the incoming dimension, for which
\(g\) has an eigenvalue different from one of multiplicity at
least \(k/2\), is at most
\[
          \exp(-c f k^2\log b)\prod_{i=1}^4 \sigma_i
\]
for an absolute \(c>0\) on sufficiently distant intervals. This
estimate allows both possible common positive dimensions in each
orbital integral.

There is an exact statement at smaller matched indices. Fix
\(A,B>0\), and on each full-dimensional edge keep only cosets with
\[
           [M_s:C(u_\gamma,a_\gamma)],\
           [M_t:u_\gamma C(u_\gamma,a_\gamma)]\le x^{2fB}.
\]
Their complete integrals may be used without an archimedean cutoff.
With the spacing of Theorem~\ref{thm:centraltraces635}, the fourth
normalized trace equals its sum over the actual rational equations
\[
                       u_1^{-1}u_2u_3^{-1}u_4=I
\]
up to error of absolute value at most \(x^{-Afk}\prod_i B_i\).
Here each \(B_i^2\) is the sum in Theorem~\ref{thm:scalarseparate636}.
\end{theorem}
\begin{proof}
For each individual fourfold term denote the relative finite ranks in
the two contractions at \(v\) by \(r_{12,v},r_{34,v}\). Their sum
is at least \(\operatorname{rank}(g_v-I)\), by writing the composite
of the two relative maps and using subadditivity of rank. This also
follows after identifying their intermediate spaces by any one of
the markings. If a nonidentity eigenvalue has multiplicity
\(m\ge k/2\), then \(m\log2>\log Y\) as above. The lemma gives
\[
      \sum_{v\mid g_1g_2}(r_{12,v}+r_{34,v})
                 \ge 2fm-\sum_v t_v\ge (1-o(1))fk .
\]
The absolute product of the two marked contractions is at most
\(q_*^{-k\sum_v(r_{12,v}+r_{34,v})/3}\) by
Theorem~\ref{thm:markedranks623}. The remaining normalized
trace is at most the product of the four operator norms, also
for the hyperplane intertwiners. By Theorem~\ref{thm:compactmass630}
those products sum to at most
\(\exp(C f k d^3\log b)\prod_i\sigma_i\). Now \(d^3=o(k)\)
and \(\log q_*\ge8\log b\). Inserting the finite bound proves
the first assertion.

In the smaller range the matched domain of an individual cycle
has \(\log Y\le8fB\log x\). In particular the preceding proof of
coordinate separation applies to each pair of scalar rotations,
using their two-edge bound. The unmarked heights are at most \(x^B\).
Therefore Theorem~\ref{thm:centraltraces635} already gives the stated
error from noncentral projective composites, with \(B_i\) computed
by the orthogonal decomposition. It remains to examine the exact
central sum. Expand the finite vectors in this sum into their actual
transports and then their individual cosets. Absolute convergence
justifies these expansions. Each resulting composite is \(\xi I\).
Since \(k\log2>\log Y\), the coefficient-norm argument of the lemma
makes \(\xi\) a root of unity, with prime-to-mark order. Both
contractions must have a common coordinate to be nonzero.

The sum of the codimensions \(t_v\) is bounded by
\(8fB\log x/\log q_*=o(f)\). Among the \(2f\) chosen factors
there is consequently one with \(t_v=0\). There the matched
domain maps onto the entire reduction of \(M\), so
\(g_v=\bar\xi I\). Both relative marked factors are unipotent
when their contractions are nonzero. If their product is
\(\bar\xi I\), the second factor is \(\bar\xi\) times the inverse
of the first. Its eigenvalues are therefore all \(\bar\xi\), and
unipotence forces \(\bar\xi=1\). Prime-to-mark reduction is
injective on the group generated by \(\xi\), so \(\xi=1\).
This leaves exactly the claimed rational equations in the
central sum.
\end{proof}

Even when that equation holds there can be several distinct paths
of rational isometries between the same rows. Their signed orbital
coefficients remain in the fourth trace. The calculation singles out
the products for which those coefficients still have to be estimated.


\subsection{Summing over all target rows in the ordinary Gram matrix}
\label{sec:rowsums637}

The Gram matrices of Section~\ref{sec:geom11} sum first over a
common target row. An absolute error for each ordered pair of
target rows can grow under this summation. For the part of the
ordinary period where the two matched indices are polynomial on
the geometric-mean scale, one can perform the sum before making
the character estimate. The resulting error is relative to the
signed sum and is uniform in the number of target rows.

Fix \(B>0\). In each block \((r,i)\) take any of the full common
positive dimension \(k\) cosets satisfying the two bounds
\([M_s:C(u_\gamma,a_\gamma)]\le x^{2fB}\) and
\([M_t:u_\gamma C(u_\gamma,a_\gamma)]\le x^{2fB}\) of
Theorem~\ref{thm:centralfour636}. Use their complete integrals.
For each actual rational isometry \(u:P_i\to P_r\) of positive
spaces put
\[
       y_{ri,u}=\sum_{\gamma:\,u_\gamma=u}
                                  c_\gamma h_{a_\gamma},\qquad
       Z_{ri}=\sum_u y_{ri,u}\otimes U_u .
\]
Take absent vectors to be zero. As in Section~\ref{sec:signedtraces635}
the unitary identifications \(U_u\) are chosen coherently for each
word by fixing endpoint models. Use the inner product conjugate-linear
in its first entry. For fixed input rows \(i,j\) of equivalent words
and a rational isometry \(g:P_j\to P_i\), define
\[
       c_{ij}(g)=\sum_r\ \sum_{u:P_i\to P_r}
                          \langle y_{ri,u},y_{rj,ug}\rangle .
\]
All sums over target rows here are finite, of arbitrary size.
The height bound gives only finitely many rational isometries
between a given pair of endpoints, by the count in
Theorem~\ref{thm:centraltraces635}. Absolute convergence within
each vector follows from Lemma~\ref{lem:massranks623}.

\begin{lemma}[Two products between fixed inputs]
\leavevmode\label{lem:uniqueproduct637}
Consider all \(g=u^{-1}v:P_j\to P_i\) for which
\(\langle y_{ri,u},y_{rj,v}\rangle\ne0\) for some target row \(r\).
Their heights satisfy \(H(g)\le x^{2B}\).
The set contains at most one actual isometry in each projective
class, even when the two products being compared pass through
different target rows. For \(i=j\), the only scalar which can
occur is the identity.
\end{lemma}
\begin{proof}
The domain calculation following Lemma~\ref{lem:indexweyl635}
bounds both indices for the twofold composite by \(x^{4fB}\).
For each nonzero contraction choose one pair of individual cosets
with nonzero marked contraction. Such a pair exists by expansion and
absolute convergence. Suppose \(g=u^{-1}v\) and \(g'={u'}^{-1}v'\)
occur and \(g=\xi g'\) for a rational norm-one scalar \(\xi\).
The four chosen cosets through the two target rows form a cycle
with composite \(g{g'}^{-1}=\xi I\) on \(P_i\). Its two integral
marked contractions are both nonzero. Each edge has the stipulated
matched index bound. The clean-factor argument of
Theorem~\ref{thm:centralfour636} applies to this individual tuple:
a marked place where the entire domain reduces surjectively
identifies the product of the two unipotent relative markings with
\(\bar\xi I\). Unipotence forces \(\bar\xi=1\); the height bound
makes \(\xi\) a root of unity of order prime to the mark. Thus
\(\xi=1\) and \(g=g'\). If \(i=j\) and \(u^{-1}v\) is a scalar,
the two-edge coordinate separation of
Theorem~\ref{thm:scalarseparate636} applies directly under these
same index bounds. This proves the last assertion without
requiring a second pair.
\end{proof}

\begin{theorem}[A relative estimate for the row-summed Gram blocks]
\leavevmode\label{thm:rowsumgram637}
Given \(A,B>0\), choose a fixed polynomial spacing in the common
word construction. For all sufficiently large \(x\) on the prescribed
intervals put \(\delta_x=x^{-Afk}\). Let
\(K_{ij}=\sum_r Z_{ri}^*Z_{rj}\) for the indicated portion of the
ordinary projected period, and let \(D\) be the common dimension
for two equivalent input words. For \(i\ne j\) put \(s=0\);
for \(i=j\) take any complex scalar \(s\). With \(sI\) omitted
in the off-diagonal case, one has
\[
 (1-\delta_x) E_{ij}(s)\ \le\
       \frac{\|K_{ij}-sI\|_{\rm HS}^2}{D}
                          \ \le\ (1+\delta_x) E_{ij}(s),
\]
where
\[
 E_{ij}(0)=\sum_g|c_{ij}(g)|^2,\qquad
 E_{ii}(s)=|c_{ii}(I)-s|^2+\sum_{g\ne I}|c_{ii}(g)|^2 .
\]
In particular \(c_{ii}(I)=\sum_{r,u}\|y_{ri,u}\|^2\).
The same \(A,B\) work for every choice of the target rows and of
the indicated cosets, including when \(E_{ij}(s)=0\).
\end{theorem}
\begin{proof}
Multiplication of blocks and unitarity give the exact identity
\[
              K_{ij}=\sum_g c_{ij}(g) U_g .
\]
Indeed each ordered pair \(u,v\) contributes its finite inner
product times \(U_u^*U_v=U_{u^{-1}v}\). Contract before taking
any norm and add all terms at the very same rational isometry.
For \(g=I\) on the diagonal one has \(v=u\), proving the formula
for its coefficient. The preceding lemma shows that the support
of the sum is projectively distinct. On the diagonal adjoin \(I\)
if needed and subtract \(s\) from its coefficient; the same
distinctness still holds.

Let \(N\) be the number of possible elements of this support with
the optional identity included. The height bound is on isometries
directly between \(P_j\) and \(P_i\), with their fixed endpoint
lattices. Thus the counting argument of
Theorem~\ref{thm:centraltraces635}, now at height \(x^{2B}\),
gives \(N\le x^{C_B f k}\) whenever the support is nonempty.
It introduces no factor for an intermediate row. For two distinct
elements \(g,g'\) in the support the relative map \(g^{-1}g'\)
is noncentral and has an integral domain of additive index at most
\(x^{8fB}\). Its height for Lemma~\ref{lem:indexweyl635} is
therefore at most \(x^{4B}\). With \(L/(2k)\ge x^K\), the lemma
bounds the absolute normalized Hilbert--Schmidt inner product of
\(U_g,U_{g'}\) by \(x^{-(K-4B) f(k-1)}\), since there are at
least \(k-1\) unequal eigenvalue pairs. Each diagonal normalized
inner product is exactly one.

Choose \(K>4B+2C_B+2A+2\), which is allowed uniformly by
Theorem~\ref{thm:common111} as in Theorem~\ref{thm:centraltraces635}.
Call the possibly modified coefficients \(d_g\). By Cauchy--Schwarz,
\((\sum_g|d_g|)^2\le N\sum_g|d_g|^2\). The absolute error from
all unequal pairs in the squared norm is hence at most
\[
            N x^{-(K-4B)f(k-1)}\sum_g|d_g|^2
                      \le\delta_x\sum_g|d_g|^2 ,
\]
using \(k-1\ge k/2\). The diagonal sum is \(E_{ij}(s)\). If
the support is empty or all \(d_g\) vanish the exact identity
gives the assertion directly. Inequivalent words have no common
nonzero full-dimensional target and give zero off-diagonal block.
\end{proof}

For \(s=1\) the theorem gives the diagonal statistic of
Corollary~\ref{cor:gramerror622} on this portion, including the
subtraction of the identity inside the relative estimate. Thus,
at these matched indices, a small diagonal defect requires both
the squared marked mass to be near one and the sum of squares of
the correlations at nonidentity rational products to be small.
Here \(fk\log x\) grows faster than \(x^{\varepsilon/8}\), so the
relative error in this measurement is smaller than the comparison
scale in Section~\ref{sec:geom11}. Establishing the required
sizes of the correlations themselves and including the remaining
matched indices and hyperplane cells are still needed for
Section~\ref{sec:missing}.


\subsection{Exact separation by the trace on the matched lattice}
\label{sec:matchedtrace638}

There is an arithmetic obstruction to two finite columns sharing
even a single raw coordinate. It uses the trace on the entire
positive space and works at both marked primes, including the
factors where the matched lattice has singular reduction.

For a rational transport \(u:P_s\to P_t\) with marking \(a\), define
\[
        h(u,a)^{2f}
          =\max\{[M_s:C(u,a)],\ [M_t:u C(u,a)]\}
\]
using the canonical matched lattices of
Section~\ref{sec:boundedcoordinates631}. Its inverse with inverse
marking has the same height. The argument below applies to any
cosets having the stated index bounds, with their complete integrals.
Here \(Q=g_1g_2\) is the product of the two original auxiliary primes.

\begin{lemma}[A trace congruence at every factor of both marks]
\leavevmode\label{lem:matchedtrace638}
Let \(g\) be a closed composite of matched rational isometries and
inverses on one definite space of dimension \(k\), and let
\(\mathfrak J\) be the index ideal of its successive matched domain
\(C\subset M\). If the composite marking at every chosen factor
has trace \(k\) in the residue field, and so does its inverse
transpose at the conjugate factor, then
\[
       \mathfrak J\,(k-\operatorname{tr}g)\subseteq Q\mathcal O_E .
\]
If \(g\ne I\) and \(\rho=\operatorname{rank}_E(g-I)\), it follows
that
\[
               Q\le2\rho\,(N\mathfrak J)^{1/(2f)}
                    \le 2\rho\prod_e h(u_e,a_e).
\]
\end{lemma}
\begin{proof}
At each finite prime choose local integral bases of \(C\) and \(M\).
Write \(T\) for the matrix of inclusion and \(S\) for that of
\(g|_C\). Both are integral, and \(\det T\) generates the local
index ideal. Thus
\[
     (\det T)(k-\operatorname{tr}g)
            = k\det T-\operatorname{tr}(S\,\operatorname{adj}T)
\]
is integral at every prime. At a marked factor there is a matrix
\(a\) for the composite marking such that \(\bar S=a\bar T\).
The displayed integer therefore reduces to
\((\det\bar T)(k-\operatorname{tr}a)\). The formula also holds
when \(\det\bar T=0\), by the adjugate identity. At the conjugate
factor the identical argument uses the inverse transpose in the
dual bases. Thus the displayed integer is divisible by the prime
at every factor over \(g_1g_2\). These primes split completely with
exponent one, giving the ideal inclusion.

The eigenvalues at every embedding are on the unit circle. If
\(k-\operatorname{tr}g=0\), the sum of their real parts is \(k\),
so all eigenvalues are one. Unitarity gives \(g=I\). Otherwise
\(\lvert k-\operatorname{tr}g\rvert\le2\rho\) at each embedding
by summing \(1-z\) over the eigenvalues different from one. The
absolute norm of the integral ideal on the left is at least
\(Q^{2f}\) and at most \((2\rho)^{2f}N\mathfrak J\).
Finally the index at each stage of the successive domain is at
most the corresponding edge index in the appropriate orientation,
as in Theorem~\ref{thm:rationalgram629}. Their product proves the
second bound.
\end{proof}

\begin{theorem}[Disjoint integral coordinates for different rational transports]
\leavevmode\label{thm:disjointtransports638}
Consider two cosets of the same ordinary block with data
\((u,a)\) and \((v,b')\), allowing either common positive dimension.
If \(u\ne v\) and the two marked columns share a coordinate of any
family, then
\[
        Q\le 2\operatorname{rank}_E(u^{-1}v-I)\,h(u,a)h(v,b').
\]
In particular, take any collection in the block of height
\(h(u,a)\le H\) with \(2kH^2<Q\). In every raw integral row family
the unions of occupied coordinates belonging to distinct rational
transports are disjoint. Each signed projected sum at a specified
transport can be recovered by coordinate restriction. On the
original intervals \(H=b^{8-\delta}\) is allowed for each fixed
\(0<\delta<8\) and all sufficiently large \(x\).
\end{theorem}
\begin{proof}
In the matrix of Lemma~\ref{lem:whittakermatrix526}, a common
coordinate implies that the quotient of the two indices is
conjugate to an upper unitriangular element at each chosen factor.
In particular \(a^{-1}b'\) is unipotent there. At the conjugate
factor its inverse transpose is also unipotent. These are the
markings of \(g=u^{-1}v\) on the successive domain obtained from
the edges \(v,u^{-1}\). Their traces are \(k\). Apply the lemma.
The resulting inequality rules out a common coordinate under the
displayed height restriction. Taking unions over all cosets for
each actual rational isometry preserves disjointness.

Coordinate restriction has entries zero and one. Thus it retrieves
each projected sum with all its signs, over any coefficient module
for the raw matrix. The assertion remains true under absolute
convergence for the orbital operators, since coordinate restriction
commutes with that sum. Finally \(k\) is polynomial in \(x\),
\(\log b=x^\varepsilon\), and \(\log Q=(16+o(1))\log b\), by
Theorem~\ref{thm:prime209} and the choice of marks in
Section~\ref{sec:markedranks623}. These estimates give
\(2k b^{16-2\delta}<Q\).
\end{proof}

\begin{corollary}[An exact diagonal Gram block]
\leavevmode\label{cor:exactdiagonal638}
Under the common height restriction \(2kH^2<Q\), let \(Z_{ri,u}\)
be the whole signed contribution at the rational transport \(u\)
in an ordinary block. With \(Z_{ri}=\sum_u Z_{ri,u}\),
\[
              Z_{ri}^* Z_{ri}=\sum_u Z_{ri,u}^*Z_{ri,u}.
\]
This identity holds for any collection of both types of cells.
On the portion of common positive dimension \(k\), write
\(y_{ri,u}=\sum_{\gamma:\,u_\gamma=u}c_\gamma h_{a_\gamma}\)
as in Section~\ref{sec:rowsums637}. For any set of target rows,
the diagonal Gram block on the input column of dimension \(D\) is
exactly
\[
       K_{ii}=m_i I,\qquad m_i=\sum_{r,u}\|y_{ri,u}\|^2,
       \qquad
       \frac{\|K_{ii}-I\|_{\rm HS}^2}{D}=|m_i-1|^2 .
\]
The identities use the given word and require no lower bound on
its compact weight gaps.
\end{corollary}
\begin{proof}
Different transports have mutually orthogonal target coordinate
spaces in the normalized marked table, by the theorem. Thus
\(Z_{ri,u}^*Z_{ri,v}=0\) for \(u\ne v\), including for operators
on the full incoming vector space. On a full-dimensional cell,
Schur's lemma identifies the real operator with a scalar times
the unitary \(U_u\). All cosets at the same \(u\) therefore give
\(Z_{ri,u}=y_{ri,u}\otimes U_u\), whose squared operator is
\(\|y_{ri,u}\|^2 I\). Sum over \(u\) and \(r\). The squared
Hilbert--Schmidt norm of \(I\) is \(D\), which proves the formula.
\end{proof}

\begin{corollary}[Two source rows with a marked identification]
\leavevmode\label{cor:alignedrows638}
Fix equivalent input words on \(P_i,P_j\). Suppose there is a
rational isometry \(w:P_j\to P_i\). Use the identity marked
matrix between their prescribed source coordinates and define
\(H_0=h(w,I)\) with those endpoint lattices. Take
full-dimensional terms of height at most \(H\) on each edge.
If \(2k H_0 H^2<Q\), then the mixed ordinary Gram block is
exactly \(c_{ij}(w)U_w\), with the correlations of
Section~\ref{sec:rowsums637} formed from these terms.
\end{corollary}
\begin{proof}
The marked domain \(C(w,I)\) is defined by the same integrality
and reduction conditions as for an orbital term. It has finite
index, by first clearing the denominators between the two
lattices and then multiplying by \(Q\). For a surviving
contraction through a target row take two individual cosets
at rational maps \(u:P_i\to P_r\), \(v:P_j\to P_r\) whose
marked columns share a coordinate. Consider the closed composite
\(w^{-1}u^{-1}v\) on \(P_j\). Its marked matrix is the relative
matrix of these two columns, as \(w\) has identity marking on its
domain. Both it and its inverse transpose are unipotent. The
trace lemma and the strict size inequality force this composite
to equal \(I\). Therefore every nonzero contraction has
\(u^{-1}v=w\). The exact expansion
\(K_{ij}=\sum_g c_{ij}(g)U_g\) now proves the claim, including the
case in which there are no such contractions.
\end{proof}

The diagonal statistic has thus reduced to the size of the signed
vectors at each actual transport throughout the stated index range.
Terms at a single transport can still cancel. Controlling that
mass, the larger indices and the remaining mixed blocks at the
normalization of the integral theta maps is part of
Section~\ref{sec:missing}.


\subsection{Reduction of the two isotropic spaces at the terminal prime}
\label{sec:terminalouter639}

An integral isometry can induce the identity on the positive column
and at both marks while changing the rational isotropic space. We
measure this change at the terminal prime already present in the
two period tests. The precision below is independent of the degree
of the field. All the operators in this section use the ordinary
sections and the pair measures of Section~\ref{sec:markedranks623}.

Fix two supported rows at a translate in Section~\ref{sec:missing}.
After transporting the second flag into the first row lattice
\(\Lambda\), denote their outer isotropic spaces by \(X_0,X_1\).
Set \(A_j=\Lambda\cap X_j\). Write \(h\) for the Hermitian form,
\(E=E_x\), \(f=[E^+:\QQ]\), \(s=\dim_E X_j\), \(d=s+1\),
and \(T_j\in\{1,T\}\) for the absolute terminal norms. Define the
integral ideal
\[
   \mathfrak c_\gamma=
       \big\langle h(v,w):v\in A_0,\ w\in A_1\big\rangle_{\mathcal O_E}.
\]
The two integral unipotent changes of representative preserve this
ideal. On the transported side they change generators or apply a
lattice isometry preserving the first isotropic space. Since the
reduced outer spaces at each mark agree and are isotropic,
\(\mathfrak c_\gamma\subseteq Q\mathcal O_E\), where \(Q=g_1g_2\).
The use of one row lattice for both flags also occurs in the source
definition of a partner at
\srcloc{Formalism/SharpRestriction/CuspSpanNormBounds/CntPartners}{199}.

\begin{lemma}[The ideal of outer pairings at a point of the orbital integral]
\leavevmode\label{lem:contentdepth639}
Suppose \(s\ge2\) and the two characters agree on their common
unipotent subgroup. At a point \(y\) of the orbital integral put
\[
   B(y)=\sqrt{T_0T_1}\,
     \big(\widetilde V_0(y)\widetilde V_1(\gamma^{-1}y)\big)^{1/f}.
\]
If \(X_0\ne X_1\), then \(\mathfrak c_\gamma\ne0\). For common
positive dimension \(k\) its absolute norm satisfies
\[
               (N\mathfrak c_\gamma)^{1/(2f)}\le B(y).
\]
For common positive dimension \(k-1\) it satisfies
\[
               (N\mathfrak c_\gamma)^{1/(2f)}
                              \le B(y)^{d/(d-2)} .
\]
\end{lemma}
\begin{proof}
Use the fractional ideals \(\mathfrak a_i\) of the ordered frame,
with \(\mathfrak a_d=\mathcal O_E\), and let
\(\mathfrak h_j=\prod_{i\le j}\mathfrak a_i\).
Write \(D_j\) for the determinant of the cross pairing of the first
\(j\) vectors. Put \(J_j=(D_j)\bar{\mathfrak h}_j\mathfrak h_j\).
The frame vectors multiplied by their indicated ideal coefficients
are integral. In local bases the determinant is therefore integral,
so \(J_j\) is an integral ideal, possibly zero. For \(D_j\ne0\)
the exterior Cauchy--Schwarz inequality in
Proposition~\ref{prop:transverseheight582} gives
\[
                   (N J_j)^{1/(2f)}\le B(y)^j .
\]
Indeed the archimedean bound on \(|\tau D_j|\) has the factor
\((N\mathfrak h_j)^{-1/f}(T_0T_1)^{j/2}\) at each place.
Multiplying over one complex embedding above each real place and
taking the \(f\)-th root cancels the norm of the two coefficient
ideals in \(J_j\).

First suppose the common positive dimension is \(k\). Put
\(t=s-\dim(X_0\cap X_1)>0\). When the intersection is nonzero,
Theorem~\ref{thm:success135} and Lemma~\ref{lem:suffixdet556}
give a nonsingular pairing of the two first \(t\)-dimensional
prefixes. Its matrix is the pairing on the quotients by the
intersection, pulled back by the two prefix isomorphisms. When the
intersection is zero, Theorem~\ref{thm:cross132} gives a common
negative vector in the extended flags. Its perpendicular at a real
place has Witt index \(s\). If any nonzero vector in \(X_1\)
annihilated \(X_0\), it would enlarge \(X_0\) there to an isotropic
space of dimension \(s+1\), because the two spaces are disjoint.
Thus \(t=s\) also gives \(D_t\ne0\). All \(t\) entries used in each
determinant monomial are outer pairings. In the integral local
bases they belong to \(\mathfrak c_\gamma\). Consequently
\(0\ne J_t\subseteq\mathfrak c_\gamma^t\). Norms of nonzero
integral ideals increase under passage to a subideal, proving the
first bound.

In common positive dimension \(k-1\), the outer intersection is
zero by Theorem~\ref{thm:success135}. The positive Gram invariant
in Theorem~\ref{thm:cross132} is
\(D_d\bar D_d-T_0T_1D_s\bar D_s\), so \(D_d\ne0\), including
when \(D_s=0\). In a monomial of \(D_d\), at most two factors
use the terminal row or the terminal column. The other \(d-2\)
factors belong to \(\mathfrak c_\gamma\) after the local frame
coefficients have been included. The two remaining factors are
integral. Hence \(0\ne J_d\subseteq\mathfrak c_\gamma^{d-2}\).
Here \(d-2>0\). The bound for \(J_d\) proves the claim.
\end{proof}

For a nonzero pairing ideal define
\[
   \delta_T(\gamma)=\frac1{2f}\sum_{\mathfrak P\mid T}
                            v_{\mathfrak P}(\mathfrak c_\gamma).
\]
The valuations assign value one to a uniformizer. Set
\(\delta_T(\gamma)=+\infty\) when \(\mathfrak c_\gamma=0\).
All \(2f\) primes over \(T\) have degree one.

\begin{theorem}[Uniform removal of distinct outer spaces with deep pairings]
\leavevmode\label{thm:terminalouter639}
For each \(L>0\) there is an integer \(a_L\ge1\) such that on
every sufficiently distant common interval, at every translate
and on every retained pair of rows \(r,i\),
\[
  \sum_{\substack{\gamma\ {\rm compatible}\\
              X_0\ne X_1,\ \delta_T(\gamma)\ge a_L}}
                \|P_\gamma^{ri}\|_{\rm HS}
              \le \sigma_{ri}\exp(-L U_x).
\]
Here \(P_\gamma^{ri}\) includes the original finite phase and
pair coefficient, and \(\sigma_{ri},U_x\) are defined in
Section~\ref{sec:markedranks623}. The sum may be restricted
further by the actual transport, common hyperplane, or marked
index.
\end{theorem}
\begin{proof}
Use the two simultaneous product cutoffs of
Lemma~\ref{lem:productcut629} for this \(L\). At a retained
point \(B(y)\le T b^{2C's}\) for a uniform \(C'\), with
\(b=\exp(x^\varepsilon)\). On the given intervals \(s\ge2\)
eventually, \(d/(d-2)\le3\), and
\(\log T=(34+o(1))d\log b\) by Theorem~\ref{thm:prime209}.
It follows that one can choose \(a_L\) independently of the rows
and of \(x\) so that \((Tb^{2C's})^3<T^{a_L}\).
For a nonzero pairing ideal the marked divisibility and the
factorization at the distinct prime \(T\) give
\[
                 (N\mathfrak c_\gamma)^{1/(2f)}
                                  \ge Q T^{\delta_T(\gamma)} .
\]
The lemma now excludes every term in the displayed sum from the
jointly retained domain. The sum of the integrals of the norms
on its complement, over all the original cosets, is bounded by
Lemma~\ref{lem:productcut629}. Integration followed by the
triangle inequality proves the result for every indicated subset.
\end{proof}

\begin{corollary}[The identity coefficient for equal terminal reductions]
\leavevmode\label{cor:terminalidentity639}
In the original expansion of \(Z_{ri}\), sum over the cosets for
which \(A_0,A_1\) have equal images in \(\Lambda/T^{a_L}\Lambda\).
Denote the resulting operator by \(Z_{ri}^{[a_L]}\), and denote
its flag-preserving part from Theorem~\ref{thm:mixedordinary541}
by \(\mathcal B_{ri}\). Then
\[
         \|Z_{ri}^{[a_L]}-\mathcal B_{ri}\|_{\rm HS}
                      \le \sigma_{ri}\exp(-LU_x).
\]
On any repeated row \(i\), contraction against its original
marked identity column gives
\[
 \big\|(h_{\iota_i}^{*}\otimes1)Z_{ii}^{[a_L]}
                 -\beta_i1_{V_i}\big\|_{\rm HS}
                \le |\beta_i|\exp(-LU_x),
           \qquad |\beta_i|=\sigma_{ii}.
\]
In particular the least singular value of this contraction is
at least \(|\beta_i|(1-\exp(-LU_x))\).
\end{corollary}
\begin{proof}
Equality of the two images is unchanged by the integral
unipotents in the double quotient. For any \(w\in A_1\) choose
\(v\in A_0\) with \(w-v\in T^{a_L}\Lambda\). Pairing with
\(A_0\) gives \(\mathfrak c_\gamma\subseteq T^{a_L}\mathcal O_E\)
by isotropy and integrality. All the compatible terms with
distinct rational outer spaces therefore satisfy the theorem.

It remains to identify the terms with equal rational outer
spaces. In this case the two saturated lattices are equal.
The common-suffix argument in Lemma~\ref{lem:suffixdet556}
identifies the whole ordered filtration. In the two ordered
bases the successive ratios are a single scalar \(\lambda\),
including the terminal ratio modulo the outer space. The two
index ideals in that lemma are the unit ideal, so
\((\lambda)^s=\mathcal O_E\) as fractional ideals.
Thus \(\lambda\) is a unit and
\(\lambda\bar\lambda=T_1/T_0\). Taking absolute norms gives
\(T_0=T_1\). All complex absolute values of \(\lambda\)
are one. Its positive powers lie in a bounded subset of the
lattice \(\mathcal O_E\) in its archimedean embedding, so it
is a root of unity. The first graded line is primitive
locally at the marks, and the two marked frames agree there
up to upper unitriangular change. Hence \(\lambda\) reduces
to one at both residue characteristics. A torsion unit reducing
to one has order a power of the residue characteristic. Indeed,
if its prime-to-residue-characteristic part \(u\ne1\) has order
\(m\), then \(1+u+\cdots+u^{m-1}=0\) would reduce to the
nonzero residue of \(m\). The two characteristics are distinct,
so \(\lambda=1\). This gives exactly the rational equations
defining the flag-preserving terms, and each such term belongs
to the selected sum.

The finite columns have unit norm. The theorem and the triangle
inequality give the first bound. On a repeated row every
flag-preserving term other than the identity has finite support
disjoint from the identity column by
Theorem~\ref{thm:flagisometry544}. Its contraction is zero.
The identity gives precisely the scalar in
Theorem~\ref{thm:identitycoefficient551}. Contraction has norm
one, and operator norm is at most Hilbert--Schmidt norm. Applying
the reverse triangle inequality to each unit input proves the
singular-value assertion.
\end{proof}

For example, on every retained diagonal row one may take the
cyclic groups of Theorem~\ref{thm:everydiagonal536} with
\((\gamma-1)\Lambda\subseteq T^{a_L}\Lambda\). Every power then
has equal outer reductions. The union of their nonidentity
terms, counting each coset once over all rational hyperplanes,
has total norm at most \(|\beta_i|\exp(-LU_x)\). Each such term
has the identity transport and the identity marking. For the
full identity coefficient one must also sum the signed terms
whose two outer reductions differ at this precision. They
occur in the original integral comparison of Section~\ref{sec:missing}.


\subsection{Recovering a rational orbital class from its terminal reduction}
\label{sec:terminalorbits640}

Fix the original source and target rows on a common interval. Write
\(T_0,T_1\in\{1,T\}\) for their terminal norms, \(f=[E^+:\QQ]\)
and \(d=s+1\). We study the terms with common positive dimension
\(k-1\). In one such term the cross matrix records the pairings of
the \(d\) ordered frame vectors at the two ends. Its rational
class \([C]\) allows lower unitriangular changes of rows and upper
unitriangular changes of columns over \(E\). Integral changes
also occur in the source definition of Gram equivalence at
\srcloc{Formalism/SharpRestriction/CuspSpanNormBounds/IsPrimitiveSub}{281}.
Several rational unipotent double classes can have the same
\([C]\). We recover the additional parameter as well.

For a positive integer \(a\), reduce an actual integral row
isometry \(\gamma:\Lambda_1\to\Lambda_0\) modulo \(T^a\) at
all the places above \(T\). Take its double class under the
images of the two local integral row unipotent groups. Denote
this finite class by \(r_a(\gamma)\). Equality means that at
each place one can compose a representative with the two local
unipotents and obtain the other representative modulo \(T^a\).
No lifting of those local changes to a simultaneous rational
change is required. This defines a partition of the terms of
the original orbital sum. We use its Haar measures and its
original finite phases throughout.

\begin{lemma}[Separation of the cross matrices at the terminal prime]
\leavevmode\label{lem:terminalpivots640}
Use the two product cutoffs of Lemma~\ref{lem:productcut629}
with exponent \(L>0\). There is a positive integer \(b_L\),
uniform on the common intervals, with the following property.
Suppose two compatible terms meet the cutoffs, possibly at
different archimedean arguments. If their reductions
\(r_{b_L d}\) agree, then either both have common positive
dimension \(k\), or both have common positive dimension \(k-1\)
and their rational cross matrices are equivalent.
\end{lemma}
\begin{proof}
Put \(\mathfrak a_I=\prod_{i\in I}\mathfrak a_i\),
\(\eta_I=(N\mathfrak a_I)^{-1/(2f)}\) and
\([j]=\{1,\ldots,j\}\) as in
Proposition~\ref{prop:pivotflags619}. At a retained point for
\(\gamma\) put \(p_{\gamma,\tau}\) equal to the product of the
local translated heights of the two rows. In particular
\[
        p_{\gamma,\tau}\ge1,\qquad
        \prod_\tau p_{\gamma,\tau}\le b^{2C'fs},
                    \qquad b=\exp(x^\varepsilon).
\]
Here \(C'\) is the constant for the cutoffs. For each leading
minor \(D_j\), exterior Cauchy--Schwarz at the single place gives
\[
        |\tau D_j|/\eta_{[j]}^2\le
                 (\sqrt{T_0T_1}\,p_{\gamma,\tau})^j
\]
by Proposition~\ref{prop:transverseheight582}. This includes
vanishing minors and both possible common dimensions.

We give the separation argument for a minor in rows \([j]\) and
columns \(K\), with \(|K|=j\), whenever its values \(z,z'\)
are invariant under the frame changes in question and satisfy
the preceding local bound with denominator
\(\eta_{[j]}\eta_K\). Put
\(\mathfrak b_{j,K}=\overline{\mathfrak a_{[j]}}\mathfrak a_K\).
Each value times this coefficient ideal is integral, by taking
the exterior pairing in local integral bases. At the terminal
places the coefficient ideals are units, since \(T\) splits
completely and is prime to the different used in the outer
ideals. Congruence of the two row maps up to integral
unipotents modulo \(T^a\) implies
\[
          (z-z')\mathfrak b_{j,K}\subseteq T^a\mathcal O_E .
\]
Indeed, after the local frame changes the reduced matrices
agree, so the determinants of the indicated submatrices agree.
The integral unipotents act by the stated triangular changes
on the ordered frame, leaving the relevant invariant unchanged.
The inclusion at the other places is just integrality.

If \(z-z'\ne0\), use the triangle inequality at each place.
Since both local height products are at least one, the
maximum of the two is at most their product. Multiplication
over the \(f\) places gives
\[
 T^a\ \le\ \bigl(N((z-z')\mathfrak b_{j,K})\bigr)^{1/(2f)}
       \ \le\ 2\bigl(\sqrt{T_0T_1}\,b^{4C's}\bigr)^j .
\]
This use of the placewise product allows the two retained
arguments to have their largest coordinates at different places.
Choose \(b_L\) so that \(T^{b_L}>2T b^{4C's}\) uniformly for
sufficiently large \(x\), using
\(\log T=(34+o(1))d\log b\). Thus already \(a\ge b_L j\)
forces equality of the two values.

Apply this first with \(K=[j]\) to every leading minor.
These are invariant under both triangular changes. Equality
of \(D_s,D_d\) gives equality of the Gram invariant
\(D_d\overline{D_d}-T_0T_1D_s\overline{D_s}\).
In common positive dimension \(k\) the concatenated frame is
dependent, so its Gram determinant vanishes. In common positive
dimension \(k-1\) the outer spaces are disjoint by relevance,
and the invariant is totally positive by
Theorems~\ref{thm:success135} and~\ref{thm:cross132}. This
distinguishes the two cases.

In the latter case the nonzero leading minors give the endpoints
of the successive reversed blocks by Lemma~\ref{lem:forms203}.
Hence both terms have the same pivot permutation \(w\).
Now put \(K_j=w([j])\), and let \(M_j\) be the minor in rows
\([j]\) and columns \(K_j\) in increasing order. These minors
are invariant by Proposition~\ref{prop:pivotflags619}. The
triangular proof of invariance works over every completion as
well. At a block endpoint \(j\), \(M_j=D_j\) and \(K_j=[j]\).
Inside a block the absolute value of
\(\tau M_j/(\eta_{[j]}\eta_{K_j})\) interpolates geometrically
between the two endpoint values, including the initial value
one at \(j=0\). The bound
\((\sqrt{T_0T_1}\,p_{\gamma,\tau})^j\) interpolates in exactly
the same way. The separation argument therefore gives equality
of all the \(M_j\), including those with \(D_j=0\). For a block
beginning after position \(u\), its first antidiagonal entry is
\(M_{u+1}/D_u\), with \(D_0=1\). This follows by taking the
first row of that block in its monomial representative. The
other entries of the block alternate in sign. Thus the two
normal forms, including the entries themselves, agree.
\end{proof}

\begin{theorem}[Recovery of the rational unipotent double class]
\leavevmode\label{thm:terminalclass640}
Fix the two projected rational hyperplanes \(H_0,H_1\), their
rational transport \(j\) with its prescribed determinant, and
the full marked index \(i\) in an ordinary hyperplane sum.
For each \(L>0\) there is an integer \(m_L\) independent of all
these choices, the rows, \(x\), and the interval, such that a
terminal reduction \(r_{m_L d}\) contains at most one rational
unipotent double class with an integral lift meeting the two
product cutoffs. The assertion concerns all the supported
rational orientations over each cross matrix.
\end{theorem}
\begin{proof}
Increase \(m_L\) so that it is at least \(b_L\). For two lifts
as in the statement the lemma supplies a single rational cross
matrix class. In the orthogonal complement of \(H_0\) put
\(L'=\operatorname{pr}_{H_0^\perp}\Lambda_0\). On its
determinant line choose as base the exterior vector \(\Omega_0\)
of one of the supported rational classes, formed from the two
ordered frames as in Theorem~\ref{thm:frameorientation611}.
The other classes over this cross matrix are determined uniquely
by their vectors \(\alpha\Omega_0\), where
\(\alpha\bar\alpha=1\). The normalizing ideal
\(\mathfrak a\) and the Gram invariant \(\mathcal G\) satisfy
\[
 Q^{-d}\mathfrak h_d^2\Omega_0=\mathfrak a\,\det L',
 \qquad
 N\mathfrak a\le X_C^2,\qquad
 X_C=Q^{-2df}(N\mathfrak h_d)^2|N_{E^+/\QQ}\mathcal G|
\]
by Proposition~\ref{prop:rotationideal611}.
Here \(\mathfrak h_d=\mathfrak a_{[d]}\).
The positivity formula
\(\tau\mathcal G\le|\tau D_d|^2\) and the leading-minor bound
in the lemma, applied to one retained point, give
\[
 X_C^{1/(2df)}\le
             Q^{-1}\sqrt{T_0T_1}\,b^{2C's}.
\]
In particular \((N\mathfrak a)^{1/(2f)}
                     \le (T b^{2C's})^{2d}\).

We check that the finite reduction compares the projected
exterior vectors themselves. For an actual integral lift take
the exterior product of the first frame and the transported
second frame and apply \(\bigwedge^{2d}\operatorname{pr}_{H_0^\perp}\).
The result is \(\alpha\Omega_0\). As in the proof of
Proposition~\ref{prop:rotationideal611}, a source unipotent
fixes its frame wedge. A target unipotent applied to a reduced
representative has, after projection on \(H_0^\perp\), block
upper triangular action there, with determinant one. This also
applies to a product of that unipotent and the one used to put
the representative in its smaller space. The assertion therefore
holds after any local integral unipotent changes of the row map.
Projection is integral with values in \(L'\) by the definition
of this lattice, even if \(L'\) is larger than
\(\Lambda_0\cap H_0^\perp\).

For two classes with parameters \(\alpha,\alpha'\), congruence
of the changed maps modulo \(T^a\) gives congruence of the two
projected wedges there by multilinearity. The factor
\(Q^{-d}\mathfrak h_d^2\) is a unit at \(T\).
Away from \(T\) both normalized wedges already lie in
\(\det L'\). Equivalently,
\[
             (\alpha-\alpha')\mathfrak a\subseteq T^a\mathcal O_E .
\]
If \(\alpha\ne\alpha'\), then
\(|\tau(\alpha-\alpha')|\le2\) at every place, since both
parameters have absolute value one. Consequently
\[
                  T^a\le 2(N\mathfrak a)^{1/(2f)}
                              \le 2(Tb^{2C's})^{2d}.
\]
By the same growth of \(T\), an integer \(m_L\) can be chosen
so that \(a=m_Ld\) violates this inequality throughout the
intervals. The parameters agree, and
Theorem~\ref{thm:frameorientation611} gives the same rational
unipotent double class.
\end{proof}

\begin{corollary}[Signed coefficients grouped at one finite precision]
\leavevmode\label{cor:terminalamplitudes640}
In the fixed hyperplane sum of the theorem let \(B_\eta\) be
the signed ordinary operator summed over all integral lifts of
the rational unipotent double class \(\eta\). Form the table of
signed operators at the finite indices \(r_{m_Ld}\), using the
original terms, including their full integration domains. Assign
each entry containing a lift meeting the cutoffs to its unique
class in the theorem, and omit the other entries. Let
\(\widehat B_\eta\) be the sum assigned to \(\eta\), with an
empty sum interpreted as zero. Then
\[
             \sum_\eta\|B_\eta-\widehat B_\eta\|_{\rm HS}
                           \le 2\sigma_{ri}\exp(-LU_x),
\]
where \(r,i\) in \(\sigma_{ri}\) denote the two rows. The same
bound holds after any selection of the original integral
terms made in both tables.
\end{corollary}
\begin{proof}
A lift whose domain meets the product cutoffs contributes its
entire signed operator once to the correct entry on both sides.
A lift with no such point contributes its norm at most twice
to the sum of the errors, even if its finite class is assigned
to a different rational class. The total of these norms is at
most \(\sigma_{ri}\exp(-LU_x)\) by
Lemma~\ref{lem:productcut629}, which bounds the integrals of
the norms on the discarded domains before grouping. Absolute
summability justifies both groupings and the triangle inequality.
\end{proof}

Thus in a fixed hyperplane and transport the signed coefficients
of different rational classes can be separated at the terminal
prime at the scale of the diagonal normalization. Integral
lifts in one such class still contribute their original signed
local factors and real tensor as in
Theorem~\ref{thm:liftphase604}. Bounding the sum of these
coefficients against the prescribed period, together with the
other common dimensions and the cohomology estimate, is required
for Section~\ref{sec:missing}.


\subsection{The positive transport at the terminal prime}
\label{sec:terminaltransport641}

We remove the fixed-hyperplane and fixed-transport assumptions
from the recovery of the rational classes in
Section~\ref{sec:terminalorbits640}. The argument also identifies
the transport on terms of full common positive dimension. It uses
an operator between the two positive spaces that has norm at most
one at every infinite place, even when the span of the two frames
has a radical.

Use the integral Witt decomposition of Lemma~\ref{lem:projectivewitt554}
at each end. Denote its middle summand by \(J_b\), the terminal
vector by \(z_b\), and its negative norm by \(-T_b\), \(b=0,1\).
Regard \(P_b=z_b^\perp\cap (J_b\otimes E)\) as a subspace of
the ambient Hermitian space \(V_b\), and denote the orthogonal
projection onto it by \(p_b\). The map \(T_bp_b\) preserves
\(\Lambda_b\). In fact projection onto \(J_b\) is integral by
the direct sum, and on \(J_b\) the formula for \(T_bp_b\) is
\(v\mapsto T_bv+z_bh(z_b,v)\).

\begin{lemma}[A projected operator in both common dimensions]
\leavevmode\label{lem:positiveadjugate641}
Let \(\gamma:\Lambda_1\to\Lambda_0\) represent a compatible
term. Work in \(V_0\), write \(X_b,Y_b\) for the outer and
extended spaces there, transporting side \(1\) by \(\gamma\).
Put \(I=X_0\cap X_1\), \(t=s-\dim I\), and \(W=Y_0+Y_1\).
Choose the following list of columns \(\mathcal Y_\gamma\) in
the target lattice.

\begin{itemize}
\item In common positive dimension \(k-1\), take both complete
ordered frames, target first.
\item In common positive dimension \(k\), take the first \(t\)
outer vectors from each frame, followed by \(z_0\). For \(t=0\)
the list consists of \(z_0\).
\end{itemize}

Let \(G_\gamma=\mathcal Y_\gamma^\dagger\mathcal Y_\gamma\) be
their Hermitian Gram matrix, and set \(\Delta_\gamma=\det G_\gamma\).
Here \(\mathcal Y_\gamma^\dagger v\) has entries \(h(y,v)\) for
the columns \(y\). In the respective cases
\[
 \Delta_\gamma=(-1)^d\mathcal G,\qquad
 \Delta_\gamma=(-1)^{t+1}T_0D_t\overline{D_t}.
\]
Both are nonzero algebraic integers of \(E^+\). Put
\[
 S_\gamma=p_0
   (1-\mathcal Y_\gamma G_\gamma^{-1}\mathcal Y_\gamma^\dagger)
                   \gamma p_1 : V_1\longrightarrow V_0 .
\]
On \(P_1\) this is the isometry of the two projected common
spaces, extended by zero on the perpendicular in \(P_1\).
It vanishes on \(P_1^\perp\). In particular it is invariant
under the rational row unipotents on both sides. The same
identity under frame changes holds over the completions at \(T\).
\end{lemma}
\begin{proof}
The frame vectors are integral since their coefficient ideals
are nonpositive powers of the different, and the terminal
vectors are integral. In the first case the two complete frames
span \(W\) nondegenerately. Their determinant is the signed Gram
invariant by \eqref{eq:cross132}. This determinant formula for
the concatenated Gram matrix also occurs at
\srcloc{Formalism/SharpRestriction/CuspSpanNormBounds/LatticeGluingRadicals}{447}.
Let \(\Pi\) be the middle factor in the formula for \(S_\gamma\).
It is orthogonal projection onto \(W^\perp\), since it kills
the displayed columns and fixes everything perpendicular to
them. The space \(W^\perp\) projects isometrically onto a
hyperplane at each end, as in Theorem~\ref{thm:smallcomplement604}.
Indeed each of its vectors on side \(b\) is a vector of that
positive hyperplane plus an outer vector. Pairing a vector of
\(P_b\) with this graph is simply its pairing with the positive
component. Thus orthogonal projection onto the graph, followed
by projection onto the target positive space, identifies the
two copies of the hyperplane and kills the source normal line.

In common dimension \(k\), \(D_t\ne0\) for \(t>0\) by the
prefix argument in Lemma~\ref{lem:contentdepth639}. Put \(D_0=1\).
The outer prefixes together span a nondegenerate hyperbolic
space with Gram determinant \((-1)^tD_t\overline{D_t}\).
The vector \(z_0\) is orthogonal to its own outer prefix.
Add a vector of that prefix to it to kill its pairings with the
opposite prefix, by nonsingularity. This leaves its norm
\(-T_0\) unchanged and is a triangular change of the displayed
list of determinant one. The second determinant formula follows.
Write \(W'\) for the span of that list. A common negative vector
in \(Y_0\cap Y_1\) exists by Theorem~\ref{thm:cross132}.
The extended spaces have intersection of dimension \(\dim I+1\),
so \(\dim W=2t+1+\dim I\). Every vector of \(I\) is orthogonal
to both extended frames. Since \(W'\) is nondegenerate and has
dimension \(2t+1\), \(W=W'\oplus I\).

Now \(\Pi\) is projection onto \((W')^\perp\). For \(v\in P_1\),
\(\gamma v\) lies in \(Y_1^\perp\). Both \(\gamma v\) and its
projection are perpendicular to \(I\), so
\(\Pi\gamma v\in W^\perp\). Their difference lies in
\(W\cap Y_1^\perp=X_1\). To check the latter identity, pass to
the nondegenerate \(W/I\) of dimension \(2t+1\). The image of
\(Y_1\) has dimension \(t+1\), and its perpendicular there
contains \(X_1/I\) of dimension \(t\), hence equals it.
Thus \(\Pi\gamma v\) represents the same vector in
\(Y_1^\perp/X_1\) as \(\gamma v\). Its image in
\(Y_0^\perp/X_0\) gives exactly the common-quotient isometry
of the two positive spaces.

In both cases the resulting map on the positive quotients
is intrinsic to the pair of frames. A row unipotent preserves
the extended flag and induces the identity on its positive
quotient. It therefore leaves \(S_\gamma\) unchanged under the
two-sided change of representative. The preceding arguments
use the defining subspaces, orthogonal projections and
nonzero determinants over \(E\). Extending these identities
to the paired factors of a split completion gives the same
invariance there.
\end{proof}

\begin{theorem}[The positive spaces and transport from terminal reduction]
\leavevmode\label{thm:allterminaltransports641}
For each \(L>0\) there is an integer \(c_L\) uniform on the
fields, rows and translates of Section~\ref{sec:missing}.
Two compatible terms of a given row pair that meet the product
cutoffs and have the same \(r_{c_Ld}\) have the same common
dimension and rational positive transport. In common dimension
\(k-1\) they have the same two rational hyperplanes. If their
original marked index is also the same, they have the same
rational unipotent double class.
\end{theorem}
\begin{proof}
Take \(c_L\ge b_L\) of Lemma~\ref{lem:terminalpivots640}. Its
proof gives equality of all leading minors and of the common
dimension. In full common dimension the rank of the outer
pairing is \(t\). For nonzero intersection this is
Theorem~\ref{thm:success135}, and for zero intersection it is
the negative-vector argument of Lemma~\ref{lem:contentdepth639}.
Thus \(t\) is the largest \(j\le s\) with \(D_j\ne0\), allowing
\(j=0\). In particular the same list positions and the same
scalar \(\Delta=\Delta_\gamma=\Delta_{\gamma'}\) occur in
the lemma for both terms.

Put \(D=T_0T_1\Delta\). The adjugate formula gives an integral
map \(\Lambda_1\to\Lambda_0\) for each \(D S_\gamma\). Explicitly
it is
\[
  (T_0p_0)\bigl(\Delta\,1-
     \mathcal Y_\gamma\operatorname{adj}(G_\gamma)
                           \mathcal Y_\gamma^\dagger\bigr)
                                      \gamma(T_1p_1).
\]
The two outside projections with their indicated factors
preserve the integral lattices. Everything in the parentheses
is integral and polynomial in the columns and their Hermitian
pairings. At the places above \(T\), use the integral
unipotents giving equality modulo \(T^a\), \(a=c_Ld\).
For these changed maps the columns in the list and their Gram
matrices agree modulo \(T^a\). This comparison uses both
conjugate factors in each Hermitian pairing. The formulas for
\(\Delta\) are invariant under those changes, and so is
\(S_\gamma\) by the lemma. Consequently
\[
      D(S_\gamma-S_{\gamma'})\Lambda_1\subseteq T^a\Lambda_0 .
\]
The inclusion at the other finite places is already given by
integrality. This calculation also applies when the two Gram
matrices become singular modulo \(T\).

Let \(B=S_\gamma-S_{\gamma'}\), and take Hermitian adjoints
in the original row spaces. The rows are self-dual integral
lattices, so \(D B/T^a\) and its adjoint are both integral
maps between them. The trace of their product is integral,
as is seen in a local lattice basis at each finite prime.
At an infinite place \(B\) is supported from \(P_1\) to
\(P_0\). Each of its two summands there is a partial isometry
on the positive spaces. Hence
\[
          \tau\operatorname{tr}(B^\dagger B)
                     =\|B_\tau|_{P_1}\|_{\rm HS}^2\le4k .
\]
If \(B\ne0\), these real numbers are strictly positive at every
embedding. The integral trace just obtained is a nonzero
algebraic integer of \(E^+\), and \(D\in E^+\). Its absolute
norm is at least one, which gives
\[
                    T^a\le 2\sqrt{k}\,
                            |N_{E^+/\QQ}D|^{1/f}.
\]

We bound the right side on the two retained domains. For every
leading minor, the local inequality in the proof of
Lemma~\ref{lem:terminalpivots640} gives
\[
 |N_{E/\QQ} D_j|^{1/f}
       \le \eta_{[j]}^4\bigl(\sqrt{T_0T_1}\,b^{2C's}\bigr)^{2j}.
\]
For a hyperplane term use
\(0<\tau\mathcal G\le|\tau D_d|^2\). For a full-dimensional
term use the second determinant formula of the lemma with
\(0\le t\le s\). The prescribed ideals give
\(\log\eta_{[j]}=O(d^2\log(2x))\), since they are products
of powers of the different and its root norm is polynomial
in \(x\), as in Section~\ref{sec:exteriormoments620}.
Together with \(T_b\le T\), polynomial growth of \(k\), and
\(\log T=(34+o(1))d\log b\), the resulting estimate is
\[
        \log\!\left(2\sqrt{k}\,|N_{E^+/\QQ}D|^{1/f}\right)
                         \le C_L d\log T .
\]
Choose \(c_L>C_L\). This rules out \(B\ne0\). The image and
the perpendicular to the kernel of the common partial isometry
in the two positive spaces give the same two hyperplanes in
common dimension \(k-1\). In common dimension \(k\) it already
gives the full transport.

In the hyperplane case the remaining action on the normal
line is determined by the determinant. The ratio of the two
ambient determinants is a unit of \(\mathcal O_E\), since
both maps identify the two fixed determinant lattices. Its
absolute value is one at all infinite places, by unitarity.
The local unipotents have determinant one, so the ratio is
one modulo \(T^a\) at all primes above \(T\). Its difference
from one, if nonzero, would have absolute norm at least
\(T^{2fa}\) and at most \(2^{2f}\), an impossibility.
The determinant-line identifications of
Lemma~\ref{lem:completedcolumn535} depend only on the two
fixed rows and their terminal frequencies. Thus the ambient
determinants give the same positive determinant map. With
the restrictions to the hyperplanes already equal, uniqueness
on the normal line gives equality of the rational transports.
Finally increase \(c_L\) to at least \(m_L\) of
Theorem~\ref{thm:terminalclass640}. At a common original marked
index that theorem now applies with the recovered hyperplanes
and transport, and gives the asserted rational unipotent class.
\end{proof}

\begin{corollary}[Simultaneous grouping on a whole ordinary block]
\leavevmode\label{cor:allterminalcoeff641}
Group the original signed terms of one row pair by their marked
index, common dimension and transport, and in common dimension
\(k-1\) also by the two hyperplanes and rational unipotent
class. Call the corresponding sums \(B_\lambda\). Alternatively
sum first at the finite indices \((i(\gamma),r_{c_Ld}(\gamma))\),
keeping the original unit marked columns. Assign each occupied
finite entry meeting the cutoffs to its recovered label, and
discard entries without a retained point. The resulting
\(\widehat B_\lambda\) satisfy
\[
       \sum_\lambda\|B_\lambda-\widehat B_\lambda\|_{\rm HS}
                             \le 2\sigma_{ri}\exp(-LU_x).
\]
The constant applies to the entire block through both common
dimensions and all its hyperplanes and transports together.
\end{corollary}
\begin{proof}
The theorem gives a single label for all the lifts meeting the
cutoffs in any finite entry. Each of their full operators is
assigned correctly. Every other misplaced or discarded lift
has its entire integration domain outside the simultaneous
cutoffs and contributes its norm at most twice. Sum the
integrals of the norms over this one row pair by
Lemma~\ref{lem:productcut629}. Tensoring with the unit marked
column preserves the Hilbert--Schmidt norm.
\end{proof}

This recovers the positive transports from the finite entries
of the ordinary period while keeping both original markings.
In the integral projected period the entries are summed with
their original phases. Bounds for that sum, its arithmetic
normalization, and the strict cohomology on the same intervals
are still required in Section~\ref{sec:missing}.


\subsection{The rational class when the common positive space has full dimension}
\label{sec:fullrationalclass645}

Use the pair of rows and the two product cutoffs of
Section~\ref{sec:terminaltransport641}. The terminal vectors
\(z_b\) have norms \(-T_b\), \(T_b\in\{1,T\}\), for \(b=0,1\).
Denote the ordered frames by \(y_b=(e_{b,1},\ldots,e_{b,s},z_b)\),
where \(d=s+1\). For an isometry \(\gamma:V_1\to V_0\) put
\(Z_{uv}=h(y_{0,u},\gamma y_{1,v})\), with \(h\) linear in its
second variable. The two outer spans, viewed in \(V_0\), are
\(X_0,X_1\); adjoining the terminal vectors gives \(Y_0,Y_1\).
Write \(I=X_0\cap X_1\), \(t=s-\dim I\), and \(D_j=\det
Z_{[j],[j]}\), \(D_0=1\). Assume first that the common positive
dimension is \(k\). Then \(D_t\ne0\) and \(Y_0\cap Y_1\)
has a negative line modulo \(I\), by
Lemma~\ref{lem:positiveadjugate641}.

\begin{lemma}[A minor determining the relative terminal vector]
\leavevmode\label{lem:borderminus645}
Put \(J=[t]\cup\{d\}\) in increasing order. The number
\[
 B_t=\det Z_{J,J}=-T_0D_t\lambda
\]
is invariant under the two rational row unipotents. Here
\(\lambda\) is determined by the images of the two prescribed
terminal vectors on the common negative line, with the source
representative equal to \(\lambda\) times the target representative
modulo \(I\). In particular \(\lambda\bar\lambda=T_1/T_0\).
For each fixed cutoff exponent \(L>0\), a common reduction
\(r_{a}\) with \(a=c_Ld\) determines \(t,\lambda\) and the
reversed blocks and their scalars on the outer prefixes for
all compatible terms meeting the cutoffs, with \(c_L\) uniform.
\end{lemma}
\begin{proof}
The two prefixes of length \(t\) have invertible mutual pairing
and complement \(I\) in their respective outer spans. Add a
vector of its own prefix to each terminal vector to make it
perpendicular to the opposite prefix. Nonsingularity gives
existence and uniqueness of the corrections. Write \(z'_0,z'_1\)
for the resulting vectors in \(V_0\). The corrections preserve
their norms, and both vectors are perpendicular to \(I\).
In \((Y_0+Y_1)/I\) the two prefixes span a nondegenerate
space of dimension \(2t\); its perpendicular there is the
one negative line. Thus \(z'_1=\lambda z'_0\bmod I\).
The scalar product is \(-T_0\lambda\). In the determinant
defining \(B_t\), the corrections add combinations of previous
rows or columns and give a block diagonal matrix. This proves
the formula, including \(t=0\). The common negative line
\((Y_0\cap Y_1)/I\) maps isomorphically to both \(Y_b/X_b\).
It defines the same terminal ratio independently of the
representatives. Since \(D_t\) is a leading minor and hence
invariant, the formula proves the invariance of \(B_t\),
including after extension to the paired factors above \(T\).

By the proof of Theorem~\ref{thm:allterminaltransports641}
the reduction at the stated order of precision determines all
\(D_j\), and \(t\) is the last nonzero leading index up to \(s\),
with \(0\) allowed. For \(t>0\) the \(t\)-square is in reversed
blocks. For \(I\ne0\) use the prefix argument of
Theorem~\ref{thm:suffix237}; at unequal terminal norms the
common terminal ratio is that of Lemma~\ref{lem:suffixdet556}.
For \(I=0\) the inversion test of Proposition~\ref{prop:class236}
acts on the two prefixes, fixing their common perpendicular
and its negative line. In both cases equality of the simple
characters for these tests imposes the same consecutive
reversals and alternating pivots, independently of \(T_0,T_1\).
The nonzero \(D_j\), \(j\le t\), give their endpoints.
For each \(j\le t\), take the minor in rows \([j]\) and
columns \(w([j])\) for the resulting permutation \(w\)
of \([t]\). Its invariance is the triangular minors argument
of Proposition~\ref{prop:pivotflags619}, applied to the two
prefixes. Inside a block its normalized absolute value at
each embedding interpolates geometrically between the two
endpoint values: the new pivot factor is constant in modulus,
and the two new outer-ideal exponents have constant sum at
each step. The proof of Lemma~\ref{lem:terminalpivots640}
therefore bounds and separates these minors up to \(t\)
at precision \(c_Ld\). The first pivot in each block is its
first new minor divided by the preceding endpoint. This
recovers the scalars also for a proper prefix.

It remains to separate \(B_t\); its rows skip the common
suffix when \(t<s\). Put \(\mathfrak h_t=\prod_{i\le t}
\mathfrak a_i\) and \(\eta_{[t]}=(N\mathfrak h_t)^{-1/(2f)}\)
with \(f=[E^+:\QQ]\). The terminal vectors are integral, so
\(B_t\bar{\mathfrak h}_t\mathfrak h_t\) is integral. A
congruence of the row maps up to the local integral unipotents
implies
\[
 (B_t-B'_t)\bar{\mathfrak h}_t\mathfrak h_t
                                  \subseteq T^a\mathcal O_E .
\]
Indeed change the maps locally to congruent representatives
and take the indicated minor, using its invariance and the
fact that the outer ideals are units above \(T\).
At a retained archimedean argument the leading minor bound
of Lemma~\ref{lem:terminalpivots640} now gives
\[
 \frac{|\tau B_t|}{\eta_{[t]}^2}
 \le \sqrt{T_0T_1}\,
       (\sqrt{T_0T_1}\,p_{\gamma,\tau})^t,\qquad
 \prod_\tau p_{\gamma,\tau}\le b^{2C'fs},\quad
 p_{\gamma,\tau}\ge1 ,
\]
because \(|\tau\lambda|=\sqrt{T_1/T_0}\).
If two values differ, take absolute ideal norms in the
inclusion and use the triangle inequality place by place.
For the two sets of local heights their maxima are bounded
by their products, hence
\[
                     T^a\le2T^{t+1}b^{4C'st}.
\]
Here \(b=\exp(x^\varepsilon)\) and \(C'\) depends only on
the cutoffs. Since \(t\le s<d\) and
\(\log T=(34+o(1))d\log b\), increasing \(c_L\) by a fixed
amount excludes this inequality uniformly for large \(x\).
Thus \(B_t\), and hence \(\lambda\), is determined.
\end{proof}

\begin{theorem}[Recovery of every rational unipotent class in full dimension]
\leavevmode\label{thm:fullrational645}
Fix an original row pair and a cutoff exponent \(L>0\).
At a uniform precision \(a=c_Ld\), two compatible retained
terms of common positive dimension \(k\) with equal
\(r_a\) have the same rational unipotent double class.
The assertion holds at both terminal frequencies and
includes \(t=0\) and \(t=s\).
\end{theorem}
\begin{proof}
We first compare the decorated frames, allowing their
stabilizers to act on the positive quotients. The lemma gives
the same prefix blocks, block scalars and terminal ratio.
When \(I\ne0\), the flags on \(I\) have a one-dimensional jump
at each of the last \(s-t\) positions. In the two ordered
suffixes the successive ratios are all \(\lambda\), by the
successor relation in Lemma~\ref{lem:suffixdet556}. Choose a
basis of \(I\) respecting these jumps, with target leading
coefficients one, and change representatives by adding earlier
vectors. The suffixes then have the form \(v_i,\lambda v_i\).
Choose the target terminal representative in the common
negative line modulo \(I\) and use its \(\lambda\)-multiple
at the other end. Both changes add only vectors of their
respective outer spans. Together the two prefixes, this basis
of \(I\), and the common terminal representative form a basis
of \(Y_0+Y_1\), with Gram matrix determined by the stated
data and \(T_0\). For \(t=s\) the same construction just omits
the suffix, as in Proposition~\ref{prop:class236}; for \(t=0\)
it omits the prefixes. The linear map between the two
configurations matching these bases preserves the form,
including its radical \(I\). Witt's extension theorem
extends it over \(E\) to the ambient space. It respects
each ordered decorated flag.

Write \(F_b\) for the rational stabilizer of the decorated
frame on side \(b\), \(N_b\) for its row unipotent group,
and \(P_b=Y_b^\perp/X_b\) for its positive quotient
(transported back to the corresponding row). The kernel of
the action \(F_b\to U(P_b)\) is \(N_b\). In a Witt basis a
frame stabilizer is unitriangular on the outer flag, fixes
the prescribed negative vector modulo the outer space, and
acts on the remaining middle summand by its action on \(P_b\).
The dual outer action is forced by the pairing. Thus trivial
action on \(P_b\) gives precisely the row unipotents.

The preceding comparison gives
\(\gamma'=s_0\gamma s_1^{-1}\) with \(s_b\in F_b\).
Theorem~\ref{thm:allterminaltransports641} gives equality of
the two positive transports \(j:P_1\to P_0\). Consequently
the actions of \(s_0,s_1\) on those spaces correspond under
\(j\). Put \(W=Y_0+Y_1\) for the configuration of \(\gamma\).
Its radical is \(I\); the common space \(W^\perp/I\) maps
isometrically onto both positive quotients. This follows
from Lemma~\ref{lem:positiveadjugate641}, or from
\(W\cap Y_b^\perp=X_b\) and the dimensions there.
Choose an \(E\)-linear complement \(A\) to \(I\) in
\(W^\perp\). Its form identifies with the nondegenerate
quotient form. Every rational isometry of \(A\) extends by
the identity on \(A^\perp\) to a rational ambient isometry
fixing \(W\) pointwise. In particular choose \(w\) this way
to induce the action of \(s_0\) on \(P_0\). Then
\(\gamma^{-1}w\gamma\) induces that of \(s_1\) on \(P_1\).
It follows that \(s_0w^{-1}\in N_0\) and
\(s_1(\gamma^{-1}w\gamma)^{-1}\in N_1\). These two elements
give the same equation for \(\gamma'\) as \(s_0,s_1\).
This proves equality in \(N_0\backslash
\operatorname{Isom}(V_1,V_0)/N_1\).
\end{proof}

\begin{corollary}[The signed contribution of each rational class]
\leavevmode\label{cor:allrational645}
For the original ordinary block, group its signed operators
by rational unipotent double class and original marked index,
in both common dimensions. Write \(B_\xi\) for the resulting
operators. Instead summing first at the finite entries
\((i(\gamma),r_{c_Ld}(\gamma))\) and assigning an entry to
the class of a retained lift if one exists gives operators
\(\widehat B_\xi\) with
\[
           \sum_\xi\|B_\xi-\widehat B_\xi\|_{\rm HS}
                       \le 2\sigma_{ri}\exp(-LU_x).
\]
This holds on the same fields and at every allowed translate.
\end{corollary}
\begin{proof}
In dimension \(k\) the theorem makes the assignment unique.
In dimension \(k-1\) the marked index and reduction determine
the class by Theorem~\ref{thm:allterminaltransports641}.
That theorem also distinguishes the two dimensions. All lifts
meeting both cutoffs anywhere in their integration domain
are therefore assigned correctly, with their original phases,
marked columns and full integrals. An omitted or incorrectly
assigned lift has its entire domain outside the simultaneous
cutoffs and contributes its norm at most twice to the sum.
Apply Lemma~\ref{lem:productcut629} for the whole row pair.
\end{proof}

This identifies the rational classes in the finite table.
The original test adds its entries with their prescribed
signs. The estimate for that sum and its arithmetic
normalization required in Section~\ref{sec:missing} remains open.


\subsection{Mixed ordinary periods at prescribed input flags}
\label{sec:partialflags653}

Consider the original marked projection on one of the common
fields. We work with the terms of full common positive dimension
\(k\) in Section~\ref{sec:rowsums637}. Fix equivalent incoming
words and coherent unitary maps \(U_u\) for their rational
transports. Write \(D\) for the common dimension of the compact
representations. We keep complete orbital integrals at matched
height \(h(u,a)\le H\), where \(H\ge1\), and allow any finite
set of target rows. The two primes have product \(Q=g_1g_2\),
\(F=E_x^+\) and \(f=[F:\QQ]\). For a set \(W\) of places of
\(F\) over \(Q\), define
\[
                Q(W)=\prod_{v\in W}q_v^{1/f}.
\]
In particular \(Q(W)=Q\) for the full set, since both marks
split completely.

Use the finite groups \(G_v=\mathrm{GL}_k(\mathbf F_{q_v})\),
their upper unitriangular subgroups \(U_v\) and their upper
triangular subgroups \(B_v\). Write \(G,U,B\) for the respective
products. A raw marked coordinate is a pair
\((aU,Uh)\) as in Lemma~\ref{lem:whittakermatrix526}.
Its input flag at \(v\) is the left coset \(B_vh_v\).
We use the fixed marked frames of the input and target lattices
to compare residue matrices. Keeping \(B_vh_v\) fixed allows
all diagonal refinements of that flag.

\begin{lemma}[A closed product at a subset of the marks]
\leavevmode\label{lem:partialcycle653}
Consider a closed sequence of \(s\) contractions between input
rows. Each contraction passes from one input to the next
through a common target row and a common raw coordinate.
Choose one contributing orbital term on each of its two
edges. Suppose that, at each \(v\in W\), all the chosen
coordinates have the same input flag. Let \(g\) be the
composite rational isometry of the initial positive space.
If \(g\ne I\), then
\[
           Q(W)\ \le\ 2\operatorname{rank}_E(g-I)
                    \prod_{\text{edges }e}h(u_e,a_e)
                    \ \le\ 2kH^{2s}.
\]
The target rows and their output flags can vary from one
contraction to the next.
\end{lemma}
\begin{proof}
Two residue matrices in the support of one raw coordinate
belong to \(aUh\). If they are \(a j_1 h,a j_2 h\) with
\(j_1,j_2\in U\), the marking of the transport from the
second input to the first is \(h^{-1}j_1^{-1}j_2h\).
At \(v\) it lies in \(h_v^{-1}U_vh_v\). The latter group
depends only on \(B_vh_v\), because \(B_v\) normalizes \(U_v\).
All the relative markings in the closed sequence thus belong
to the same unitriangular group at each \(v\in W\). Their
product has trace \(k\), as does its inverse transpose at the
conjugate factor.

Let \(\mathfrak J\) be the index ideal of the successive
matched domain in its source lattice. As in
Lemma~\ref{lem:matchedtrace638}, at a finite place write
\(T,S\) for inclusion of that domain and for \(g\) on it,
respectively. Then
\[
  (\det T)(k-\operatorname{tr}g)
       =k\det T-\operatorname{tr}(S\operatorname{adj}T).
\]
This is integral everywhere. At each selected factor
\(\bar S=a_{\mathrm{comp}}\bar T\), with \(a_{\mathrm{comp}}\)
the composite marking. The reduction of the right side is
\((\det\bar T)(k-\operatorname{tr}a_{\mathrm{comp}})=0\),
including when the reduction of \(T\) is singular. The same
holds at the conjugate factor. Thus
\(\mathfrak J(k-\operatorname{tr}g)\) is divisible by both
primes of \(E\) over every \(v\in W\).

If \(g\ne I\), definiteness implies \(k-\operatorname{tr}g\ne0\).
At every embedding its absolute value is at most
\(2\operatorname{rank}_E(g-I)\), by the unitary eigenvalue
decomposition. The norm of the product of the selected prime
ideals is \(\prod_{v\in W}q_v^2\). Taking \(2f\)-th roots
in the norm inequality proves the first bound, using
\((N\mathfrak J)^{1/(2f)}\le\prod_e h(u_e,a_e)\). For this
last inequality impose the matched domain at each step. Its
additional index is at most the index for the individual
edge in the chosen orientation. There are \(2s\) edges.
\end{proof}

\begin{theorem}[The mixed Gram block on a prescribed tuple of flags]
\leavevmode\label{thm:mixedflags653}
Fix \(W\) with \(2kH^4<Q(W)\) and fix just the input flags
\(\varphi=(B_vh_v)_{v\in W}\). Sum over all their refinements,
over the flags at the other places, over all output coordinates
and over all the chosen target rows. For any two of the incoming
rows \(i,j\), the contribution \(K_{ij,\varphi}\) to the mixed
Gram block is zero or is a scalar times a single rational
transport \(U_{g_{ij,\varphi}}\) from \(j\) to \(i\).
This includes rows with no specified rational identification.

More precisely, let \(o\) run over the retained target rows
with their raw coordinates having the prescribed flags.
There are signed scalars \(z_{oi}\), obtained from the original
orbital coefficients and marked matrix entries, such that
\[
 c_{ij,\varphi}=\sum_o\overline{z_{oi}}z_{oj},\qquad
       \frac{\|K_{ij,\varphi}\|_{\mathrm{HS}}^2}{D}
                    =|c_{ij,\varphi}|^2 .
\]
On the diagonal one has \(K_{ii,\varphi}=c_{ii,\varphi}I\).
If \(z_{oi}z_{oj}\ne0\) for at least one \(o\), the rational
transport is unique among all the contributing contractions,
including when \(c_{ij,\varphi}=0\).
\end{theorem}
\begin{proof}
First fix \(o,i\). Since \(2kH^2<Q\),
Theorem~\ref{thm:disjointtransports638} shows that all terms
meeting this individual coordinate have the same rational
transport \(u_{oi}\). By Schur's lemma in full common
dimension their sum is \(z_{oi}U_{u_{oi}}\), where \(z_{oi}\)
retains the original signs, the Whittaker phases and the
normalization of the raw matrix. Entries with zero scalar
give zero operators. Absolute convergence within each entry
justifies forming its sum.

If \(o,o'\) both contribute, compare \(u_{oi}^{-1}u_{oj}\)
and \(u_{o'i}^{-1}u_{o'j}\). Composing the first with the
inverse of the second gives a closed sequence of two
contractions. Choose one original term on each edge with the
specified coordinate. The lemma forces the closed rational
isometry to be the identity. This proves uniqueness even for
different target rows and different refinements outside \(W\).
Coherence of the compact identifications gives
\(U_{u_{oi}}^*U_{u_{oj}}=U_{u_{oi}^{-1}u_{oj}}\).
Consequently the entire Gram block for these flags is the
displayed signed scalar times this one unitary. Its squared
Hilbert--Schmidt norm is \(D|c_{ij,\varphi}|^2\).
For \(i=j\) the relative transport is the identity in
every individual contraction.
\end{proof}

For example, with all places in \(W\) one may take
\(H=b^{4-\delta}\) for any fixed \(0<\delta<4\). With just
the places over the larger of \(g_1,g_2\) one may take
\(H=b^{2-\delta}\) for \(0<\delta<2\), and sum the other
mark completely. Indeed
\(\log Q=(16+o(1))\log b\), \(Q(W)\ge\sqrt Q\) in the second
case, and \(\log(2k)=o(\log b)\) uniformly on the interval.
More generally two contractions yielding different rational
transports can have equal input flags only at a set of
places \(W'\) satisfying \(Q(W')\le2kH^4\), by the same
closed sequence.

\begin{corollary}[Trace moments of the signed scalar matrix]
\leavevmode\label{cor:flagmoments653}
Let \(a\ge1\) be an integer with \(2kH^{2a}<Q(W)\), and keep
the same tuple of flags on \(W\). Write \(Z_\varphi\) for the
rectangular matrix of compact operators on the indicated
rows and coordinates, and \(z_\varphi=(z_{oi})\) for its
scalar entries as constructed in the proof. For each
\(1\le l\le a\),
\[
 \operatorname{tr}\bigl((Z_\varphi^* Z_\varphi)^l\bigr)
       =D\,\operatorname{tr}\bigl((z_\varphi^* z_\varphi)^l\bigr).
\]
\end{corollary}
\begin{proof}
The assumed inequality already gives \(2kH^2<Q\), so each
nonzero entry has the unique transport needed to define
\(z_{oi}\). Expand the trace on the left over closed walks
alternating \(l\) input rows and \(l\) target coordinates.
Each nonzero summand is its scalar coefficient times the
trace of the compact operator of the closed rational product.
The lemma with \(s=l\) forces this product to be the
identity. Its compact trace is \(D\). The scalar coefficients
with their signs sum to the trace on the right.
\end{proof}

For the full mixed block within the chosen height and common
dimension one must add the operators over \(\varphi\) before
taking a norm. In particular its squared norm still contains
terms from distinct flag tuples, and the sizes of the signed
scalars on a fixed tuple are arithmetic quantities. The local
separation above does not bound these quantities at the
normalization of the period comparison \eqref{eq:gap}.


\subsection{Relative input flags in the fourth period moment}
\label{sec:pairflags658}

The two contractions in the fourth moment of the original marked
projection have separate input coordinates. We compute their
weights on pairs of input flags and estimate the contribution of
the pairs far from opposition. Fix \(x\) on one of the common
intervals, with \(F=E_x^+\), \(f=[F:\QQ]\), and the two marked
primes \(g_1,g_2\). Use \(k,b,d,\sigma_{ri},U_x\) and the actual
blocks \(Z_{ri}\) of Section~\ref{sec:markedranks623}. There are
\(2f\) marked places. At one of them write \(q=q_v\),
\(G=\mathrm{GL}_k(\mathbb F_q)\), and let \(U,B\) be the upper
unitriangular and upper triangular groups. A flag is a left
coset \(C=Bh\), whose unitriangular subgroup is \(U_C=h^{-1}Uh\).

For \(s\in U_C\) let \(m_C(s)\) count the nonzero entries on the
first superdiagonal of \(hsh^{-1}\). This is independent of the
representative. Conjugation by \(U\) preserves these entries,
and diagonal conjugation scales them by nonzero elements. Define
\[
 b_C(s)=
 \begin{cases}
       (-1/(q-1))^{m_C(s)},&s\in U_C,\\
       0,&s\notin U_C .
 \end{cases}
\]
For a pair of flags let \(w\in\mathfrak S_k\) denote their relative
position, with the reverse permutation for opposite flags.
Put \(N=k(k-1)/2\), and let \(\ell(w)\) count inversions.

\begin{lemma}[The signed marked weight at prescribed flags]
\leavevmode\label{lem:flagsign658}
Fix four orbital terms with marking indices \(a_1,a_2,a_3,a_4\)
in the order of Theorem~\ref{thm:fourkernels622}. Put
\(s_v=a_{1,v}^{-1}a_{2,v}\), \(z_v=a_{3,v}^{-1}a_{4,v}\) in
the input frames of that expansion. For any set \(\mathcal D\) of
ordered pairs of full flag tuples the marked part of their
fourfold contraction on \(\mathcal D\) is exactly
\[
   \mathbb E_{C,C'}
   \left[{\bf1}_{\mathcal D}(C,C')\,
         \prod_{v\mid g_1g_2} b_{C_v}(s_v)b_{C'_v}(z_v)\right],
\]
where the two tuples and all their local entries are independent
uniform flags. In particular its absolute value is at most
\(\Pr(\mathcal D)\).
\end{lemma}
\begin{proof}
Work first at one place. In the raw matrix of
Lemma~\ref{lem:whittakermatrix526}, fix an input refinement
\(Uh\). An index \(a\) occurs at the unique output \(ah^{-1}U\)
with an entry of modulus one. Two indices occur at the same output
if \(ha_1^{-1}a_2h^{-1}\in U\). If \(a_j=g u_j h\) there,
their coefficient is \(\overline{\theta(u_1)}\theta(u_2)
=\theta(ha_1^{-1}a_2h^{-1})\). Dividing by the column norm
squared \([G:U]\) makes summation over \(Uh\) a uniform average.

Above \(Bh\), the refinements are precisely \(Uth\) for
diagonal \(t\). The ratios \(t_i/t_{i+1}\) are independent
uniform elements of \(\mathbb F_q^\times\), since every tuple of
ratios has \(q-1\) preimages. The prescribed Whittaker character
is a product of nontrivial additive characters on the simple
root entries. If such an entry is nonzero, its average over
the corresponding ratio is \(-1/(q-1)\), obtained by removing
the term at zero from a full additive character sum. A zero
entry contributes one. This gives \(b_C(s)\). The second pair
has its own summed coordinate even when the two target rows
agree, as follows by expanding the product of the two inner
products in the trace. The local groups form a direct product,
so all the stated independent averages follow. The bound uses
\(|b_C|\le1\). It also holds before averaging diagonal
refinements if one sums their absolute contributions instead.
\end{proof}

Define the total deficit from the opposite position by
\[
               J(C,C')=\sum_{v\mid g_1g_2}\bigl(N-\ell(w_v)\bigr),
               \qquad q_*=\min(g_1,g_2).
\]
\begin{lemma}[Distribution and common unitriangular subgroups]
\leavevmode\label{lem:flagdistribution658}
For all \(j\ge0\) one has
\[
       \Pr(J\ge j)\le
             q_*^{-j/2}(1-q_*^{-1/2})^{-2fk}.
\]
At one place the intersection \(U_C\cap U_{C'}\) has order
\(q^{N-\ell(w)}\). It is trivial for opposite flags.
\end{lemma}
\begin{proof}
Choose a basis adapted to both flags, so that their steps
are initial segments in the ordinary and the permuted orders.
Such a basis is obtained by echelon reduction of a basis of
the second flag with respect to the first, adding earlier
columns to later ones and applying upper triangular row
operations. In this basis the common unitriangular matrices have
one free entry per pair ordered increasingly in both orders.
There are \(N-\ell(w)\) such pairs. The stabilizer of the second
flag within the triangular stabilizer of the first also allows
the \(k\) diagonal entries to be arbitrary units. Therefore the
number of flags of position \(w\) to a fixed first flag is
\(|B|/((q-1)^kq^{N-\ell(w)})=q^{\ell(w)}\). In particular the
total number of flags is at least \(q^N\).

The expectation of \(q_*^{(N-\ell(w))/2}\) at this place is
thus at most \(\sum_w q^{-(N-\ell(w))/2}\). Reversing a
permutation interchanges length and deficit. Inserting the
largest letter into a permutation of \(i-1\) letters adds
any of \(0,\ldots,i-1\) inversions once. The sum is consequently
at most \(\prod_{i=1}^k(1-q^{-1/2})^{-1}\).
Multiply the \(2f\) bounds and apply Markov's inequality to
\(q_*^{J/2}\).
\end{proof}

\begin{theorem}[An absolute bound away from opposite input flags]
\leavevmode\label{thm:flagfourth658}
For every fixed \(L>0\) there is \(C_L>0\) such that if
\[
   j_x=\left\lceil C_L f k
             \left(d^3+\frac{k\log(2x)}{\log b}\right)\right\rceil ,
                         \qquad j_x/(fk^2)\longrightarrow0 ,
\]
the part of \(\operatorname{tr}(Z_{ri}^*Z_{rj}Z_{tj}^*Z_{ti})\)
with \(J\ge j_x\) has total absolute contribution at most
\[
              e^{-L U_x}\,\sigma_{ri}\sigma_{rj}\sigma_{tj}\sigma_{ti}.
\]
This includes the full range of orbital heights and both terminal
frequencies. For the squared norm of the mixed Gram block on
any finite set of target rows the corresponding bound is
\(e^{-L U_x}(\sum_r\sigma_{ri}\sigma_{rj})^2\).
\end{theorem}
\begin{proof}
For each orbital quartet apply the preceding probability bound.
The absolute value of the remaining operator trace is at most
the product of the four Hilbert--Schmidt norms, by
Cauchy--Schwarz. Lemma~\ref{lem:massranks623} bounds each
sum of norms over the original orbital terms by
\(\sigma_{ri}\exp(CU_x)\), also with absolute values taken
inside the real integrals. Hence the total over any subset
of these quartets and flag pairs with \(J\ge j_x\) is bounded by
the product of the four \(\sigma\)'s times
\[
    \exp(4CU_x)\,(1-q_*^{-1/2})^{-2fk}q_*^{-j_x/2}.
\]
Here \(U_x=fk(d^3\log b+k\log(2x))\).
The prefactor \((1-q_*^{-1/2})^{-2fk}\) tends to one, since
\(f,k\) are polynomially bounded and \(q_*\ge(kb)^8\),
with \(b=\exp(x^\varepsilon)\). Furthermore
\(j_x\log q_*/2\ge4C_L U_x\). Increasing \(C_L\) proves
the estimate. The asserted limit follows from \(d^3/k\to0\)
and \(\log(2x)/\log b\to0\) at the prescribed scales.
Summation of the bound over \(r,t\) gives the mixed Gram bound.
\end{proof}

If the flags coincide, or lie on a common panel, at a fixed
positive fraction of the \(2f\) places, then \(J\gg fk^2\)
as \(k\to\infty\). Indeed each such place has \(\ell(w)\le1\).
Thus those pairs are covered at every height. Theorem~\ref{thm:markedranks623}
also bounds the quartets with large total rank of the two
relative markings. Both omissions can be made together because
its proof uses the probability of unipotent collisions before
taking any Whittaker cancellations. The residual fourth moment
includes pairs with small \(J\) and small comparison rank.
Its comparison with the designated integral period is still
required in Section~\ref{sec:missing}.


\subsection{The first contributing flag position for a paired unipotent}
\label{sec:pairedunipotent659}

We next compute a signed contribution inside the range of small
deficit left by Theorem~\ref{thm:flagfourth658}. Work at one
actual marked place, of order \(q\). Keep the notation
\(C,C',w,N,\ell(w),b_C\) of Section~\ref{sec:pairflags658} and
put \(\delta=N-\ell(w)\). Let the two relative markings be
\(s=I+T\), \(z=I+P(T)\), where \(T\) is nilpotent and \(P\)
is a polynomial over \(\mathbb F_q\) with \(P(0)=0,\ P'(0)\ne0\).
This includes \(z=s^{-1}\), since the geometric series for the
inverse terminates. Such pairs occur in the repeated terms of a
squared Gram norm. Write \(\lambda_1,\ldots,\lambda_l\) for the
Jordan block sizes of \(T\), and let \(a_j\) count the blocks
of size \(j\). Put
\[
   r=k-l,\qquad m_\lambda=\sum_{i=1}^l\binom{\lambda_i}{2},
   \qquad A_\lambda=\frac{l!}{\prod_j a_j!}.
\]
Let \(X\) denote the set of full flags and \(\mathcal O_\lambda\)
the conjugacy class of nilpotent matrices of this type.

\begin{theorem}[Least possible deficit and its full signed contribution]
\leavevmode\label{thm:firstflag659}
If \(b_C(s)b_{C'}(z)\ne0\), then \(\delta\ge m_\lambda\).
The lower bound is attained. At equality the two flag orders
are obtained by reversing the order of consecutive blocks of
sizes \(\lambda_i\), preserving the order within each block.
For independent uniform flags the exact contribution at equality is
\[
   \mathbb E\!\left[{\bf1}_{\delta=m_\lambda} b_C(s)b_{C'}(z)\right]
     =\frac{A_\lambda q^{N-r}}{|X|(q-1)^r|\mathcal O_\lambda|}.
\]
\end{theorem}
\begin{proof}
The strictly triangular matrices in either flag form a nilpotent
algebra. The polynomial \(P\) has a composition inverse modulo
\(t^{k+1}\), obtained by equating successive coefficients using
its nonzero linear coefficient. Thus \(P(T)\) lies in one of
these algebras exactly when \(T\) does. If the product of the
two weights is nonzero, take a basis adapted to both flags.
Order its indices by \(i\prec j\) when \(i\) precedes \(j\) in
both flag orders. This partial order has \(\delta\) comparable
pairs, and \(T_{ij}\) can be nonzero only for \(i\prec j\).

For a column \(j\) let \(d_j\) be the number of its predecessors
and \(h_j\) the maximum number of strict steps in a chain ending
at \(j\). The vector \(T^n e_j\) is zero for \(n>h_j\), by the
matrix product formula. Hence
\[
 \sum_{n\ge1}\operatorname{rank}T^n
    \le\sum_{n\ge1}\#\{j:h_j\ge n\}
    =\sum_j h_j\le\sum_j d_j=\delta .
\]
On a Jordan block of size \(b\), the sum on the left contributes
\(\sum_{n\ge1}\max(b-n,0)=\binom b2\). This proves the lower
bound. For equality one must have \(h_j=d_j\) for every \(j\),
so all predecessors of each element form a chain. The argument
applied to \(T^{\mathsf t}\) in the reversed orders gives the
same property for successors. In the diagram of immediate
predecessors and successors, every element therefore has at most
one of each. Its components are chains.

Each such chain is consecutive in each flag order. Indeed an
element of a different component lying between two chain elements
in one order would have both of its comparisons with them reversed
in the other order. That would reverse the comparison of the
two chain elements themselves. Distinct chain blocks must thus
appear in opposite order. If the chain lengths are \(b_i\),
\(T\) is block diagonal with a strictly upper triangular matrix
on each block. At power \(n\) the rank on such a block is at
most \(\max(b_i-n,0)\). Equality of the total sum of ranks with
\(\delta=\sum_i\binom{b_i}{2}\) forces equality on every block
at every power. In particular the rank on a block of size \(b_i\)
is \(b_i-1\), so it is one Jordan block. Conversely, put the
Jordan chains of \(T\) in consecutive blocks and reverse the
order of the blocks in the second flag. Then both weights are
nonzero and there are exactly \(m_\lambda\) comparable pairs.

There are \(A_\lambda\) relative permutations of the required
form, one for each distinct ordering of the block sizes. For a
fixed one, the number of common strictly triangular matrices
of type \(\lambda\) is
\((q-1)^r q^{m_\lambda-r}\). Within each block all entries just
above its diagonal must be nonzero and suffice for regular
nilpotence, since their product is the extreme entry in the
highest possible power. Every entry farther above the diagonal
is free. A uniformly conjugated flag pair of the specified
position sees the fixed \(T\) uniformly over \(\mathcal O_\lambda\),
by transitivity and equal stabilizer sizes. For each surviving
matrix the two weights have product \((q-1)^{-2r}\). In fact
there are \(r\) nonzero simple entries in each flag, and replacing
\(T\) by \(P(T)\) scales these entries by \(P'(0)\).
The probability of each relative permutation is
\(q^{N-m_\lambda}/|X|\) by Lemma~\ref{lem:flagdistribution658}.
Multiplication of these factors proves the formula.
\end{proof}

\begin{theorem}[Cancellation between the first position and the later ones]
\leavevmode\label{thm:flagcancel659}
Put \(D_j=\prod_{i=1}^j(q^i-1)\) and
\(\lambda'_h=\#\{i:\lambda_i\ge h\}\).
The unrestricted signed weight is \((D_l/D_k)^2\).
The ratio of the preceding theorem's contribution to the
unrestricted weight is
\[
 A_\lambda q^{e_\lambda}
 \frac{(1-q^{-1})^l
       \prod_j\prod_{i=1}^{a_j}(1-q^{-i})}
      {\prod_{i=1}^{l}(1-q^{-i})^2},\qquad
 e_\lambda=\sum_{h\ge2}\lambda'_h(\lambda'_h-1).
\]
The factor multiplying \(A_\lambda q^{e_\lambda}\) is
\(\exp(O(k/q))\) with an absolute implied constant for \(q\ge2\).
\end{theorem}
\begin{proof}
The two uniform flags are independent. Averaging a single weight
gives the normalized induced Whittaker character by
Lemma~\ref{lem:flagsign658}. Theorem~\ref{thm:char125} gives
\((-1)^r D_l/D_k\) for each, since \(P(T)\) has the same kernel
dimension as \(T\). Now \(|X|=D_k/(q-1)^k\) by counting ordered
bases and their triangular stabilizers. Let \(Z(T)\) be the
invertible centralizer of \(T\). The orbit has size
\(q^N D_k/|Z(T)|\). Substitution in Theorem~\ref{thm:firstflag659}
gives the ratio \(A_\lambda q^{-r}(q-1)^l|Z(T)|/D_l^2\).

To evaluate the centralizer, view the space as the module
\(\bigoplus_j(\mathbb F_q[t]/t^j)^{a_j}\), with \(t\) acting
as \(T\). Homomorphisms from a single length-\(u\) summand to a
length-\(v\) summand have dimension \(\min(u,v)\), by choosing
the image of its generator. The full endomorphism algebra
therefore has dimension \(\sum_h(\lambda'_h)^2\).
Modulo \(t\), the induced matrix is triangular by groups of
equal size, since a map from a shorter to a longer summand has
image divisible by \(t\). Its diagonal blocks are unrestricted
independent matrices of sizes \(a_j\). An endomorphism is
invertible precisely when each such block is invertible. Indeed
surjectivity modulo \(t\) gives surjectivity on the whole module
by repeatedly substituting the image plus \(t\) times the module
and using nilpotence. Thus
\[
 |Z(T)|=q^{\sum_h(\lambda'_h)^2}
                     \prod_j\prod_{i=1}^{a_j}(1-q^{-i}).
\]
Using \(\lambda'_1=l\) and \(\sum_{h\ge2}\lambda'_h=r\)
now gives the displayed exponent. Taking logarithms of the
remaining factors gives \(O(k/q)\), since at most \(4k\)
terms occur and \(|\log(1-q^{-i})|\le2/q\) for \(q\ge2\).
\end{proof}

\begin{corollary}[Two small ranks on the marked fields]
\leavevmode\label{cor:smalljordan659}
For a transvection paired with its inverse the opposite flags
give zero. The first nonzero contribution, at deficit one, is
\[
       \frac{(k-1)q^{N-1}}
       {|X|(q^k-1)(q^{k-1}-1)} .
\]
Its ratio to the whole signed weight is
\((k-1)\exp(O(k/q))\). For two blocks of size two and the
remaining blocks of size one, the first contribution occurs at
deficit two and its ratio is \(\binom{k-2}{2}q^2\exp(O(k/q))\).
\end{corollary}
\begin{proof}
For a transvection, there are \(k-1\) possible positions of
the block of size two. The orbit of its nilpotent part consists
of \(v\otimes\phi\), with \(v,\phi\ne0\), \(\phi(v)=0\).
Each such matrix has \(q-1\) presentations, giving orbit size
\((q^k-1)(q^{k-1}-1)/(q-1)\). Apply the two theorems. For the
second partition, \(l=k-2\), \(\lambda'_2=2\), and there are
\(\binom{k-2}{2}\) possible placements of the two larger blocks.
\end{proof}

All formulas multiply over the actual marked places for prescribed
local positions. In particular, when every pair has the indicated
polynomial relation, the least possible total deficit is
\(\sum_v m_{\lambda_v}\) and its signed weight is the product
of the local values in Theorem~\ref{thm:firstflag659}.
At the marked primes \(q\ge(kb)^8\), the error in the logarithm
of the product of the ratio factors is \(O(fk/q_*)=o(1)\).
For the two-block example, the signed sum above the first
position must cancel almost all of its positive contribution.
The remaining orbital coefficients and the mixed pairs with
independent relative markings still enter the period comparison.


\subsection{Transvections on opposite and neighboring flags}
\label{sec:oppositetrans660}

We compute the second flag average when its relative marking
is a transvection. The first relative marking can be arbitrary.
All the flags, places and signed weights are those of
Lemma~\ref{lem:flagsign658}. At one place write \(V=\mathbb F_q^k\)
with \(k\ge2\) and \(z=I+v\otimes\phi\), where \(v,\phi\ne0\)
and \(\phi(v)=0\). Only the incident pair consisting of the
line \(L=\langle v\rangle\) and the hyperplane \(H=\ker\phi\)
will enter the weights. At a fixed full flag \(C\) define
\[
        a=\min\{i:L\subset C_i\},\qquad
        b=\min\{i:C_i\not\subset H\}.
\]
In an adapted basis \(a\) is the last nonzero vector coordinate,
and \(b\) is the first nonzero covector coordinate. They are
distinct because equality would give \(\phi(v)\ne0\).
The matrix belongs to \(U_C\) if and only if \(a<b\).
In this case it has one nonzero simple entry if \(b=a+1\),
and none if \(b\ge a+2\).

\begin{theorem}[Averaging a transvection on opposite flags]
\leavevmode\label{thm:oppositeaverage660}
Fix \(C\) and choose \(C'\) uniformly from its opposite flags.
The entire signed average at \(C'\) is
\[
 \mathbb E\big[b_{C'}(z)\mid C,\ C'\text{ opposite to }C\big]
    =\begin{cases}
      0,&a<b,\\[2pt]
      -1/((q-1)q^{k-1}),&a=b+1,\\[2pt]
      -q^{-k},&a\ge b+2 .
    \end{cases}
\]
\end{theorem}
\begin{proof}
Take \(C\) to be the standard flag. A triangular basis change
depending on \(L,H\) sends these to the coordinate line \(e_a\)
and coordinate hyperplane \(e_b^*=0\). One can see this by
choosing the basis in \(H\) at every step except \(b\), and
choosing the vector spanning \(L\) at step \(a\). Each choice
extends the preceding basis since these are the pivot steps.
A triangular change preserves the uniform measure on opposite
flags. Such flags are the translates of the reversed flag by
the upper unitriangular group, acting simply transitively by
Lemma~\ref{lem:flagdistribution658}. Keeping the reversed flag
fixed and moving the incident pair by the inverse element
therefore makes it uniform on the set of pairs of the given
pivots. Indeed the triangular group acts transitively there by
the basis construction and its diagonal subgroup fixes the
coordinate pair. This proves the same transitivity for the
unitriangular subgroup.

Normalize the vector's coordinate at \(a\) and the covector's
coordinate at \(b\) both to one. Their scalings have no effect
on flag weights. If \(a<b\), a vector supported in positions
at most \(a\) and a covector supported in positions at least
\(b\) cannot give a strictly upper triangular transvection in
the reversed order. If \(a>b\), there are \(q^{a+k-b-2}\)
normalized incident pairs. There are initially \(a-1+k-b\)
free entries, and \(\phi(v)=0\) determines the vector entry
at \(b\) uniquely.

Write \(A\) for the first nonzero vector coordinate and \(B\)
for the last nonzero covector coordinate in the ordinary order.
The surviving pairs satisfy \(B<A\). Incidence is automatic
there. Write \(n=a-b,\ u=a-A,\ v'=B-b\).
For fixed \(u,v'\ge0\), \(u+v'\le n-1\), their count is
\(f(u)f(v')\), with \(f(0)=1\) and \(f(j)=(q-1)q^{j-1}\)
for \(j\ge1\). Indeed the two ends of each support are
specified and all its interior coordinates are free. The
simple-entry weight is one if \(u+v'\le n-2\) and
\(-1/(q-1)\) if \(u+v'=n-1\).

Let \(S_n\) be the resulting signed sum. In the ring of formal
power series put \(F(X)=\sum_{j\ge0}f(j)X^j=(1-X)/(1-qX)\).
The counting rule just proved gives
\[
  \sum_{n\ge1}S_n X^{n-1}
    =\left(\frac{X}{1-X}-\frac1{q-1}\right)F(X)^2
    =-\frac{F(X)}{q-1}.
\]
Hence \(S_1=-1/(q-1)\) and \(S_n=-q^{n-2}\) for \(n\ge2\).
Division by \(q^{k+n-2}\), the number of incident pairs of
these pivots, gives the assertion.
\end{proof}

\begin{theorem}[A recurrence for every relative position]
\leavevmode\label{thm:pivotrecursion660}
For \(w\in\mathfrak S_k\) put
\(F_w(a,b)=\mathbb E[b_{C'}(z)\mid C,\operatorname{pos}(C,C')=w]\),
where \(w\) lists the reference coordinates in the order of the
second flag in a simultaneous basis. All \(k(k-1)\) values of
\(F_w\) can be obtained from the opposite-flag formula in
\(N-\ell(w)\) successive steps. If the current permutation is
\(s_i w\) of length \(\ell(w)+1\), the step is
\[
                   F_w=\big(qA_i-(q-1)I\big)F_{s_i w}.
\]
Here \(A_i\) is the following operator on the pivot pairs. Form a
word with the first symbol at position \(a\), the last symbol at
\(b\) and the middle symbol elsewhere, ordering the three symbols
as first, middle, last. If symbols \(i,i+1\) are equal, keep
the current pair. If they increase, exchange them. If they
decrease, take the exchanged pair with probability \(1/q\) and
the original with probability \((q-1)/q\). Apply this rule to
average the function.
\end{theorem}
\begin{proof}
Changing only the \(i\)-dimensional space in a full flag amounts
to choosing one of the other \(q\) lines in
\(C_{i+1}/C_{i-1}\). The partial flag \(L\subseteq H\) induces
a filtration on this quotient, by projecting its intersections
with \(C_{i+1}\). The two symbols specify the two graded
dimensions. For distinct symbols there is exactly one line
in the earlier grade and \(q\) others. Thus the unique line
gives the increasing order, while all other lines give the
decreasing order. Equal symbols are unchanged. This proves
the stated rule, also when \(k=2\) and \(L=H\).

For a function on flags let \(T_w\) average it over the flags
of relative position \(w\) from the argument, so \(A_i=T_{s_i}\)
on functions of these pivots. If \(\ell(s_i w)=\ell(w)+1\), then
\(A_iT_w=T_{s_i w}\). To verify this count, an endpoint of
position \(s_i w\) determines a unique intermediate flag in
the \(i\)-panel of the starting flag which has shorter position
\(w\) to that endpoint. The filtration on the panel quotient
from the endpoint full flag distinguishes a unique earlier
line. Taking it exchanges the inverted adjacent reference
coordinates. Conversely, from the shortened flag all \(q\)
other lines lengthen the position. The number of endpoints
is \(q^{\ell(w)+1}\), by Lemma~\ref{lem:flagdistribution658},
which also verifies the normalization. On each \(q+1\) element
panel the neighbor average is \((\mathbf1\mathbf1^{\mathsf t}-I)/q\),
so its inverse is \(qA_i-(q-1)I\). Inverting the composition
gives the asserted step. Every permutation can be lengthened
to the reverse order by successive left multiplications by
adjacent exchanges. Reversing that sequence takes exactly
\(N-\ell(w)\) steps. This proves the recurrence.
\end{proof}

\begin{corollary}[Absolute bounds for signed rank-one contractions]
\leavevmode\label{cor:weightedtrans660}
For an arbitrary first relative marking \(s\), its local paired
coefficient conditional on position \(w\) equals
\(\mathbb E_C[b_C(s)F_w(a_z(C),b_z(C))]\). If \(s\) also is a
transvection, put \(D=q^k-1\). For every set \(\mathcal E\) of
pairs of flags one has
\[
 \left|\mathbb E[{\bf1}_{\mathcal E} b_C(s)b_{C'}(z)]\right|
          \le (2k-3)^2D^{-2}.
\]
Conditional on the opposite position, the absolute value is at
most \((2k-3)qD^{-2}/(q-1)\), uniformly in the two incident
pairs. No relation between the markings is required.
\end{corollary}
\begin{proof}
For one transvection the mean signed value is \(-1/D\), by
Theorem~\ref{thm:char125}. In its flag average the negative
values occur exactly at \(b=a+1\). There are
\(M=(q^k-1)(q^{k-1}-1)/(q-1)^2\) incident line--hyperplane
pairs, by choosing a line and then a hyperplane in its quotient.
For fixed \(a,b=a+1\), the two normalized coordinate lists
have disjoint supports and contain \(q^{k-2}\) possibilities.
Thus the mean of the magnitude of the negative part is
\[
  B_-=\frac{(k-1)q^{k-2}}{(q-1)M}
     =\frac{k-1}{D}\frac{q^{k-2}(q-1)}{q^{k-1}-1}
     \le\frac{k-1}{D}.
\]
Consequently \(\mathbb E|b_C(s)|=2B_--1/D\le(2k-3)/D\).
Drop the indicator and use the independence of \(C,C'\) for
the first bound. For the second, average over the opposite flags
first using Theorem~\ref{thm:oppositeaverage660}. Its absolute
conditional mean is at most \(q^{1-k}/(q-1)\), and
\(D q^{1-k}<q\). The conditional expression for general \(s,w\)
follows by the same order of averaging.
\end{proof}

For an orbital quartet on the original fields these calculations
apply independently at the marked places. For example, at \(h\)
places with two transvections, any restriction on the full pair
of flag tuples costs at most
\(\prod_{v}(2k-3)^2/(q_v^k-1)^2\), with the product over these
\(h\) places and trivial upper bound one elsewhere. This bound
uses the signs after averaging the diagonal refinements on each
flag. With specified relative positions, the recurrence instead
keeps their signed coefficients. It leaves the archimedean
orbital weights in the same fourth moment to be estimated.


\subsection{Real flag angles in the untruncated signed period}
\label{sec:realflag664}

In a retained row put \(d=s+1\), \(n=2s+k+1\),
\(f=[E_x^+:\QQ]\) and \(b=\exp(x^\varepsilon)\).
Keep the physical Bessel section \(W_l\), the original finite
projection and \(T_l\in\{1,T\}\) at either end of each orbital
integral in Theorem~\ref{thm:fourkernels622}. Its normalized
squared row measure on \(N_l\backslash G_\infty\) is
\(d\mu_l=w_l^{-2}\|W_l(y)\|^2\,dy\).
For \(1\le i\le d\) let \(a_{l,i,\tau}(y)\) be the successive
Gram--Schmidt lengths of the physical outer frame followed by
its terminal vector, in the reference positive norm after applying
\(y_\tau^{-1}\). Adding earlier vectors leaves these lengths
unchanged. Set \(V=bn^4\). Write \(\Omega_l\) for the region
where \(a_{l,i+1,\tau}<V a_{l,i,\tau}\) for every \(i<d\)
and every \(\tau\).

\begin{lemma}[Transfer from a row measure to all its orbital integrals]
\leavevmode\label{lem:rowmeasure664}
Let \(A_0,A_1\) be measurable subsets of the two integral row
quotients. In every supported coset take the union of the
conditions \(y\in A_0\) and \(\gamma^{-1}y\in A_1\).
Let \(I_{\mathrm{bad}}\) be the sum over these regions of the
integrals of
\(|c_{01}|\,\|W_0(y)\|\,\|W_1(\gamma^{-1}y)\|\).
Divide by \(\lambda_{01}=|c_{01}|w_0w_1\) for a nonzero pair.
With \(D=\sum_{i=1}^d i(n-i+1)\), \(\nu=4D+1\),
\(q_0=4n^2\), and \(P_\nu\) as in
Theorem~\ref{thm:twoenergies592}, one has
\[
 \frac{I_{\mathrm{bad}}}{\lambda_{01}}
 \le P_\nu e^{Cfs\nu\log b}
       \big\{\mu_0(A_0)^{\vartheta}+\mu_1(A_1)^{\vartheta}\big\},
 \qquad \vartheta=\tfrac12(1-\nu/q_0)\ge\tfrac38
\]
for all sufficiently large \(x\) on the original intervals.
The same bound holds on any measurable subregion of this union
and for any subset of the supported cosets.
\end{lemma}
\begin{proof}
Write \(v_0,v_1\) for the two translated heights \(V'_l\) of
Proposition~\ref{prop:fulltail621} at their respective row
arguments. In its unit-frequency coordinates the same proposition
and Theorem~\ref{thm:moment286} give
\(\int (V'_l)^{q_0}\,d\mu_l\le e^{Cfsq_0\log b}\).
The change of majorant costs \(R^{fsq_0}\), where
\(\log R=O(\log x)\), which is included in this bound.
H\"older's inequality gives
\[
 \int_{A_l}(V'_l)^\nu\,d\mu_l
     \le \mu_l(A_l)^{1-\nu/q_0}e^{Cfs\nu\log b}.
\]
The exponent is positive since \(\nu\le n^2\) for sufficiently
large \(x\), as \(d^2/n\to0\).

First impose only \(y\in A_0\). In the disjoint union of all
orbital domains, apply Cauchy--Schwarz after multiplying the
first section norm by \(v_0^{\nu/2}v_1^{-\nu/2}\) and the second
by its reciprocal. Keep the indicator in the first factor.
Unfolding the first square counts all translates from the second
row with weight \(v_1^{-\nu}\). Their sum is at most \(P_\nu\)
pointwise by Theorem~\ref{thm:twoenergies592}. Reversing rows
bounds the second square by the same count times the unrestricted
weighted row integral. The resulting bound is
\[
 P_\nu |c_{01}|w_0w_1
 \left(\int_{A_0}(V'_0)^\nu\,d\mu_0
             \int(V'_1)^\nu\,d\mu_1\right)^{1/2}.
\]
Apply the moment estimates, interchange the rows for the second
condition, and add. All unfolded quantities are nonnegative,
so the estimate applies to the stated restrictions.
\end{proof}

\begin{theorem}[Simultaneous real localization of the signed moments]
\leavevmode\label{thm:conemoments664}
Replace every original orbital domain by its intersection with
\(y\in\Omega_0,\ \gamma^{-1}y\in\Omega_1\). Write
\(Z_{ri}^{\Omega}\) for the resulting marked block and
\(\sigma_{ri}\) for the normalization of
Lemma~\ref{lem:massranks623}. There is a uniform \(c_1>0\)
such that with \(\delta=e^{-c_1bn^2}\),
\begin{align*}
 \big|\|Z_{ri}\|_{\rm HS}^2-\|Z_{ri}^{\Omega}\|_{\rm HS}^2\big|
   &\le\delta\sigma_{ri}^2,\\
 \left|\operatorname{tr}(Z_{ri}^*Z_{rj}Z_{tj}^*Z_{ti})
 -\operatorname{tr}\big((Z_{ri}^{\Omega})^*Z_{rj}^{\Omega}
                       (Z_{tj}^{\Omega})^*Z_{ti}^{\Omega}\big)\right|
   &\le\delta\sigma_{ri}\sigma_{rj}\sigma_{tj}\sigma_{ti}.
\end{align*}
These errors also bound the absolute integral contributions
where at least one of the incident row arguments is outside its
region. For the squared norm of the mixed Gram block summed over
any finite set of target rows the error is at most
\(\delta(\sum_r\sigma_{ri}\sigma_{rj})^2\).
\end{theorem}
\begin{proof}
By Theorem~\ref{thm:localmass156},
\(\mu_l(\Omega_l^c)\le e^{-c b n^2}\) at prime frequency.
The normalized squared measure on the integral unipotent quotient
gives the same probabilities for these conditions as the measure
on the real Bessel quotient. Indeed the pointwise norm and the
conditions are invariant under the real decorated stabilizer,
and its lattice fibre has the constant volume used in \(w_l^2\).
The estimate holds also at unit frequency. Under the change
of real section of Proposition~\ref{prop:scale151}, the
corresponding prime-frequency measure is its pushforward and
all \(d\) physical Gram--Schmidt lengths acquire the common
factor \(\sqrt T\). The successive ratios are preserved. The
argument applies to every unit vector of the prescribed product
lowest compact type.

Apply Lemma~\ref{lem:rowmeasure664} with \(A_l=\Omega_l^c\).
The logarithms of \(P_\nu\), of \(e^{Cfs\nu\log b}\), and of
\(\lambda_{ri}/\sigma_{ri}=\sqrt{D_rD_i}\) for a nonzero
row pair are bounded by a fixed polynomial in \(x\).
For \(P_\nu\) this follows from its formula in
Theorem~\ref{thm:twoenergies592}, since \(\log T=O(d\log b)\).
For the positive column dimensions it follows from the Weyl
product in Lemma~\ref{lem:massranks623}. As \(b=\exp(x^\varepsilon)\),
each logarithm is \(o(bn^2)\). It follows that the sum of
absolute integrals discarded in a single block is at most
\(\sigma_{ri}e^{-c'bn^2}\) for some \(c'>0\).

In a product of two or four orbital integrals, bound the norm of
each finite column by one and the absolute value of the trace
of the real rank-one factors by the product of their norms.
Use the preceding estimate for any one failed incident condition.
For the other blocks use the unrestricted absolute bound
\(\sigma_{ri}e^{C U_x}\) of Lemma~\ref{lem:massranks623},
where \(U_x=fk(d^3\log b+k\log(2x))=o(bn^2)\).
Adding the finitely many choices of a failed condition and
decreasing \(c_1\) gives both inequalities and the stated
absolute form. Summation over the two target indices gives the
mixed Gram estimate. Structurally zero blocks contribute zero.
\end{proof}

\begin{corollary}[The cross pairing of the surviving real flags]
\leavevmode\label{cor:crossangles664}
Fix a supported coset and \(1\le j\le d\). Use the cross
determinant \(D_j\) of Proposition~\ref{prop:transverseheight582}.
At a real place \(\tau\), transport the first \(j\) vectors from
both physical frames to the common reference space at the two
orbital arguments. Let \(G_{l,j,\tau}\) be their positive
Hermitian Gram matrices, and \(H_{01,j,\tau}\) their cross matrix
for the original indefinite form. Put
\[
 A_{j,\tau}=
 \frac{|\det H_{01,j,\tau}|}
 {\sqrt{\det G_{0,j,\tau}\det G_{1,j,\tau}}}.
\]
Then \(0\le A_{j,\tau}\le1\). Write
\(A_j=(\prod_{\tau}A_{j,\tau})^{1/(fj)}\). On the common region
of the theorem, simultaneously for every prefix,
\[
    A_j \le
       |N_{E_x/\QQ}D_j|^{1/(2fj)}
               \frac{V^{\,2d-j-1}}{\sqrt{T_0T_1}} .
\]
In particular, fix any \(B>0\) and put \(Q=g_1g_2\). At a pair
with at least one prime frequency, the part on which, for some \(j\),
\[
       |N_{E_x/\QQ}D_j|^{1/(2fj)}\le Q^B,\qquad A_j>b^{-14d}
\]
has the absolute error bound of the theorem whenever it occurs
in a second or fourth moment. If both terminal norms are
\(-T\), the same holds with \(b^{-31d}\). All prefix lengths
and cosets in the indicated union cost one error bound.
\end{corollary}
\begin{proof}
In the reference positive metric the original Hermitian form
is represented by a unitary involution. After orthonormalizing
the two lists, their cross matrix represents an orthogonal
projection composed with this involution on one of the two
\(j\)-planes. All its singular values are at most one, and
the modulus of its determinant is \(A_{j,\tau}\). Thus these
numbers measure the product of the cosines of the principal
angles with the image of the other plane under the involution.
The cross determinant before orthonormalization has absolute
value \(|\tau D_j|\), since simultaneous transport is an isometry
for the indefinite form. These statements also hold when that
determinant vanishes.

The last Gram--Schmidt length is at least \(\sqrt{T_l}\).
Subtracting any vector in the preceding isotropic span from the
terminal vector preserves its Hermitian norm \(-T_l\), whose
absolute value is bounded by its squared positive norm.
Telescoping the defining ratios on \(\Omega_l\) gives
\(a_{l,i,\tau}\ge\sqrt{T_l}V^{-(d-i)}\) for all \(i\le d\).
The positive \(j\)-volume is the product of the first \(j\)
lengths. After multiplying the two volumes, the exponent of
\(V^{-1}\) is \(2\sum_{i=1}^j(d-i)=j(2d-j-1)\). This proves
the bound at each place with \(|\tau D_j|^{1/j}\) on the right
after taking \(j\)-th roots. The product of \(|\tau D_j|\)
over the real places of the maximal real field is
\(|N_{E_x/\QQ}D_j|^{1/2}\), giving the assertion.

Finally \(\log V=(1+o(1))\log b\),
\(\log Q=(16+o(1))\log b\), and
\(\log T=(34+o(1))d\log b\) on these intervals by
Theorem~\ref{thm:prime209}. When \(\sqrt{T_0T_1}\ge\sqrt T\)
and the norm test holds, the bound on \(A_j\) is at most
\(b^{-(15-o(1))d}\), uniformly in \(j\). With both frequencies
prime the exponent becomes \(-(32-o(1))d\).
Thus each indicated region lies in the single complement
estimated in Theorem~\ref{thm:conemoments664}.
\end{proof}

The calculation retains the signed ordinary integrals on the
original fields. It localizes low cross determinants at prime
frequency to small normalized real pairings. The integrals
on that remaining region and their comparison with the
prescribed integral theta map in Section~\ref{sec:missing}
still require bounds.

\clearpage
\section{References}\label{nav:bibliography-XXII}
\begingroup\def\section*#1{}\begin{thebibliography}{99}\small
\bibitem{XXII-XVII-koornwinder610} T. H. Koornwinder, \emph{The addition formula for Jacobi polynomials II. The Laplace type integral representation and the product formula}, Mathematisch Centrum, first printing 1972, second printing 1976. \href{https://ir.cwi.nl/pub/12597}{Author's report}.
\bibitem{XXII-XVII-IV-serre211} J.-P. Serre, \emph{Bounds for the orders of the finite subgroups of \(G(k)\)}, in \emph{Group Representation Theory}, EPFL Press, 2007. \href{https://arxiv.org/pdf/1011.0346}{Lecture I, Section 1.1, Theorem 1}.
\end{thebibliography}\endgroup
\end{document}
